% appendix-video-compression.tex
% Standalone appendix — include from retrospective.tex via:
%   \appendix
%   % appendix-info-theory.tex
% Standalone appendix — include from retrospective.tex via:
%   \appendix
%   % appendix-info-theory.tex
% Standalone appendix — include from retrospective.tex via:
%   \appendix
%   % appendix-info-theory.tex
% Standalone appendix — include from retrospective.tex via:
%   \appendix
%   \input{appendix-info-theory}
%
% Requires in the preamble: amsmath, amssymb, tabularx, booktabs

\section{Information-Theoretic View}
\label{app:info-theory}

This appendix applies information-theoretic measures to the real-time
web communication stack described in this document. We quantify channel
capacity, source coding efficiency, and effective information rates
across three time horizons: today (March~2026), the near term
(2027--2028), and the medium term (2029--2031). The analysis draws on
the standards, codec parameters, and infrastructure projections
developed in the main text.

\subsection{Channel Capacity of the Transport Layer}

The Shannon-Hartley theorem gives the theoretical maximum information
rate for a channel with additive white Gaussian noise:
\begin{equation}
  C = B \log_2\!\left(1 + \frac{S}{N}\right) \quad \text{[bits/s]}
  \label{eq:shannon}
\end{equation}
where $B$ is the channel bandwidth in Hz and $S/N$ is the
signal-to-noise ratio.

While the physical-layer capacity of a broadband link is governed by
(\ref{eq:shannon}), the \emph{effective} capacity available to the
application is reduced by protocol overhead, congestion control, and
transport inefficiencies. We define the effective application-layer
throughput as:
\begin{equation}
  R_{\text{eff}} = C_{\text{link}} \cdot \eta_{\text{transport}}
    \cdot (1 - \rho_{\text{overhead}})
  \label{eq:reff}
\end{equation}
where $C_{\text{link}}$ is the link-layer capacity,
$\eta_{\text{transport}} \in (0,1]$ is the transport efficiency factor
(capturing congestion control, retransmission, and multiplexing
behavior), and $\rho_{\text{overhead}} \in [0,1)$ is the fractional
protocol overhead (headers, encryption, framing).

\subsubsection{Head-of-Line Blocking and Multiplexed Capacity}

Consider $n$ independent streams multiplexed over a single connection.
Under TCP (HTTP/2), a single lost packet blocks all $n$ streams. The
expected throughput per stream under loss rate $p$ is bounded by:
\begin{equation}
  R_{\text{TCP}}^{(i)} \leq \frac{R_{\text{eff}}}{n}
    \cdot (1 - p)^{n-1}
  \label{eq:hol-tcp}
\end{equation}
because a loss in \emph{any} of the $n$ streams stalls all others. Under
QUIC (RFC~9000), streams are independent at the transport layer, so:
\begin{equation}
  R_{\text{QUIC}}^{(i)} \leq \frac{R_{\text{eff}}}{n}
    \cdot (1 - p)
  \label{eq:hol-quic}
\end{equation}
The ratio of effective per-stream throughput is:
\begin{equation}
  \frac{R_{\text{QUIC}}^{(i)}}{R_{\text{TCP}}^{(i)}}
    = \frac{1}{(1-p)^{n-2}}
  \label{eq:hol-ratio}
\end{equation}
\excursus{Excursus: Why $n = 8$?}{In a typical multi-participant video
conference, each peer contributes multiple concurrent streams: a camera video
track, an audio track, and often a screen-share or data channel. A four-person
call with two media tracks each produces $n = 8$ multiplexed streams over a
single connection --- the reference value used throughout this section. Larger
meetings (e.g.\ 10 participants with simulcast layers) push $n$ higher,
amplifying the HOL blocking effect further.}

For $n = 8$ streams (a typical multi-track video conference) at
$p = 0.01$ (1\% loss), this gives a QUIC advantage factor of
$1/(0.99)^6 \approx 1.062$, a modest 6.2\% improvement. At $p = 0.05$
(5\% loss, common on mobile or congested links), the factor is
$1/(0.95)^6 \approx 1.34$, a 34\% throughput advantage per stream.

\begin{figure}[ht]
\centering
\begin{tikzpicture}[>=Stealth]

% --- Shaded region (drawn FIRST so curves appear on top) ---
\fill[orange!8] (2, 0) rectangle (10, 5.3);
\node[font=\tiny\itshape, text=black!40, anchor=north east] at (9.8, 5.2)
  {mobile / congested};
\node[font=\tiny\itshape, text=black!40, anchor=north west] at (0.2, 5.2)
  {fixed broadband};

% --- Grid lines ---
\foreach \y in {1, 2, 3, 4, 5} {
  \draw[gray!20] (0, \y) -- (10, \y);
}

% --- Axes ---
\draw[->, thick] (0, 0) -- (10.5, 0);
\draw[->, thick] (0, 0) -- (0, 5.5);

% --- Axis titles (positioned to clear tick labels) ---
\node[font=\small] at (5, -0.85) {Packet loss rate $p$};
\node[font=\small, rotate=90] at (-0.9, 2.75) {QUIC advantage factor};

% --- X-axis labels ---
\foreach \x/\lab in {0/0, 2/0.02, 4/0.04, 6/0.06, 8/0.08, 10/0.10} {
  \draw (\x, -0.1) -- (\x, 0.1);
  \node[font=\tiny, anchor=north] at (\x, -0.15) {\lab};
}

% --- Y-axis labels ---
\foreach \y/\lab in {0/1.0, 1/1.1, 2/1.2, 3/1.3, 4/1.4, 5/1.5} {
  \draw (-0.1, \y) -- (0.1, \y);
  \node[font=\tiny, anchor=east] at (-0.15, \y) {\lab};
}

% --- n=4 curve: 1/(1-p)^2 ---
\draw[very thick, blue!70]
  (0, 0) .. controls (2.5, 0.5) and (5, 1.1) ..
  (7.5, 2.0) .. controls (8.5, 2.4) and (9.5, 2.9) ..
  (10, 3.1);
\node[font=\scriptsize, blue!70, anchor=south west] at (9.0, 3.0)
  {$n = 4$};

% --- n=8 curve: 1/(1-p)^6 ---
\draw[very thick, orange!80!black]
  (0, 0) .. controls (1.5, 0.5) and (3, 1.7) ..
  (4, 2.5) .. controls (5, 3.5) and (5.5, 4.0) ..
  (6, 4.7);
\node[font=\scriptsize, orange!80!black, anchor=south] at (5.5, 4.8)
  {$n = 8$};

% --- Reference points ---
\fill[orange!80!black] (1.0, 0.62) circle (2.5pt);
\node[font=\tiny, orange!80!black, anchor=south west] at (1.1, 0.62)
  {6.2\%};

\fill[orange!80!black] (5.0, 3.7) circle (2.5pt);
\node[font=\tiny, orange!80!black, anchor=west] at (5.1, 3.7)
  {34\%};

\fill[blue!70] (1.0, 0.20) circle (2.5pt);

% --- Annotation ---
\node[font=\scriptsize\itshape, text=black!60, align=center,
  anchor=north] at (5.0, -1.0)
  {QUIC recovers per-stream throughput closer to $C/n$ (Shannon fair share)\\
   by eliminating head-of-line blocking (eq.~\ref{eq:hol-ratio}).};

\end{tikzpicture}
\caption{QUIC advantage factor $1/(1-p)^{n-2}$ over TCP for multiplexed
  streams. At 5\% loss with $n=8$ streams, QUIC delivers 34\% more
  per-stream throughput by maintaining independence at the transport layer.}
\label{fig:hol-blocking}
\end{figure}

\subsubsection{Connection Establishment Latency}

Let $\tau$ denote the one-way latency (half the round-trip time). The
information-carrying capacity during connection establishment is zero;
the handshake represents pure overhead. The time to first useful byte
is:
\begin{align}
  T_{\text{TCP+TLS\,1.3}} &= 2\tau_{\text{TCP}} + \tau_{\text{TLS}}
    = 3\tau \quad \text{(1-RTT TLS~1.3)} \label{eq:tcp-setup} \\
  T_{\text{QUIC}} &= \tau \quad \text{(1-RTT)}
    \label{eq:quic-setup} \\
  T_{\text{QUIC,\,0-RTT}} &= 0 \quad \text{(resumption)}
    \label{eq:quic-0rtt}
\end{align}
The QUIC 0-RTT case (\ref{eq:quic-0rtt}) eliminates the handshake
entirely for resumed connections. The information lost during
handshake dead time, relative to steady-state capacity, is:
\begin{equation}
  I_{\text{lost}} = C_{\text{link}} \cdot T_{\text{setup}}
  \label{eq:ilost}
\end{equation}
For a 120~Mbps link with $\tau = 25$~ms (typical fixed broadband in
2026), the TCP+TLS~1.3 handshake wastes $120 \times 10^6 \times 0.075
= 9 \times 10^6$ bits (1.125~MB), while QUIC~1-RTT wastes only
$3 \times 10^6$ bits (375~KB).

\subsection{Source Coding Efficiency}

The rate-distortion function $R(D)$ defines the minimum bit rate
required to represent a source at distortion level $D$. For a
Gaussian source with variance $\sigma^2$ under mean squared error
distortion:
\begin{equation}
  R(D) = \frac{1}{2} \log_2 \frac{\sigma^2}{D}
    \quad \text{[bits/sample]}
  \label{eq:rate-distortion}
\end{equation}

\excursus{Excursus: Why Gaussian source?}{The Gaussian memoryless model is used
because it yields the highest (worst-case) $R(D)$ among all sources of equal
variance. Any codec that approaches the Gaussian bound on real video is
therefore doing at least as well as the bound predicts --- the true $R(D)$ for
spatially and temporally correlated video lies below the Gaussian curve, making
it a conservative benchmark.}

No practical codec achieves $R(D)$;\footnote{The $R(D)$ curve in
Figure~\ref{fig:rate-distortion} is the Gaussian memoryless source bound
(equation~\ref{eq:rate-distortion}), which is an upper bound on the true
$R(D)$ for non-Gaussian sources such as video.  Additionally, $R(D)$ assumes
zero excess-distortion probability, whereas practical codecs permit a
frame-level excess-distortion probability on the order of $10^{-6}$ to
$10^{-8}$.  Both factors allow codec operating points measured on real content
to approach or marginally cross the Gaussian bound.} real codecs operate at
some rate $R_{\text{codec}} > R(D)$. We define the coding efficiency gap as:
\begin{equation}
  \Delta = \frac{R_{\text{codec}} - R(D)}{R(D)}
  \label{eq:coding-gap}
\end{equation}

The codecs in the WebRTC stack can be ordered by their empirical
proximity to $R(D)$ for typical video content at conversational
resolutions (720p, 30~fps). Using published Bj\o ntegaard Delta (BD)
rate measurements:
\begin{equation}
  R_{\text{H.264}} > R_{\text{VP9}} > R_{\text{AV1}}
    > R(D)
  \label{eq:codec-ordering}
\end{equation}

Specifically, the empirical rate savings at equivalent PSNR are
approximately:
\begin{align}
  R_{\text{VP9}}  &\approx 0.65 \cdot R_{\text{H.264}}
    \quad (\sim\!35\% \text{ savings}) \label{eq:vp9-savings} \\
  R_{\text{AV1}}  &\approx 0.70 \cdot R_{\text{VP9}}
    \approx 0.46 \cdot R_{\text{H.264}}
    \quad (\sim\!30\% \text{ over VP9})
  \label{eq:av1-savings}
\end{align}

\begin{figure}[ht]
\centering
\begin{tikzpicture}[>=Stealth]

% --- Axes ---
\draw[->, thick] (0, 0) -- (10.5, 0);
\draw[->, thick] (0, 0) -- (0, 6.5);

% --- Axis titles ---
\node[font=\small] at (5, -0.85) {Distortion $D$};
\node[font=\small, rotate=90] at (-0.9, 3.25) {Rate [Mbps]};

% --- Y-axis rate labels ---
\foreach \y/\lab in {1.0/0.5, 2.0/1.0, 3.0/1.5, 4.0/2.0, 5.0/2.5, 6.0/3.0} {
  \draw (-0.1, \y) -- (0.1, \y);
  \node[font=\tiny, anchor=east] at (-0.15, \y) {\lab};
}

% --- R(D) Shannon bound curve ---
\draw[very thick, green!60!black]
  (1.0, 5.8) .. controls (2.0, 3.5) and (3.5, 2.0) ..
  (5.0, 1.3) .. controls (6.5, 0.8) and (8.0, 0.55) ..
  (9.5, 0.4);
\node[font=\scriptsize, green!60!black, anchor=north west] at (8.5, 0.55)
  {$R(D)$ bound};

% --- Reference distortion line (720p/30fps operating point) ---
\draw[thin, gray, dashed] (4.0, 0) -- (4.0, 6.2);
\node[font=\tiny, gray, anchor=south] at (4.0, 6.2) {720p/30fps};

% --- R(D) value at reference distortion ---
\fill[green!60!black] (4.0, 1.55) circle (2pt);

% --- Codec operating points at reference distortion ---
% H.264: ~2.5 Mbps
\draw[thick, dashed, red!70!black] (0.3, 4.8) -- (9.5, 4.8);
\fill[red!70!black] (4.0, 4.8) circle (3pt);
\node[font=\scriptsize, red!70!black, anchor=west] at (9.6, 4.8) {H.264};

% VP9: ~1.6 Mbps (0.65 * H.264)
\draw[thick, dashed, blue!70] (0.3, 3.3) -- (9.5, 3.3);
\fill[blue!70] (4.0, 3.3) circle (3pt);
\node[font=\scriptsize, blue!70, anchor=west] at (9.6, 3.3) {VP9};

% AV1: ~1.15 Mbps (0.46 * H.264)
\draw[thick, dashed, orange!80!black] (0.3, 2.3) -- (9.5, 2.3);
\fill[orange!80!black] (4.0, 2.3) circle (3pt);
\node[font=\scriptsize, orange!80!black, anchor=west] at (9.6, 2.3) {AV1};

% --- Gap annotations (braces) ---
% H.264 gap
\draw[decorate, decoration={brace, amplitude=4pt, mirror},
  thick, red!50!black]
  (3.6, 1.55) -- (3.6, 4.8)
  node[midway, anchor=east, font=\scriptsize, xshift=-5pt]
  {$\Delta_{\text{H.264}}$};

% AV1 gap
\draw[decorate, decoration={brace, amplitude=4pt},
  thick, orange!60!black]
  (4.4, 1.55) -- (4.4, 2.3)
  node[midway, anchor=west, font=\scriptsize, xshift=5pt]
  {$\Delta_{\text{AV1}}$};

% --- Annotation ---
\node[font=\scriptsize\itshape, text=black!60, align=center,
  anchor=north] at (5.0, -0.5)
  {The Shannon rate-distortion bound $R(D)$ is the theoretical minimum.\\
   AV1 narrows the gap $\Delta$ to ${\sim}0.4$ by 2030 (eq.~\ref{eq:coding-gap}).};

\end{tikzpicture}
\caption{Codec operating rates vs.\ the Shannon rate-distortion bound $R(D)$
  at 720p/30fps. Each codec operates above the theoretical minimum; the gap
  $\Delta$ (eq.~\ref{eq:coding-gap}) shrinks from H.264 through VP9 to AV1,
  approaching but never reaching the bound.}
\label{fig:rate-distortion}
\end{figure}

\subsubsection{Codec Entropy and WebCodecs}

The WebCodecs API (W3C Working Draft) exposes per-frame codec access,
allowing measurement and control of the instantaneous encoding rate.
For a video frame $f_t$ at time $t$, the encoder output has empirical
entropy:
\begin{equation}
  \hat{H}(f_t) = -\sum_{s \in \mathcal{S}}
    p(s \mid f_t) \log_2 p(s \mid f_t)
  \label{eq:frame-entropy}
\end{equation}
where $\mathcal{S}$ is the set of symbols emitted by the entropy coder
(CABAC for H.264/AV1, arithmetic coding for VP9). WebCodecs enables
applications to observe $|E(f_t)|$ (the encoded frame size in bits)
directly, which serves as a practical upper bound on
$\hat{H}(f_t)$:
\begin{equation}
  \hat{H}(f_t) \leq |E(f_t)| \leq R_{\text{codec}} / \text{fps}
  \label{eq:entropy-bound}
\end{equation}

\subsection{Effective Information Rate by Time Horizon}

We now combine transport and source coding measures to estimate the
effective information rate for a representative real-time video
session: a single 720p/30fps video call with one audio channel.

\excursus{Excursus: Why 720p?}{Throughout this appendix, 720p/30fps is used as
a fixed reference for cross-era comparison, not as a prediction that 720p
remains the dominant conferencing resolution. The bandwidth surplus quantified
below is precisely what enables migration to 1080p (feasible by ${\sim}$2028
with AV1 hardware encode at today's 720p bit rates) and 4K (feasible by
${\sim}$2030). The 720p baseline makes the efficiency gains legible: the same
task gets cheaper, freeing capacity for higher quality or more streams.}

Define the \emph{net information rate} as:
\begin{equation}
  I_{\text{net}} = R_{\text{codec}} \cdot \eta_{\text{transport}}
    \cdot (1 - \rho_{\text{overhead}})
    \cdot (1 - p_{\text{loss}})
  \label{eq:inet}
\end{equation}

\begin{table}[ht]
\centering
\small
\renewcommand{\arraystretch}{1.4}
\begin{tabularx}{\textwidth}{lXXX}
\toprule
\textbf{Parameter} & \textbf{March 2026} & \textbf{2027--2028} &
  \textbf{2029--2031} \\
\midrule
Median link capacity &
  120~Mbps &
  $\sim$150~Mbps &
  200--300~Mbps \\
Dominant codec (WebRTC) &
  VP9 / H.264 &
  VP9 / AV1 (HW) &
  AV1 (HW default) \\
Codec rate (720p/30fps) &
  1.5--2.5~Mbps &
  1.0--1.8~Mbps &
  0.7--1.2~Mbps \\
$\eta_{\text{transport}}$ &
  0.92 (QUIC) &
  0.93 &
  0.94 \\
$\rho_{\text{overhead}}$ (SRTP+QUIC) &
  0.06 &
  0.05 &
  0.04 \\
Typical loss $p$ &
  0.01--0.03 &
  0.01--0.02 &
  0.005--0.01 \\
$I_{\text{net}}$ (720p video) &
  1.3--2.2~Mbps &
  0.9--1.6~Mbps &
  0.6--1.1~Mbps \\
$I_{\text{net}}$ as \% of $C_{\text{link}}$ &
  1.1--1.8\% &
  0.6--1.1\% &
  0.2--0.6\% \\
\bottomrule
\end{tabularx}
\caption{Effective information rate for a 720p/30fps WebRTC video stream.
  The declining ratio $I_{\text{net}} / C_{\text{link}}$ reflects both
  improved codec efficiency and growing link capacity, freeing bandwidth
  for higher resolutions, additional streams, or concurrent data
  channels.}
\label{tab:inet}
\end{table}

The declining ratio $I_{\text{net}} / C_{\text{link}}$ in
Table~\ref{tab:inet} is the central quantitative observation: by 2030, a
720p video call consumes less than 0.5\% of available link capacity. The
surplus capacity enables qualitative changes --- 4K video conferencing,
multi-stream layouts, or concurrent data channels via WebTransport ---
that are quantitatively infeasible in 2026.

\begin{figure}[ht]
\centering
\begin{tikzpicture}[>=Stealth]

% --- Axes ---
\draw[->, thick] (0, 0) -- (11, 0);
\draw[->, thick] (0, 0) -- (0, 5.5);

% --- Axis title ---
\node[font=\small, rotate=90] at (-0.9, 2.75) {Mbps};

% --- Y-axis labels ---
\foreach \y/\lab in {0/0, 1/50, 2/100, 3/150, 4/200, 5/250} {
  \draw (-0.1, \y) -- (0.1, \y);
  \node[font=\tiny, anchor=east] at (-0.15, \y) {\lab};
}

% --- Group: March 2026 ---
% C_link = 120 Mbps -> y = 2.4
\fill[blue!12] (0.8, 0) rectangle (2.4, 2.4);
\draw[blue!50, thick] (0.8, 0) rectangle (2.4, 2.4);
% I_net = 1.75 Mbps (midpoint) -> y = 0.035
\fill[orange!60] (0.8, 0) rectangle (2.4, 0.07);
\draw[orange!80!black, thick] (0.8, 0) rectangle (2.4, 0.07);
\node[font=\footnotesize, anchor=south] at (1.6, 2.5) {120};
\node[font=\tiny, orange!80!black] at (1.6, -0.4) {1.5\%};
\node[font=\small, anchor=north] at (1.6, -0.7) {2026};

% --- Group: 2027--2028 ---
% C_link = 150 Mbps -> y = 3.0
\fill[blue!12] (4.0, 0) rectangle (5.6, 3.0);
\draw[blue!50, thick] (4.0, 0) rectangle (5.6, 3.0);
% I_net = 1.25 Mbps (midpoint) -> y = 0.025
\fill[orange!60] (4.0, 0) rectangle (5.6, 0.05);
\draw[orange!80!black, thick] (4.0, 0) rectangle (5.6, 0.05);
\node[font=\footnotesize, anchor=south] at (4.8, 3.1) {150};
\node[font=\tiny, orange!80!black] at (4.8, -0.4) {0.8\%};
\node[font=\small, anchor=north] at (4.8, -0.7) {2027--28};

% --- Group: 2029--2031 ---
% C_link = 250 Mbps -> y = 5.0
\fill[blue!12] (7.2, 0) rectangle (8.8, 5.0);
\draw[blue!50, thick] (7.2, 0) rectangle (8.8, 5.0);
% I_net = 0.85 Mbps (midpoint) -> y = 0.017
\fill[orange!60] (7.2, 0) rectangle (8.8, 0.034);
\draw[orange!80!black, thick] (7.2, 0) rectangle (8.8, 0.034);
\node[font=\footnotesize, anchor=south] at (8.0, 5.1) {250};
\node[font=\tiny, orange!80!black] at (8.0, -0.4) {0.3\%};
\node[font=\small, anchor=north] at (8.0, -0.7) {2029--31};

% --- Legend ---
\node[anchor=west, font=\footnotesize] at (0.5, 5.8) {%
  \tikz[baseline=-0.5ex]{\fill[blue!12] (0,0) rectangle (0.3,0.3);
    \draw[blue!50, thick] (0,0) rectangle (0.3,0.3);}
  ~$C_{\text{link}}$ (channel capacity) \qquad
  \tikz[baseline=-0.5ex]{\fill[orange!60] (0,0) rectangle (0.3,0.3);
    \draw[orange!80!black, thick] (0,0) rectangle (0.3,0.3);}
  ~$I_{\text{net}}$ (720p stream)};

% --- Annotation ---
\node[font=\scriptsize\itshape, text=black!60, align=center,
  anchor=north] at (5.5, -1.2)
  {A single 720p call uses a vanishing fraction of Shannon-derived\\
   channel capacity, freeing bandwidth for 4K, multi-stream, and data channels.};

\end{tikzpicture}
\caption{Channel capacity $C_{\text{link}}$ vs.\ effective information rate
  $I_{\text{net}}$ for a single 720p/30fps stream across time horizons. The
  orange bar (stream rate) is barely visible against the blue bar (capacity),
  illustrating the growing surplus as codecs improve and bandwidth increases.}
\label{fig:inet-capacity}
\end{figure}

\subsection{Security Entropy and Key Space}

The DTLS-SRTP security architecture (RFCs~8827, 3711) protects all
WebRTC media. Information-theoretic security requires that the key
entropy exceed the information content of the protected material.

For a DTLS~1.3 handshake with ECDHE on P-256, the key agreement
provides:
\begin{equation}
  H_{\text{key}} = 128 \text{~bits of security}
  \label{eq:key-entropy}
\end{equation}

The SRTP master key is 128 bits (AES-128-GCM), giving an exhaustive
search cost of $2^{128}$ operations. The total information content of a
one-hour 720p video call at 2~Mbps is:
\begin{equation}
  I_{\text{call}} = 2 \times 10^6 \times 3600
    = 7.2 \times 10^9 \text{~bits}
    \approx 2^{32.7} \text{~bits}
  \label{eq:call-content}
\end{equation}

The ratio of key entropy to protected content:
\begin{equation}
  \frac{H_{\text{key}}}{I_{\text{call}}}
    = \frac{128}{32.7} \approx 3.91
  \label{eq:security-margin}
\end{equation}

This ratio remains comfortably above~1 even for much longer sessions.
For a continuous 24-hour 4K AV1 stream at 15~Mbps:
\begin{equation}
  I_{\text{stream}} = 15 \times 10^6 \times 86400
    = 1.296 \times 10^{12} \text{~bits}
    \approx 2^{40.2} \text{~bits}
\end{equation}
yielding $H_{\text{key}} / I_{\text{stream}} \approx 3.18$, still
well above the threshold.

\subsubsection{End-to-End Encryption via Encoded Transforms}

The RTCRtpScriptTransform API (Encoded Transforms) enables
application-layer E2EE within the WebRTC pipeline. From an
information-theoretic perspective, this introduces a second encryption
layer with its own key, creating a \emph{nested channel}:
\begin{equation}
  X \xrightarrow{\;K_{\text{E2EE}}\;}
  E_1(X) \xrightarrow{\;K_{\text{SRTP}}\;}
  E_2(E_1(X))
  \label{eq:nested}
\end{equation}

The mutual information between the ciphertext and plaintext, given
neither key, satisfies:
\begin{equation}
  I\!\left(X;\, E_2(E_1(X))\right) \leq
    \min\!\big(H(K_{\text{E2EE}})^{-1},\;
    H(K_{\text{SRTP}})^{-1}\big)
    \cdot I_{\text{call}}
  \label{eq:mutual-e2ee}
\end{equation}

In practice, with $H(K_{\text{E2EE}}) = H(K_{\text{SRTP}}) = 128$
bits, this quantity is negligible ($\ll 1$ bit for any practical
session length), confirming that the nested encryption provides
information-theoretic confidentiality well beyond computational
security guarantees.

Today (March~2026), Encoded Transforms remain Chrome-only. By
2027--2028 (Section~8.3), cross-browser standardization will make
this nested channel universally available. By 2029--2031, SFrame
for group calls will extend this to $n$-party sessions, where the
mutual information between any non-member and the group plaintext
remains negligible:
\begin{equation}
  I\!\left(X;\, E(X) \;\middle|\; K_i \notin
    \{K_1, \ldots, K_n\}\right) \approx 0
  \label{eq:sframe}
\end{equation}

\subsection{WebTransport: Unreliable Datagrams and Channel Capacity}

WebTransport (RFC~9297) introduces unreliable datagrams to browser-server
communication. For latency-sensitive applications (game state, sensor
telemetry), the relevant measure is not throughput but
\emph{freshness-weighted information rate}. Define the value of a
datagram as exponentially decaying with delivery latency $\delta$:
\begin{equation}
  V(\delta) = e^{-\lambda \delta}
  \label{eq:freshness}
\end{equation}
where $\lambda > 0$ is an application-specific decay rate. The
effective information rate under unreliable delivery is:
\begin{equation}
  I_{\text{WT}} = R_{\text{send}} \cdot (1 - p_{\text{loss}})
    \cdot \mathbb{E}\!\left[e^{-\lambda \delta}\right]
  \label{eq:iwt}
\end{equation}

Under reliable delivery (TCP/WebSocket), a retransmission delays
subsequent data, so:
\begin{equation}
  I_{\text{WS}} = R_{\text{send}}
    \cdot \mathbb{E}\!\left[e^{-\lambda(\delta + p \cdot 2\tau)}\right]
  \label{eq:iws}
\end{equation}

The ratio $I_{\text{WT}} / I_{\text{WS}}$ exceeds 1 when:
\begin{equation}
  (1 - p) \cdot e^{\lambda \cdot p \cdot 2\tau} > 1
  \quad \Longleftrightarrow \quad
  \lambda > \frac{-\ln(1 - p)}{2 p \tau}
  \label{eq:wt-advantage}
\end{equation}

For $p = 0.02$ and $\tau = 25$~ms, this gives $\lambda > 20.2$~s$^{-1}$
(half-life $< 34$~ms).

\excursus{Excursus: Why $p = 0.02$ and $\tau = 25$\,ms?}{These values represent
median fixed-broadband conditions in 2026 (Table~\ref{tab:inet}). On mobile or
congested links ($p \approx 0.05$, $\tau \approx 50$\,ms), the threshold
$\lambda^*$ drops to ${\sim}10$\,s$^{-1}$, broadening WebTransport's advantage
to applications with data half-lives up to ${\sim}70$\,ms. The fixed-broadband
reference is thus conservative: mobile scenarios favor unreliable delivery even
more strongly.}

Any application where data older than
$\sim$30~ms is useless benefits from WebTransport's unreliable
mode over WebSocket --- precisely the regime of interactive gaming, live
cursor tracking, and real-time sensor feeds.

As of March~2026, Safari does not support WebTransport
(Section~5), limiting coverage to $\sim$82\%. By 2027--2028
(Section~8.2), universal support will allow applications to rely on
(\ref{eq:iwt}) without fallback paths.

\begin{figure}[ht]
\centering
\begin{tikzpicture}[>=Stealth]

% --- Axes ---
\draw[->, thick] (0, 0) -- (10.5, 0);
\draw[->, thick] (0, 0) -- (0, 5.5);

% --- Axis titles (positioned to clear tick labels) ---
\node[font=\small] at (5, -0.75)
  {Freshness decay rate $\lambda$ [s$^{-1}$]};
\node[font=\small, rotate=90] at (-1.1, 2.75)
  {$I_{\text{WT}} / I_{\text{WS}}$};

% --- X-axis labels ---
\foreach \x/\lab in {0/0, 1.67/10, 3.33/20, 5.0/30, 6.67/40, 8.33/50, 10/60} {
  \draw (\x, -0.1) -- (\x, 0.1);
  \node[font=\tiny, anchor=north] at (\x, -0.15) {\lab};
}

% --- Y-axis labels (0.96 to 1.06, each 0.02 = 1 unit) ---
\foreach \y/\lab in {0/0.96, 1/0.98, 2/1.00, 3/1.02, 4/1.04, 5/1.06} {
  \draw (-0.1, \y) -- (0.1, \y);
  \node[font=\tiny, anchor=east] at (-0.15, \y) {\lab};
}

% --- Ratio = 1.0 line (y=2 maps to ratio 1.0) ---
\draw[thin, gray, dashed] (0, 2) -- (10, 2);

% --- Shade WebTransport advantage region ---
% lambda* = 20.2 -> x = 3.37
\fill[green!10] (3.37, 2) rectangle (10, 5.3);

% --- I_WT/I_WS curve: 0.98*exp(lambda*0.001) ---
% Y range: 0.96-1.06 (each 0.02 = 1 unit). lambda=0: y=1, lambda*=20.2: y=2, lambda=60: y=4.05
\draw[very thick, orange!80!black]
  (0, 1.0) .. controls (1.5, 1.3) and (2.5, 1.7) ..
  (3.37, 2.0) .. controls (4.5, 2.4) and (6, 2.8) ..
  (8, 3.4) .. controls (9, 3.7) and (9.5, 3.85) ..
  (10, 4.05);

% --- Threshold line ---
\draw[thick, dashed, red!60!black] (3.37, 0) -- (3.37, 5.2);
\node[font=\scriptsize, red!60!black, anchor=south, rotate=90] at (3.17, 3.5)
  {$\lambda^* = 20.2$ s$^{-1}$};

% --- Application regime labels ---
\node[font=\scriptsize\itshape, text=blue!60, anchor=north] at (1.5, 5.3)
  {video conf};
\node[font=\scriptsize\itshape, text=green!50!black, anchor=north] at (7.0, 5.3)
  {gaming / cursor / sensors};

% --- Mark advantage/disadvantage ---
\node[font=\tiny, text=red!50!black, anchor=north east] at (3.2, 1.8)
  {WebSocket wins};
\node[font=\tiny, text=green!50!black, anchor=north west] at (3.5, 4.5)
  {WebTransport wins};

% --- Annotation ---
\node[font=\scriptsize\itshape, text=black!60, align=center,
  anchor=north] at (5.0, -1.0)
  {For $p = 0.02$, $\tau = 25$~ms. Applications where data older than\\
   ${\sim}$30~ms is stale gain from unreliable delivery (eq.~\ref{eq:wt-advantage}).};

\end{tikzpicture}
\caption{Freshness-weighted information rate ratio $I_{\text{WT}} / I_{\text{WS}}$
  as a function of decay rate $\lambda$. Above the threshold $\lambda^*$,
  WebTransport's unreliable datagrams deliver more useful information per
  unit time than WebSocket's reliable delivery, because retransmission
  wastes capacity on stale data.}
\label{fig:freshness}
\end{figure}

\subsection{Projections Summary and Practical Implications}

\begin{table}[ht]
\centering
\small
\renewcommand{\arraystretch}{1.4}
\begin{tabularx}{\textwidth}{lXXX}
\toprule
\textbf{Measure} & \textbf{March 2026} & \textbf{2027--2028} &
  \textbf{2029--2031} \\
\midrule
Channel capacity $C_{\text{link}}$ &
  120~Mbps &
  150~Mbps &
  250~Mbps \\
HOL blocking factor $(1\!-\!p)^{n-2}$ at $n\!=\!8$ &
  0.94 ($p\!=\!0.01$) &
  0.94 &
  0.97 ($p\!=\!0.005$) \\
Source coding gap $\Delta$ &
  VP9: $\sim$0.8 &
  AV1 HW: $\sim$0.5 &
  AV1 opt: $\sim$0.4 \\
$I_{\text{net}}$ per 720p stream &
  1.3--2.2~Mbps &
  0.9--1.6~Mbps &
  0.6--1.1~Mbps \\
Max 720p streams in $C_{\text{link}}$ &
  $\sim$55--90 &
  $\sim$95--165 &
  $\sim$230--415 \\
E2EE availability &
  Chrome only &
  Cross-browser &
  SFrame $n$-party \\
WebTransport coverage &
  82\% &
  $\sim$100\% &
  100\% \\
Freshness threshold $\lambda^*$ &
  20.2~s$^{-1}$ &
  $\sim$18~s$^{-1}$ &
  $\sim$15~s$^{-1}$ \\
\bottomrule
\end{tabularx}
\caption{Information-theoretic measures across time horizons. $\Delta$
  is the coding efficiency gap (\ref{eq:coding-gap}), $I_{\text{net}}$
  is the net information rate (\ref{eq:inet}), and $\lambda^*$ is the
  freshness decay threshold above which WebTransport dominates WebSocket
  (\ref{eq:wt-advantage}).}
\label{tab:projections}
\end{table}

The trajectory is clear: the combination of improving source coding
(AV1 hardware), expanding channel capacity (broadband growth), and
maturing transport protocols (QUIC universal, WebTransport universal)
drives the \emph{information cost per unit of communicated experience}
steadily downward. The remaining question is not whether the
Shannon capacity suffices --- it already does, by orders of magnitude
--- but how efficiently the protocol stack converts that capacity into
delivered, fresh, secure information at the application layer. By
2030, the gap between theoretical capacity and realized throughput will
be the narrowest it has ever been for browser-based real-time
communication.

\subsubsection{What the Theory Means for Applications}

The equations and figures in this appendix are not purely academic. Each
theoretical result maps to a concrete design decision that application
developers face today and over the next five years.

\textbf{Codec selection determines bandwidth cost, not quality.}
The rate-distortion analysis (Section~A.2, Figure~\ref{fig:rate-distortion})
shows that AV1 operates at roughly 46\% of H.264's bit rate for equivalent
perceptual quality at 720p. In practice, this means a WebRTC application
switching from H.264 to AV1 can either halve its bandwidth consumption at
the same quality or substantially increase resolution (e.g., 720p to 1080p)
within the same bandwidth envelope. By 2030, with AV1 hardware encoding
ubiquitous and the coding gap $\Delta$ approaching 0.4, further codec
improvements yield diminishing returns --- the stack will be operating close
enough to the Shannon rate-distortion bound that the next generation of
efficiency gains will come from transport and architecture, not codecs alone.

\textbf{QUIC's advantage is situational but decisive on bad networks.}
The head-of-line blocking analysis (Section~A.1.1,
Figure~\ref{fig:hol-blocking}) shows that QUIC's per-stream throughput
advantage over TCP is modest on good networks (6\% at 1\% loss) but
substantial on degraded ones (34\% at 5\% loss). The practical implication
is that applications targeting mobile users, emerging markets, or congested
enterprise networks --- precisely the environments where real-time
communication is hardest --- benefit disproportionately from QUIC-based
transports. For a video conference with 8 participants on a cellular
connection experiencing 3--5\% packet loss, the difference between QUIC and
TCP is the difference between usable and degraded video.

\textbf{Connection latency is a solved problem for returning users.}
The handshake analysis (Section~A.1.2) shows that QUIC 0-RTT eliminates
connection establishment overhead entirely for resumed connections. For a
WebTransport application that a user opens daily --- a collaboration tool,
a live dashboard, a recurring meeting client --- the time to first useful
byte drops to zero. The 1.125~MB of wasted capacity during a TCP+TLS
handshake on a 120~Mbps link may seem small, but for latency-sensitive
applications where the first 100~ms of perceived responsiveness determines
user experience, the difference is measurable.

\textbf{Bandwidth surplus enables architectural ambition.}
The effective information rate analysis (Section~A.3,
Figure~\ref{fig:inet-capacity}) reveals the most consequential practical
finding: a single 720p video stream consumes a vanishing fraction of
available link capacity (under 0.5\% by 2030). This surplus is not merely
``headroom'' --- it is the resource that makes new application architectures
feasible. A 250~Mbps link can theoretically carry 230--415 simultaneous 720p
streams. In practice, this means applications can afford to run multiple
concurrent media flows (video, screen share, spatial audio), data channels
(collaborative editing, file transfer), and telemetry streams without
contending for bandwidth. The constraint shifts from ``can the network carry
this?''\ to ``can the application manage this complexity?''

\textbf{End-to-end encryption adds negligible overhead.}
The security analysis (Section~A.4) confirms that the nested DTLS-SRTP plus
E2EE architecture imposes no meaningful information-theoretic cost. The key
entropy ($2^{128}$) exceeds the information content of any practical session
by a factor of at least 3. For application developers, this means E2EE can
be treated as a default rather than a premium feature --- there is no
bandwidth, latency, or capacity penalty for encrypting at the application
layer on top of the transport-layer encryption that WebRTC already mandates.
The only constraint is API availability, which moves from Chrome-only (2026)
to universal (2028).

\textbf{WebTransport unlocks real-time data that WebSocket cannot serve.}
The freshness analysis (Section~A.5, Figure~\ref{fig:freshness}) quantifies
when unreliable delivery outperforms reliable delivery: any data with a
useful half-life below approximately 34~ms benefits from WebTransport's
unreliable datagrams. This covers a wide class of interactive applications
--- game state synchronization, live cursor and pointer sharing, real-time
sensor feeds, collaborative whiteboard strokes --- where retransmitting a
stale update is worse than dropping it. Before WebTransport, these
applications had to choose between WebSocket (reliable but stale under loss)
and WebRTC DataChannel (unreliable mode available but requiring the full
peer-to-peer ICE/DTLS machinery). WebTransport provides the unreliable
mode with a simple client-server model over HTTP/3, removing the
architectural overhead.

\textbf{The overall picture: the stack converges on efficiency.}
Taken together, the information-theoretic measures trace a consistent
trajectory. The protocol stack is converging toward the theoretical limits
established by Shannon: codecs approach the rate-distortion bound, transports
eliminate unnecessary blocking and handshake overhead, and encryption adds
negligible cost. By 2030, the practical distance between what the standards
deliver and what information theory permits will be the smallest it has ever
been for browser-based real-time communication. The consequence for
application developers is that the infrastructure --- standards, hardware,
and networks --- will no longer be the binding constraint. The remaining
challenge is purely architectural: how to compose these efficient, universal
primitives into applications that fully exploit the capacity and low latency
they provide.

%
% Requires in the preamble: amsmath, amssymb, tabularx, booktabs

\section{Information-Theoretic View}
\label{app:info-theory}

This appendix applies information-theoretic measures to the real-time
web communication stack described in this document. We quantify channel
capacity, source coding efficiency, and effective information rates
across three time horizons: today (March~2026), the near term
(2027--2028), and the medium term (2029--2031). The analysis draws on
the standards, codec parameters, and infrastructure projections
developed in the main text.

\subsection{Channel Capacity of the Transport Layer}

The Shannon-Hartley theorem gives the theoretical maximum information
rate for a channel with additive white Gaussian noise:
\begin{equation}
  C = B \log_2\!\left(1 + \frac{S}{N}\right) \quad \text{[bits/s]}
  \label{eq:shannon}
\end{equation}
where $B$ is the channel bandwidth in Hz and $S/N$ is the
signal-to-noise ratio.

While the physical-layer capacity of a broadband link is governed by
(\ref{eq:shannon}), the \emph{effective} capacity available to the
application is reduced by protocol overhead, congestion control, and
transport inefficiencies. We define the effective application-layer
throughput as:
\begin{equation}
  R_{\text{eff}} = C_{\text{link}} \cdot \eta_{\text{transport}}
    \cdot (1 - \rho_{\text{overhead}})
  \label{eq:reff}
\end{equation}
where $C_{\text{link}}$ is the link-layer capacity,
$\eta_{\text{transport}} \in (0,1]$ is the transport efficiency factor
(capturing congestion control, retransmission, and multiplexing
behavior), and $\rho_{\text{overhead}} \in [0,1)$ is the fractional
protocol overhead (headers, encryption, framing).

\subsubsection{Head-of-Line Blocking and Multiplexed Capacity}

Consider $n$ independent streams multiplexed over a single connection.
Under TCP (HTTP/2), a single lost packet blocks all $n$ streams. The
expected throughput per stream under loss rate $p$ is bounded by:
\begin{equation}
  R_{\text{TCP}}^{(i)} \leq \frac{R_{\text{eff}}}{n}
    \cdot (1 - p)^{n-1}
  \label{eq:hol-tcp}
\end{equation}
because a loss in \emph{any} of the $n$ streams stalls all others. Under
QUIC (RFC~9000), streams are independent at the transport layer, so:
\begin{equation}
  R_{\text{QUIC}}^{(i)} \leq \frac{R_{\text{eff}}}{n}
    \cdot (1 - p)
  \label{eq:hol-quic}
\end{equation}
The ratio of effective per-stream throughput is:
\begin{equation}
  \frac{R_{\text{QUIC}}^{(i)}}{R_{\text{TCP}}^{(i)}}
    = \frac{1}{(1-p)^{n-2}}
  \label{eq:hol-ratio}
\end{equation}
\excursus{Excursus: Why $n = 8$?}{In a typical multi-participant video
conference, each peer contributes multiple concurrent streams: a camera video
track, an audio track, and often a screen-share or data channel. A four-person
call with two media tracks each produces $n = 8$ multiplexed streams over a
single connection --- the reference value used throughout this section. Larger
meetings (e.g.\ 10 participants with simulcast layers) push $n$ higher,
amplifying the HOL blocking effect further.}

For $n = 8$ streams (a typical multi-track video conference) at
$p = 0.01$ (1\% loss), this gives a QUIC advantage factor of
$1/(0.99)^6 \approx 1.062$, a modest 6.2\% improvement. At $p = 0.05$
(5\% loss, common on mobile or congested links), the factor is
$1/(0.95)^6 \approx 1.34$, a 34\% throughput advantage per stream.

\begin{figure}[ht]
\centering
\begin{tikzpicture}[>=Stealth]

% --- Shaded region (drawn FIRST so curves appear on top) ---
\fill[orange!8] (2, 0) rectangle (10, 5.3);
\node[font=\tiny\itshape, text=black!40, anchor=north east] at (9.8, 5.2)
  {mobile / congested};
\node[font=\tiny\itshape, text=black!40, anchor=north west] at (0.2, 5.2)
  {fixed broadband};

% --- Grid lines ---
\foreach \y in {1, 2, 3, 4, 5} {
  \draw[gray!20] (0, \y) -- (10, \y);
}

% --- Axes ---
\draw[->, thick] (0, 0) -- (10.5, 0);
\draw[->, thick] (0, 0) -- (0, 5.5);

% --- Axis titles (positioned to clear tick labels) ---
\node[font=\small] at (5, -0.85) {Packet loss rate $p$};
\node[font=\small, rotate=90] at (-0.9, 2.75) {QUIC advantage factor};

% --- X-axis labels ---
\foreach \x/\lab in {0/0, 2/0.02, 4/0.04, 6/0.06, 8/0.08, 10/0.10} {
  \draw (\x, -0.1) -- (\x, 0.1);
  \node[font=\tiny, anchor=north] at (\x, -0.15) {\lab};
}

% --- Y-axis labels ---
\foreach \y/\lab in {0/1.0, 1/1.1, 2/1.2, 3/1.3, 4/1.4, 5/1.5} {
  \draw (-0.1, \y) -- (0.1, \y);
  \node[font=\tiny, anchor=east] at (-0.15, \y) {\lab};
}

% --- n=4 curve: 1/(1-p)^2 ---
\draw[very thick, blue!70]
  (0, 0) .. controls (2.5, 0.5) and (5, 1.1) ..
  (7.5, 2.0) .. controls (8.5, 2.4) and (9.5, 2.9) ..
  (10, 3.1);
\node[font=\scriptsize, blue!70, anchor=south west] at (9.0, 3.0)
  {$n = 4$};

% --- n=8 curve: 1/(1-p)^6 ---
\draw[very thick, orange!80!black]
  (0, 0) .. controls (1.5, 0.5) and (3, 1.7) ..
  (4, 2.5) .. controls (5, 3.5) and (5.5, 4.0) ..
  (6, 4.7);
\node[font=\scriptsize, orange!80!black, anchor=south] at (5.5, 4.8)
  {$n = 8$};

% --- Reference points ---
\fill[orange!80!black] (1.0, 0.62) circle (2.5pt);
\node[font=\tiny, orange!80!black, anchor=south west] at (1.1, 0.62)
  {6.2\%};

\fill[orange!80!black] (5.0, 3.7) circle (2.5pt);
\node[font=\tiny, orange!80!black, anchor=west] at (5.1, 3.7)
  {34\%};

\fill[blue!70] (1.0, 0.20) circle (2.5pt);

% --- Annotation ---
\node[font=\scriptsize\itshape, text=black!60, align=center,
  anchor=north] at (5.0, -1.0)
  {QUIC recovers per-stream throughput closer to $C/n$ (Shannon fair share)\\
   by eliminating head-of-line blocking (eq.~\ref{eq:hol-ratio}).};

\end{tikzpicture}
\caption{QUIC advantage factor $1/(1-p)^{n-2}$ over TCP for multiplexed
  streams. At 5\% loss with $n=8$ streams, QUIC delivers 34\% more
  per-stream throughput by maintaining independence at the transport layer.}
\label{fig:hol-blocking}
\end{figure}

\subsubsection{Connection Establishment Latency}

Let $\tau$ denote the one-way latency (half the round-trip time). The
information-carrying capacity during connection establishment is zero;
the handshake represents pure overhead. The time to first useful byte
is:
\begin{align}
  T_{\text{TCP+TLS\,1.3}} &= 2\tau_{\text{TCP}} + \tau_{\text{TLS}}
    = 3\tau \quad \text{(1-RTT TLS~1.3)} \label{eq:tcp-setup} \\
  T_{\text{QUIC}} &= \tau \quad \text{(1-RTT)}
    \label{eq:quic-setup} \\
  T_{\text{QUIC,\,0-RTT}} &= 0 \quad \text{(resumption)}
    \label{eq:quic-0rtt}
\end{align}
The QUIC 0-RTT case (\ref{eq:quic-0rtt}) eliminates the handshake
entirely for resumed connections. The information lost during
handshake dead time, relative to steady-state capacity, is:
\begin{equation}
  I_{\text{lost}} = C_{\text{link}} \cdot T_{\text{setup}}
  \label{eq:ilost}
\end{equation}
For a 120~Mbps link with $\tau = 25$~ms (typical fixed broadband in
2026), the TCP+TLS~1.3 handshake wastes $120 \times 10^6 \times 0.075
= 9 \times 10^6$ bits (1.125~MB), while QUIC~1-RTT wastes only
$3 \times 10^6$ bits (375~KB).

\subsection{Source Coding Efficiency}

The rate-distortion function $R(D)$ defines the minimum bit rate
required to represent a source at distortion level $D$. For a
Gaussian source with variance $\sigma^2$ under mean squared error
distortion:
\begin{equation}
  R(D) = \frac{1}{2} \log_2 \frac{\sigma^2}{D}
    \quad \text{[bits/sample]}
  \label{eq:rate-distortion}
\end{equation}

\excursus{Excursus: Why Gaussian source?}{The Gaussian memoryless model is used
because it yields the highest (worst-case) $R(D)$ among all sources of equal
variance. Any codec that approaches the Gaussian bound on real video is
therefore doing at least as well as the bound predicts --- the true $R(D)$ for
spatially and temporally correlated video lies below the Gaussian curve, making
it a conservative benchmark.}

No practical codec achieves $R(D)$;\footnote{The $R(D)$ curve in
Figure~\ref{fig:rate-distortion} is the Gaussian memoryless source bound
(equation~\ref{eq:rate-distortion}), which is an upper bound on the true
$R(D)$ for non-Gaussian sources such as video.  Additionally, $R(D)$ assumes
zero excess-distortion probability, whereas practical codecs permit a
frame-level excess-distortion probability on the order of $10^{-6}$ to
$10^{-8}$.  Both factors allow codec operating points measured on real content
to approach or marginally cross the Gaussian bound.} real codecs operate at
some rate $R_{\text{codec}} > R(D)$. We define the coding efficiency gap as:
\begin{equation}
  \Delta = \frac{R_{\text{codec}} - R(D)}{R(D)}
  \label{eq:coding-gap}
\end{equation}

The codecs in the WebRTC stack can be ordered by their empirical
proximity to $R(D)$ for typical video content at conversational
resolutions (720p, 30~fps). Using published Bj\o ntegaard Delta (BD)
rate measurements:
\begin{equation}
  R_{\text{H.264}} > R_{\text{VP9}} > R_{\text{AV1}}
    > R(D)
  \label{eq:codec-ordering}
\end{equation}

Specifically, the empirical rate savings at equivalent PSNR are
approximately:
\begin{align}
  R_{\text{VP9}}  &\approx 0.65 \cdot R_{\text{H.264}}
    \quad (\sim\!35\% \text{ savings}) \label{eq:vp9-savings} \\
  R_{\text{AV1}}  &\approx 0.70 \cdot R_{\text{VP9}}
    \approx 0.46 \cdot R_{\text{H.264}}
    \quad (\sim\!30\% \text{ over VP9})
  \label{eq:av1-savings}
\end{align}

\begin{figure}[ht]
\centering
\begin{tikzpicture}[>=Stealth]

% --- Axes ---
\draw[->, thick] (0, 0) -- (10.5, 0);
\draw[->, thick] (0, 0) -- (0, 6.5);

% --- Axis titles ---
\node[font=\small] at (5, -0.85) {Distortion $D$};
\node[font=\small, rotate=90] at (-0.9, 3.25) {Rate [Mbps]};

% --- Y-axis rate labels ---
\foreach \y/\lab in {1.0/0.5, 2.0/1.0, 3.0/1.5, 4.0/2.0, 5.0/2.5, 6.0/3.0} {
  \draw (-0.1, \y) -- (0.1, \y);
  \node[font=\tiny, anchor=east] at (-0.15, \y) {\lab};
}

% --- R(D) Shannon bound curve ---
\draw[very thick, green!60!black]
  (1.0, 5.8) .. controls (2.0, 3.5) and (3.5, 2.0) ..
  (5.0, 1.3) .. controls (6.5, 0.8) and (8.0, 0.55) ..
  (9.5, 0.4);
\node[font=\scriptsize, green!60!black, anchor=north west] at (8.5, 0.55)
  {$R(D)$ bound};

% --- Reference distortion line (720p/30fps operating point) ---
\draw[thin, gray, dashed] (4.0, 0) -- (4.0, 6.2);
\node[font=\tiny, gray, anchor=south] at (4.0, 6.2) {720p/30fps};

% --- R(D) value at reference distortion ---
\fill[green!60!black] (4.0, 1.55) circle (2pt);

% --- Codec operating points at reference distortion ---
% H.264: ~2.5 Mbps
\draw[thick, dashed, red!70!black] (0.3, 4.8) -- (9.5, 4.8);
\fill[red!70!black] (4.0, 4.8) circle (3pt);
\node[font=\scriptsize, red!70!black, anchor=west] at (9.6, 4.8) {H.264};

% VP9: ~1.6 Mbps (0.65 * H.264)
\draw[thick, dashed, blue!70] (0.3, 3.3) -- (9.5, 3.3);
\fill[blue!70] (4.0, 3.3) circle (3pt);
\node[font=\scriptsize, blue!70, anchor=west] at (9.6, 3.3) {VP9};

% AV1: ~1.15 Mbps (0.46 * H.264)
\draw[thick, dashed, orange!80!black] (0.3, 2.3) -- (9.5, 2.3);
\fill[orange!80!black] (4.0, 2.3) circle (3pt);
\node[font=\scriptsize, orange!80!black, anchor=west] at (9.6, 2.3) {AV1};

% --- Gap annotations (braces) ---
% H.264 gap
\draw[decorate, decoration={brace, amplitude=4pt, mirror},
  thick, red!50!black]
  (3.6, 1.55) -- (3.6, 4.8)
  node[midway, anchor=east, font=\scriptsize, xshift=-5pt]
  {$\Delta_{\text{H.264}}$};

% AV1 gap
\draw[decorate, decoration={brace, amplitude=4pt},
  thick, orange!60!black]
  (4.4, 1.55) -- (4.4, 2.3)
  node[midway, anchor=west, font=\scriptsize, xshift=5pt]
  {$\Delta_{\text{AV1}}$};

% --- Annotation ---
\node[font=\scriptsize\itshape, text=black!60, align=center,
  anchor=north] at (5.0, -0.5)
  {The Shannon rate-distortion bound $R(D)$ is the theoretical minimum.\\
   AV1 narrows the gap $\Delta$ to ${\sim}0.4$ by 2030 (eq.~\ref{eq:coding-gap}).};

\end{tikzpicture}
\caption{Codec operating rates vs.\ the Shannon rate-distortion bound $R(D)$
  at 720p/30fps. Each codec operates above the theoretical minimum; the gap
  $\Delta$ (eq.~\ref{eq:coding-gap}) shrinks from H.264 through VP9 to AV1,
  approaching but never reaching the bound.}
\label{fig:rate-distortion}
\end{figure}

\subsubsection{Codec Entropy and WebCodecs}

The WebCodecs API (W3C Working Draft) exposes per-frame codec access,
allowing measurement and control of the instantaneous encoding rate.
For a video frame $f_t$ at time $t$, the encoder output has empirical
entropy:
\begin{equation}
  \hat{H}(f_t) = -\sum_{s \in \mathcal{S}}
    p(s \mid f_t) \log_2 p(s \mid f_t)
  \label{eq:frame-entropy}
\end{equation}
where $\mathcal{S}$ is the set of symbols emitted by the entropy coder
(CABAC for H.264/AV1, arithmetic coding for VP9). WebCodecs enables
applications to observe $|E(f_t)|$ (the encoded frame size in bits)
directly, which serves as a practical upper bound on
$\hat{H}(f_t)$:
\begin{equation}
  \hat{H}(f_t) \leq |E(f_t)| \leq R_{\text{codec}} / \text{fps}
  \label{eq:entropy-bound}
\end{equation}

\subsection{Effective Information Rate by Time Horizon}

We now combine transport and source coding measures to estimate the
effective information rate for a representative real-time video
session: a single 720p/30fps video call with one audio channel.

\excursus{Excursus: Why 720p?}{Throughout this appendix, 720p/30fps is used as
a fixed reference for cross-era comparison, not as a prediction that 720p
remains the dominant conferencing resolution. The bandwidth surplus quantified
below is precisely what enables migration to 1080p (feasible by ${\sim}$2028
with AV1 hardware encode at today's 720p bit rates) and 4K (feasible by
${\sim}$2030). The 720p baseline makes the efficiency gains legible: the same
task gets cheaper, freeing capacity for higher quality or more streams.}

Define the \emph{net information rate} as:
\begin{equation}
  I_{\text{net}} = R_{\text{codec}} \cdot \eta_{\text{transport}}
    \cdot (1 - \rho_{\text{overhead}})
    \cdot (1 - p_{\text{loss}})
  \label{eq:inet}
\end{equation}

\begin{table}[ht]
\centering
\small
\renewcommand{\arraystretch}{1.4}
\begin{tabularx}{\textwidth}{lXXX}
\toprule
\textbf{Parameter} & \textbf{March 2026} & \textbf{2027--2028} &
  \textbf{2029--2031} \\
\midrule
Median link capacity &
  120~Mbps &
  $\sim$150~Mbps &
  200--300~Mbps \\
Dominant codec (WebRTC) &
  VP9 / H.264 &
  VP9 / AV1 (HW) &
  AV1 (HW default) \\
Codec rate (720p/30fps) &
  1.5--2.5~Mbps &
  1.0--1.8~Mbps &
  0.7--1.2~Mbps \\
$\eta_{\text{transport}}$ &
  0.92 (QUIC) &
  0.93 &
  0.94 \\
$\rho_{\text{overhead}}$ (SRTP+QUIC) &
  0.06 &
  0.05 &
  0.04 \\
Typical loss $p$ &
  0.01--0.03 &
  0.01--0.02 &
  0.005--0.01 \\
$I_{\text{net}}$ (720p video) &
  1.3--2.2~Mbps &
  0.9--1.6~Mbps &
  0.6--1.1~Mbps \\
$I_{\text{net}}$ as \% of $C_{\text{link}}$ &
  1.1--1.8\% &
  0.6--1.1\% &
  0.2--0.6\% \\
\bottomrule
\end{tabularx}
\caption{Effective information rate for a 720p/30fps WebRTC video stream.
  The declining ratio $I_{\text{net}} / C_{\text{link}}$ reflects both
  improved codec efficiency and growing link capacity, freeing bandwidth
  for higher resolutions, additional streams, or concurrent data
  channels.}
\label{tab:inet}
\end{table}

The declining ratio $I_{\text{net}} / C_{\text{link}}$ in
Table~\ref{tab:inet} is the central quantitative observation: by 2030, a
720p video call consumes less than 0.5\% of available link capacity. The
surplus capacity enables qualitative changes --- 4K video conferencing,
multi-stream layouts, or concurrent data channels via WebTransport ---
that are quantitatively infeasible in 2026.

\begin{figure}[ht]
\centering
\begin{tikzpicture}[>=Stealth]

% --- Axes ---
\draw[->, thick] (0, 0) -- (11, 0);
\draw[->, thick] (0, 0) -- (0, 5.5);

% --- Axis title ---
\node[font=\small, rotate=90] at (-0.9, 2.75) {Mbps};

% --- Y-axis labels ---
\foreach \y/\lab in {0/0, 1/50, 2/100, 3/150, 4/200, 5/250} {
  \draw (-0.1, \y) -- (0.1, \y);
  \node[font=\tiny, anchor=east] at (-0.15, \y) {\lab};
}

% --- Group: March 2026 ---
% C_link = 120 Mbps -> y = 2.4
\fill[blue!12] (0.8, 0) rectangle (2.4, 2.4);
\draw[blue!50, thick] (0.8, 0) rectangle (2.4, 2.4);
% I_net = 1.75 Mbps (midpoint) -> y = 0.035
\fill[orange!60] (0.8, 0) rectangle (2.4, 0.07);
\draw[orange!80!black, thick] (0.8, 0) rectangle (2.4, 0.07);
\node[font=\footnotesize, anchor=south] at (1.6, 2.5) {120};
\node[font=\tiny, orange!80!black] at (1.6, -0.4) {1.5\%};
\node[font=\small, anchor=north] at (1.6, -0.7) {2026};

% --- Group: 2027--2028 ---
% C_link = 150 Mbps -> y = 3.0
\fill[blue!12] (4.0, 0) rectangle (5.6, 3.0);
\draw[blue!50, thick] (4.0, 0) rectangle (5.6, 3.0);
% I_net = 1.25 Mbps (midpoint) -> y = 0.025
\fill[orange!60] (4.0, 0) rectangle (5.6, 0.05);
\draw[orange!80!black, thick] (4.0, 0) rectangle (5.6, 0.05);
\node[font=\footnotesize, anchor=south] at (4.8, 3.1) {150};
\node[font=\tiny, orange!80!black] at (4.8, -0.4) {0.8\%};
\node[font=\small, anchor=north] at (4.8, -0.7) {2027--28};

% --- Group: 2029--2031 ---
% C_link = 250 Mbps -> y = 5.0
\fill[blue!12] (7.2, 0) rectangle (8.8, 5.0);
\draw[blue!50, thick] (7.2, 0) rectangle (8.8, 5.0);
% I_net = 0.85 Mbps (midpoint) -> y = 0.017
\fill[orange!60] (7.2, 0) rectangle (8.8, 0.034);
\draw[orange!80!black, thick] (7.2, 0) rectangle (8.8, 0.034);
\node[font=\footnotesize, anchor=south] at (8.0, 5.1) {250};
\node[font=\tiny, orange!80!black] at (8.0, -0.4) {0.3\%};
\node[font=\small, anchor=north] at (8.0, -0.7) {2029--31};

% --- Legend ---
\node[anchor=west, font=\footnotesize] at (0.5, 5.8) {%
  \tikz[baseline=-0.5ex]{\fill[blue!12] (0,0) rectangle (0.3,0.3);
    \draw[blue!50, thick] (0,0) rectangle (0.3,0.3);}
  ~$C_{\text{link}}$ (channel capacity) \qquad
  \tikz[baseline=-0.5ex]{\fill[orange!60] (0,0) rectangle (0.3,0.3);
    \draw[orange!80!black, thick] (0,0) rectangle (0.3,0.3);}
  ~$I_{\text{net}}$ (720p stream)};

% --- Annotation ---
\node[font=\scriptsize\itshape, text=black!60, align=center,
  anchor=north] at (5.5, -1.2)
  {A single 720p call uses a vanishing fraction of Shannon-derived\\
   channel capacity, freeing bandwidth for 4K, multi-stream, and data channels.};

\end{tikzpicture}
\caption{Channel capacity $C_{\text{link}}$ vs.\ effective information rate
  $I_{\text{net}}$ for a single 720p/30fps stream across time horizons. The
  orange bar (stream rate) is barely visible against the blue bar (capacity),
  illustrating the growing surplus as codecs improve and bandwidth increases.}
\label{fig:inet-capacity}
\end{figure}

\subsection{Security Entropy and Key Space}

The DTLS-SRTP security architecture (RFCs~8827, 3711) protects all
WebRTC media. Information-theoretic security requires that the key
entropy exceed the information content of the protected material.

For a DTLS~1.3 handshake with ECDHE on P-256, the key agreement
provides:
\begin{equation}
  H_{\text{key}} = 128 \text{~bits of security}
  \label{eq:key-entropy}
\end{equation}

The SRTP master key is 128 bits (AES-128-GCM), giving an exhaustive
search cost of $2^{128}$ operations. The total information content of a
one-hour 720p video call at 2~Mbps is:
\begin{equation}
  I_{\text{call}} = 2 \times 10^6 \times 3600
    = 7.2 \times 10^9 \text{~bits}
    \approx 2^{32.7} \text{~bits}
  \label{eq:call-content}
\end{equation}

The ratio of key entropy to protected content:
\begin{equation}
  \frac{H_{\text{key}}}{I_{\text{call}}}
    = \frac{128}{32.7} \approx 3.91
  \label{eq:security-margin}
\end{equation}

This ratio remains comfortably above~1 even for much longer sessions.
For a continuous 24-hour 4K AV1 stream at 15~Mbps:
\begin{equation}
  I_{\text{stream}} = 15 \times 10^6 \times 86400
    = 1.296 \times 10^{12} \text{~bits}
    \approx 2^{40.2} \text{~bits}
\end{equation}
yielding $H_{\text{key}} / I_{\text{stream}} \approx 3.18$, still
well above the threshold.

\subsubsection{End-to-End Encryption via Encoded Transforms}

The RTCRtpScriptTransform API (Encoded Transforms) enables
application-layer E2EE within the WebRTC pipeline. From an
information-theoretic perspective, this introduces a second encryption
layer with its own key, creating a \emph{nested channel}:
\begin{equation}
  X \xrightarrow{\;K_{\text{E2EE}}\;}
  E_1(X) \xrightarrow{\;K_{\text{SRTP}}\;}
  E_2(E_1(X))
  \label{eq:nested}
\end{equation}

The mutual information between the ciphertext and plaintext, given
neither key, satisfies:
\begin{equation}
  I\!\left(X;\, E_2(E_1(X))\right) \leq
    \min\!\big(H(K_{\text{E2EE}})^{-1},\;
    H(K_{\text{SRTP}})^{-1}\big)
    \cdot I_{\text{call}}
  \label{eq:mutual-e2ee}
\end{equation}

In practice, with $H(K_{\text{E2EE}}) = H(K_{\text{SRTP}}) = 128$
bits, this quantity is negligible ($\ll 1$ bit for any practical
session length), confirming that the nested encryption provides
information-theoretic confidentiality well beyond computational
security guarantees.

Today (March~2026), Encoded Transforms remain Chrome-only. By
2027--2028 (Section~8.3), cross-browser standardization will make
this nested channel universally available. By 2029--2031, SFrame
for group calls will extend this to $n$-party sessions, where the
mutual information between any non-member and the group plaintext
remains negligible:
\begin{equation}
  I\!\left(X;\, E(X) \;\middle|\; K_i \notin
    \{K_1, \ldots, K_n\}\right) \approx 0
  \label{eq:sframe}
\end{equation}

\subsection{WebTransport: Unreliable Datagrams and Channel Capacity}

WebTransport (RFC~9297) introduces unreliable datagrams to browser-server
communication. For latency-sensitive applications (game state, sensor
telemetry), the relevant measure is not throughput but
\emph{freshness-weighted information rate}. Define the value of a
datagram as exponentially decaying with delivery latency $\delta$:
\begin{equation}
  V(\delta) = e^{-\lambda \delta}
  \label{eq:freshness}
\end{equation}
where $\lambda > 0$ is an application-specific decay rate. The
effective information rate under unreliable delivery is:
\begin{equation}
  I_{\text{WT}} = R_{\text{send}} \cdot (1 - p_{\text{loss}})
    \cdot \mathbb{E}\!\left[e^{-\lambda \delta}\right]
  \label{eq:iwt}
\end{equation}

Under reliable delivery (TCP/WebSocket), a retransmission delays
subsequent data, so:
\begin{equation}
  I_{\text{WS}} = R_{\text{send}}
    \cdot \mathbb{E}\!\left[e^{-\lambda(\delta + p \cdot 2\tau)}\right]
  \label{eq:iws}
\end{equation}

The ratio $I_{\text{WT}} / I_{\text{WS}}$ exceeds 1 when:
\begin{equation}
  (1 - p) \cdot e^{\lambda \cdot p \cdot 2\tau} > 1
  \quad \Longleftrightarrow \quad
  \lambda > \frac{-\ln(1 - p)}{2 p \tau}
  \label{eq:wt-advantage}
\end{equation}

For $p = 0.02$ and $\tau = 25$~ms, this gives $\lambda > 20.2$~s$^{-1}$
(half-life $< 34$~ms).

\excursus{Excursus: Why $p = 0.02$ and $\tau = 25$\,ms?}{These values represent
median fixed-broadband conditions in 2026 (Table~\ref{tab:inet}). On mobile or
congested links ($p \approx 0.05$, $\tau \approx 50$\,ms), the threshold
$\lambda^*$ drops to ${\sim}10$\,s$^{-1}$, broadening WebTransport's advantage
to applications with data half-lives up to ${\sim}70$\,ms. The fixed-broadband
reference is thus conservative: mobile scenarios favor unreliable delivery even
more strongly.}

Any application where data older than
$\sim$30~ms is useless benefits from WebTransport's unreliable
mode over WebSocket --- precisely the regime of interactive gaming, live
cursor tracking, and real-time sensor feeds.

As of March~2026, Safari does not support WebTransport
(Section~5), limiting coverage to $\sim$82\%. By 2027--2028
(Section~8.2), universal support will allow applications to rely on
(\ref{eq:iwt}) without fallback paths.

\begin{figure}[ht]
\centering
\begin{tikzpicture}[>=Stealth]

% --- Axes ---
\draw[->, thick] (0, 0) -- (10.5, 0);
\draw[->, thick] (0, 0) -- (0, 5.5);

% --- Axis titles (positioned to clear tick labels) ---
\node[font=\small] at (5, -0.75)
  {Freshness decay rate $\lambda$ [s$^{-1}$]};
\node[font=\small, rotate=90] at (-1.1, 2.75)
  {$I_{\text{WT}} / I_{\text{WS}}$};

% --- X-axis labels ---
\foreach \x/\lab in {0/0, 1.67/10, 3.33/20, 5.0/30, 6.67/40, 8.33/50, 10/60} {
  \draw (\x, -0.1) -- (\x, 0.1);
  \node[font=\tiny, anchor=north] at (\x, -0.15) {\lab};
}

% --- Y-axis labels (0.96 to 1.06, each 0.02 = 1 unit) ---
\foreach \y/\lab in {0/0.96, 1/0.98, 2/1.00, 3/1.02, 4/1.04, 5/1.06} {
  \draw (-0.1, \y) -- (0.1, \y);
  \node[font=\tiny, anchor=east] at (-0.15, \y) {\lab};
}

% --- Ratio = 1.0 line (y=2 maps to ratio 1.0) ---
\draw[thin, gray, dashed] (0, 2) -- (10, 2);

% --- Shade WebTransport advantage region ---
% lambda* = 20.2 -> x = 3.37
\fill[green!10] (3.37, 2) rectangle (10, 5.3);

% --- I_WT/I_WS curve: 0.98*exp(lambda*0.001) ---
% Y range: 0.96-1.06 (each 0.02 = 1 unit). lambda=0: y=1, lambda*=20.2: y=2, lambda=60: y=4.05
\draw[very thick, orange!80!black]
  (0, 1.0) .. controls (1.5, 1.3) and (2.5, 1.7) ..
  (3.37, 2.0) .. controls (4.5, 2.4) and (6, 2.8) ..
  (8, 3.4) .. controls (9, 3.7) and (9.5, 3.85) ..
  (10, 4.05);

% --- Threshold line ---
\draw[thick, dashed, red!60!black] (3.37, 0) -- (3.37, 5.2);
\node[font=\scriptsize, red!60!black, anchor=south, rotate=90] at (3.17, 3.5)
  {$\lambda^* = 20.2$ s$^{-1}$};

% --- Application regime labels ---
\node[font=\scriptsize\itshape, text=blue!60, anchor=north] at (1.5, 5.3)
  {video conf};
\node[font=\scriptsize\itshape, text=green!50!black, anchor=north] at (7.0, 5.3)
  {gaming / cursor / sensors};

% --- Mark advantage/disadvantage ---
\node[font=\tiny, text=red!50!black, anchor=north east] at (3.2, 1.8)
  {WebSocket wins};
\node[font=\tiny, text=green!50!black, anchor=north west] at (3.5, 4.5)
  {WebTransport wins};

% --- Annotation ---
\node[font=\scriptsize\itshape, text=black!60, align=center,
  anchor=north] at (5.0, -1.0)
  {For $p = 0.02$, $\tau = 25$~ms. Applications where data older than\\
   ${\sim}$30~ms is stale gain from unreliable delivery (eq.~\ref{eq:wt-advantage}).};

\end{tikzpicture}
\caption{Freshness-weighted information rate ratio $I_{\text{WT}} / I_{\text{WS}}$
  as a function of decay rate $\lambda$. Above the threshold $\lambda^*$,
  WebTransport's unreliable datagrams deliver more useful information per
  unit time than WebSocket's reliable delivery, because retransmission
  wastes capacity on stale data.}
\label{fig:freshness}
\end{figure}

\subsection{Projections Summary and Practical Implications}

\begin{table}[ht]
\centering
\small
\renewcommand{\arraystretch}{1.4}
\begin{tabularx}{\textwidth}{lXXX}
\toprule
\textbf{Measure} & \textbf{March 2026} & \textbf{2027--2028} &
  \textbf{2029--2031} \\
\midrule
Channel capacity $C_{\text{link}}$ &
  120~Mbps &
  150~Mbps &
  250~Mbps \\
HOL blocking factor $(1\!-\!p)^{n-2}$ at $n\!=\!8$ &
  0.94 ($p\!=\!0.01$) &
  0.94 &
  0.97 ($p\!=\!0.005$) \\
Source coding gap $\Delta$ &
  VP9: $\sim$0.8 &
  AV1 HW: $\sim$0.5 &
  AV1 opt: $\sim$0.4 \\
$I_{\text{net}}$ per 720p stream &
  1.3--2.2~Mbps &
  0.9--1.6~Mbps &
  0.6--1.1~Mbps \\
Max 720p streams in $C_{\text{link}}$ &
  $\sim$55--90 &
  $\sim$95--165 &
  $\sim$230--415 \\
E2EE availability &
  Chrome only &
  Cross-browser &
  SFrame $n$-party \\
WebTransport coverage &
  82\% &
  $\sim$100\% &
  100\% \\
Freshness threshold $\lambda^*$ &
  20.2~s$^{-1}$ &
  $\sim$18~s$^{-1}$ &
  $\sim$15~s$^{-1}$ \\
\bottomrule
\end{tabularx}
\caption{Information-theoretic measures across time horizons. $\Delta$
  is the coding efficiency gap (\ref{eq:coding-gap}), $I_{\text{net}}$
  is the net information rate (\ref{eq:inet}), and $\lambda^*$ is the
  freshness decay threshold above which WebTransport dominates WebSocket
  (\ref{eq:wt-advantage}).}
\label{tab:projections}
\end{table}

The trajectory is clear: the combination of improving source coding
(AV1 hardware), expanding channel capacity (broadband growth), and
maturing transport protocols (QUIC universal, WebTransport universal)
drives the \emph{information cost per unit of communicated experience}
steadily downward. The remaining question is not whether the
Shannon capacity suffices --- it already does, by orders of magnitude
--- but how efficiently the protocol stack converts that capacity into
delivered, fresh, secure information at the application layer. By
2030, the gap between theoretical capacity and realized throughput will
be the narrowest it has ever been for browser-based real-time
communication.

\subsubsection{What the Theory Means for Applications}

The equations and figures in this appendix are not purely academic. Each
theoretical result maps to a concrete design decision that application
developers face today and over the next five years.

\textbf{Codec selection determines bandwidth cost, not quality.}
The rate-distortion analysis (Section~A.2, Figure~\ref{fig:rate-distortion})
shows that AV1 operates at roughly 46\% of H.264's bit rate for equivalent
perceptual quality at 720p. In practice, this means a WebRTC application
switching from H.264 to AV1 can either halve its bandwidth consumption at
the same quality or substantially increase resolution (e.g., 720p to 1080p)
within the same bandwidth envelope. By 2030, with AV1 hardware encoding
ubiquitous and the coding gap $\Delta$ approaching 0.4, further codec
improvements yield diminishing returns --- the stack will be operating close
enough to the Shannon rate-distortion bound that the next generation of
efficiency gains will come from transport and architecture, not codecs alone.

\textbf{QUIC's advantage is situational but decisive on bad networks.}
The head-of-line blocking analysis (Section~A.1.1,
Figure~\ref{fig:hol-blocking}) shows that QUIC's per-stream throughput
advantage over TCP is modest on good networks (6\% at 1\% loss) but
substantial on degraded ones (34\% at 5\% loss). The practical implication
is that applications targeting mobile users, emerging markets, or congested
enterprise networks --- precisely the environments where real-time
communication is hardest --- benefit disproportionately from QUIC-based
transports. For a video conference with 8 participants on a cellular
connection experiencing 3--5\% packet loss, the difference between QUIC and
TCP is the difference between usable and degraded video.

\textbf{Connection latency is a solved problem for returning users.}
The handshake analysis (Section~A.1.2) shows that QUIC 0-RTT eliminates
connection establishment overhead entirely for resumed connections. For a
WebTransport application that a user opens daily --- a collaboration tool,
a live dashboard, a recurring meeting client --- the time to first useful
byte drops to zero. The 1.125~MB of wasted capacity during a TCP+TLS
handshake on a 120~Mbps link may seem small, but for latency-sensitive
applications where the first 100~ms of perceived responsiveness determines
user experience, the difference is measurable.

\textbf{Bandwidth surplus enables architectural ambition.}
The effective information rate analysis (Section~A.3,
Figure~\ref{fig:inet-capacity}) reveals the most consequential practical
finding: a single 720p video stream consumes a vanishing fraction of
available link capacity (under 0.5\% by 2030). This surplus is not merely
``headroom'' --- it is the resource that makes new application architectures
feasible. A 250~Mbps link can theoretically carry 230--415 simultaneous 720p
streams. In practice, this means applications can afford to run multiple
concurrent media flows (video, screen share, spatial audio), data channels
(collaborative editing, file transfer), and telemetry streams without
contending for bandwidth. The constraint shifts from ``can the network carry
this?''\ to ``can the application manage this complexity?''

\textbf{End-to-end encryption adds negligible overhead.}
The security analysis (Section~A.4) confirms that the nested DTLS-SRTP plus
E2EE architecture imposes no meaningful information-theoretic cost. The key
entropy ($2^{128}$) exceeds the information content of any practical session
by a factor of at least 3. For application developers, this means E2EE can
be treated as a default rather than a premium feature --- there is no
bandwidth, latency, or capacity penalty for encrypting at the application
layer on top of the transport-layer encryption that WebRTC already mandates.
The only constraint is API availability, which moves from Chrome-only (2026)
to universal (2028).

\textbf{WebTransport unlocks real-time data that WebSocket cannot serve.}
The freshness analysis (Section~A.5, Figure~\ref{fig:freshness}) quantifies
when unreliable delivery outperforms reliable delivery: any data with a
useful half-life below approximately 34~ms benefits from WebTransport's
unreliable datagrams. This covers a wide class of interactive applications
--- game state synchronization, live cursor and pointer sharing, real-time
sensor feeds, collaborative whiteboard strokes --- where retransmitting a
stale update is worse than dropping it. Before WebTransport, these
applications had to choose between WebSocket (reliable but stale under loss)
and WebRTC DataChannel (unreliable mode available but requiring the full
peer-to-peer ICE/DTLS machinery). WebTransport provides the unreliable
mode with a simple client-server model over HTTP/3, removing the
architectural overhead.

\textbf{The overall picture: the stack converges on efficiency.}
Taken together, the information-theoretic measures trace a consistent
trajectory. The protocol stack is converging toward the theoretical limits
established by Shannon: codecs approach the rate-distortion bound, transports
eliminate unnecessary blocking and handshake overhead, and encryption adds
negligible cost. By 2030, the practical distance between what the standards
deliver and what information theory permits will be the smallest it has ever
been for browser-based real-time communication. The consequence for
application developers is that the infrastructure --- standards, hardware,
and networks --- will no longer be the binding constraint. The remaining
challenge is purely architectural: how to compose these efficient, universal
primitives into applications that fully exploit the capacity and low latency
they provide.

%
% Requires in the preamble: amsmath, amssymb, tabularx, booktabs

\section{Information-Theoretic View}
\label{app:info-theory}

This appendix applies information-theoretic measures to the real-time
web communication stack described in this document. We quantify channel
capacity, source coding efficiency, and effective information rates
across three time horizons: today (March~2026), the near term
(2027--2028), and the medium term (2029--2031). The analysis draws on
the standards, codec parameters, and infrastructure projections
developed in the main text.

\subsection{Channel Capacity of the Transport Layer}

The Shannon-Hartley theorem gives the theoretical maximum information
rate for a channel with additive white Gaussian noise:
\begin{equation}
  C = B \log_2\!\left(1 + \frac{S}{N}\right) \quad \text{[bits/s]}
  \label{eq:shannon}
\end{equation}
where $B$ is the channel bandwidth in Hz and $S/N$ is the
signal-to-noise ratio.

While the physical-layer capacity of a broadband link is governed by
(\ref{eq:shannon}), the \emph{effective} capacity available to the
application is reduced by protocol overhead, congestion control, and
transport inefficiencies. We define the effective application-layer
throughput as:
\begin{equation}
  R_{\text{eff}} = C_{\text{link}} \cdot \eta_{\text{transport}}
    \cdot (1 - \rho_{\text{overhead}})
  \label{eq:reff}
\end{equation}
where $C_{\text{link}}$ is the link-layer capacity,
$\eta_{\text{transport}} \in (0,1]$ is the transport efficiency factor
(capturing congestion control, retransmission, and multiplexing
behavior), and $\rho_{\text{overhead}} \in [0,1)$ is the fractional
protocol overhead (headers, encryption, framing).

\subsubsection{Head-of-Line Blocking and Multiplexed Capacity}

Consider $n$ independent streams multiplexed over a single connection.
Under TCP (HTTP/2), a single lost packet blocks all $n$ streams. The
expected throughput per stream under loss rate $p$ is bounded by:
\begin{equation}
  R_{\text{TCP}}^{(i)} \leq \frac{R_{\text{eff}}}{n}
    \cdot (1 - p)^{n-1}
  \label{eq:hol-tcp}
\end{equation}
because a loss in \emph{any} of the $n$ streams stalls all others. Under
QUIC (RFC~9000), streams are independent at the transport layer, so:
\begin{equation}
  R_{\text{QUIC}}^{(i)} \leq \frac{R_{\text{eff}}}{n}
    \cdot (1 - p)
  \label{eq:hol-quic}
\end{equation}
The ratio of effective per-stream throughput is:
\begin{equation}
  \frac{R_{\text{QUIC}}^{(i)}}{R_{\text{TCP}}^{(i)}}
    = \frac{1}{(1-p)^{n-2}}
  \label{eq:hol-ratio}
\end{equation}
\excursus{Excursus: Why $n = 8$?}{In a typical multi-participant video
conference, each peer contributes multiple concurrent streams: a camera video
track, an audio track, and often a screen-share or data channel. A four-person
call with two media tracks each produces $n = 8$ multiplexed streams over a
single connection --- the reference value used throughout this section. Larger
meetings (e.g.\ 10 participants with simulcast layers) push $n$ higher,
amplifying the HOL blocking effect further.}

For $n = 8$ streams (a typical multi-track video conference) at
$p = 0.01$ (1\% loss), this gives a QUIC advantage factor of
$1/(0.99)^6 \approx 1.062$, a modest 6.2\% improvement. At $p = 0.05$
(5\% loss, common on mobile or congested links), the factor is
$1/(0.95)^6 \approx 1.34$, a 34\% throughput advantage per stream.

\begin{figure}[ht]
\centering
\begin{tikzpicture}[>=Stealth]

% --- Shaded region (drawn FIRST so curves appear on top) ---
\fill[orange!8] (2, 0) rectangle (10, 5.3);
\node[font=\tiny\itshape, text=black!40, anchor=north east] at (9.8, 5.2)
  {mobile / congested};
\node[font=\tiny\itshape, text=black!40, anchor=north west] at (0.2, 5.2)
  {fixed broadband};

% --- Grid lines ---
\foreach \y in {1, 2, 3, 4, 5} {
  \draw[gray!20] (0, \y) -- (10, \y);
}

% --- Axes ---
\draw[->, thick] (0, 0) -- (10.5, 0);
\draw[->, thick] (0, 0) -- (0, 5.5);

% --- Axis titles (positioned to clear tick labels) ---
\node[font=\small] at (5, -0.85) {Packet loss rate $p$};
\node[font=\small, rotate=90] at (-0.9, 2.75) {QUIC advantage factor};

% --- X-axis labels ---
\foreach \x/\lab in {0/0, 2/0.02, 4/0.04, 6/0.06, 8/0.08, 10/0.10} {
  \draw (\x, -0.1) -- (\x, 0.1);
  \node[font=\tiny, anchor=north] at (\x, -0.15) {\lab};
}

% --- Y-axis labels ---
\foreach \y/\lab in {0/1.0, 1/1.1, 2/1.2, 3/1.3, 4/1.4, 5/1.5} {
  \draw (-0.1, \y) -- (0.1, \y);
  \node[font=\tiny, anchor=east] at (-0.15, \y) {\lab};
}

% --- n=4 curve: 1/(1-p)^2 ---
\draw[very thick, blue!70]
  (0, 0) .. controls (2.5, 0.5) and (5, 1.1) ..
  (7.5, 2.0) .. controls (8.5, 2.4) and (9.5, 2.9) ..
  (10, 3.1);
\node[font=\scriptsize, blue!70, anchor=south west] at (9.0, 3.0)
  {$n = 4$};

% --- n=8 curve: 1/(1-p)^6 ---
\draw[very thick, orange!80!black]
  (0, 0) .. controls (1.5, 0.5) and (3, 1.7) ..
  (4, 2.5) .. controls (5, 3.5) and (5.5, 4.0) ..
  (6, 4.7);
\node[font=\scriptsize, orange!80!black, anchor=south] at (5.5, 4.8)
  {$n = 8$};

% --- Reference points ---
\fill[orange!80!black] (1.0, 0.62) circle (2.5pt);
\node[font=\tiny, orange!80!black, anchor=south west] at (1.1, 0.62)
  {6.2\%};

\fill[orange!80!black] (5.0, 3.7) circle (2.5pt);
\node[font=\tiny, orange!80!black, anchor=west] at (5.1, 3.7)
  {34\%};

\fill[blue!70] (1.0, 0.20) circle (2.5pt);

% --- Annotation ---
\node[font=\scriptsize\itshape, text=black!60, align=center,
  anchor=north] at (5.0, -1.0)
  {QUIC recovers per-stream throughput closer to $C/n$ (Shannon fair share)\\
   by eliminating head-of-line blocking (eq.~\ref{eq:hol-ratio}).};

\end{tikzpicture}
\caption{QUIC advantage factor $1/(1-p)^{n-2}$ over TCP for multiplexed
  streams. At 5\% loss with $n=8$ streams, QUIC delivers 34\% more
  per-stream throughput by maintaining independence at the transport layer.}
\label{fig:hol-blocking}
\end{figure}

\subsubsection{Connection Establishment Latency}

Let $\tau$ denote the one-way latency (half the round-trip time). The
information-carrying capacity during connection establishment is zero;
the handshake represents pure overhead. The time to first useful byte
is:
\begin{align}
  T_{\text{TCP+TLS\,1.3}} &= 2\tau_{\text{TCP}} + \tau_{\text{TLS}}
    = 3\tau \quad \text{(1-RTT TLS~1.3)} \label{eq:tcp-setup} \\
  T_{\text{QUIC}} &= \tau \quad \text{(1-RTT)}
    \label{eq:quic-setup} \\
  T_{\text{QUIC,\,0-RTT}} &= 0 \quad \text{(resumption)}
    \label{eq:quic-0rtt}
\end{align}
The QUIC 0-RTT case (\ref{eq:quic-0rtt}) eliminates the handshake
entirely for resumed connections. The information lost during
handshake dead time, relative to steady-state capacity, is:
\begin{equation}
  I_{\text{lost}} = C_{\text{link}} \cdot T_{\text{setup}}
  \label{eq:ilost}
\end{equation}
For a 120~Mbps link with $\tau = 25$~ms (typical fixed broadband in
2026), the TCP+TLS~1.3 handshake wastes $120 \times 10^6 \times 0.075
= 9 \times 10^6$ bits (1.125~MB), while QUIC~1-RTT wastes only
$3 \times 10^6$ bits (375~KB).

\subsection{Source Coding Efficiency}

The rate-distortion function $R(D)$ defines the minimum bit rate
required to represent a source at distortion level $D$. For a
Gaussian source with variance $\sigma^2$ under mean squared error
distortion:
\begin{equation}
  R(D) = \frac{1}{2} \log_2 \frac{\sigma^2}{D}
    \quad \text{[bits/sample]}
  \label{eq:rate-distortion}
\end{equation}

\excursus{Excursus: Why Gaussian source?}{The Gaussian memoryless model is used
because it yields the highest (worst-case) $R(D)$ among all sources of equal
variance. Any codec that approaches the Gaussian bound on real video is
therefore doing at least as well as the bound predicts --- the true $R(D)$ for
spatially and temporally correlated video lies below the Gaussian curve, making
it a conservative benchmark.}

No practical codec achieves $R(D)$;\footnote{The $R(D)$ curve in
Figure~\ref{fig:rate-distortion} is the Gaussian memoryless source bound
(equation~\ref{eq:rate-distortion}), which is an upper bound on the true
$R(D)$ for non-Gaussian sources such as video.  Additionally, $R(D)$ assumes
zero excess-distortion probability, whereas practical codecs permit a
frame-level excess-distortion probability on the order of $10^{-6}$ to
$10^{-8}$.  Both factors allow codec operating points measured on real content
to approach or marginally cross the Gaussian bound.} real codecs operate at
some rate $R_{\text{codec}} > R(D)$. We define the coding efficiency gap as:
\begin{equation}
  \Delta = \frac{R_{\text{codec}} - R(D)}{R(D)}
  \label{eq:coding-gap}
\end{equation}

The codecs in the WebRTC stack can be ordered by their empirical
proximity to $R(D)$ for typical video content at conversational
resolutions (720p, 30~fps). Using published Bj\o ntegaard Delta (BD)
rate measurements:
\begin{equation}
  R_{\text{H.264}} > R_{\text{VP9}} > R_{\text{AV1}}
    > R(D)
  \label{eq:codec-ordering}
\end{equation}

Specifically, the empirical rate savings at equivalent PSNR are
approximately:
\begin{align}
  R_{\text{VP9}}  &\approx 0.65 \cdot R_{\text{H.264}}
    \quad (\sim\!35\% \text{ savings}) \label{eq:vp9-savings} \\
  R_{\text{AV1}}  &\approx 0.70 \cdot R_{\text{VP9}}
    \approx 0.46 \cdot R_{\text{H.264}}
    \quad (\sim\!30\% \text{ over VP9})
  \label{eq:av1-savings}
\end{align}

\begin{figure}[ht]
\centering
\begin{tikzpicture}[>=Stealth]

% --- Axes ---
\draw[->, thick] (0, 0) -- (10.5, 0);
\draw[->, thick] (0, 0) -- (0, 6.5);

% --- Axis titles ---
\node[font=\small] at (5, -0.85) {Distortion $D$};
\node[font=\small, rotate=90] at (-0.9, 3.25) {Rate [Mbps]};

% --- Y-axis rate labels ---
\foreach \y/\lab in {1.0/0.5, 2.0/1.0, 3.0/1.5, 4.0/2.0, 5.0/2.5, 6.0/3.0} {
  \draw (-0.1, \y) -- (0.1, \y);
  \node[font=\tiny, anchor=east] at (-0.15, \y) {\lab};
}

% --- R(D) Shannon bound curve ---
\draw[very thick, green!60!black]
  (1.0, 5.8) .. controls (2.0, 3.5) and (3.5, 2.0) ..
  (5.0, 1.3) .. controls (6.5, 0.8) and (8.0, 0.55) ..
  (9.5, 0.4);
\node[font=\scriptsize, green!60!black, anchor=north west] at (8.5, 0.55)
  {$R(D)$ bound};

% --- Reference distortion line (720p/30fps operating point) ---
\draw[thin, gray, dashed] (4.0, 0) -- (4.0, 6.2);
\node[font=\tiny, gray, anchor=south] at (4.0, 6.2) {720p/30fps};

% --- R(D) value at reference distortion ---
\fill[green!60!black] (4.0, 1.55) circle (2pt);

% --- Codec operating points at reference distortion ---
% H.264: ~2.5 Mbps
\draw[thick, dashed, red!70!black] (0.3, 4.8) -- (9.5, 4.8);
\fill[red!70!black] (4.0, 4.8) circle (3pt);
\node[font=\scriptsize, red!70!black, anchor=west] at (9.6, 4.8) {H.264};

% VP9: ~1.6 Mbps (0.65 * H.264)
\draw[thick, dashed, blue!70] (0.3, 3.3) -- (9.5, 3.3);
\fill[blue!70] (4.0, 3.3) circle (3pt);
\node[font=\scriptsize, blue!70, anchor=west] at (9.6, 3.3) {VP9};

% AV1: ~1.15 Mbps (0.46 * H.264)
\draw[thick, dashed, orange!80!black] (0.3, 2.3) -- (9.5, 2.3);
\fill[orange!80!black] (4.0, 2.3) circle (3pt);
\node[font=\scriptsize, orange!80!black, anchor=west] at (9.6, 2.3) {AV1};

% --- Gap annotations (braces) ---
% H.264 gap
\draw[decorate, decoration={brace, amplitude=4pt, mirror},
  thick, red!50!black]
  (3.6, 1.55) -- (3.6, 4.8)
  node[midway, anchor=east, font=\scriptsize, xshift=-5pt]
  {$\Delta_{\text{H.264}}$};

% AV1 gap
\draw[decorate, decoration={brace, amplitude=4pt},
  thick, orange!60!black]
  (4.4, 1.55) -- (4.4, 2.3)
  node[midway, anchor=west, font=\scriptsize, xshift=5pt]
  {$\Delta_{\text{AV1}}$};

% --- Annotation ---
\node[font=\scriptsize\itshape, text=black!60, align=center,
  anchor=north] at (5.0, -0.5)
  {The Shannon rate-distortion bound $R(D)$ is the theoretical minimum.\\
   AV1 narrows the gap $\Delta$ to ${\sim}0.4$ by 2030 (eq.~\ref{eq:coding-gap}).};

\end{tikzpicture}
\caption{Codec operating rates vs.\ the Shannon rate-distortion bound $R(D)$
  at 720p/30fps. Each codec operates above the theoretical minimum; the gap
  $\Delta$ (eq.~\ref{eq:coding-gap}) shrinks from H.264 through VP9 to AV1,
  approaching but never reaching the bound.}
\label{fig:rate-distortion}
\end{figure}

\subsubsection{Codec Entropy and WebCodecs}

The WebCodecs API (W3C Working Draft) exposes per-frame codec access,
allowing measurement and control of the instantaneous encoding rate.
For a video frame $f_t$ at time $t$, the encoder output has empirical
entropy:
\begin{equation}
  \hat{H}(f_t) = -\sum_{s \in \mathcal{S}}
    p(s \mid f_t) \log_2 p(s \mid f_t)
  \label{eq:frame-entropy}
\end{equation}
where $\mathcal{S}$ is the set of symbols emitted by the entropy coder
(CABAC for H.264/AV1, arithmetic coding for VP9). WebCodecs enables
applications to observe $|E(f_t)|$ (the encoded frame size in bits)
directly, which serves as a practical upper bound on
$\hat{H}(f_t)$:
\begin{equation}
  \hat{H}(f_t) \leq |E(f_t)| \leq R_{\text{codec}} / \text{fps}
  \label{eq:entropy-bound}
\end{equation}

\subsection{Effective Information Rate by Time Horizon}

We now combine transport and source coding measures to estimate the
effective information rate for a representative real-time video
session: a single 720p/30fps video call with one audio channel.

\excursus{Excursus: Why 720p?}{Throughout this appendix, 720p/30fps is used as
a fixed reference for cross-era comparison, not as a prediction that 720p
remains the dominant conferencing resolution. The bandwidth surplus quantified
below is precisely what enables migration to 1080p (feasible by ${\sim}$2028
with AV1 hardware encode at today's 720p bit rates) and 4K (feasible by
${\sim}$2030). The 720p baseline makes the efficiency gains legible: the same
task gets cheaper, freeing capacity for higher quality or more streams.}

Define the \emph{net information rate} as:
\begin{equation}
  I_{\text{net}} = R_{\text{codec}} \cdot \eta_{\text{transport}}
    \cdot (1 - \rho_{\text{overhead}})
    \cdot (1 - p_{\text{loss}})
  \label{eq:inet}
\end{equation}

\begin{table}[ht]
\centering
\small
\renewcommand{\arraystretch}{1.4}
\begin{tabularx}{\textwidth}{lXXX}
\toprule
\textbf{Parameter} & \textbf{March 2026} & \textbf{2027--2028} &
  \textbf{2029--2031} \\
\midrule
Median link capacity &
  120~Mbps &
  $\sim$150~Mbps &
  200--300~Mbps \\
Dominant codec (WebRTC) &
  VP9 / H.264 &
  VP9 / AV1 (HW) &
  AV1 (HW default) \\
Codec rate (720p/30fps) &
  1.5--2.5~Mbps &
  1.0--1.8~Mbps &
  0.7--1.2~Mbps \\
$\eta_{\text{transport}}$ &
  0.92 (QUIC) &
  0.93 &
  0.94 \\
$\rho_{\text{overhead}}$ (SRTP+QUIC) &
  0.06 &
  0.05 &
  0.04 \\
Typical loss $p$ &
  0.01--0.03 &
  0.01--0.02 &
  0.005--0.01 \\
$I_{\text{net}}$ (720p video) &
  1.3--2.2~Mbps &
  0.9--1.6~Mbps &
  0.6--1.1~Mbps \\
$I_{\text{net}}$ as \% of $C_{\text{link}}$ &
  1.1--1.8\% &
  0.6--1.1\% &
  0.2--0.6\% \\
\bottomrule
\end{tabularx}
\caption{Effective information rate for a 720p/30fps WebRTC video stream.
  The declining ratio $I_{\text{net}} / C_{\text{link}}$ reflects both
  improved codec efficiency and growing link capacity, freeing bandwidth
  for higher resolutions, additional streams, or concurrent data
  channels.}
\label{tab:inet}
\end{table}

The declining ratio $I_{\text{net}} / C_{\text{link}}$ in
Table~\ref{tab:inet} is the central quantitative observation: by 2030, a
720p video call consumes less than 0.5\% of available link capacity. The
surplus capacity enables qualitative changes --- 4K video conferencing,
multi-stream layouts, or concurrent data channels via WebTransport ---
that are quantitatively infeasible in 2026.

\begin{figure}[ht]
\centering
\begin{tikzpicture}[>=Stealth]

% --- Axes ---
\draw[->, thick] (0, 0) -- (11, 0);
\draw[->, thick] (0, 0) -- (0, 5.5);

% --- Axis title ---
\node[font=\small, rotate=90] at (-0.9, 2.75) {Mbps};

% --- Y-axis labels ---
\foreach \y/\lab in {0/0, 1/50, 2/100, 3/150, 4/200, 5/250} {
  \draw (-0.1, \y) -- (0.1, \y);
  \node[font=\tiny, anchor=east] at (-0.15, \y) {\lab};
}

% --- Group: March 2026 ---
% C_link = 120 Mbps -> y = 2.4
\fill[blue!12] (0.8, 0) rectangle (2.4, 2.4);
\draw[blue!50, thick] (0.8, 0) rectangle (2.4, 2.4);
% I_net = 1.75 Mbps (midpoint) -> y = 0.035
\fill[orange!60] (0.8, 0) rectangle (2.4, 0.07);
\draw[orange!80!black, thick] (0.8, 0) rectangle (2.4, 0.07);
\node[font=\footnotesize, anchor=south] at (1.6, 2.5) {120};
\node[font=\tiny, orange!80!black] at (1.6, -0.4) {1.5\%};
\node[font=\small, anchor=north] at (1.6, -0.7) {2026};

% --- Group: 2027--2028 ---
% C_link = 150 Mbps -> y = 3.0
\fill[blue!12] (4.0, 0) rectangle (5.6, 3.0);
\draw[blue!50, thick] (4.0, 0) rectangle (5.6, 3.0);
% I_net = 1.25 Mbps (midpoint) -> y = 0.025
\fill[orange!60] (4.0, 0) rectangle (5.6, 0.05);
\draw[orange!80!black, thick] (4.0, 0) rectangle (5.6, 0.05);
\node[font=\footnotesize, anchor=south] at (4.8, 3.1) {150};
\node[font=\tiny, orange!80!black] at (4.8, -0.4) {0.8\%};
\node[font=\small, anchor=north] at (4.8, -0.7) {2027--28};

% --- Group: 2029--2031 ---
% C_link = 250 Mbps -> y = 5.0
\fill[blue!12] (7.2, 0) rectangle (8.8, 5.0);
\draw[blue!50, thick] (7.2, 0) rectangle (8.8, 5.0);
% I_net = 0.85 Mbps (midpoint) -> y = 0.017
\fill[orange!60] (7.2, 0) rectangle (8.8, 0.034);
\draw[orange!80!black, thick] (7.2, 0) rectangle (8.8, 0.034);
\node[font=\footnotesize, anchor=south] at (8.0, 5.1) {250};
\node[font=\tiny, orange!80!black] at (8.0, -0.4) {0.3\%};
\node[font=\small, anchor=north] at (8.0, -0.7) {2029--31};

% --- Legend ---
\node[anchor=west, font=\footnotesize] at (0.5, 5.8) {%
  \tikz[baseline=-0.5ex]{\fill[blue!12] (0,0) rectangle (0.3,0.3);
    \draw[blue!50, thick] (0,0) rectangle (0.3,0.3);}
  ~$C_{\text{link}}$ (channel capacity) \qquad
  \tikz[baseline=-0.5ex]{\fill[orange!60] (0,0) rectangle (0.3,0.3);
    \draw[orange!80!black, thick] (0,0) rectangle (0.3,0.3);}
  ~$I_{\text{net}}$ (720p stream)};

% --- Annotation ---
\node[font=\scriptsize\itshape, text=black!60, align=center,
  anchor=north] at (5.5, -1.2)
  {A single 720p call uses a vanishing fraction of Shannon-derived\\
   channel capacity, freeing bandwidth for 4K, multi-stream, and data channels.};

\end{tikzpicture}
\caption{Channel capacity $C_{\text{link}}$ vs.\ effective information rate
  $I_{\text{net}}$ for a single 720p/30fps stream across time horizons. The
  orange bar (stream rate) is barely visible against the blue bar (capacity),
  illustrating the growing surplus as codecs improve and bandwidth increases.}
\label{fig:inet-capacity}
\end{figure}

\subsection{Security Entropy and Key Space}

The DTLS-SRTP security architecture (RFCs~8827, 3711) protects all
WebRTC media. Information-theoretic security requires that the key
entropy exceed the information content of the protected material.

For a DTLS~1.3 handshake with ECDHE on P-256, the key agreement
provides:
\begin{equation}
  H_{\text{key}} = 128 \text{~bits of security}
  \label{eq:key-entropy}
\end{equation}

The SRTP master key is 128 bits (AES-128-GCM), giving an exhaustive
search cost of $2^{128}$ operations. The total information content of a
one-hour 720p video call at 2~Mbps is:
\begin{equation}
  I_{\text{call}} = 2 \times 10^6 \times 3600
    = 7.2 \times 10^9 \text{~bits}
    \approx 2^{32.7} \text{~bits}
  \label{eq:call-content}
\end{equation}

The ratio of key entropy to protected content:
\begin{equation}
  \frac{H_{\text{key}}}{I_{\text{call}}}
    = \frac{128}{32.7} \approx 3.91
  \label{eq:security-margin}
\end{equation}

This ratio remains comfortably above~1 even for much longer sessions.
For a continuous 24-hour 4K AV1 stream at 15~Mbps:
\begin{equation}
  I_{\text{stream}} = 15 \times 10^6 \times 86400
    = 1.296 \times 10^{12} \text{~bits}
    \approx 2^{40.2} \text{~bits}
\end{equation}
yielding $H_{\text{key}} / I_{\text{stream}} \approx 3.18$, still
well above the threshold.

\subsubsection{End-to-End Encryption via Encoded Transforms}

The RTCRtpScriptTransform API (Encoded Transforms) enables
application-layer E2EE within the WebRTC pipeline. From an
information-theoretic perspective, this introduces a second encryption
layer with its own key, creating a \emph{nested channel}:
\begin{equation}
  X \xrightarrow{\;K_{\text{E2EE}}\;}
  E_1(X) \xrightarrow{\;K_{\text{SRTP}}\;}
  E_2(E_1(X))
  \label{eq:nested}
\end{equation}

The mutual information between the ciphertext and plaintext, given
neither key, satisfies:
\begin{equation}
  I\!\left(X;\, E_2(E_1(X))\right) \leq
    \min\!\big(H(K_{\text{E2EE}})^{-1},\;
    H(K_{\text{SRTP}})^{-1}\big)
    \cdot I_{\text{call}}
  \label{eq:mutual-e2ee}
\end{equation}

In practice, with $H(K_{\text{E2EE}}) = H(K_{\text{SRTP}}) = 128$
bits, this quantity is negligible ($\ll 1$ bit for any practical
session length), confirming that the nested encryption provides
information-theoretic confidentiality well beyond computational
security guarantees.

Today (March~2026), Encoded Transforms remain Chrome-only. By
2027--2028 (Section~8.3), cross-browser standardization will make
this nested channel universally available. By 2029--2031, SFrame
for group calls will extend this to $n$-party sessions, where the
mutual information between any non-member and the group plaintext
remains negligible:
\begin{equation}
  I\!\left(X;\, E(X) \;\middle|\; K_i \notin
    \{K_1, \ldots, K_n\}\right) \approx 0
  \label{eq:sframe}
\end{equation}

\subsection{WebTransport: Unreliable Datagrams and Channel Capacity}

WebTransport (RFC~9297) introduces unreliable datagrams to browser-server
communication. For latency-sensitive applications (game state, sensor
telemetry), the relevant measure is not throughput but
\emph{freshness-weighted information rate}. Define the value of a
datagram as exponentially decaying with delivery latency $\delta$:
\begin{equation}
  V(\delta) = e^{-\lambda \delta}
  \label{eq:freshness}
\end{equation}
where $\lambda > 0$ is an application-specific decay rate. The
effective information rate under unreliable delivery is:
\begin{equation}
  I_{\text{WT}} = R_{\text{send}} \cdot (1 - p_{\text{loss}})
    \cdot \mathbb{E}\!\left[e^{-\lambda \delta}\right]
  \label{eq:iwt}
\end{equation}

Under reliable delivery (TCP/WebSocket), a retransmission delays
subsequent data, so:
\begin{equation}
  I_{\text{WS}} = R_{\text{send}}
    \cdot \mathbb{E}\!\left[e^{-\lambda(\delta + p \cdot 2\tau)}\right]
  \label{eq:iws}
\end{equation}

The ratio $I_{\text{WT}} / I_{\text{WS}}$ exceeds 1 when:
\begin{equation}
  (1 - p) \cdot e^{\lambda \cdot p \cdot 2\tau} > 1
  \quad \Longleftrightarrow \quad
  \lambda > \frac{-\ln(1 - p)}{2 p \tau}
  \label{eq:wt-advantage}
\end{equation}

For $p = 0.02$ and $\tau = 25$~ms, this gives $\lambda > 20.2$~s$^{-1}$
(half-life $< 34$~ms).

\excursus{Excursus: Why $p = 0.02$ and $\tau = 25$\,ms?}{These values represent
median fixed-broadband conditions in 2026 (Table~\ref{tab:inet}). On mobile or
congested links ($p \approx 0.05$, $\tau \approx 50$\,ms), the threshold
$\lambda^*$ drops to ${\sim}10$\,s$^{-1}$, broadening WebTransport's advantage
to applications with data half-lives up to ${\sim}70$\,ms. The fixed-broadband
reference is thus conservative: mobile scenarios favor unreliable delivery even
more strongly.}

Any application where data older than
$\sim$30~ms is useless benefits from WebTransport's unreliable
mode over WebSocket --- precisely the regime of interactive gaming, live
cursor tracking, and real-time sensor feeds.

As of March~2026, Safari does not support WebTransport
(Section~5), limiting coverage to $\sim$82\%. By 2027--2028
(Section~8.2), universal support will allow applications to rely on
(\ref{eq:iwt}) without fallback paths.

\begin{figure}[ht]
\centering
\begin{tikzpicture}[>=Stealth]

% --- Axes ---
\draw[->, thick] (0, 0) -- (10.5, 0);
\draw[->, thick] (0, 0) -- (0, 5.5);

% --- Axis titles (positioned to clear tick labels) ---
\node[font=\small] at (5, -0.75)
  {Freshness decay rate $\lambda$ [s$^{-1}$]};
\node[font=\small, rotate=90] at (-1.1, 2.75)
  {$I_{\text{WT}} / I_{\text{WS}}$};

% --- X-axis labels ---
\foreach \x/\lab in {0/0, 1.67/10, 3.33/20, 5.0/30, 6.67/40, 8.33/50, 10/60} {
  \draw (\x, -0.1) -- (\x, 0.1);
  \node[font=\tiny, anchor=north] at (\x, -0.15) {\lab};
}

% --- Y-axis labels (0.96 to 1.06, each 0.02 = 1 unit) ---
\foreach \y/\lab in {0/0.96, 1/0.98, 2/1.00, 3/1.02, 4/1.04, 5/1.06} {
  \draw (-0.1, \y) -- (0.1, \y);
  \node[font=\tiny, anchor=east] at (-0.15, \y) {\lab};
}

% --- Ratio = 1.0 line (y=2 maps to ratio 1.0) ---
\draw[thin, gray, dashed] (0, 2) -- (10, 2);

% --- Shade WebTransport advantage region ---
% lambda* = 20.2 -> x = 3.37
\fill[green!10] (3.37, 2) rectangle (10, 5.3);

% --- I_WT/I_WS curve: 0.98*exp(lambda*0.001) ---
% Y range: 0.96-1.06 (each 0.02 = 1 unit). lambda=0: y=1, lambda*=20.2: y=2, lambda=60: y=4.05
\draw[very thick, orange!80!black]
  (0, 1.0) .. controls (1.5, 1.3) and (2.5, 1.7) ..
  (3.37, 2.0) .. controls (4.5, 2.4) and (6, 2.8) ..
  (8, 3.4) .. controls (9, 3.7) and (9.5, 3.85) ..
  (10, 4.05);

% --- Threshold line ---
\draw[thick, dashed, red!60!black] (3.37, 0) -- (3.37, 5.2);
\node[font=\scriptsize, red!60!black, anchor=south, rotate=90] at (3.17, 3.5)
  {$\lambda^* = 20.2$ s$^{-1}$};

% --- Application regime labels ---
\node[font=\scriptsize\itshape, text=blue!60, anchor=north] at (1.5, 5.3)
  {video conf};
\node[font=\scriptsize\itshape, text=green!50!black, anchor=north] at (7.0, 5.3)
  {gaming / cursor / sensors};

% --- Mark advantage/disadvantage ---
\node[font=\tiny, text=red!50!black, anchor=north east] at (3.2, 1.8)
  {WebSocket wins};
\node[font=\tiny, text=green!50!black, anchor=north west] at (3.5, 4.5)
  {WebTransport wins};

% --- Annotation ---
\node[font=\scriptsize\itshape, text=black!60, align=center,
  anchor=north] at (5.0, -1.0)
  {For $p = 0.02$, $\tau = 25$~ms. Applications where data older than\\
   ${\sim}$30~ms is stale gain from unreliable delivery (eq.~\ref{eq:wt-advantage}).};

\end{tikzpicture}
\caption{Freshness-weighted information rate ratio $I_{\text{WT}} / I_{\text{WS}}$
  as a function of decay rate $\lambda$. Above the threshold $\lambda^*$,
  WebTransport's unreliable datagrams deliver more useful information per
  unit time than WebSocket's reliable delivery, because retransmission
  wastes capacity on stale data.}
\label{fig:freshness}
\end{figure}

\subsection{Projections Summary and Practical Implications}

\begin{table}[ht]
\centering
\small
\renewcommand{\arraystretch}{1.4}
\begin{tabularx}{\textwidth}{lXXX}
\toprule
\textbf{Measure} & \textbf{March 2026} & \textbf{2027--2028} &
  \textbf{2029--2031} \\
\midrule
Channel capacity $C_{\text{link}}$ &
  120~Mbps &
  150~Mbps &
  250~Mbps \\
HOL blocking factor $(1\!-\!p)^{n-2}$ at $n\!=\!8$ &
  0.94 ($p\!=\!0.01$) &
  0.94 &
  0.97 ($p\!=\!0.005$) \\
Source coding gap $\Delta$ &
  VP9: $\sim$0.8 &
  AV1 HW: $\sim$0.5 &
  AV1 opt: $\sim$0.4 \\
$I_{\text{net}}$ per 720p stream &
  1.3--2.2~Mbps &
  0.9--1.6~Mbps &
  0.6--1.1~Mbps \\
Max 720p streams in $C_{\text{link}}$ &
  $\sim$55--90 &
  $\sim$95--165 &
  $\sim$230--415 \\
E2EE availability &
  Chrome only &
  Cross-browser &
  SFrame $n$-party \\
WebTransport coverage &
  82\% &
  $\sim$100\% &
  100\% \\
Freshness threshold $\lambda^*$ &
  20.2~s$^{-1}$ &
  $\sim$18~s$^{-1}$ &
  $\sim$15~s$^{-1}$ \\
\bottomrule
\end{tabularx}
\caption{Information-theoretic measures across time horizons. $\Delta$
  is the coding efficiency gap (\ref{eq:coding-gap}), $I_{\text{net}}$
  is the net information rate (\ref{eq:inet}), and $\lambda^*$ is the
  freshness decay threshold above which WebTransport dominates WebSocket
  (\ref{eq:wt-advantage}).}
\label{tab:projections}
\end{table}

The trajectory is clear: the combination of improving source coding
(AV1 hardware), expanding channel capacity (broadband growth), and
maturing transport protocols (QUIC universal, WebTransport universal)
drives the \emph{information cost per unit of communicated experience}
steadily downward. The remaining question is not whether the
Shannon capacity suffices --- it already does, by orders of magnitude
--- but how efficiently the protocol stack converts that capacity into
delivered, fresh, secure information at the application layer. By
2030, the gap between theoretical capacity and realized throughput will
be the narrowest it has ever been for browser-based real-time
communication.

\subsubsection{What the Theory Means for Applications}

The equations and figures in this appendix are not purely academic. Each
theoretical result maps to a concrete design decision that application
developers face today and over the next five years.

\textbf{Codec selection determines bandwidth cost, not quality.}
The rate-distortion analysis (Section~A.2, Figure~\ref{fig:rate-distortion})
shows that AV1 operates at roughly 46\% of H.264's bit rate for equivalent
perceptual quality at 720p. In practice, this means a WebRTC application
switching from H.264 to AV1 can either halve its bandwidth consumption at
the same quality or substantially increase resolution (e.g., 720p to 1080p)
within the same bandwidth envelope. By 2030, with AV1 hardware encoding
ubiquitous and the coding gap $\Delta$ approaching 0.4, further codec
improvements yield diminishing returns --- the stack will be operating close
enough to the Shannon rate-distortion bound that the next generation of
efficiency gains will come from transport and architecture, not codecs alone.

\textbf{QUIC's advantage is situational but decisive on bad networks.}
The head-of-line blocking analysis (Section~A.1.1,
Figure~\ref{fig:hol-blocking}) shows that QUIC's per-stream throughput
advantage over TCP is modest on good networks (6\% at 1\% loss) but
substantial on degraded ones (34\% at 5\% loss). The practical implication
is that applications targeting mobile users, emerging markets, or congested
enterprise networks --- precisely the environments where real-time
communication is hardest --- benefit disproportionately from QUIC-based
transports. For a video conference with 8 participants on a cellular
connection experiencing 3--5\% packet loss, the difference between QUIC and
TCP is the difference between usable and degraded video.

\textbf{Connection latency is a solved problem for returning users.}
The handshake analysis (Section~A.1.2) shows that QUIC 0-RTT eliminates
connection establishment overhead entirely for resumed connections. For a
WebTransport application that a user opens daily --- a collaboration tool,
a live dashboard, a recurring meeting client --- the time to first useful
byte drops to zero. The 1.125~MB of wasted capacity during a TCP+TLS
handshake on a 120~Mbps link may seem small, but for latency-sensitive
applications where the first 100~ms of perceived responsiveness determines
user experience, the difference is measurable.

\textbf{Bandwidth surplus enables architectural ambition.}
The effective information rate analysis (Section~A.3,
Figure~\ref{fig:inet-capacity}) reveals the most consequential practical
finding: a single 720p video stream consumes a vanishing fraction of
available link capacity (under 0.5\% by 2030). This surplus is not merely
``headroom'' --- it is the resource that makes new application architectures
feasible. A 250~Mbps link can theoretically carry 230--415 simultaneous 720p
streams. In practice, this means applications can afford to run multiple
concurrent media flows (video, screen share, spatial audio), data channels
(collaborative editing, file transfer), and telemetry streams without
contending for bandwidth. The constraint shifts from ``can the network carry
this?''\ to ``can the application manage this complexity?''

\textbf{End-to-end encryption adds negligible overhead.}
The security analysis (Section~A.4) confirms that the nested DTLS-SRTP plus
E2EE architecture imposes no meaningful information-theoretic cost. The key
entropy ($2^{128}$) exceeds the information content of any practical session
by a factor of at least 3. For application developers, this means E2EE can
be treated as a default rather than a premium feature --- there is no
bandwidth, latency, or capacity penalty for encrypting at the application
layer on top of the transport-layer encryption that WebRTC already mandates.
The only constraint is API availability, which moves from Chrome-only (2026)
to universal (2028).

\textbf{WebTransport unlocks real-time data that WebSocket cannot serve.}
The freshness analysis (Section~A.5, Figure~\ref{fig:freshness}) quantifies
when unreliable delivery outperforms reliable delivery: any data with a
useful half-life below approximately 34~ms benefits from WebTransport's
unreliable datagrams. This covers a wide class of interactive applications
--- game state synchronization, live cursor and pointer sharing, real-time
sensor feeds, collaborative whiteboard strokes --- where retransmitting a
stale update is worse than dropping it. Before WebTransport, these
applications had to choose between WebSocket (reliable but stale under loss)
and WebRTC DataChannel (unreliable mode available but requiring the full
peer-to-peer ICE/DTLS machinery). WebTransport provides the unreliable
mode with a simple client-server model over HTTP/3, removing the
architectural overhead.

\textbf{The overall picture: the stack converges on efficiency.}
Taken together, the information-theoretic measures trace a consistent
trajectory. The protocol stack is converging toward the theoretical limits
established by Shannon: codecs approach the rate-distortion bound, transports
eliminate unnecessary blocking and handshake overhead, and encryption adds
negligible cost. By 2030, the practical distance between what the standards
deliver and what information theory permits will be the smallest it has ever
been for browser-based real-time communication. The consequence for
application developers is that the infrastructure --- standards, hardware,
and networks --- will no longer be the binding constraint. The remaining
challenge is purely architectural: how to compose these efficient, universal
primitives into applications that fully exploit the capacity and low latency
they provide.

%   % appendix-video-compression.tex
% Standalone appendix — include from retrospective.tex via:
%   \appendix
%   % appendix-info-theory.tex
% Standalone appendix — include from retrospective.tex via:
%   \appendix
%   % appendix-info-theory.tex
% Standalone appendix — include from retrospective.tex via:
%   \appendix
%   \input{appendix-info-theory}
%
% Requires in the preamble: amsmath, amssymb, tabularx, booktabs

\section{Information-Theoretic View}
\label{app:info-theory}

This appendix applies information-theoretic measures to the real-time
web communication stack described in this document. We quantify channel
capacity, source coding efficiency, and effective information rates
across three time horizons: today (March~2026), the near term
(2027--2028), and the medium term (2029--2031). The analysis draws on
the standards, codec parameters, and infrastructure projections
developed in the main text.

\subsection{Channel Capacity of the Transport Layer}

The Shannon-Hartley theorem gives the theoretical maximum information
rate for a channel with additive white Gaussian noise:
\begin{equation}
  C = B \log_2\!\left(1 + \frac{S}{N}\right) \quad \text{[bits/s]}
  \label{eq:shannon}
\end{equation}
where $B$ is the channel bandwidth in Hz and $S/N$ is the
signal-to-noise ratio.

While the physical-layer capacity of a broadband link is governed by
(\ref{eq:shannon}), the \emph{effective} capacity available to the
application is reduced by protocol overhead, congestion control, and
transport inefficiencies. We define the effective application-layer
throughput as:
\begin{equation}
  R_{\text{eff}} = C_{\text{link}} \cdot \eta_{\text{transport}}
    \cdot (1 - \rho_{\text{overhead}})
  \label{eq:reff}
\end{equation}
where $C_{\text{link}}$ is the link-layer capacity,
$\eta_{\text{transport}} \in (0,1]$ is the transport efficiency factor
(capturing congestion control, retransmission, and multiplexing
behavior), and $\rho_{\text{overhead}} \in [0,1)$ is the fractional
protocol overhead (headers, encryption, framing).

\subsubsection{Head-of-Line Blocking and Multiplexed Capacity}

Consider $n$ independent streams multiplexed over a single connection.
Under TCP (HTTP/2), a single lost packet blocks all $n$ streams. The
expected throughput per stream under loss rate $p$ is bounded by:
\begin{equation}
  R_{\text{TCP}}^{(i)} \leq \frac{R_{\text{eff}}}{n}
    \cdot (1 - p)^{n-1}
  \label{eq:hol-tcp}
\end{equation}
because a loss in \emph{any} of the $n$ streams stalls all others. Under
QUIC (RFC~9000), streams are independent at the transport layer, so:
\begin{equation}
  R_{\text{QUIC}}^{(i)} \leq \frac{R_{\text{eff}}}{n}
    \cdot (1 - p)
  \label{eq:hol-quic}
\end{equation}
The ratio of effective per-stream throughput is:
\begin{equation}
  \frac{R_{\text{QUIC}}^{(i)}}{R_{\text{TCP}}^{(i)}}
    = \frac{1}{(1-p)^{n-2}}
  \label{eq:hol-ratio}
\end{equation}
\excursus{Excursus: Why $n = 8$?}{In a typical multi-participant video
conference, each peer contributes multiple concurrent streams: a camera video
track, an audio track, and often a screen-share or data channel. A four-person
call with two media tracks each produces $n = 8$ multiplexed streams over a
single connection --- the reference value used throughout this section. Larger
meetings (e.g.\ 10 participants with simulcast layers) push $n$ higher,
amplifying the HOL blocking effect further.}

For $n = 8$ streams (a typical multi-track video conference) at
$p = 0.01$ (1\% loss), this gives a QUIC advantage factor of
$1/(0.99)^6 \approx 1.062$, a modest 6.2\% improvement. At $p = 0.05$
(5\% loss, common on mobile or congested links), the factor is
$1/(0.95)^6 \approx 1.34$, a 34\% throughput advantage per stream.

\begin{figure}[ht]
\centering
\begin{tikzpicture}[>=Stealth]

% --- Shaded region (drawn FIRST so curves appear on top) ---
\fill[orange!8] (2, 0) rectangle (10, 5.3);
\node[font=\tiny\itshape, text=black!40, anchor=north east] at (9.8, 5.2)
  {mobile / congested};
\node[font=\tiny\itshape, text=black!40, anchor=north west] at (0.2, 5.2)
  {fixed broadband};

% --- Grid lines ---
\foreach \y in {1, 2, 3, 4, 5} {
  \draw[gray!20] (0, \y) -- (10, \y);
}

% --- Axes ---
\draw[->, thick] (0, 0) -- (10.5, 0);
\draw[->, thick] (0, 0) -- (0, 5.5);

% --- Axis titles (positioned to clear tick labels) ---
\node[font=\small] at (5, -0.85) {Packet loss rate $p$};
\node[font=\small, rotate=90] at (-0.9, 2.75) {QUIC advantage factor};

% --- X-axis labels ---
\foreach \x/\lab in {0/0, 2/0.02, 4/0.04, 6/0.06, 8/0.08, 10/0.10} {
  \draw (\x, -0.1) -- (\x, 0.1);
  \node[font=\tiny, anchor=north] at (\x, -0.15) {\lab};
}

% --- Y-axis labels ---
\foreach \y/\lab in {0/1.0, 1/1.1, 2/1.2, 3/1.3, 4/1.4, 5/1.5} {
  \draw (-0.1, \y) -- (0.1, \y);
  \node[font=\tiny, anchor=east] at (-0.15, \y) {\lab};
}

% --- n=4 curve: 1/(1-p)^2 ---
\draw[very thick, blue!70]
  (0, 0) .. controls (2.5, 0.5) and (5, 1.1) ..
  (7.5, 2.0) .. controls (8.5, 2.4) and (9.5, 2.9) ..
  (10, 3.1);
\node[font=\scriptsize, blue!70, anchor=south west] at (9.0, 3.0)
  {$n = 4$};

% --- n=8 curve: 1/(1-p)^6 ---
\draw[very thick, orange!80!black]
  (0, 0) .. controls (1.5, 0.5) and (3, 1.7) ..
  (4, 2.5) .. controls (5, 3.5) and (5.5, 4.0) ..
  (6, 4.7);
\node[font=\scriptsize, orange!80!black, anchor=south] at (5.5, 4.8)
  {$n = 8$};

% --- Reference points ---
\fill[orange!80!black] (1.0, 0.62) circle (2.5pt);
\node[font=\tiny, orange!80!black, anchor=south west] at (1.1, 0.62)
  {6.2\%};

\fill[orange!80!black] (5.0, 3.7) circle (2.5pt);
\node[font=\tiny, orange!80!black, anchor=west] at (5.1, 3.7)
  {34\%};

\fill[blue!70] (1.0, 0.20) circle (2.5pt);

% --- Annotation ---
\node[font=\scriptsize\itshape, text=black!60, align=center,
  anchor=north] at (5.0, -1.0)
  {QUIC recovers per-stream throughput closer to $C/n$ (Shannon fair share)\\
   by eliminating head-of-line blocking (eq.~\ref{eq:hol-ratio}).};

\end{tikzpicture}
\caption{QUIC advantage factor $1/(1-p)^{n-2}$ over TCP for multiplexed
  streams. At 5\% loss with $n=8$ streams, QUIC delivers 34\% more
  per-stream throughput by maintaining independence at the transport layer.}
\label{fig:hol-blocking}
\end{figure}

\subsubsection{Connection Establishment Latency}

Let $\tau$ denote the one-way latency (half the round-trip time). The
information-carrying capacity during connection establishment is zero;
the handshake represents pure overhead. The time to first useful byte
is:
\begin{align}
  T_{\text{TCP+TLS\,1.3}} &= 2\tau_{\text{TCP}} + \tau_{\text{TLS}}
    = 3\tau \quad \text{(1-RTT TLS~1.3)} \label{eq:tcp-setup} \\
  T_{\text{QUIC}} &= \tau \quad \text{(1-RTT)}
    \label{eq:quic-setup} \\
  T_{\text{QUIC,\,0-RTT}} &= 0 \quad \text{(resumption)}
    \label{eq:quic-0rtt}
\end{align}
The QUIC 0-RTT case (\ref{eq:quic-0rtt}) eliminates the handshake
entirely for resumed connections. The information lost during
handshake dead time, relative to steady-state capacity, is:
\begin{equation}
  I_{\text{lost}} = C_{\text{link}} \cdot T_{\text{setup}}
  \label{eq:ilost}
\end{equation}
For a 120~Mbps link with $\tau = 25$~ms (typical fixed broadband in
2026), the TCP+TLS~1.3 handshake wastes $120 \times 10^6 \times 0.075
= 9 \times 10^6$ bits (1.125~MB), while QUIC~1-RTT wastes only
$3 \times 10^6$ bits (375~KB).

\subsection{Source Coding Efficiency}

The rate-distortion function $R(D)$ defines the minimum bit rate
required to represent a source at distortion level $D$. For a
Gaussian source with variance $\sigma^2$ under mean squared error
distortion:
\begin{equation}
  R(D) = \frac{1}{2} \log_2 \frac{\sigma^2}{D}
    \quad \text{[bits/sample]}
  \label{eq:rate-distortion}
\end{equation}

\excursus{Excursus: Why Gaussian source?}{The Gaussian memoryless model is used
because it yields the highest (worst-case) $R(D)$ among all sources of equal
variance. Any codec that approaches the Gaussian bound on real video is
therefore doing at least as well as the bound predicts --- the true $R(D)$ for
spatially and temporally correlated video lies below the Gaussian curve, making
it a conservative benchmark.}

No practical codec achieves $R(D)$;\footnote{The $R(D)$ curve in
Figure~\ref{fig:rate-distortion} is the Gaussian memoryless source bound
(equation~\ref{eq:rate-distortion}), which is an upper bound on the true
$R(D)$ for non-Gaussian sources such as video.  Additionally, $R(D)$ assumes
zero excess-distortion probability, whereas practical codecs permit a
frame-level excess-distortion probability on the order of $10^{-6}$ to
$10^{-8}$.  Both factors allow codec operating points measured on real content
to approach or marginally cross the Gaussian bound.} real codecs operate at
some rate $R_{\text{codec}} > R(D)$. We define the coding efficiency gap as:
\begin{equation}
  \Delta = \frac{R_{\text{codec}} - R(D)}{R(D)}
  \label{eq:coding-gap}
\end{equation}

The codecs in the WebRTC stack can be ordered by their empirical
proximity to $R(D)$ for typical video content at conversational
resolutions (720p, 30~fps). Using published Bj\o ntegaard Delta (BD)
rate measurements:
\begin{equation}
  R_{\text{H.264}} > R_{\text{VP9}} > R_{\text{AV1}}
    > R(D)
  \label{eq:codec-ordering}
\end{equation}

Specifically, the empirical rate savings at equivalent PSNR are
approximately:
\begin{align}
  R_{\text{VP9}}  &\approx 0.65 \cdot R_{\text{H.264}}
    \quad (\sim\!35\% \text{ savings}) \label{eq:vp9-savings} \\
  R_{\text{AV1}}  &\approx 0.70 \cdot R_{\text{VP9}}
    \approx 0.46 \cdot R_{\text{H.264}}
    \quad (\sim\!30\% \text{ over VP9})
  \label{eq:av1-savings}
\end{align}

\begin{figure}[ht]
\centering
\begin{tikzpicture}[>=Stealth]

% --- Axes ---
\draw[->, thick] (0, 0) -- (10.5, 0);
\draw[->, thick] (0, 0) -- (0, 6.5);

% --- Axis titles ---
\node[font=\small] at (5, -0.85) {Distortion $D$};
\node[font=\small, rotate=90] at (-0.9, 3.25) {Rate [Mbps]};

% --- Y-axis rate labels ---
\foreach \y/\lab in {1.0/0.5, 2.0/1.0, 3.0/1.5, 4.0/2.0, 5.0/2.5, 6.0/3.0} {
  \draw (-0.1, \y) -- (0.1, \y);
  \node[font=\tiny, anchor=east] at (-0.15, \y) {\lab};
}

% --- R(D) Shannon bound curve ---
\draw[very thick, green!60!black]
  (1.0, 5.8) .. controls (2.0, 3.5) and (3.5, 2.0) ..
  (5.0, 1.3) .. controls (6.5, 0.8) and (8.0, 0.55) ..
  (9.5, 0.4);
\node[font=\scriptsize, green!60!black, anchor=north west] at (8.5, 0.55)
  {$R(D)$ bound};

% --- Reference distortion line (720p/30fps operating point) ---
\draw[thin, gray, dashed] (4.0, 0) -- (4.0, 6.2);
\node[font=\tiny, gray, anchor=south] at (4.0, 6.2) {720p/30fps};

% --- R(D) value at reference distortion ---
\fill[green!60!black] (4.0, 1.55) circle (2pt);

% --- Codec operating points at reference distortion ---
% H.264: ~2.5 Mbps
\draw[thick, dashed, red!70!black] (0.3, 4.8) -- (9.5, 4.8);
\fill[red!70!black] (4.0, 4.8) circle (3pt);
\node[font=\scriptsize, red!70!black, anchor=west] at (9.6, 4.8) {H.264};

% VP9: ~1.6 Mbps (0.65 * H.264)
\draw[thick, dashed, blue!70] (0.3, 3.3) -- (9.5, 3.3);
\fill[blue!70] (4.0, 3.3) circle (3pt);
\node[font=\scriptsize, blue!70, anchor=west] at (9.6, 3.3) {VP9};

% AV1: ~1.15 Mbps (0.46 * H.264)
\draw[thick, dashed, orange!80!black] (0.3, 2.3) -- (9.5, 2.3);
\fill[orange!80!black] (4.0, 2.3) circle (3pt);
\node[font=\scriptsize, orange!80!black, anchor=west] at (9.6, 2.3) {AV1};

% --- Gap annotations (braces) ---
% H.264 gap
\draw[decorate, decoration={brace, amplitude=4pt, mirror},
  thick, red!50!black]
  (3.6, 1.55) -- (3.6, 4.8)
  node[midway, anchor=east, font=\scriptsize, xshift=-5pt]
  {$\Delta_{\text{H.264}}$};

% AV1 gap
\draw[decorate, decoration={brace, amplitude=4pt},
  thick, orange!60!black]
  (4.4, 1.55) -- (4.4, 2.3)
  node[midway, anchor=west, font=\scriptsize, xshift=5pt]
  {$\Delta_{\text{AV1}}$};

% --- Annotation ---
\node[font=\scriptsize\itshape, text=black!60, align=center,
  anchor=north] at (5.0, -0.5)
  {The Shannon rate-distortion bound $R(D)$ is the theoretical minimum.\\
   AV1 narrows the gap $\Delta$ to ${\sim}0.4$ by 2030 (eq.~\ref{eq:coding-gap}).};

\end{tikzpicture}
\caption{Codec operating rates vs.\ the Shannon rate-distortion bound $R(D)$
  at 720p/30fps. Each codec operates above the theoretical minimum; the gap
  $\Delta$ (eq.~\ref{eq:coding-gap}) shrinks from H.264 through VP9 to AV1,
  approaching but never reaching the bound.}
\label{fig:rate-distortion}
\end{figure}

\subsubsection{Codec Entropy and WebCodecs}

The WebCodecs API (W3C Working Draft) exposes per-frame codec access,
allowing measurement and control of the instantaneous encoding rate.
For a video frame $f_t$ at time $t$, the encoder output has empirical
entropy:
\begin{equation}
  \hat{H}(f_t) = -\sum_{s \in \mathcal{S}}
    p(s \mid f_t) \log_2 p(s \mid f_t)
  \label{eq:frame-entropy}
\end{equation}
where $\mathcal{S}$ is the set of symbols emitted by the entropy coder
(CABAC for H.264/AV1, arithmetic coding for VP9). WebCodecs enables
applications to observe $|E(f_t)|$ (the encoded frame size in bits)
directly, which serves as a practical upper bound on
$\hat{H}(f_t)$:
\begin{equation}
  \hat{H}(f_t) \leq |E(f_t)| \leq R_{\text{codec}} / \text{fps}
  \label{eq:entropy-bound}
\end{equation}

\subsection{Effective Information Rate by Time Horizon}

We now combine transport and source coding measures to estimate the
effective information rate for a representative real-time video
session: a single 720p/30fps video call with one audio channel.

\excursus{Excursus: Why 720p?}{Throughout this appendix, 720p/30fps is used as
a fixed reference for cross-era comparison, not as a prediction that 720p
remains the dominant conferencing resolution. The bandwidth surplus quantified
below is precisely what enables migration to 1080p (feasible by ${\sim}$2028
with AV1 hardware encode at today's 720p bit rates) and 4K (feasible by
${\sim}$2030). The 720p baseline makes the efficiency gains legible: the same
task gets cheaper, freeing capacity for higher quality or more streams.}

Define the \emph{net information rate} as:
\begin{equation}
  I_{\text{net}} = R_{\text{codec}} \cdot \eta_{\text{transport}}
    \cdot (1 - \rho_{\text{overhead}})
    \cdot (1 - p_{\text{loss}})
  \label{eq:inet}
\end{equation}

\begin{table}[ht]
\centering
\small
\renewcommand{\arraystretch}{1.4}
\begin{tabularx}{\textwidth}{lXXX}
\toprule
\textbf{Parameter} & \textbf{March 2026} & \textbf{2027--2028} &
  \textbf{2029--2031} \\
\midrule
Median link capacity &
  120~Mbps &
  $\sim$150~Mbps &
  200--300~Mbps \\
Dominant codec (WebRTC) &
  VP9 / H.264 &
  VP9 / AV1 (HW) &
  AV1 (HW default) \\
Codec rate (720p/30fps) &
  1.5--2.5~Mbps &
  1.0--1.8~Mbps &
  0.7--1.2~Mbps \\
$\eta_{\text{transport}}$ &
  0.92 (QUIC) &
  0.93 &
  0.94 \\
$\rho_{\text{overhead}}$ (SRTP+QUIC) &
  0.06 &
  0.05 &
  0.04 \\
Typical loss $p$ &
  0.01--0.03 &
  0.01--0.02 &
  0.005--0.01 \\
$I_{\text{net}}$ (720p video) &
  1.3--2.2~Mbps &
  0.9--1.6~Mbps &
  0.6--1.1~Mbps \\
$I_{\text{net}}$ as \% of $C_{\text{link}}$ &
  1.1--1.8\% &
  0.6--1.1\% &
  0.2--0.6\% \\
\bottomrule
\end{tabularx}
\caption{Effective information rate for a 720p/30fps WebRTC video stream.
  The declining ratio $I_{\text{net}} / C_{\text{link}}$ reflects both
  improved codec efficiency and growing link capacity, freeing bandwidth
  for higher resolutions, additional streams, or concurrent data
  channels.}
\label{tab:inet}
\end{table}

The declining ratio $I_{\text{net}} / C_{\text{link}}$ in
Table~\ref{tab:inet} is the central quantitative observation: by 2030, a
720p video call consumes less than 0.5\% of available link capacity. The
surplus capacity enables qualitative changes --- 4K video conferencing,
multi-stream layouts, or concurrent data channels via WebTransport ---
that are quantitatively infeasible in 2026.

\begin{figure}[ht]
\centering
\begin{tikzpicture}[>=Stealth]

% --- Axes ---
\draw[->, thick] (0, 0) -- (11, 0);
\draw[->, thick] (0, 0) -- (0, 5.5);

% --- Axis title ---
\node[font=\small, rotate=90] at (-0.9, 2.75) {Mbps};

% --- Y-axis labels ---
\foreach \y/\lab in {0/0, 1/50, 2/100, 3/150, 4/200, 5/250} {
  \draw (-0.1, \y) -- (0.1, \y);
  \node[font=\tiny, anchor=east] at (-0.15, \y) {\lab};
}

% --- Group: March 2026 ---
% C_link = 120 Mbps -> y = 2.4
\fill[blue!12] (0.8, 0) rectangle (2.4, 2.4);
\draw[blue!50, thick] (0.8, 0) rectangle (2.4, 2.4);
% I_net = 1.75 Mbps (midpoint) -> y = 0.035
\fill[orange!60] (0.8, 0) rectangle (2.4, 0.07);
\draw[orange!80!black, thick] (0.8, 0) rectangle (2.4, 0.07);
\node[font=\footnotesize, anchor=south] at (1.6, 2.5) {120};
\node[font=\tiny, orange!80!black] at (1.6, -0.4) {1.5\%};
\node[font=\small, anchor=north] at (1.6, -0.7) {2026};

% --- Group: 2027--2028 ---
% C_link = 150 Mbps -> y = 3.0
\fill[blue!12] (4.0, 0) rectangle (5.6, 3.0);
\draw[blue!50, thick] (4.0, 0) rectangle (5.6, 3.0);
% I_net = 1.25 Mbps (midpoint) -> y = 0.025
\fill[orange!60] (4.0, 0) rectangle (5.6, 0.05);
\draw[orange!80!black, thick] (4.0, 0) rectangle (5.6, 0.05);
\node[font=\footnotesize, anchor=south] at (4.8, 3.1) {150};
\node[font=\tiny, orange!80!black] at (4.8, -0.4) {0.8\%};
\node[font=\small, anchor=north] at (4.8, -0.7) {2027--28};

% --- Group: 2029--2031 ---
% C_link = 250 Mbps -> y = 5.0
\fill[blue!12] (7.2, 0) rectangle (8.8, 5.0);
\draw[blue!50, thick] (7.2, 0) rectangle (8.8, 5.0);
% I_net = 0.85 Mbps (midpoint) -> y = 0.017
\fill[orange!60] (7.2, 0) rectangle (8.8, 0.034);
\draw[orange!80!black, thick] (7.2, 0) rectangle (8.8, 0.034);
\node[font=\footnotesize, anchor=south] at (8.0, 5.1) {250};
\node[font=\tiny, orange!80!black] at (8.0, -0.4) {0.3\%};
\node[font=\small, anchor=north] at (8.0, -0.7) {2029--31};

% --- Legend ---
\node[anchor=west, font=\footnotesize] at (0.5, 5.8) {%
  \tikz[baseline=-0.5ex]{\fill[blue!12] (0,0) rectangle (0.3,0.3);
    \draw[blue!50, thick] (0,0) rectangle (0.3,0.3);}
  ~$C_{\text{link}}$ (channel capacity) \qquad
  \tikz[baseline=-0.5ex]{\fill[orange!60] (0,0) rectangle (0.3,0.3);
    \draw[orange!80!black, thick] (0,0) rectangle (0.3,0.3);}
  ~$I_{\text{net}}$ (720p stream)};

% --- Annotation ---
\node[font=\scriptsize\itshape, text=black!60, align=center,
  anchor=north] at (5.5, -1.2)
  {A single 720p call uses a vanishing fraction of Shannon-derived\\
   channel capacity, freeing bandwidth for 4K, multi-stream, and data channels.};

\end{tikzpicture}
\caption{Channel capacity $C_{\text{link}}$ vs.\ effective information rate
  $I_{\text{net}}$ for a single 720p/30fps stream across time horizons. The
  orange bar (stream rate) is barely visible against the blue bar (capacity),
  illustrating the growing surplus as codecs improve and bandwidth increases.}
\label{fig:inet-capacity}
\end{figure}

\subsection{Security Entropy and Key Space}

The DTLS-SRTP security architecture (RFCs~8827, 3711) protects all
WebRTC media. Information-theoretic security requires that the key
entropy exceed the information content of the protected material.

For a DTLS~1.3 handshake with ECDHE on P-256, the key agreement
provides:
\begin{equation}
  H_{\text{key}} = 128 \text{~bits of security}
  \label{eq:key-entropy}
\end{equation}

The SRTP master key is 128 bits (AES-128-GCM), giving an exhaustive
search cost of $2^{128}$ operations. The total information content of a
one-hour 720p video call at 2~Mbps is:
\begin{equation}
  I_{\text{call}} = 2 \times 10^6 \times 3600
    = 7.2 \times 10^9 \text{~bits}
    \approx 2^{32.7} \text{~bits}
  \label{eq:call-content}
\end{equation}

The ratio of key entropy to protected content:
\begin{equation}
  \frac{H_{\text{key}}}{I_{\text{call}}}
    = \frac{128}{32.7} \approx 3.91
  \label{eq:security-margin}
\end{equation}

This ratio remains comfortably above~1 even for much longer sessions.
For a continuous 24-hour 4K AV1 stream at 15~Mbps:
\begin{equation}
  I_{\text{stream}} = 15 \times 10^6 \times 86400
    = 1.296 \times 10^{12} \text{~bits}
    \approx 2^{40.2} \text{~bits}
\end{equation}
yielding $H_{\text{key}} / I_{\text{stream}} \approx 3.18$, still
well above the threshold.

\subsubsection{End-to-End Encryption via Encoded Transforms}

The RTCRtpScriptTransform API (Encoded Transforms) enables
application-layer E2EE within the WebRTC pipeline. From an
information-theoretic perspective, this introduces a second encryption
layer with its own key, creating a \emph{nested channel}:
\begin{equation}
  X \xrightarrow{\;K_{\text{E2EE}}\;}
  E_1(X) \xrightarrow{\;K_{\text{SRTP}}\;}
  E_2(E_1(X))
  \label{eq:nested}
\end{equation}

The mutual information between the ciphertext and plaintext, given
neither key, satisfies:
\begin{equation}
  I\!\left(X;\, E_2(E_1(X))\right) \leq
    \min\!\big(H(K_{\text{E2EE}})^{-1},\;
    H(K_{\text{SRTP}})^{-1}\big)
    \cdot I_{\text{call}}
  \label{eq:mutual-e2ee}
\end{equation}

In practice, with $H(K_{\text{E2EE}}) = H(K_{\text{SRTP}}) = 128$
bits, this quantity is negligible ($\ll 1$ bit for any practical
session length), confirming that the nested encryption provides
information-theoretic confidentiality well beyond computational
security guarantees.

Today (March~2026), Encoded Transforms remain Chrome-only. By
2027--2028 (Section~8.3), cross-browser standardization will make
this nested channel universally available. By 2029--2031, SFrame
for group calls will extend this to $n$-party sessions, where the
mutual information between any non-member and the group plaintext
remains negligible:
\begin{equation}
  I\!\left(X;\, E(X) \;\middle|\; K_i \notin
    \{K_1, \ldots, K_n\}\right) \approx 0
  \label{eq:sframe}
\end{equation}

\subsection{WebTransport: Unreliable Datagrams and Channel Capacity}

WebTransport (RFC~9297) introduces unreliable datagrams to browser-server
communication. For latency-sensitive applications (game state, sensor
telemetry), the relevant measure is not throughput but
\emph{freshness-weighted information rate}. Define the value of a
datagram as exponentially decaying with delivery latency $\delta$:
\begin{equation}
  V(\delta) = e^{-\lambda \delta}
  \label{eq:freshness}
\end{equation}
where $\lambda > 0$ is an application-specific decay rate. The
effective information rate under unreliable delivery is:
\begin{equation}
  I_{\text{WT}} = R_{\text{send}} \cdot (1 - p_{\text{loss}})
    \cdot \mathbb{E}\!\left[e^{-\lambda \delta}\right]
  \label{eq:iwt}
\end{equation}

Under reliable delivery (TCP/WebSocket), a retransmission delays
subsequent data, so:
\begin{equation}
  I_{\text{WS}} = R_{\text{send}}
    \cdot \mathbb{E}\!\left[e^{-\lambda(\delta + p \cdot 2\tau)}\right]
  \label{eq:iws}
\end{equation}

The ratio $I_{\text{WT}} / I_{\text{WS}}$ exceeds 1 when:
\begin{equation}
  (1 - p) \cdot e^{\lambda \cdot p \cdot 2\tau} > 1
  \quad \Longleftrightarrow \quad
  \lambda > \frac{-\ln(1 - p)}{2 p \tau}
  \label{eq:wt-advantage}
\end{equation}

For $p = 0.02$ and $\tau = 25$~ms, this gives $\lambda > 20.2$~s$^{-1}$
(half-life $< 34$~ms).

\excursus{Excursus: Why $p = 0.02$ and $\tau = 25$\,ms?}{These values represent
median fixed-broadband conditions in 2026 (Table~\ref{tab:inet}). On mobile or
congested links ($p \approx 0.05$, $\tau \approx 50$\,ms), the threshold
$\lambda^*$ drops to ${\sim}10$\,s$^{-1}$, broadening WebTransport's advantage
to applications with data half-lives up to ${\sim}70$\,ms. The fixed-broadband
reference is thus conservative: mobile scenarios favor unreliable delivery even
more strongly.}

Any application where data older than
$\sim$30~ms is useless benefits from WebTransport's unreliable
mode over WebSocket --- precisely the regime of interactive gaming, live
cursor tracking, and real-time sensor feeds.

As of March~2026, Safari does not support WebTransport
(Section~5), limiting coverage to $\sim$82\%. By 2027--2028
(Section~8.2), universal support will allow applications to rely on
(\ref{eq:iwt}) without fallback paths.

\begin{figure}[ht]
\centering
\begin{tikzpicture}[>=Stealth]

% --- Axes ---
\draw[->, thick] (0, 0) -- (10.5, 0);
\draw[->, thick] (0, 0) -- (0, 5.5);

% --- Axis titles (positioned to clear tick labels) ---
\node[font=\small] at (5, -0.75)
  {Freshness decay rate $\lambda$ [s$^{-1}$]};
\node[font=\small, rotate=90] at (-1.1, 2.75)
  {$I_{\text{WT}} / I_{\text{WS}}$};

% --- X-axis labels ---
\foreach \x/\lab in {0/0, 1.67/10, 3.33/20, 5.0/30, 6.67/40, 8.33/50, 10/60} {
  \draw (\x, -0.1) -- (\x, 0.1);
  \node[font=\tiny, anchor=north] at (\x, -0.15) {\lab};
}

% --- Y-axis labels (0.96 to 1.06, each 0.02 = 1 unit) ---
\foreach \y/\lab in {0/0.96, 1/0.98, 2/1.00, 3/1.02, 4/1.04, 5/1.06} {
  \draw (-0.1, \y) -- (0.1, \y);
  \node[font=\tiny, anchor=east] at (-0.15, \y) {\lab};
}

% --- Ratio = 1.0 line (y=2 maps to ratio 1.0) ---
\draw[thin, gray, dashed] (0, 2) -- (10, 2);

% --- Shade WebTransport advantage region ---
% lambda* = 20.2 -> x = 3.37
\fill[green!10] (3.37, 2) rectangle (10, 5.3);

% --- I_WT/I_WS curve: 0.98*exp(lambda*0.001) ---
% Y range: 0.96-1.06 (each 0.02 = 1 unit). lambda=0: y=1, lambda*=20.2: y=2, lambda=60: y=4.05
\draw[very thick, orange!80!black]
  (0, 1.0) .. controls (1.5, 1.3) and (2.5, 1.7) ..
  (3.37, 2.0) .. controls (4.5, 2.4) and (6, 2.8) ..
  (8, 3.4) .. controls (9, 3.7) and (9.5, 3.85) ..
  (10, 4.05);

% --- Threshold line ---
\draw[thick, dashed, red!60!black] (3.37, 0) -- (3.37, 5.2);
\node[font=\scriptsize, red!60!black, anchor=south, rotate=90] at (3.17, 3.5)
  {$\lambda^* = 20.2$ s$^{-1}$};

% --- Application regime labels ---
\node[font=\scriptsize\itshape, text=blue!60, anchor=north] at (1.5, 5.3)
  {video conf};
\node[font=\scriptsize\itshape, text=green!50!black, anchor=north] at (7.0, 5.3)
  {gaming / cursor / sensors};

% --- Mark advantage/disadvantage ---
\node[font=\tiny, text=red!50!black, anchor=north east] at (3.2, 1.8)
  {WebSocket wins};
\node[font=\tiny, text=green!50!black, anchor=north west] at (3.5, 4.5)
  {WebTransport wins};

% --- Annotation ---
\node[font=\scriptsize\itshape, text=black!60, align=center,
  anchor=north] at (5.0, -1.0)
  {For $p = 0.02$, $\tau = 25$~ms. Applications where data older than\\
   ${\sim}$30~ms is stale gain from unreliable delivery (eq.~\ref{eq:wt-advantage}).};

\end{tikzpicture}
\caption{Freshness-weighted information rate ratio $I_{\text{WT}} / I_{\text{WS}}$
  as a function of decay rate $\lambda$. Above the threshold $\lambda^*$,
  WebTransport's unreliable datagrams deliver more useful information per
  unit time than WebSocket's reliable delivery, because retransmission
  wastes capacity on stale data.}
\label{fig:freshness}
\end{figure}

\subsection{Projections Summary and Practical Implications}

\begin{table}[ht]
\centering
\small
\renewcommand{\arraystretch}{1.4}
\begin{tabularx}{\textwidth}{lXXX}
\toprule
\textbf{Measure} & \textbf{March 2026} & \textbf{2027--2028} &
  \textbf{2029--2031} \\
\midrule
Channel capacity $C_{\text{link}}$ &
  120~Mbps &
  150~Mbps &
  250~Mbps \\
HOL blocking factor $(1\!-\!p)^{n-2}$ at $n\!=\!8$ &
  0.94 ($p\!=\!0.01$) &
  0.94 &
  0.97 ($p\!=\!0.005$) \\
Source coding gap $\Delta$ &
  VP9: $\sim$0.8 &
  AV1 HW: $\sim$0.5 &
  AV1 opt: $\sim$0.4 \\
$I_{\text{net}}$ per 720p stream &
  1.3--2.2~Mbps &
  0.9--1.6~Mbps &
  0.6--1.1~Mbps \\
Max 720p streams in $C_{\text{link}}$ &
  $\sim$55--90 &
  $\sim$95--165 &
  $\sim$230--415 \\
E2EE availability &
  Chrome only &
  Cross-browser &
  SFrame $n$-party \\
WebTransport coverage &
  82\% &
  $\sim$100\% &
  100\% \\
Freshness threshold $\lambda^*$ &
  20.2~s$^{-1}$ &
  $\sim$18~s$^{-1}$ &
  $\sim$15~s$^{-1}$ \\
\bottomrule
\end{tabularx}
\caption{Information-theoretic measures across time horizons. $\Delta$
  is the coding efficiency gap (\ref{eq:coding-gap}), $I_{\text{net}}$
  is the net information rate (\ref{eq:inet}), and $\lambda^*$ is the
  freshness decay threshold above which WebTransport dominates WebSocket
  (\ref{eq:wt-advantage}).}
\label{tab:projections}
\end{table}

The trajectory is clear: the combination of improving source coding
(AV1 hardware), expanding channel capacity (broadband growth), and
maturing transport protocols (QUIC universal, WebTransport universal)
drives the \emph{information cost per unit of communicated experience}
steadily downward. The remaining question is not whether the
Shannon capacity suffices --- it already does, by orders of magnitude
--- but how efficiently the protocol stack converts that capacity into
delivered, fresh, secure information at the application layer. By
2030, the gap between theoretical capacity and realized throughput will
be the narrowest it has ever been for browser-based real-time
communication.

\subsubsection{What the Theory Means for Applications}

The equations and figures in this appendix are not purely academic. Each
theoretical result maps to a concrete design decision that application
developers face today and over the next five years.

\textbf{Codec selection determines bandwidth cost, not quality.}
The rate-distortion analysis (Section~A.2, Figure~\ref{fig:rate-distortion})
shows that AV1 operates at roughly 46\% of H.264's bit rate for equivalent
perceptual quality at 720p. In practice, this means a WebRTC application
switching from H.264 to AV1 can either halve its bandwidth consumption at
the same quality or substantially increase resolution (e.g., 720p to 1080p)
within the same bandwidth envelope. By 2030, with AV1 hardware encoding
ubiquitous and the coding gap $\Delta$ approaching 0.4, further codec
improvements yield diminishing returns --- the stack will be operating close
enough to the Shannon rate-distortion bound that the next generation of
efficiency gains will come from transport and architecture, not codecs alone.

\textbf{QUIC's advantage is situational but decisive on bad networks.}
The head-of-line blocking analysis (Section~A.1.1,
Figure~\ref{fig:hol-blocking}) shows that QUIC's per-stream throughput
advantage over TCP is modest on good networks (6\% at 1\% loss) but
substantial on degraded ones (34\% at 5\% loss). The practical implication
is that applications targeting mobile users, emerging markets, or congested
enterprise networks --- precisely the environments where real-time
communication is hardest --- benefit disproportionately from QUIC-based
transports. For a video conference with 8 participants on a cellular
connection experiencing 3--5\% packet loss, the difference between QUIC and
TCP is the difference between usable and degraded video.

\textbf{Connection latency is a solved problem for returning users.}
The handshake analysis (Section~A.1.2) shows that QUIC 0-RTT eliminates
connection establishment overhead entirely for resumed connections. For a
WebTransport application that a user opens daily --- a collaboration tool,
a live dashboard, a recurring meeting client --- the time to first useful
byte drops to zero. The 1.125~MB of wasted capacity during a TCP+TLS
handshake on a 120~Mbps link may seem small, but for latency-sensitive
applications where the first 100~ms of perceived responsiveness determines
user experience, the difference is measurable.

\textbf{Bandwidth surplus enables architectural ambition.}
The effective information rate analysis (Section~A.3,
Figure~\ref{fig:inet-capacity}) reveals the most consequential practical
finding: a single 720p video stream consumes a vanishing fraction of
available link capacity (under 0.5\% by 2030). This surplus is not merely
``headroom'' --- it is the resource that makes new application architectures
feasible. A 250~Mbps link can theoretically carry 230--415 simultaneous 720p
streams. In practice, this means applications can afford to run multiple
concurrent media flows (video, screen share, spatial audio), data channels
(collaborative editing, file transfer), and telemetry streams without
contending for bandwidth. The constraint shifts from ``can the network carry
this?''\ to ``can the application manage this complexity?''

\textbf{End-to-end encryption adds negligible overhead.}
The security analysis (Section~A.4) confirms that the nested DTLS-SRTP plus
E2EE architecture imposes no meaningful information-theoretic cost. The key
entropy ($2^{128}$) exceeds the information content of any practical session
by a factor of at least 3. For application developers, this means E2EE can
be treated as a default rather than a premium feature --- there is no
bandwidth, latency, or capacity penalty for encrypting at the application
layer on top of the transport-layer encryption that WebRTC already mandates.
The only constraint is API availability, which moves from Chrome-only (2026)
to universal (2028).

\textbf{WebTransport unlocks real-time data that WebSocket cannot serve.}
The freshness analysis (Section~A.5, Figure~\ref{fig:freshness}) quantifies
when unreliable delivery outperforms reliable delivery: any data with a
useful half-life below approximately 34~ms benefits from WebTransport's
unreliable datagrams. This covers a wide class of interactive applications
--- game state synchronization, live cursor and pointer sharing, real-time
sensor feeds, collaborative whiteboard strokes --- where retransmitting a
stale update is worse than dropping it. Before WebTransport, these
applications had to choose between WebSocket (reliable but stale under loss)
and WebRTC DataChannel (unreliable mode available but requiring the full
peer-to-peer ICE/DTLS machinery). WebTransport provides the unreliable
mode with a simple client-server model over HTTP/3, removing the
architectural overhead.

\textbf{The overall picture: the stack converges on efficiency.}
Taken together, the information-theoretic measures trace a consistent
trajectory. The protocol stack is converging toward the theoretical limits
established by Shannon: codecs approach the rate-distortion bound, transports
eliminate unnecessary blocking and handshake overhead, and encryption adds
negligible cost. By 2030, the practical distance between what the standards
deliver and what information theory permits will be the smallest it has ever
been for browser-based real-time communication. The consequence for
application developers is that the infrastructure --- standards, hardware,
and networks --- will no longer be the binding constraint. The remaining
challenge is purely architectural: how to compose these efficient, universal
primitives into applications that fully exploit the capacity and low latency
they provide.

%
% Requires in the preamble: amsmath, amssymb, tabularx, booktabs

\section{Information-Theoretic View}
\label{app:info-theory}

This appendix applies information-theoretic measures to the real-time
web communication stack described in this document. We quantify channel
capacity, source coding efficiency, and effective information rates
across three time horizons: today (March~2026), the near term
(2027--2028), and the medium term (2029--2031). The analysis draws on
the standards, codec parameters, and infrastructure projections
developed in the main text.

\subsection{Channel Capacity of the Transport Layer}

The Shannon-Hartley theorem gives the theoretical maximum information
rate for a channel with additive white Gaussian noise:
\begin{equation}
  C = B \log_2\!\left(1 + \frac{S}{N}\right) \quad \text{[bits/s]}
  \label{eq:shannon}
\end{equation}
where $B$ is the channel bandwidth in Hz and $S/N$ is the
signal-to-noise ratio.

While the physical-layer capacity of a broadband link is governed by
(\ref{eq:shannon}), the \emph{effective} capacity available to the
application is reduced by protocol overhead, congestion control, and
transport inefficiencies. We define the effective application-layer
throughput as:
\begin{equation}
  R_{\text{eff}} = C_{\text{link}} \cdot \eta_{\text{transport}}
    \cdot (1 - \rho_{\text{overhead}})
  \label{eq:reff}
\end{equation}
where $C_{\text{link}}$ is the link-layer capacity,
$\eta_{\text{transport}} \in (0,1]$ is the transport efficiency factor
(capturing congestion control, retransmission, and multiplexing
behavior), and $\rho_{\text{overhead}} \in [0,1)$ is the fractional
protocol overhead (headers, encryption, framing).

\subsubsection{Head-of-Line Blocking and Multiplexed Capacity}

Consider $n$ independent streams multiplexed over a single connection.
Under TCP (HTTP/2), a single lost packet blocks all $n$ streams. The
expected throughput per stream under loss rate $p$ is bounded by:
\begin{equation}
  R_{\text{TCP}}^{(i)} \leq \frac{R_{\text{eff}}}{n}
    \cdot (1 - p)^{n-1}
  \label{eq:hol-tcp}
\end{equation}
because a loss in \emph{any} of the $n$ streams stalls all others. Under
QUIC (RFC~9000), streams are independent at the transport layer, so:
\begin{equation}
  R_{\text{QUIC}}^{(i)} \leq \frac{R_{\text{eff}}}{n}
    \cdot (1 - p)
  \label{eq:hol-quic}
\end{equation}
The ratio of effective per-stream throughput is:
\begin{equation}
  \frac{R_{\text{QUIC}}^{(i)}}{R_{\text{TCP}}^{(i)}}
    = \frac{1}{(1-p)^{n-2}}
  \label{eq:hol-ratio}
\end{equation}
\excursus{Excursus: Why $n = 8$?}{In a typical multi-participant video
conference, each peer contributes multiple concurrent streams: a camera video
track, an audio track, and often a screen-share or data channel. A four-person
call with two media tracks each produces $n = 8$ multiplexed streams over a
single connection --- the reference value used throughout this section. Larger
meetings (e.g.\ 10 participants with simulcast layers) push $n$ higher,
amplifying the HOL blocking effect further.}

For $n = 8$ streams (a typical multi-track video conference) at
$p = 0.01$ (1\% loss), this gives a QUIC advantage factor of
$1/(0.99)^6 \approx 1.062$, a modest 6.2\% improvement. At $p = 0.05$
(5\% loss, common on mobile or congested links), the factor is
$1/(0.95)^6 \approx 1.34$, a 34\% throughput advantage per stream.

\begin{figure}[ht]
\centering
\begin{tikzpicture}[>=Stealth]

% --- Shaded region (drawn FIRST so curves appear on top) ---
\fill[orange!8] (2, 0) rectangle (10, 5.3);
\node[font=\tiny\itshape, text=black!40, anchor=north east] at (9.8, 5.2)
  {mobile / congested};
\node[font=\tiny\itshape, text=black!40, anchor=north west] at (0.2, 5.2)
  {fixed broadband};

% --- Grid lines ---
\foreach \y in {1, 2, 3, 4, 5} {
  \draw[gray!20] (0, \y) -- (10, \y);
}

% --- Axes ---
\draw[->, thick] (0, 0) -- (10.5, 0);
\draw[->, thick] (0, 0) -- (0, 5.5);

% --- Axis titles (positioned to clear tick labels) ---
\node[font=\small] at (5, -0.85) {Packet loss rate $p$};
\node[font=\small, rotate=90] at (-0.9, 2.75) {QUIC advantage factor};

% --- X-axis labels ---
\foreach \x/\lab in {0/0, 2/0.02, 4/0.04, 6/0.06, 8/0.08, 10/0.10} {
  \draw (\x, -0.1) -- (\x, 0.1);
  \node[font=\tiny, anchor=north] at (\x, -0.15) {\lab};
}

% --- Y-axis labels ---
\foreach \y/\lab in {0/1.0, 1/1.1, 2/1.2, 3/1.3, 4/1.4, 5/1.5} {
  \draw (-0.1, \y) -- (0.1, \y);
  \node[font=\tiny, anchor=east] at (-0.15, \y) {\lab};
}

% --- n=4 curve: 1/(1-p)^2 ---
\draw[very thick, blue!70]
  (0, 0) .. controls (2.5, 0.5) and (5, 1.1) ..
  (7.5, 2.0) .. controls (8.5, 2.4) and (9.5, 2.9) ..
  (10, 3.1);
\node[font=\scriptsize, blue!70, anchor=south west] at (9.0, 3.0)
  {$n = 4$};

% --- n=8 curve: 1/(1-p)^6 ---
\draw[very thick, orange!80!black]
  (0, 0) .. controls (1.5, 0.5) and (3, 1.7) ..
  (4, 2.5) .. controls (5, 3.5) and (5.5, 4.0) ..
  (6, 4.7);
\node[font=\scriptsize, orange!80!black, anchor=south] at (5.5, 4.8)
  {$n = 8$};

% --- Reference points ---
\fill[orange!80!black] (1.0, 0.62) circle (2.5pt);
\node[font=\tiny, orange!80!black, anchor=south west] at (1.1, 0.62)
  {6.2\%};

\fill[orange!80!black] (5.0, 3.7) circle (2.5pt);
\node[font=\tiny, orange!80!black, anchor=west] at (5.1, 3.7)
  {34\%};

\fill[blue!70] (1.0, 0.20) circle (2.5pt);

% --- Annotation ---
\node[font=\scriptsize\itshape, text=black!60, align=center,
  anchor=north] at (5.0, -1.0)
  {QUIC recovers per-stream throughput closer to $C/n$ (Shannon fair share)\\
   by eliminating head-of-line blocking (eq.~\ref{eq:hol-ratio}).};

\end{tikzpicture}
\caption{QUIC advantage factor $1/(1-p)^{n-2}$ over TCP for multiplexed
  streams. At 5\% loss with $n=8$ streams, QUIC delivers 34\% more
  per-stream throughput by maintaining independence at the transport layer.}
\label{fig:hol-blocking}
\end{figure}

\subsubsection{Connection Establishment Latency}

Let $\tau$ denote the one-way latency (half the round-trip time). The
information-carrying capacity during connection establishment is zero;
the handshake represents pure overhead. The time to first useful byte
is:
\begin{align}
  T_{\text{TCP+TLS\,1.3}} &= 2\tau_{\text{TCP}} + \tau_{\text{TLS}}
    = 3\tau \quad \text{(1-RTT TLS~1.3)} \label{eq:tcp-setup} \\
  T_{\text{QUIC}} &= \tau \quad \text{(1-RTT)}
    \label{eq:quic-setup} \\
  T_{\text{QUIC,\,0-RTT}} &= 0 \quad \text{(resumption)}
    \label{eq:quic-0rtt}
\end{align}
The QUIC 0-RTT case (\ref{eq:quic-0rtt}) eliminates the handshake
entirely for resumed connections. The information lost during
handshake dead time, relative to steady-state capacity, is:
\begin{equation}
  I_{\text{lost}} = C_{\text{link}} \cdot T_{\text{setup}}
  \label{eq:ilost}
\end{equation}
For a 120~Mbps link with $\tau = 25$~ms (typical fixed broadband in
2026), the TCP+TLS~1.3 handshake wastes $120 \times 10^6 \times 0.075
= 9 \times 10^6$ bits (1.125~MB), while QUIC~1-RTT wastes only
$3 \times 10^6$ bits (375~KB).

\subsection{Source Coding Efficiency}

The rate-distortion function $R(D)$ defines the minimum bit rate
required to represent a source at distortion level $D$. For a
Gaussian source with variance $\sigma^2$ under mean squared error
distortion:
\begin{equation}
  R(D) = \frac{1}{2} \log_2 \frac{\sigma^2}{D}
    \quad \text{[bits/sample]}
  \label{eq:rate-distortion}
\end{equation}

\excursus{Excursus: Why Gaussian source?}{The Gaussian memoryless model is used
because it yields the highest (worst-case) $R(D)$ among all sources of equal
variance. Any codec that approaches the Gaussian bound on real video is
therefore doing at least as well as the bound predicts --- the true $R(D)$ for
spatially and temporally correlated video lies below the Gaussian curve, making
it a conservative benchmark.}

No practical codec achieves $R(D)$;\footnote{The $R(D)$ curve in
Figure~\ref{fig:rate-distortion} is the Gaussian memoryless source bound
(equation~\ref{eq:rate-distortion}), which is an upper bound on the true
$R(D)$ for non-Gaussian sources such as video.  Additionally, $R(D)$ assumes
zero excess-distortion probability, whereas practical codecs permit a
frame-level excess-distortion probability on the order of $10^{-6}$ to
$10^{-8}$.  Both factors allow codec operating points measured on real content
to approach or marginally cross the Gaussian bound.} real codecs operate at
some rate $R_{\text{codec}} > R(D)$. We define the coding efficiency gap as:
\begin{equation}
  \Delta = \frac{R_{\text{codec}} - R(D)}{R(D)}
  \label{eq:coding-gap}
\end{equation}

The codecs in the WebRTC stack can be ordered by their empirical
proximity to $R(D)$ for typical video content at conversational
resolutions (720p, 30~fps). Using published Bj\o ntegaard Delta (BD)
rate measurements:
\begin{equation}
  R_{\text{H.264}} > R_{\text{VP9}} > R_{\text{AV1}}
    > R(D)
  \label{eq:codec-ordering}
\end{equation}

Specifically, the empirical rate savings at equivalent PSNR are
approximately:
\begin{align}
  R_{\text{VP9}}  &\approx 0.65 \cdot R_{\text{H.264}}
    \quad (\sim\!35\% \text{ savings}) \label{eq:vp9-savings} \\
  R_{\text{AV1}}  &\approx 0.70 \cdot R_{\text{VP9}}
    \approx 0.46 \cdot R_{\text{H.264}}
    \quad (\sim\!30\% \text{ over VP9})
  \label{eq:av1-savings}
\end{align}

\begin{figure}[ht]
\centering
\begin{tikzpicture}[>=Stealth]

% --- Axes ---
\draw[->, thick] (0, 0) -- (10.5, 0);
\draw[->, thick] (0, 0) -- (0, 6.5);

% --- Axis titles ---
\node[font=\small] at (5, -0.85) {Distortion $D$};
\node[font=\small, rotate=90] at (-0.9, 3.25) {Rate [Mbps]};

% --- Y-axis rate labels ---
\foreach \y/\lab in {1.0/0.5, 2.0/1.0, 3.0/1.5, 4.0/2.0, 5.0/2.5, 6.0/3.0} {
  \draw (-0.1, \y) -- (0.1, \y);
  \node[font=\tiny, anchor=east] at (-0.15, \y) {\lab};
}

% --- R(D) Shannon bound curve ---
\draw[very thick, green!60!black]
  (1.0, 5.8) .. controls (2.0, 3.5) and (3.5, 2.0) ..
  (5.0, 1.3) .. controls (6.5, 0.8) and (8.0, 0.55) ..
  (9.5, 0.4);
\node[font=\scriptsize, green!60!black, anchor=north west] at (8.5, 0.55)
  {$R(D)$ bound};

% --- Reference distortion line (720p/30fps operating point) ---
\draw[thin, gray, dashed] (4.0, 0) -- (4.0, 6.2);
\node[font=\tiny, gray, anchor=south] at (4.0, 6.2) {720p/30fps};

% --- R(D) value at reference distortion ---
\fill[green!60!black] (4.0, 1.55) circle (2pt);

% --- Codec operating points at reference distortion ---
% H.264: ~2.5 Mbps
\draw[thick, dashed, red!70!black] (0.3, 4.8) -- (9.5, 4.8);
\fill[red!70!black] (4.0, 4.8) circle (3pt);
\node[font=\scriptsize, red!70!black, anchor=west] at (9.6, 4.8) {H.264};

% VP9: ~1.6 Mbps (0.65 * H.264)
\draw[thick, dashed, blue!70] (0.3, 3.3) -- (9.5, 3.3);
\fill[blue!70] (4.0, 3.3) circle (3pt);
\node[font=\scriptsize, blue!70, anchor=west] at (9.6, 3.3) {VP9};

% AV1: ~1.15 Mbps (0.46 * H.264)
\draw[thick, dashed, orange!80!black] (0.3, 2.3) -- (9.5, 2.3);
\fill[orange!80!black] (4.0, 2.3) circle (3pt);
\node[font=\scriptsize, orange!80!black, anchor=west] at (9.6, 2.3) {AV1};

% --- Gap annotations (braces) ---
% H.264 gap
\draw[decorate, decoration={brace, amplitude=4pt, mirror},
  thick, red!50!black]
  (3.6, 1.55) -- (3.6, 4.8)
  node[midway, anchor=east, font=\scriptsize, xshift=-5pt]
  {$\Delta_{\text{H.264}}$};

% AV1 gap
\draw[decorate, decoration={brace, amplitude=4pt},
  thick, orange!60!black]
  (4.4, 1.55) -- (4.4, 2.3)
  node[midway, anchor=west, font=\scriptsize, xshift=5pt]
  {$\Delta_{\text{AV1}}$};

% --- Annotation ---
\node[font=\scriptsize\itshape, text=black!60, align=center,
  anchor=north] at (5.0, -0.5)
  {The Shannon rate-distortion bound $R(D)$ is the theoretical minimum.\\
   AV1 narrows the gap $\Delta$ to ${\sim}0.4$ by 2030 (eq.~\ref{eq:coding-gap}).};

\end{tikzpicture}
\caption{Codec operating rates vs.\ the Shannon rate-distortion bound $R(D)$
  at 720p/30fps. Each codec operates above the theoretical minimum; the gap
  $\Delta$ (eq.~\ref{eq:coding-gap}) shrinks from H.264 through VP9 to AV1,
  approaching but never reaching the bound.}
\label{fig:rate-distortion}
\end{figure}

\subsubsection{Codec Entropy and WebCodecs}

The WebCodecs API (W3C Working Draft) exposes per-frame codec access,
allowing measurement and control of the instantaneous encoding rate.
For a video frame $f_t$ at time $t$, the encoder output has empirical
entropy:
\begin{equation}
  \hat{H}(f_t) = -\sum_{s \in \mathcal{S}}
    p(s \mid f_t) \log_2 p(s \mid f_t)
  \label{eq:frame-entropy}
\end{equation}
where $\mathcal{S}$ is the set of symbols emitted by the entropy coder
(CABAC for H.264/AV1, arithmetic coding for VP9). WebCodecs enables
applications to observe $|E(f_t)|$ (the encoded frame size in bits)
directly, which serves as a practical upper bound on
$\hat{H}(f_t)$:
\begin{equation}
  \hat{H}(f_t) \leq |E(f_t)| \leq R_{\text{codec}} / \text{fps}
  \label{eq:entropy-bound}
\end{equation}

\subsection{Effective Information Rate by Time Horizon}

We now combine transport and source coding measures to estimate the
effective information rate for a representative real-time video
session: a single 720p/30fps video call with one audio channel.

\excursus{Excursus: Why 720p?}{Throughout this appendix, 720p/30fps is used as
a fixed reference for cross-era comparison, not as a prediction that 720p
remains the dominant conferencing resolution. The bandwidth surplus quantified
below is precisely what enables migration to 1080p (feasible by ${\sim}$2028
with AV1 hardware encode at today's 720p bit rates) and 4K (feasible by
${\sim}$2030). The 720p baseline makes the efficiency gains legible: the same
task gets cheaper, freeing capacity for higher quality or more streams.}

Define the \emph{net information rate} as:
\begin{equation}
  I_{\text{net}} = R_{\text{codec}} \cdot \eta_{\text{transport}}
    \cdot (1 - \rho_{\text{overhead}})
    \cdot (1 - p_{\text{loss}})
  \label{eq:inet}
\end{equation}

\begin{table}[ht]
\centering
\small
\renewcommand{\arraystretch}{1.4}
\begin{tabularx}{\textwidth}{lXXX}
\toprule
\textbf{Parameter} & \textbf{March 2026} & \textbf{2027--2028} &
  \textbf{2029--2031} \\
\midrule
Median link capacity &
  120~Mbps &
  $\sim$150~Mbps &
  200--300~Mbps \\
Dominant codec (WebRTC) &
  VP9 / H.264 &
  VP9 / AV1 (HW) &
  AV1 (HW default) \\
Codec rate (720p/30fps) &
  1.5--2.5~Mbps &
  1.0--1.8~Mbps &
  0.7--1.2~Mbps \\
$\eta_{\text{transport}}$ &
  0.92 (QUIC) &
  0.93 &
  0.94 \\
$\rho_{\text{overhead}}$ (SRTP+QUIC) &
  0.06 &
  0.05 &
  0.04 \\
Typical loss $p$ &
  0.01--0.03 &
  0.01--0.02 &
  0.005--0.01 \\
$I_{\text{net}}$ (720p video) &
  1.3--2.2~Mbps &
  0.9--1.6~Mbps &
  0.6--1.1~Mbps \\
$I_{\text{net}}$ as \% of $C_{\text{link}}$ &
  1.1--1.8\% &
  0.6--1.1\% &
  0.2--0.6\% \\
\bottomrule
\end{tabularx}
\caption{Effective information rate for a 720p/30fps WebRTC video stream.
  The declining ratio $I_{\text{net}} / C_{\text{link}}$ reflects both
  improved codec efficiency and growing link capacity, freeing bandwidth
  for higher resolutions, additional streams, or concurrent data
  channels.}
\label{tab:inet}
\end{table}

The declining ratio $I_{\text{net}} / C_{\text{link}}$ in
Table~\ref{tab:inet} is the central quantitative observation: by 2030, a
720p video call consumes less than 0.5\% of available link capacity. The
surplus capacity enables qualitative changes --- 4K video conferencing,
multi-stream layouts, or concurrent data channels via WebTransport ---
that are quantitatively infeasible in 2026.

\begin{figure}[ht]
\centering
\begin{tikzpicture}[>=Stealth]

% --- Axes ---
\draw[->, thick] (0, 0) -- (11, 0);
\draw[->, thick] (0, 0) -- (0, 5.5);

% --- Axis title ---
\node[font=\small, rotate=90] at (-0.9, 2.75) {Mbps};

% --- Y-axis labels ---
\foreach \y/\lab in {0/0, 1/50, 2/100, 3/150, 4/200, 5/250} {
  \draw (-0.1, \y) -- (0.1, \y);
  \node[font=\tiny, anchor=east] at (-0.15, \y) {\lab};
}

% --- Group: March 2026 ---
% C_link = 120 Mbps -> y = 2.4
\fill[blue!12] (0.8, 0) rectangle (2.4, 2.4);
\draw[blue!50, thick] (0.8, 0) rectangle (2.4, 2.4);
% I_net = 1.75 Mbps (midpoint) -> y = 0.035
\fill[orange!60] (0.8, 0) rectangle (2.4, 0.07);
\draw[orange!80!black, thick] (0.8, 0) rectangle (2.4, 0.07);
\node[font=\footnotesize, anchor=south] at (1.6, 2.5) {120};
\node[font=\tiny, orange!80!black] at (1.6, -0.4) {1.5\%};
\node[font=\small, anchor=north] at (1.6, -0.7) {2026};

% --- Group: 2027--2028 ---
% C_link = 150 Mbps -> y = 3.0
\fill[blue!12] (4.0, 0) rectangle (5.6, 3.0);
\draw[blue!50, thick] (4.0, 0) rectangle (5.6, 3.0);
% I_net = 1.25 Mbps (midpoint) -> y = 0.025
\fill[orange!60] (4.0, 0) rectangle (5.6, 0.05);
\draw[orange!80!black, thick] (4.0, 0) rectangle (5.6, 0.05);
\node[font=\footnotesize, anchor=south] at (4.8, 3.1) {150};
\node[font=\tiny, orange!80!black] at (4.8, -0.4) {0.8\%};
\node[font=\small, anchor=north] at (4.8, -0.7) {2027--28};

% --- Group: 2029--2031 ---
% C_link = 250 Mbps -> y = 5.0
\fill[blue!12] (7.2, 0) rectangle (8.8, 5.0);
\draw[blue!50, thick] (7.2, 0) rectangle (8.8, 5.0);
% I_net = 0.85 Mbps (midpoint) -> y = 0.017
\fill[orange!60] (7.2, 0) rectangle (8.8, 0.034);
\draw[orange!80!black, thick] (7.2, 0) rectangle (8.8, 0.034);
\node[font=\footnotesize, anchor=south] at (8.0, 5.1) {250};
\node[font=\tiny, orange!80!black] at (8.0, -0.4) {0.3\%};
\node[font=\small, anchor=north] at (8.0, -0.7) {2029--31};

% --- Legend ---
\node[anchor=west, font=\footnotesize] at (0.5, 5.8) {%
  \tikz[baseline=-0.5ex]{\fill[blue!12] (0,0) rectangle (0.3,0.3);
    \draw[blue!50, thick] (0,0) rectangle (0.3,0.3);}
  ~$C_{\text{link}}$ (channel capacity) \qquad
  \tikz[baseline=-0.5ex]{\fill[orange!60] (0,0) rectangle (0.3,0.3);
    \draw[orange!80!black, thick] (0,0) rectangle (0.3,0.3);}
  ~$I_{\text{net}}$ (720p stream)};

% --- Annotation ---
\node[font=\scriptsize\itshape, text=black!60, align=center,
  anchor=north] at (5.5, -1.2)
  {A single 720p call uses a vanishing fraction of Shannon-derived\\
   channel capacity, freeing bandwidth for 4K, multi-stream, and data channels.};

\end{tikzpicture}
\caption{Channel capacity $C_{\text{link}}$ vs.\ effective information rate
  $I_{\text{net}}$ for a single 720p/30fps stream across time horizons. The
  orange bar (stream rate) is barely visible against the blue bar (capacity),
  illustrating the growing surplus as codecs improve and bandwidth increases.}
\label{fig:inet-capacity}
\end{figure}

\subsection{Security Entropy and Key Space}

The DTLS-SRTP security architecture (RFCs~8827, 3711) protects all
WebRTC media. Information-theoretic security requires that the key
entropy exceed the information content of the protected material.

For a DTLS~1.3 handshake with ECDHE on P-256, the key agreement
provides:
\begin{equation}
  H_{\text{key}} = 128 \text{~bits of security}
  \label{eq:key-entropy}
\end{equation}

The SRTP master key is 128 bits (AES-128-GCM), giving an exhaustive
search cost of $2^{128}$ operations. The total information content of a
one-hour 720p video call at 2~Mbps is:
\begin{equation}
  I_{\text{call}} = 2 \times 10^6 \times 3600
    = 7.2 \times 10^9 \text{~bits}
    \approx 2^{32.7} \text{~bits}
  \label{eq:call-content}
\end{equation}

The ratio of key entropy to protected content:
\begin{equation}
  \frac{H_{\text{key}}}{I_{\text{call}}}
    = \frac{128}{32.7} \approx 3.91
  \label{eq:security-margin}
\end{equation}

This ratio remains comfortably above~1 even for much longer sessions.
For a continuous 24-hour 4K AV1 stream at 15~Mbps:
\begin{equation}
  I_{\text{stream}} = 15 \times 10^6 \times 86400
    = 1.296 \times 10^{12} \text{~bits}
    \approx 2^{40.2} \text{~bits}
\end{equation}
yielding $H_{\text{key}} / I_{\text{stream}} \approx 3.18$, still
well above the threshold.

\subsubsection{End-to-End Encryption via Encoded Transforms}

The RTCRtpScriptTransform API (Encoded Transforms) enables
application-layer E2EE within the WebRTC pipeline. From an
information-theoretic perspective, this introduces a second encryption
layer with its own key, creating a \emph{nested channel}:
\begin{equation}
  X \xrightarrow{\;K_{\text{E2EE}}\;}
  E_1(X) \xrightarrow{\;K_{\text{SRTP}}\;}
  E_2(E_1(X))
  \label{eq:nested}
\end{equation}

The mutual information between the ciphertext and plaintext, given
neither key, satisfies:
\begin{equation}
  I\!\left(X;\, E_2(E_1(X))\right) \leq
    \min\!\big(H(K_{\text{E2EE}})^{-1},\;
    H(K_{\text{SRTP}})^{-1}\big)
    \cdot I_{\text{call}}
  \label{eq:mutual-e2ee}
\end{equation}

In practice, with $H(K_{\text{E2EE}}) = H(K_{\text{SRTP}}) = 128$
bits, this quantity is negligible ($\ll 1$ bit for any practical
session length), confirming that the nested encryption provides
information-theoretic confidentiality well beyond computational
security guarantees.

Today (March~2026), Encoded Transforms remain Chrome-only. By
2027--2028 (Section~8.3), cross-browser standardization will make
this nested channel universally available. By 2029--2031, SFrame
for group calls will extend this to $n$-party sessions, where the
mutual information between any non-member and the group plaintext
remains negligible:
\begin{equation}
  I\!\left(X;\, E(X) \;\middle|\; K_i \notin
    \{K_1, \ldots, K_n\}\right) \approx 0
  \label{eq:sframe}
\end{equation}

\subsection{WebTransport: Unreliable Datagrams and Channel Capacity}

WebTransport (RFC~9297) introduces unreliable datagrams to browser-server
communication. For latency-sensitive applications (game state, sensor
telemetry), the relevant measure is not throughput but
\emph{freshness-weighted information rate}. Define the value of a
datagram as exponentially decaying with delivery latency $\delta$:
\begin{equation}
  V(\delta) = e^{-\lambda \delta}
  \label{eq:freshness}
\end{equation}
where $\lambda > 0$ is an application-specific decay rate. The
effective information rate under unreliable delivery is:
\begin{equation}
  I_{\text{WT}} = R_{\text{send}} \cdot (1 - p_{\text{loss}})
    \cdot \mathbb{E}\!\left[e^{-\lambda \delta}\right]
  \label{eq:iwt}
\end{equation}

Under reliable delivery (TCP/WebSocket), a retransmission delays
subsequent data, so:
\begin{equation}
  I_{\text{WS}} = R_{\text{send}}
    \cdot \mathbb{E}\!\left[e^{-\lambda(\delta + p \cdot 2\tau)}\right]
  \label{eq:iws}
\end{equation}

The ratio $I_{\text{WT}} / I_{\text{WS}}$ exceeds 1 when:
\begin{equation}
  (1 - p) \cdot e^{\lambda \cdot p \cdot 2\tau} > 1
  \quad \Longleftrightarrow \quad
  \lambda > \frac{-\ln(1 - p)}{2 p \tau}
  \label{eq:wt-advantage}
\end{equation}

For $p = 0.02$ and $\tau = 25$~ms, this gives $\lambda > 20.2$~s$^{-1}$
(half-life $< 34$~ms).

\excursus{Excursus: Why $p = 0.02$ and $\tau = 25$\,ms?}{These values represent
median fixed-broadband conditions in 2026 (Table~\ref{tab:inet}). On mobile or
congested links ($p \approx 0.05$, $\tau \approx 50$\,ms), the threshold
$\lambda^*$ drops to ${\sim}10$\,s$^{-1}$, broadening WebTransport's advantage
to applications with data half-lives up to ${\sim}70$\,ms. The fixed-broadband
reference is thus conservative: mobile scenarios favor unreliable delivery even
more strongly.}

Any application where data older than
$\sim$30~ms is useless benefits from WebTransport's unreliable
mode over WebSocket --- precisely the regime of interactive gaming, live
cursor tracking, and real-time sensor feeds.

As of March~2026, Safari does not support WebTransport
(Section~5), limiting coverage to $\sim$82\%. By 2027--2028
(Section~8.2), universal support will allow applications to rely on
(\ref{eq:iwt}) without fallback paths.

\begin{figure}[ht]
\centering
\begin{tikzpicture}[>=Stealth]

% --- Axes ---
\draw[->, thick] (0, 0) -- (10.5, 0);
\draw[->, thick] (0, 0) -- (0, 5.5);

% --- Axis titles (positioned to clear tick labels) ---
\node[font=\small] at (5, -0.75)
  {Freshness decay rate $\lambda$ [s$^{-1}$]};
\node[font=\small, rotate=90] at (-1.1, 2.75)
  {$I_{\text{WT}} / I_{\text{WS}}$};

% --- X-axis labels ---
\foreach \x/\lab in {0/0, 1.67/10, 3.33/20, 5.0/30, 6.67/40, 8.33/50, 10/60} {
  \draw (\x, -0.1) -- (\x, 0.1);
  \node[font=\tiny, anchor=north] at (\x, -0.15) {\lab};
}

% --- Y-axis labels (0.96 to 1.06, each 0.02 = 1 unit) ---
\foreach \y/\lab in {0/0.96, 1/0.98, 2/1.00, 3/1.02, 4/1.04, 5/1.06} {
  \draw (-0.1, \y) -- (0.1, \y);
  \node[font=\tiny, anchor=east] at (-0.15, \y) {\lab};
}

% --- Ratio = 1.0 line (y=2 maps to ratio 1.0) ---
\draw[thin, gray, dashed] (0, 2) -- (10, 2);

% --- Shade WebTransport advantage region ---
% lambda* = 20.2 -> x = 3.37
\fill[green!10] (3.37, 2) rectangle (10, 5.3);

% --- I_WT/I_WS curve: 0.98*exp(lambda*0.001) ---
% Y range: 0.96-1.06 (each 0.02 = 1 unit). lambda=0: y=1, lambda*=20.2: y=2, lambda=60: y=4.05
\draw[very thick, orange!80!black]
  (0, 1.0) .. controls (1.5, 1.3) and (2.5, 1.7) ..
  (3.37, 2.0) .. controls (4.5, 2.4) and (6, 2.8) ..
  (8, 3.4) .. controls (9, 3.7) and (9.5, 3.85) ..
  (10, 4.05);

% --- Threshold line ---
\draw[thick, dashed, red!60!black] (3.37, 0) -- (3.37, 5.2);
\node[font=\scriptsize, red!60!black, anchor=south, rotate=90] at (3.17, 3.5)
  {$\lambda^* = 20.2$ s$^{-1}$};

% --- Application regime labels ---
\node[font=\scriptsize\itshape, text=blue!60, anchor=north] at (1.5, 5.3)
  {video conf};
\node[font=\scriptsize\itshape, text=green!50!black, anchor=north] at (7.0, 5.3)
  {gaming / cursor / sensors};

% --- Mark advantage/disadvantage ---
\node[font=\tiny, text=red!50!black, anchor=north east] at (3.2, 1.8)
  {WebSocket wins};
\node[font=\tiny, text=green!50!black, anchor=north west] at (3.5, 4.5)
  {WebTransport wins};

% --- Annotation ---
\node[font=\scriptsize\itshape, text=black!60, align=center,
  anchor=north] at (5.0, -1.0)
  {For $p = 0.02$, $\tau = 25$~ms. Applications where data older than\\
   ${\sim}$30~ms is stale gain from unreliable delivery (eq.~\ref{eq:wt-advantage}).};

\end{tikzpicture}
\caption{Freshness-weighted information rate ratio $I_{\text{WT}} / I_{\text{WS}}$
  as a function of decay rate $\lambda$. Above the threshold $\lambda^*$,
  WebTransport's unreliable datagrams deliver more useful information per
  unit time than WebSocket's reliable delivery, because retransmission
  wastes capacity on stale data.}
\label{fig:freshness}
\end{figure}

\subsection{Projections Summary and Practical Implications}

\begin{table}[ht]
\centering
\small
\renewcommand{\arraystretch}{1.4}
\begin{tabularx}{\textwidth}{lXXX}
\toprule
\textbf{Measure} & \textbf{March 2026} & \textbf{2027--2028} &
  \textbf{2029--2031} \\
\midrule
Channel capacity $C_{\text{link}}$ &
  120~Mbps &
  150~Mbps &
  250~Mbps \\
HOL blocking factor $(1\!-\!p)^{n-2}$ at $n\!=\!8$ &
  0.94 ($p\!=\!0.01$) &
  0.94 &
  0.97 ($p\!=\!0.005$) \\
Source coding gap $\Delta$ &
  VP9: $\sim$0.8 &
  AV1 HW: $\sim$0.5 &
  AV1 opt: $\sim$0.4 \\
$I_{\text{net}}$ per 720p stream &
  1.3--2.2~Mbps &
  0.9--1.6~Mbps &
  0.6--1.1~Mbps \\
Max 720p streams in $C_{\text{link}}$ &
  $\sim$55--90 &
  $\sim$95--165 &
  $\sim$230--415 \\
E2EE availability &
  Chrome only &
  Cross-browser &
  SFrame $n$-party \\
WebTransport coverage &
  82\% &
  $\sim$100\% &
  100\% \\
Freshness threshold $\lambda^*$ &
  20.2~s$^{-1}$ &
  $\sim$18~s$^{-1}$ &
  $\sim$15~s$^{-1}$ \\
\bottomrule
\end{tabularx}
\caption{Information-theoretic measures across time horizons. $\Delta$
  is the coding efficiency gap (\ref{eq:coding-gap}), $I_{\text{net}}$
  is the net information rate (\ref{eq:inet}), and $\lambda^*$ is the
  freshness decay threshold above which WebTransport dominates WebSocket
  (\ref{eq:wt-advantage}).}
\label{tab:projections}
\end{table}

The trajectory is clear: the combination of improving source coding
(AV1 hardware), expanding channel capacity (broadband growth), and
maturing transport protocols (QUIC universal, WebTransport universal)
drives the \emph{information cost per unit of communicated experience}
steadily downward. The remaining question is not whether the
Shannon capacity suffices --- it already does, by orders of magnitude
--- but how efficiently the protocol stack converts that capacity into
delivered, fresh, secure information at the application layer. By
2030, the gap between theoretical capacity and realized throughput will
be the narrowest it has ever been for browser-based real-time
communication.

\subsubsection{What the Theory Means for Applications}

The equations and figures in this appendix are not purely academic. Each
theoretical result maps to a concrete design decision that application
developers face today and over the next five years.

\textbf{Codec selection determines bandwidth cost, not quality.}
The rate-distortion analysis (Section~A.2, Figure~\ref{fig:rate-distortion})
shows that AV1 operates at roughly 46\% of H.264's bit rate for equivalent
perceptual quality at 720p. In practice, this means a WebRTC application
switching from H.264 to AV1 can either halve its bandwidth consumption at
the same quality or substantially increase resolution (e.g., 720p to 1080p)
within the same bandwidth envelope. By 2030, with AV1 hardware encoding
ubiquitous and the coding gap $\Delta$ approaching 0.4, further codec
improvements yield diminishing returns --- the stack will be operating close
enough to the Shannon rate-distortion bound that the next generation of
efficiency gains will come from transport and architecture, not codecs alone.

\textbf{QUIC's advantage is situational but decisive on bad networks.}
The head-of-line blocking analysis (Section~A.1.1,
Figure~\ref{fig:hol-blocking}) shows that QUIC's per-stream throughput
advantage over TCP is modest on good networks (6\% at 1\% loss) but
substantial on degraded ones (34\% at 5\% loss). The practical implication
is that applications targeting mobile users, emerging markets, or congested
enterprise networks --- precisely the environments where real-time
communication is hardest --- benefit disproportionately from QUIC-based
transports. For a video conference with 8 participants on a cellular
connection experiencing 3--5\% packet loss, the difference between QUIC and
TCP is the difference between usable and degraded video.

\textbf{Connection latency is a solved problem for returning users.}
The handshake analysis (Section~A.1.2) shows that QUIC 0-RTT eliminates
connection establishment overhead entirely for resumed connections. For a
WebTransport application that a user opens daily --- a collaboration tool,
a live dashboard, a recurring meeting client --- the time to first useful
byte drops to zero. The 1.125~MB of wasted capacity during a TCP+TLS
handshake on a 120~Mbps link may seem small, but for latency-sensitive
applications where the first 100~ms of perceived responsiveness determines
user experience, the difference is measurable.

\textbf{Bandwidth surplus enables architectural ambition.}
The effective information rate analysis (Section~A.3,
Figure~\ref{fig:inet-capacity}) reveals the most consequential practical
finding: a single 720p video stream consumes a vanishing fraction of
available link capacity (under 0.5\% by 2030). This surplus is not merely
``headroom'' --- it is the resource that makes new application architectures
feasible. A 250~Mbps link can theoretically carry 230--415 simultaneous 720p
streams. In practice, this means applications can afford to run multiple
concurrent media flows (video, screen share, spatial audio), data channels
(collaborative editing, file transfer), and telemetry streams without
contending for bandwidth. The constraint shifts from ``can the network carry
this?''\ to ``can the application manage this complexity?''

\textbf{End-to-end encryption adds negligible overhead.}
The security analysis (Section~A.4) confirms that the nested DTLS-SRTP plus
E2EE architecture imposes no meaningful information-theoretic cost. The key
entropy ($2^{128}$) exceeds the information content of any practical session
by a factor of at least 3. For application developers, this means E2EE can
be treated as a default rather than a premium feature --- there is no
bandwidth, latency, or capacity penalty for encrypting at the application
layer on top of the transport-layer encryption that WebRTC already mandates.
The only constraint is API availability, which moves from Chrome-only (2026)
to universal (2028).

\textbf{WebTransport unlocks real-time data that WebSocket cannot serve.}
The freshness analysis (Section~A.5, Figure~\ref{fig:freshness}) quantifies
when unreliable delivery outperforms reliable delivery: any data with a
useful half-life below approximately 34~ms benefits from WebTransport's
unreliable datagrams. This covers a wide class of interactive applications
--- game state synchronization, live cursor and pointer sharing, real-time
sensor feeds, collaborative whiteboard strokes --- where retransmitting a
stale update is worse than dropping it. Before WebTransport, these
applications had to choose between WebSocket (reliable but stale under loss)
and WebRTC DataChannel (unreliable mode available but requiring the full
peer-to-peer ICE/DTLS machinery). WebTransport provides the unreliable
mode with a simple client-server model over HTTP/3, removing the
architectural overhead.

\textbf{The overall picture: the stack converges on efficiency.}
Taken together, the information-theoretic measures trace a consistent
trajectory. The protocol stack is converging toward the theoretical limits
established by Shannon: codecs approach the rate-distortion bound, transports
eliminate unnecessary blocking and handshake overhead, and encryption adds
negligible cost. By 2030, the practical distance between what the standards
deliver and what information theory permits will be the smallest it has ever
been for browser-based real-time communication. The consequence for
application developers is that the infrastructure --- standards, hardware,
and networks --- will no longer be the binding constraint. The remaining
challenge is purely architectural: how to compose these efficient, universal
primitives into applications that fully exploit the capacity and low latency
they provide.

%   % appendix-video-compression.tex
% Standalone appendix — include from retrospective.tex via:
%   \appendix
%   % appendix-info-theory.tex
% Standalone appendix — include from retrospective.tex via:
%   \appendix
%   \input{appendix-info-theory}
%
% Requires in the preamble: amsmath, amssymb, tabularx, booktabs

\section{Information-Theoretic View}
\label{app:info-theory}

This appendix applies information-theoretic measures to the real-time
web communication stack described in this document. We quantify channel
capacity, source coding efficiency, and effective information rates
across three time horizons: today (March~2026), the near term
(2027--2028), and the medium term (2029--2031). The analysis draws on
the standards, codec parameters, and infrastructure projections
developed in the main text.

\subsection{Channel Capacity of the Transport Layer}

The Shannon-Hartley theorem gives the theoretical maximum information
rate for a channel with additive white Gaussian noise:
\begin{equation}
  C = B \log_2\!\left(1 + \frac{S}{N}\right) \quad \text{[bits/s]}
  \label{eq:shannon}
\end{equation}
where $B$ is the channel bandwidth in Hz and $S/N$ is the
signal-to-noise ratio.

While the physical-layer capacity of a broadband link is governed by
(\ref{eq:shannon}), the \emph{effective} capacity available to the
application is reduced by protocol overhead, congestion control, and
transport inefficiencies. We define the effective application-layer
throughput as:
\begin{equation}
  R_{\text{eff}} = C_{\text{link}} \cdot \eta_{\text{transport}}
    \cdot (1 - \rho_{\text{overhead}})
  \label{eq:reff}
\end{equation}
where $C_{\text{link}}$ is the link-layer capacity,
$\eta_{\text{transport}} \in (0,1]$ is the transport efficiency factor
(capturing congestion control, retransmission, and multiplexing
behavior), and $\rho_{\text{overhead}} \in [0,1)$ is the fractional
protocol overhead (headers, encryption, framing).

\subsubsection{Head-of-Line Blocking and Multiplexed Capacity}

Consider $n$ independent streams multiplexed over a single connection.
Under TCP (HTTP/2), a single lost packet blocks all $n$ streams. The
expected throughput per stream under loss rate $p$ is bounded by:
\begin{equation}
  R_{\text{TCP}}^{(i)} \leq \frac{R_{\text{eff}}}{n}
    \cdot (1 - p)^{n-1}
  \label{eq:hol-tcp}
\end{equation}
because a loss in \emph{any} of the $n$ streams stalls all others. Under
QUIC (RFC~9000), streams are independent at the transport layer, so:
\begin{equation}
  R_{\text{QUIC}}^{(i)} \leq \frac{R_{\text{eff}}}{n}
    \cdot (1 - p)
  \label{eq:hol-quic}
\end{equation}
The ratio of effective per-stream throughput is:
\begin{equation}
  \frac{R_{\text{QUIC}}^{(i)}}{R_{\text{TCP}}^{(i)}}
    = \frac{1}{(1-p)^{n-2}}
  \label{eq:hol-ratio}
\end{equation}
\excursus{Excursus: Why $n = 8$?}{In a typical multi-participant video
conference, each peer contributes multiple concurrent streams: a camera video
track, an audio track, and often a screen-share or data channel. A four-person
call with two media tracks each produces $n = 8$ multiplexed streams over a
single connection --- the reference value used throughout this section. Larger
meetings (e.g.\ 10 participants with simulcast layers) push $n$ higher,
amplifying the HOL blocking effect further.}

For $n = 8$ streams (a typical multi-track video conference) at
$p = 0.01$ (1\% loss), this gives a QUIC advantage factor of
$1/(0.99)^6 \approx 1.062$, a modest 6.2\% improvement. At $p = 0.05$
(5\% loss, common on mobile or congested links), the factor is
$1/(0.95)^6 \approx 1.34$, a 34\% throughput advantage per stream.

\begin{figure}[ht]
\centering
\begin{tikzpicture}[>=Stealth]

% --- Shaded region (drawn FIRST so curves appear on top) ---
\fill[orange!8] (2, 0) rectangle (10, 5.3);
\node[font=\tiny\itshape, text=black!40, anchor=north east] at (9.8, 5.2)
  {mobile / congested};
\node[font=\tiny\itshape, text=black!40, anchor=north west] at (0.2, 5.2)
  {fixed broadband};

% --- Grid lines ---
\foreach \y in {1, 2, 3, 4, 5} {
  \draw[gray!20] (0, \y) -- (10, \y);
}

% --- Axes ---
\draw[->, thick] (0, 0) -- (10.5, 0);
\draw[->, thick] (0, 0) -- (0, 5.5);

% --- Axis titles (positioned to clear tick labels) ---
\node[font=\small] at (5, -0.85) {Packet loss rate $p$};
\node[font=\small, rotate=90] at (-0.9, 2.75) {QUIC advantage factor};

% --- X-axis labels ---
\foreach \x/\lab in {0/0, 2/0.02, 4/0.04, 6/0.06, 8/0.08, 10/0.10} {
  \draw (\x, -0.1) -- (\x, 0.1);
  \node[font=\tiny, anchor=north] at (\x, -0.15) {\lab};
}

% --- Y-axis labels ---
\foreach \y/\lab in {0/1.0, 1/1.1, 2/1.2, 3/1.3, 4/1.4, 5/1.5} {
  \draw (-0.1, \y) -- (0.1, \y);
  \node[font=\tiny, anchor=east] at (-0.15, \y) {\lab};
}

% --- n=4 curve: 1/(1-p)^2 ---
\draw[very thick, blue!70]
  (0, 0) .. controls (2.5, 0.5) and (5, 1.1) ..
  (7.5, 2.0) .. controls (8.5, 2.4) and (9.5, 2.9) ..
  (10, 3.1);
\node[font=\scriptsize, blue!70, anchor=south west] at (9.0, 3.0)
  {$n = 4$};

% --- n=8 curve: 1/(1-p)^6 ---
\draw[very thick, orange!80!black]
  (0, 0) .. controls (1.5, 0.5) and (3, 1.7) ..
  (4, 2.5) .. controls (5, 3.5) and (5.5, 4.0) ..
  (6, 4.7);
\node[font=\scriptsize, orange!80!black, anchor=south] at (5.5, 4.8)
  {$n = 8$};

% --- Reference points ---
\fill[orange!80!black] (1.0, 0.62) circle (2.5pt);
\node[font=\tiny, orange!80!black, anchor=south west] at (1.1, 0.62)
  {6.2\%};

\fill[orange!80!black] (5.0, 3.7) circle (2.5pt);
\node[font=\tiny, orange!80!black, anchor=west] at (5.1, 3.7)
  {34\%};

\fill[blue!70] (1.0, 0.20) circle (2.5pt);

% --- Annotation ---
\node[font=\scriptsize\itshape, text=black!60, align=center,
  anchor=north] at (5.0, -1.0)
  {QUIC recovers per-stream throughput closer to $C/n$ (Shannon fair share)\\
   by eliminating head-of-line blocking (eq.~\ref{eq:hol-ratio}).};

\end{tikzpicture}
\caption{QUIC advantage factor $1/(1-p)^{n-2}$ over TCP for multiplexed
  streams. At 5\% loss with $n=8$ streams, QUIC delivers 34\% more
  per-stream throughput by maintaining independence at the transport layer.}
\label{fig:hol-blocking}
\end{figure}

\subsubsection{Connection Establishment Latency}

Let $\tau$ denote the one-way latency (half the round-trip time). The
information-carrying capacity during connection establishment is zero;
the handshake represents pure overhead. The time to first useful byte
is:
\begin{align}
  T_{\text{TCP+TLS\,1.3}} &= 2\tau_{\text{TCP}} + \tau_{\text{TLS}}
    = 3\tau \quad \text{(1-RTT TLS~1.3)} \label{eq:tcp-setup} \\
  T_{\text{QUIC}} &= \tau \quad \text{(1-RTT)}
    \label{eq:quic-setup} \\
  T_{\text{QUIC,\,0-RTT}} &= 0 \quad \text{(resumption)}
    \label{eq:quic-0rtt}
\end{align}
The QUIC 0-RTT case (\ref{eq:quic-0rtt}) eliminates the handshake
entirely for resumed connections. The information lost during
handshake dead time, relative to steady-state capacity, is:
\begin{equation}
  I_{\text{lost}} = C_{\text{link}} \cdot T_{\text{setup}}
  \label{eq:ilost}
\end{equation}
For a 120~Mbps link with $\tau = 25$~ms (typical fixed broadband in
2026), the TCP+TLS~1.3 handshake wastes $120 \times 10^6 \times 0.075
= 9 \times 10^6$ bits (1.125~MB), while QUIC~1-RTT wastes only
$3 \times 10^6$ bits (375~KB).

\subsection{Source Coding Efficiency}

The rate-distortion function $R(D)$ defines the minimum bit rate
required to represent a source at distortion level $D$. For a
Gaussian source with variance $\sigma^2$ under mean squared error
distortion:
\begin{equation}
  R(D) = \frac{1}{2} \log_2 \frac{\sigma^2}{D}
    \quad \text{[bits/sample]}
  \label{eq:rate-distortion}
\end{equation}

\excursus{Excursus: Why Gaussian source?}{The Gaussian memoryless model is used
because it yields the highest (worst-case) $R(D)$ among all sources of equal
variance. Any codec that approaches the Gaussian bound on real video is
therefore doing at least as well as the bound predicts --- the true $R(D)$ for
spatially and temporally correlated video lies below the Gaussian curve, making
it a conservative benchmark.}

No practical codec achieves $R(D)$;\footnote{The $R(D)$ curve in
Figure~\ref{fig:rate-distortion} is the Gaussian memoryless source bound
(equation~\ref{eq:rate-distortion}), which is an upper bound on the true
$R(D)$ for non-Gaussian sources such as video.  Additionally, $R(D)$ assumes
zero excess-distortion probability, whereas practical codecs permit a
frame-level excess-distortion probability on the order of $10^{-6}$ to
$10^{-8}$.  Both factors allow codec operating points measured on real content
to approach or marginally cross the Gaussian bound.} real codecs operate at
some rate $R_{\text{codec}} > R(D)$. We define the coding efficiency gap as:
\begin{equation}
  \Delta = \frac{R_{\text{codec}} - R(D)}{R(D)}
  \label{eq:coding-gap}
\end{equation}

The codecs in the WebRTC stack can be ordered by their empirical
proximity to $R(D)$ for typical video content at conversational
resolutions (720p, 30~fps). Using published Bj\o ntegaard Delta (BD)
rate measurements:
\begin{equation}
  R_{\text{H.264}} > R_{\text{VP9}} > R_{\text{AV1}}
    > R(D)
  \label{eq:codec-ordering}
\end{equation}

Specifically, the empirical rate savings at equivalent PSNR are
approximately:
\begin{align}
  R_{\text{VP9}}  &\approx 0.65 \cdot R_{\text{H.264}}
    \quad (\sim\!35\% \text{ savings}) \label{eq:vp9-savings} \\
  R_{\text{AV1}}  &\approx 0.70 \cdot R_{\text{VP9}}
    \approx 0.46 \cdot R_{\text{H.264}}
    \quad (\sim\!30\% \text{ over VP9})
  \label{eq:av1-savings}
\end{align}

\begin{figure}[ht]
\centering
\begin{tikzpicture}[>=Stealth]

% --- Axes ---
\draw[->, thick] (0, 0) -- (10.5, 0);
\draw[->, thick] (0, 0) -- (0, 6.5);

% --- Axis titles ---
\node[font=\small] at (5, -0.85) {Distortion $D$};
\node[font=\small, rotate=90] at (-0.9, 3.25) {Rate [Mbps]};

% --- Y-axis rate labels ---
\foreach \y/\lab in {1.0/0.5, 2.0/1.0, 3.0/1.5, 4.0/2.0, 5.0/2.5, 6.0/3.0} {
  \draw (-0.1, \y) -- (0.1, \y);
  \node[font=\tiny, anchor=east] at (-0.15, \y) {\lab};
}

% --- R(D) Shannon bound curve ---
\draw[very thick, green!60!black]
  (1.0, 5.8) .. controls (2.0, 3.5) and (3.5, 2.0) ..
  (5.0, 1.3) .. controls (6.5, 0.8) and (8.0, 0.55) ..
  (9.5, 0.4);
\node[font=\scriptsize, green!60!black, anchor=north west] at (8.5, 0.55)
  {$R(D)$ bound};

% --- Reference distortion line (720p/30fps operating point) ---
\draw[thin, gray, dashed] (4.0, 0) -- (4.0, 6.2);
\node[font=\tiny, gray, anchor=south] at (4.0, 6.2) {720p/30fps};

% --- R(D) value at reference distortion ---
\fill[green!60!black] (4.0, 1.55) circle (2pt);

% --- Codec operating points at reference distortion ---
% H.264: ~2.5 Mbps
\draw[thick, dashed, red!70!black] (0.3, 4.8) -- (9.5, 4.8);
\fill[red!70!black] (4.0, 4.8) circle (3pt);
\node[font=\scriptsize, red!70!black, anchor=west] at (9.6, 4.8) {H.264};

% VP9: ~1.6 Mbps (0.65 * H.264)
\draw[thick, dashed, blue!70] (0.3, 3.3) -- (9.5, 3.3);
\fill[blue!70] (4.0, 3.3) circle (3pt);
\node[font=\scriptsize, blue!70, anchor=west] at (9.6, 3.3) {VP9};

% AV1: ~1.15 Mbps (0.46 * H.264)
\draw[thick, dashed, orange!80!black] (0.3, 2.3) -- (9.5, 2.3);
\fill[orange!80!black] (4.0, 2.3) circle (3pt);
\node[font=\scriptsize, orange!80!black, anchor=west] at (9.6, 2.3) {AV1};

% --- Gap annotations (braces) ---
% H.264 gap
\draw[decorate, decoration={brace, amplitude=4pt, mirror},
  thick, red!50!black]
  (3.6, 1.55) -- (3.6, 4.8)
  node[midway, anchor=east, font=\scriptsize, xshift=-5pt]
  {$\Delta_{\text{H.264}}$};

% AV1 gap
\draw[decorate, decoration={brace, amplitude=4pt},
  thick, orange!60!black]
  (4.4, 1.55) -- (4.4, 2.3)
  node[midway, anchor=west, font=\scriptsize, xshift=5pt]
  {$\Delta_{\text{AV1}}$};

% --- Annotation ---
\node[font=\scriptsize\itshape, text=black!60, align=center,
  anchor=north] at (5.0, -0.5)
  {The Shannon rate-distortion bound $R(D)$ is the theoretical minimum.\\
   AV1 narrows the gap $\Delta$ to ${\sim}0.4$ by 2030 (eq.~\ref{eq:coding-gap}).};

\end{tikzpicture}
\caption{Codec operating rates vs.\ the Shannon rate-distortion bound $R(D)$
  at 720p/30fps. Each codec operates above the theoretical minimum; the gap
  $\Delta$ (eq.~\ref{eq:coding-gap}) shrinks from H.264 through VP9 to AV1,
  approaching but never reaching the bound.}
\label{fig:rate-distortion}
\end{figure}

\subsubsection{Codec Entropy and WebCodecs}

The WebCodecs API (W3C Working Draft) exposes per-frame codec access,
allowing measurement and control of the instantaneous encoding rate.
For a video frame $f_t$ at time $t$, the encoder output has empirical
entropy:
\begin{equation}
  \hat{H}(f_t) = -\sum_{s \in \mathcal{S}}
    p(s \mid f_t) \log_2 p(s \mid f_t)
  \label{eq:frame-entropy}
\end{equation}
where $\mathcal{S}$ is the set of symbols emitted by the entropy coder
(CABAC for H.264/AV1, arithmetic coding for VP9). WebCodecs enables
applications to observe $|E(f_t)|$ (the encoded frame size in bits)
directly, which serves as a practical upper bound on
$\hat{H}(f_t)$:
\begin{equation}
  \hat{H}(f_t) \leq |E(f_t)| \leq R_{\text{codec}} / \text{fps}
  \label{eq:entropy-bound}
\end{equation}

\subsection{Effective Information Rate by Time Horizon}

We now combine transport and source coding measures to estimate the
effective information rate for a representative real-time video
session: a single 720p/30fps video call with one audio channel.

\excursus{Excursus: Why 720p?}{Throughout this appendix, 720p/30fps is used as
a fixed reference for cross-era comparison, not as a prediction that 720p
remains the dominant conferencing resolution. The bandwidth surplus quantified
below is precisely what enables migration to 1080p (feasible by ${\sim}$2028
with AV1 hardware encode at today's 720p bit rates) and 4K (feasible by
${\sim}$2030). The 720p baseline makes the efficiency gains legible: the same
task gets cheaper, freeing capacity for higher quality or more streams.}

Define the \emph{net information rate} as:
\begin{equation}
  I_{\text{net}} = R_{\text{codec}} \cdot \eta_{\text{transport}}
    \cdot (1 - \rho_{\text{overhead}})
    \cdot (1 - p_{\text{loss}})
  \label{eq:inet}
\end{equation}

\begin{table}[ht]
\centering
\small
\renewcommand{\arraystretch}{1.4}
\begin{tabularx}{\textwidth}{lXXX}
\toprule
\textbf{Parameter} & \textbf{March 2026} & \textbf{2027--2028} &
  \textbf{2029--2031} \\
\midrule
Median link capacity &
  120~Mbps &
  $\sim$150~Mbps &
  200--300~Mbps \\
Dominant codec (WebRTC) &
  VP9 / H.264 &
  VP9 / AV1 (HW) &
  AV1 (HW default) \\
Codec rate (720p/30fps) &
  1.5--2.5~Mbps &
  1.0--1.8~Mbps &
  0.7--1.2~Mbps \\
$\eta_{\text{transport}}$ &
  0.92 (QUIC) &
  0.93 &
  0.94 \\
$\rho_{\text{overhead}}$ (SRTP+QUIC) &
  0.06 &
  0.05 &
  0.04 \\
Typical loss $p$ &
  0.01--0.03 &
  0.01--0.02 &
  0.005--0.01 \\
$I_{\text{net}}$ (720p video) &
  1.3--2.2~Mbps &
  0.9--1.6~Mbps &
  0.6--1.1~Mbps \\
$I_{\text{net}}$ as \% of $C_{\text{link}}$ &
  1.1--1.8\% &
  0.6--1.1\% &
  0.2--0.6\% \\
\bottomrule
\end{tabularx}
\caption{Effective information rate for a 720p/30fps WebRTC video stream.
  The declining ratio $I_{\text{net}} / C_{\text{link}}$ reflects both
  improved codec efficiency and growing link capacity, freeing bandwidth
  for higher resolutions, additional streams, or concurrent data
  channels.}
\label{tab:inet}
\end{table}

The declining ratio $I_{\text{net}} / C_{\text{link}}$ in
Table~\ref{tab:inet} is the central quantitative observation: by 2030, a
720p video call consumes less than 0.5\% of available link capacity. The
surplus capacity enables qualitative changes --- 4K video conferencing,
multi-stream layouts, or concurrent data channels via WebTransport ---
that are quantitatively infeasible in 2026.

\begin{figure}[ht]
\centering
\begin{tikzpicture}[>=Stealth]

% --- Axes ---
\draw[->, thick] (0, 0) -- (11, 0);
\draw[->, thick] (0, 0) -- (0, 5.5);

% --- Axis title ---
\node[font=\small, rotate=90] at (-0.9, 2.75) {Mbps};

% --- Y-axis labels ---
\foreach \y/\lab in {0/0, 1/50, 2/100, 3/150, 4/200, 5/250} {
  \draw (-0.1, \y) -- (0.1, \y);
  \node[font=\tiny, anchor=east] at (-0.15, \y) {\lab};
}

% --- Group: March 2026 ---
% C_link = 120 Mbps -> y = 2.4
\fill[blue!12] (0.8, 0) rectangle (2.4, 2.4);
\draw[blue!50, thick] (0.8, 0) rectangle (2.4, 2.4);
% I_net = 1.75 Mbps (midpoint) -> y = 0.035
\fill[orange!60] (0.8, 0) rectangle (2.4, 0.07);
\draw[orange!80!black, thick] (0.8, 0) rectangle (2.4, 0.07);
\node[font=\footnotesize, anchor=south] at (1.6, 2.5) {120};
\node[font=\tiny, orange!80!black] at (1.6, -0.4) {1.5\%};
\node[font=\small, anchor=north] at (1.6, -0.7) {2026};

% --- Group: 2027--2028 ---
% C_link = 150 Mbps -> y = 3.0
\fill[blue!12] (4.0, 0) rectangle (5.6, 3.0);
\draw[blue!50, thick] (4.0, 0) rectangle (5.6, 3.0);
% I_net = 1.25 Mbps (midpoint) -> y = 0.025
\fill[orange!60] (4.0, 0) rectangle (5.6, 0.05);
\draw[orange!80!black, thick] (4.0, 0) rectangle (5.6, 0.05);
\node[font=\footnotesize, anchor=south] at (4.8, 3.1) {150};
\node[font=\tiny, orange!80!black] at (4.8, -0.4) {0.8\%};
\node[font=\small, anchor=north] at (4.8, -0.7) {2027--28};

% --- Group: 2029--2031 ---
% C_link = 250 Mbps -> y = 5.0
\fill[blue!12] (7.2, 0) rectangle (8.8, 5.0);
\draw[blue!50, thick] (7.2, 0) rectangle (8.8, 5.0);
% I_net = 0.85 Mbps (midpoint) -> y = 0.017
\fill[orange!60] (7.2, 0) rectangle (8.8, 0.034);
\draw[orange!80!black, thick] (7.2, 0) rectangle (8.8, 0.034);
\node[font=\footnotesize, anchor=south] at (8.0, 5.1) {250};
\node[font=\tiny, orange!80!black] at (8.0, -0.4) {0.3\%};
\node[font=\small, anchor=north] at (8.0, -0.7) {2029--31};

% --- Legend ---
\node[anchor=west, font=\footnotesize] at (0.5, 5.8) {%
  \tikz[baseline=-0.5ex]{\fill[blue!12] (0,0) rectangle (0.3,0.3);
    \draw[blue!50, thick] (0,0) rectangle (0.3,0.3);}
  ~$C_{\text{link}}$ (channel capacity) \qquad
  \tikz[baseline=-0.5ex]{\fill[orange!60] (0,0) rectangle (0.3,0.3);
    \draw[orange!80!black, thick] (0,0) rectangle (0.3,0.3);}
  ~$I_{\text{net}}$ (720p stream)};

% --- Annotation ---
\node[font=\scriptsize\itshape, text=black!60, align=center,
  anchor=north] at (5.5, -1.2)
  {A single 720p call uses a vanishing fraction of Shannon-derived\\
   channel capacity, freeing bandwidth for 4K, multi-stream, and data channels.};

\end{tikzpicture}
\caption{Channel capacity $C_{\text{link}}$ vs.\ effective information rate
  $I_{\text{net}}$ for a single 720p/30fps stream across time horizons. The
  orange bar (stream rate) is barely visible against the blue bar (capacity),
  illustrating the growing surplus as codecs improve and bandwidth increases.}
\label{fig:inet-capacity}
\end{figure}

\subsection{Security Entropy and Key Space}

The DTLS-SRTP security architecture (RFCs~8827, 3711) protects all
WebRTC media. Information-theoretic security requires that the key
entropy exceed the information content of the protected material.

For a DTLS~1.3 handshake with ECDHE on P-256, the key agreement
provides:
\begin{equation}
  H_{\text{key}} = 128 \text{~bits of security}
  \label{eq:key-entropy}
\end{equation}

The SRTP master key is 128 bits (AES-128-GCM), giving an exhaustive
search cost of $2^{128}$ operations. The total information content of a
one-hour 720p video call at 2~Mbps is:
\begin{equation}
  I_{\text{call}} = 2 \times 10^6 \times 3600
    = 7.2 \times 10^9 \text{~bits}
    \approx 2^{32.7} \text{~bits}
  \label{eq:call-content}
\end{equation}

The ratio of key entropy to protected content:
\begin{equation}
  \frac{H_{\text{key}}}{I_{\text{call}}}
    = \frac{128}{32.7} \approx 3.91
  \label{eq:security-margin}
\end{equation}

This ratio remains comfortably above~1 even for much longer sessions.
For a continuous 24-hour 4K AV1 stream at 15~Mbps:
\begin{equation}
  I_{\text{stream}} = 15 \times 10^6 \times 86400
    = 1.296 \times 10^{12} \text{~bits}
    \approx 2^{40.2} \text{~bits}
\end{equation}
yielding $H_{\text{key}} / I_{\text{stream}} \approx 3.18$, still
well above the threshold.

\subsubsection{End-to-End Encryption via Encoded Transforms}

The RTCRtpScriptTransform API (Encoded Transforms) enables
application-layer E2EE within the WebRTC pipeline. From an
information-theoretic perspective, this introduces a second encryption
layer with its own key, creating a \emph{nested channel}:
\begin{equation}
  X \xrightarrow{\;K_{\text{E2EE}}\;}
  E_1(X) \xrightarrow{\;K_{\text{SRTP}}\;}
  E_2(E_1(X))
  \label{eq:nested}
\end{equation}

The mutual information between the ciphertext and plaintext, given
neither key, satisfies:
\begin{equation}
  I\!\left(X;\, E_2(E_1(X))\right) \leq
    \min\!\big(H(K_{\text{E2EE}})^{-1},\;
    H(K_{\text{SRTP}})^{-1}\big)
    \cdot I_{\text{call}}
  \label{eq:mutual-e2ee}
\end{equation}

In practice, with $H(K_{\text{E2EE}}) = H(K_{\text{SRTP}}) = 128$
bits, this quantity is negligible ($\ll 1$ bit for any practical
session length), confirming that the nested encryption provides
information-theoretic confidentiality well beyond computational
security guarantees.

Today (March~2026), Encoded Transforms remain Chrome-only. By
2027--2028 (Section~8.3), cross-browser standardization will make
this nested channel universally available. By 2029--2031, SFrame
for group calls will extend this to $n$-party sessions, where the
mutual information between any non-member and the group plaintext
remains negligible:
\begin{equation}
  I\!\left(X;\, E(X) \;\middle|\; K_i \notin
    \{K_1, \ldots, K_n\}\right) \approx 0
  \label{eq:sframe}
\end{equation}

\subsection{WebTransport: Unreliable Datagrams and Channel Capacity}

WebTransport (RFC~9297) introduces unreliable datagrams to browser-server
communication. For latency-sensitive applications (game state, sensor
telemetry), the relevant measure is not throughput but
\emph{freshness-weighted information rate}. Define the value of a
datagram as exponentially decaying with delivery latency $\delta$:
\begin{equation}
  V(\delta) = e^{-\lambda \delta}
  \label{eq:freshness}
\end{equation}
where $\lambda > 0$ is an application-specific decay rate. The
effective information rate under unreliable delivery is:
\begin{equation}
  I_{\text{WT}} = R_{\text{send}} \cdot (1 - p_{\text{loss}})
    \cdot \mathbb{E}\!\left[e^{-\lambda \delta}\right]
  \label{eq:iwt}
\end{equation}

Under reliable delivery (TCP/WebSocket), a retransmission delays
subsequent data, so:
\begin{equation}
  I_{\text{WS}} = R_{\text{send}}
    \cdot \mathbb{E}\!\left[e^{-\lambda(\delta + p \cdot 2\tau)}\right]
  \label{eq:iws}
\end{equation}

The ratio $I_{\text{WT}} / I_{\text{WS}}$ exceeds 1 when:
\begin{equation}
  (1 - p) \cdot e^{\lambda \cdot p \cdot 2\tau} > 1
  \quad \Longleftrightarrow \quad
  \lambda > \frac{-\ln(1 - p)}{2 p \tau}
  \label{eq:wt-advantage}
\end{equation}

For $p = 0.02$ and $\tau = 25$~ms, this gives $\lambda > 20.2$~s$^{-1}$
(half-life $< 34$~ms).

\excursus{Excursus: Why $p = 0.02$ and $\tau = 25$\,ms?}{These values represent
median fixed-broadband conditions in 2026 (Table~\ref{tab:inet}). On mobile or
congested links ($p \approx 0.05$, $\tau \approx 50$\,ms), the threshold
$\lambda^*$ drops to ${\sim}10$\,s$^{-1}$, broadening WebTransport's advantage
to applications with data half-lives up to ${\sim}70$\,ms. The fixed-broadband
reference is thus conservative: mobile scenarios favor unreliable delivery even
more strongly.}

Any application where data older than
$\sim$30~ms is useless benefits from WebTransport's unreliable
mode over WebSocket --- precisely the regime of interactive gaming, live
cursor tracking, and real-time sensor feeds.

As of March~2026, Safari does not support WebTransport
(Section~5), limiting coverage to $\sim$82\%. By 2027--2028
(Section~8.2), universal support will allow applications to rely on
(\ref{eq:iwt}) without fallback paths.

\begin{figure}[ht]
\centering
\begin{tikzpicture}[>=Stealth]

% --- Axes ---
\draw[->, thick] (0, 0) -- (10.5, 0);
\draw[->, thick] (0, 0) -- (0, 5.5);

% --- Axis titles (positioned to clear tick labels) ---
\node[font=\small] at (5, -0.75)
  {Freshness decay rate $\lambda$ [s$^{-1}$]};
\node[font=\small, rotate=90] at (-1.1, 2.75)
  {$I_{\text{WT}} / I_{\text{WS}}$};

% --- X-axis labels ---
\foreach \x/\lab in {0/0, 1.67/10, 3.33/20, 5.0/30, 6.67/40, 8.33/50, 10/60} {
  \draw (\x, -0.1) -- (\x, 0.1);
  \node[font=\tiny, anchor=north] at (\x, -0.15) {\lab};
}

% --- Y-axis labels (0.96 to 1.06, each 0.02 = 1 unit) ---
\foreach \y/\lab in {0/0.96, 1/0.98, 2/1.00, 3/1.02, 4/1.04, 5/1.06} {
  \draw (-0.1, \y) -- (0.1, \y);
  \node[font=\tiny, anchor=east] at (-0.15, \y) {\lab};
}

% --- Ratio = 1.0 line (y=2 maps to ratio 1.0) ---
\draw[thin, gray, dashed] (0, 2) -- (10, 2);

% --- Shade WebTransport advantage region ---
% lambda* = 20.2 -> x = 3.37
\fill[green!10] (3.37, 2) rectangle (10, 5.3);

% --- I_WT/I_WS curve: 0.98*exp(lambda*0.001) ---
% Y range: 0.96-1.06 (each 0.02 = 1 unit). lambda=0: y=1, lambda*=20.2: y=2, lambda=60: y=4.05
\draw[very thick, orange!80!black]
  (0, 1.0) .. controls (1.5, 1.3) and (2.5, 1.7) ..
  (3.37, 2.0) .. controls (4.5, 2.4) and (6, 2.8) ..
  (8, 3.4) .. controls (9, 3.7) and (9.5, 3.85) ..
  (10, 4.05);

% --- Threshold line ---
\draw[thick, dashed, red!60!black] (3.37, 0) -- (3.37, 5.2);
\node[font=\scriptsize, red!60!black, anchor=south, rotate=90] at (3.17, 3.5)
  {$\lambda^* = 20.2$ s$^{-1}$};

% --- Application regime labels ---
\node[font=\scriptsize\itshape, text=blue!60, anchor=north] at (1.5, 5.3)
  {video conf};
\node[font=\scriptsize\itshape, text=green!50!black, anchor=north] at (7.0, 5.3)
  {gaming / cursor / sensors};

% --- Mark advantage/disadvantage ---
\node[font=\tiny, text=red!50!black, anchor=north east] at (3.2, 1.8)
  {WebSocket wins};
\node[font=\tiny, text=green!50!black, anchor=north west] at (3.5, 4.5)
  {WebTransport wins};

% --- Annotation ---
\node[font=\scriptsize\itshape, text=black!60, align=center,
  anchor=north] at (5.0, -1.0)
  {For $p = 0.02$, $\tau = 25$~ms. Applications where data older than\\
   ${\sim}$30~ms is stale gain from unreliable delivery (eq.~\ref{eq:wt-advantage}).};

\end{tikzpicture}
\caption{Freshness-weighted information rate ratio $I_{\text{WT}} / I_{\text{WS}}$
  as a function of decay rate $\lambda$. Above the threshold $\lambda^*$,
  WebTransport's unreliable datagrams deliver more useful information per
  unit time than WebSocket's reliable delivery, because retransmission
  wastes capacity on stale data.}
\label{fig:freshness}
\end{figure}

\subsection{Projections Summary and Practical Implications}

\begin{table}[ht]
\centering
\small
\renewcommand{\arraystretch}{1.4}
\begin{tabularx}{\textwidth}{lXXX}
\toprule
\textbf{Measure} & \textbf{March 2026} & \textbf{2027--2028} &
  \textbf{2029--2031} \\
\midrule
Channel capacity $C_{\text{link}}$ &
  120~Mbps &
  150~Mbps &
  250~Mbps \\
HOL blocking factor $(1\!-\!p)^{n-2}$ at $n\!=\!8$ &
  0.94 ($p\!=\!0.01$) &
  0.94 &
  0.97 ($p\!=\!0.005$) \\
Source coding gap $\Delta$ &
  VP9: $\sim$0.8 &
  AV1 HW: $\sim$0.5 &
  AV1 opt: $\sim$0.4 \\
$I_{\text{net}}$ per 720p stream &
  1.3--2.2~Mbps &
  0.9--1.6~Mbps &
  0.6--1.1~Mbps \\
Max 720p streams in $C_{\text{link}}$ &
  $\sim$55--90 &
  $\sim$95--165 &
  $\sim$230--415 \\
E2EE availability &
  Chrome only &
  Cross-browser &
  SFrame $n$-party \\
WebTransport coverage &
  82\% &
  $\sim$100\% &
  100\% \\
Freshness threshold $\lambda^*$ &
  20.2~s$^{-1}$ &
  $\sim$18~s$^{-1}$ &
  $\sim$15~s$^{-1}$ \\
\bottomrule
\end{tabularx}
\caption{Information-theoretic measures across time horizons. $\Delta$
  is the coding efficiency gap (\ref{eq:coding-gap}), $I_{\text{net}}$
  is the net information rate (\ref{eq:inet}), and $\lambda^*$ is the
  freshness decay threshold above which WebTransport dominates WebSocket
  (\ref{eq:wt-advantage}).}
\label{tab:projections}
\end{table}

The trajectory is clear: the combination of improving source coding
(AV1 hardware), expanding channel capacity (broadband growth), and
maturing transport protocols (QUIC universal, WebTransport universal)
drives the \emph{information cost per unit of communicated experience}
steadily downward. The remaining question is not whether the
Shannon capacity suffices --- it already does, by orders of magnitude
--- but how efficiently the protocol stack converts that capacity into
delivered, fresh, secure information at the application layer. By
2030, the gap between theoretical capacity and realized throughput will
be the narrowest it has ever been for browser-based real-time
communication.

\subsubsection{What the Theory Means for Applications}

The equations and figures in this appendix are not purely academic. Each
theoretical result maps to a concrete design decision that application
developers face today and over the next five years.

\textbf{Codec selection determines bandwidth cost, not quality.}
The rate-distortion analysis (Section~A.2, Figure~\ref{fig:rate-distortion})
shows that AV1 operates at roughly 46\% of H.264's bit rate for equivalent
perceptual quality at 720p. In practice, this means a WebRTC application
switching from H.264 to AV1 can either halve its bandwidth consumption at
the same quality or substantially increase resolution (e.g., 720p to 1080p)
within the same bandwidth envelope. By 2030, with AV1 hardware encoding
ubiquitous and the coding gap $\Delta$ approaching 0.4, further codec
improvements yield diminishing returns --- the stack will be operating close
enough to the Shannon rate-distortion bound that the next generation of
efficiency gains will come from transport and architecture, not codecs alone.

\textbf{QUIC's advantage is situational but decisive on bad networks.}
The head-of-line blocking analysis (Section~A.1.1,
Figure~\ref{fig:hol-blocking}) shows that QUIC's per-stream throughput
advantage over TCP is modest on good networks (6\% at 1\% loss) but
substantial on degraded ones (34\% at 5\% loss). The practical implication
is that applications targeting mobile users, emerging markets, or congested
enterprise networks --- precisely the environments where real-time
communication is hardest --- benefit disproportionately from QUIC-based
transports. For a video conference with 8 participants on a cellular
connection experiencing 3--5\% packet loss, the difference between QUIC and
TCP is the difference between usable and degraded video.

\textbf{Connection latency is a solved problem for returning users.}
The handshake analysis (Section~A.1.2) shows that QUIC 0-RTT eliminates
connection establishment overhead entirely for resumed connections. For a
WebTransport application that a user opens daily --- a collaboration tool,
a live dashboard, a recurring meeting client --- the time to first useful
byte drops to zero. The 1.125~MB of wasted capacity during a TCP+TLS
handshake on a 120~Mbps link may seem small, but for latency-sensitive
applications where the first 100~ms of perceived responsiveness determines
user experience, the difference is measurable.

\textbf{Bandwidth surplus enables architectural ambition.}
The effective information rate analysis (Section~A.3,
Figure~\ref{fig:inet-capacity}) reveals the most consequential practical
finding: a single 720p video stream consumes a vanishing fraction of
available link capacity (under 0.5\% by 2030). This surplus is not merely
``headroom'' --- it is the resource that makes new application architectures
feasible. A 250~Mbps link can theoretically carry 230--415 simultaneous 720p
streams. In practice, this means applications can afford to run multiple
concurrent media flows (video, screen share, spatial audio), data channels
(collaborative editing, file transfer), and telemetry streams without
contending for bandwidth. The constraint shifts from ``can the network carry
this?''\ to ``can the application manage this complexity?''

\textbf{End-to-end encryption adds negligible overhead.}
The security analysis (Section~A.4) confirms that the nested DTLS-SRTP plus
E2EE architecture imposes no meaningful information-theoretic cost. The key
entropy ($2^{128}$) exceeds the information content of any practical session
by a factor of at least 3. For application developers, this means E2EE can
be treated as a default rather than a premium feature --- there is no
bandwidth, latency, or capacity penalty for encrypting at the application
layer on top of the transport-layer encryption that WebRTC already mandates.
The only constraint is API availability, which moves from Chrome-only (2026)
to universal (2028).

\textbf{WebTransport unlocks real-time data that WebSocket cannot serve.}
The freshness analysis (Section~A.5, Figure~\ref{fig:freshness}) quantifies
when unreliable delivery outperforms reliable delivery: any data with a
useful half-life below approximately 34~ms benefits from WebTransport's
unreliable datagrams. This covers a wide class of interactive applications
--- game state synchronization, live cursor and pointer sharing, real-time
sensor feeds, collaborative whiteboard strokes --- where retransmitting a
stale update is worse than dropping it. Before WebTransport, these
applications had to choose between WebSocket (reliable but stale under loss)
and WebRTC DataChannel (unreliable mode available but requiring the full
peer-to-peer ICE/DTLS machinery). WebTransport provides the unreliable
mode with a simple client-server model over HTTP/3, removing the
architectural overhead.

\textbf{The overall picture: the stack converges on efficiency.}
Taken together, the information-theoretic measures trace a consistent
trajectory. The protocol stack is converging toward the theoretical limits
established by Shannon: codecs approach the rate-distortion bound, transports
eliminate unnecessary blocking and handshake overhead, and encryption adds
negligible cost. By 2030, the practical distance between what the standards
deliver and what information theory permits will be the smallest it has ever
been for browser-based real-time communication. The consequence for
application developers is that the infrastructure --- standards, hardware,
and networks --- will no longer be the binding constraint. The remaining
challenge is purely architectural: how to compose these efficient, universal
primitives into applications that fully exploit the capacity and low latency
they provide.

%   % appendix-video-compression.tex
% Standalone appendix — include from retrospective.tex via:
%   \appendix
%   \input{appendix-info-theory}
%   \input{appendix-video-compression}
%
% Requires in the preamble: amsmath, amssymb, tabularx, booktabs

\section{Video Compression Fundamentals}
\label{app:video-compression}

Appendix~\ref{app:info-theory} quantified the coding efficiency gap $\Delta$
(eq.~\ref{eq:coding-gap}) for each codec generation, treating codecs as black
boxes characterized by their distance from the Shannon rate-distortion bound.
This appendix opens those black boxes. We examine the signal processing and
coding techniques that determine $\Delta$: how video codecs exploit spatial,
temporal, and perceptual redundancy, what distinguishes one codec generation
from the next, and why computational complexity---not algorithmic
capability---remains the binding constraint for real-time web communication.

%% ---------------------------------------------------------------------------
\subsection{Why Video Is Compressible}
\label{sec:compressible}

A raw digital video signal carries far more data than information. Three
fundamental redundancies make compression ratios of 200:1 to 500:1 achievable
at perceptually acceptable quality.

\subsubsection{Spatial Redundancy}

Natural images exhibit strong local correlation. The autocorrelation function
for a typical image row can be modeled as a first-order Markov process:
\begin{equation}
  R_{xx}(k) = \sigma^2 \rho^{|k|}
  \label{eq:spatial-autocorr}
\end{equation}
where $\sigma^2$ is the pixel variance and $\rho \in [0.90, 0.97]$ for
natural content. The separable two-dimensional extension is:
\begin{equation}
  R_{xx}(k, l) = \sigma^2 \,\rho_h^{|k|}\, \rho_v^{|l|}
  \label{eq:spatial-autocorr-2d}
\end{equation}
The power spectral density---the Fourier transform of
(\ref{eq:spatial-autocorr-2d})---is concentrated at low spatial frequencies,
meaning most of the signal energy occupies a small fraction of the frequency
domain. This concentration is the foundation of transform coding
(Section~\ref{sec:transforms}).

\subsubsection{Temporal Redundancy}

Consecutive video frames are highly correlated. For typical conversational
video at 30~fps, the inter-frame correlation coefficient is:
\begin{equation}
  \text{Corr}(f_t, f_{t-1}) \approx 0.90\text{--}0.98
  \label{eq:temporal-corr}
\end{equation}
Most pixels change little between frames; only regions containing motion,
lighting changes, or scene cuts carry significant new information. Motion
estimation (Section~\ref{sec:motion}) exploits this by predicting each block
from a displaced version of a reference frame, transmitting only the
prediction residual.

\subsubsection{Perceptual Irrelevance}

The human visual system (HVS) has lower spatial acuity for chrominance than
luminance. This asymmetry motivates \emph{chroma subsampling}: representing
color difference channels at reduced resolution relative to the luma channel.

\excursus{Excursus: Why 4:2:0?}{All WebRTC codecs operate on YCbCr 4:2:0
video, where each chroma channel is sampled at half the horizontal and half
the vertical resolution of luma. This reduces the sample count by a factor of
$1 + 0.25 + 0.25 = 1.5$ relative to $1 + 1 + 1 = 3$ for full-resolution
4:4:4 color, a 2$\times$ saving before any coding begins. Perceptual
experiments confirm that 4:2:0 is essentially indistinguishable from 4:4:4 at
typical viewing distances for video conferencing---the HVS cannot resolve
the missing chroma detail.}

\subsubsection{The Compression Imperative}

A raw 720p/30fps 4:2:0 video stream at 8~bits per sample requires:
\begin{equation}
  R_{\text{raw}} = 1280 \times 720 \times 1.5 \times 8 \times 30
    = 331.8 \;\text{Mbps}
  \label{eq:raw-rate}
\end{equation}
This exceeds the 120~Mbps median link capacity cited in
Table~\ref{tab:inet}. The Shannon rate-distortion bound
(eq.~\ref{eq:rate-distortion}) yields approximately 0.7--1.5~Mbps at
perceptually acceptable distortion---a compression ratio of 220:1 to 475:1.
Video compression is not optional for real-time communication; it is a
prerequisite.

The first-order entropy of raw pixel values $H(X)$ greatly exceeds the
conditional entropy given neighboring samples:
\begin{equation}
  H(X) \gg H(X \mid X_{\text{neighbors}})
  \label{eq:conditional-entropy}
\end{equation}
This gap---the mutual information between a pixel and its
neighbors---represents the exploitable redundancy. The entire hybrid coding
framework (Section~\ref{sec:hybrid}) is designed to extract this mutual
information through prediction, leaving only the irreducible residual for
transform coding and quantization.

%% ---------------------------------------------------------------------------
\subsection{The Hybrid Coding Framework}
\label{sec:hybrid}

Every major video codec since H.261 (1988)---including every codec in the
WebRTC stack---is a variation on the \emph{block-based hybrid coding}
framework. The encoder partitions each frame into blocks, predicts each block
from previously coded data, transforms and quantizes the prediction residual,
and entropy codes the result.

\subsubsection{Encoder Signal Flow}

For a block at position $(x,y)$ in frame $f_t$, the encoder computes a
prediction $\hat{f}_{\text{pred}}(x,y)$ from either the current frame
(intra prediction) or a reference frame (inter prediction). The prediction
residual is:
\begin{equation}
  e(x,y) = f_t(x,y) - \hat{f}_{\text{pred}}(x,y)
  \label{eq:residual}
\end{equation}
The residual is transformed, quantized, and entropy coded to produce the
bitstream. The decoder reconstructs the frame as:
\begin{equation}
  \hat{f}(x,y) = \hat{e}(x,y) + \hat{f}_{\text{pred}}(x,y)
  \label{eq:reconstruction}
\end{equation}
where $\hat{e}$ is the inverse-transformed, dequantized residual.

\subsubsection{Rate-Distortion Optimization}

Every coding decision---block partition, prediction mode, motion vector,
transform type, quantization level---is made by minimizing a Lagrangian
rate-distortion cost:
\begin{equation}
  J(m) = D(m) + \lambda \cdot R(m)
  \label{eq:rd-cost}
\end{equation}
where $m \in \mathcal{M}$ is a coding mode, $D(m)$ is the distortion
(typically sum of squared errors), $R(m)$ is the bit cost, and $\lambda$ is a
Lagrange multiplier controlling the rate-distortion trade-off. The optimal
mode satisfies:
\begin{equation}
  m^* = \arg\min_{m \in \mathcal{M}} \left[ D(m) + \lambda \cdot R(m) \right]
  \label{eq:rd-opt}
\end{equation}
The Lagrange multiplier relates to the quantization parameter (QP) via an
empirically validated exponential relationship:
\begin{equation}
  \lambda = c \cdot 2^{(QP - 12)/3}
  \label{eq:lambda-qp}
\end{equation}
where $c$ is a codec-specific constant. This relationship ensures that
finer quantization (lower QP) assigns higher cost to bits, favoring modes
that reduce distortion, while coarser quantization (higher QP) favors
modes that reduce rate.

\begin{figure}[ht]
\centering
\begin{tikzpicture}[>=Stealth, scale=0.85, every node/.style={scale=0.85},
  block/.style={draw, rounded corners=3pt, minimum width=1.6cm,
    minimum height=0.8cm, font=\scriptsize, align=center, thick},
  arrow/.style={->, thick},
  label/.style={font=\tiny, text=black!60}]

% --- Input ---
\node[block, fill=blue!15] (input) at (0, 0) {Input\\frame $f_t$};

% --- Subtract ---
\node[draw, circle, inner sep=1.5pt, thick] (sub) at (2.5, 0) {$-$};

% --- Transform ---
\node[block, fill=orange!20] (transform) at (5, 0) {Transform\\$T$};

% --- Quantize ---
\node[block, fill=orange!20] (quant) at (7.5, 0) {Quantize\\$Q$};

% --- Entropy ---
\node[block, fill=purple!20] (entropy) at (10.5, 0) {Entropy\\coder};

% --- Bitstream ---
\node[font=\scriptsize, right=0.4cm of entropy] (bitstream) {Bitstream};

% --- Inverse Quantize ---
\node[block, fill=orange!20] (iquant) at (7.5, -2) {Inverse\\$Q^{-1}$};

% --- Inverse Transform ---
\node[block, fill=orange!20] (itransform) at (5, -2) {Inverse\\$T^{-1}$};

% --- Add ---
\node[draw, circle, inner sep=1.5pt, thick] (add) at (2.5, -2) {$+$};

% --- In-loop filter ---
\node[block, fill=red!15] (filter) at (2.5, -3.8) {In-loop\\filter};

% --- Reference buffer ---
\node[block, fill=red!15] (refbuf) at (5.5, -3.8) {Reference\\buffer};

% --- Mode decision ---
\node[block, fill=green!15] (mode) at (8.5, -3.8) {Mode\\decision};

% --- Prediction ---
\node[block, fill=green!15] (pred) at (5.5, -5.5) {Prediction\\(intra/inter)};

% --- Forward path ---
\draw[arrow] (input) -- (sub);
\draw[arrow] (sub) -- node[above, label] {$e$} (transform);
\draw[arrow] (transform) -- (quant);
\draw[arrow] (quant) -- (entropy);
\draw[arrow] (entropy) -- (bitstream);

% --- Reconstruction path ---
\draw[arrow] (quant) -- (iquant);
\draw[arrow] (iquant) -- (itransform);
\draw[arrow] (itransform) -- node[above, label] {$\hat{e}$} (add);

% --- Filter and buffer ---
\draw[arrow] (add) -- (filter);
\draw[arrow] (filter) -- (refbuf);

% --- Prediction feedback ---
\draw[arrow] (refbuf) -- (mode);
\draw[arrow] (mode) -- (pred);
\draw[arrow] (pred) -| (sub);
\draw[arrow] (pred) -| (add);

% --- Labels ---
\node[label, anchor=west] at (0, -5.5) {$\hat{f}_{\text{pred}}$};

\end{tikzpicture}
\caption{Block-based hybrid encoder architecture. Every codec from H.264
  through VVC follows this structure: predict each block from previously
  coded data (green), transform and quantize the residual (orange), entropy
  code the result (purple), then reconstruct and filter for future
  reference (red). The coding efficiency gap $\Delta$
  (eq.~\ref{eq:coding-gap}) arises from sub-optimal prediction, imperfect
  energy compaction, quantization beyond the rate-distortion bound, and
  entropy coding overhead.}
\label{fig:hybrid-encoder}
\end{figure}

The four sub-optimal components of Figure~\ref{fig:hybrid-encoder}---prediction,
transform, quantization, and entropy coding---each contribute to $\Delta$:
\begin{equation}
  \Delta \approx \Delta_{\text{pred}} + \Delta_{\text{transform}}
    + \Delta_{\text{quant}} + \Delta_{\text{entropy}}
  \label{eq:delta-decomp}
\end{equation}
Generational codec improvements reduce one or more of these terms. The codec
lineages in Sections~\ref{sec:mpeg-lineage}--\ref{sec:aom-lineage} trace
exactly which tools target which terms.

%% ---------------------------------------------------------------------------
\subsection{Transforms}
\label{sec:transforms}

The transform stage decorrelates the prediction residual, concentrating
energy into a few coefficients that can be efficiently quantized and coded.

\subsubsection{The KLT and Its DCT Approximation}
\label{sec:klt-dct}

The optimal decorrelating transform for a stationary source is the
\emph{Karhunen-Lo\`eve Transform} (KLT), whose basis vectors are the
eigenvectors of the input autocorrelation matrix $\mathbf{R}_{xx}$. For a
first-order Markov source with correlation $\rho$ (eq.~\ref{eq:spatial-autocorr}),
the KLT basis converges to the Type-II DCT as block size grows:
\begin{equation}
  C_k(n) = \alpha_k \cos\!\left[\frac{\pi k (2n+1)}{2N}\right],
  \quad k,n = 0, \ldots, N-1
  \label{eq:dct-ii}
\end{equation}
where $\alpha_0 = \sqrt{1/N}$ and $\alpha_k = \sqrt{2/N}$ for $k > 0$.

The energy compaction efficiency---the fraction of total energy captured by
the first $K$ of $N$ coefficients---is:
\begin{equation}
  \eta(K, N) = \frac{\sum_{k=0}^{K-1} \sigma_k^2}
    {\sum_{k=0}^{N-1} \sigma_k^2}
  \label{eq:energy-compaction}
\end{equation}
where $\sigma_k^2$ is the variance of the $k$-th transform coefficient.
For a typical intra residual with $\rho = 0.95$ and $N = 8$, the first 3
DCT coefficients capture approximately 95\% of the energy. This
concentration enables aggressive quantization of high-frequency coefficients
with minimal perceptual impact.

\subsubsection{Integer Transforms in Practice}
\label{sec:integer-transforms}

Practical codecs use integer approximations of the DCT to ensure bit-exact
encoder-decoder agreement. H.264 introduced a $4 \times 4$ integer
transform:
\begin{equation}
  \mathbf{H}_4 = \begin{pmatrix}
    1 &  1 &  1 &  1 \\
    2 &  1 & -1 & -2 \\
    1 & -1 & -1 &  1 \\
    1 & -2 &  2 & -1
  \end{pmatrix}
  \label{eq:h264-transform}
\end{equation}
The non-unity scaling factors are absorbed into the quantization step,
preserving exact integer arithmetic throughout the encode-decode loop.

Each codec generation expands the transform toolkit:

\begin{table}[ht]
\centering
\small
\renewcommand{\arraystretch}{1.4}
\begin{tabularx}{\textwidth}{lXXX}
\toprule
\textbf{Codec} & \textbf{Transform sizes} & \textbf{Transform types}
  & \textbf{Selection method} \\
\midrule
H.264 &
  $4\!\times\!4$, $8\!\times\!8$ &
  Integer DCT &
  Fixed per block size \\
VP9 &
  $4\!\times\!4$ to $32\!\times\!32$ &
  DCT, ADST &
  Row/column adaptive \\
HEVC &
  $4\!\times\!4$ to $32\!\times\!32$ &
  Integer DCT, DST ($4\!\times\!4$ intra) &
  Mode-dependent \\
AV1 &
  $4\!\times\!4$ to $64\!\times\!64$ &
  DCT, ADST, FlipADST, Identity &
  RD-optimized per block \\
VVC &
  $4\!\times\!4$ to $64\!\times\!64$ &
  DCT-II, DST-VII $+$ LFNST &
  RD-optimized $+$ secondary \\
\bottomrule
\end{tabularx}
\caption{Transform sizes and types across codec generations. AV1 applies
  four 1D transform types independently to rows and columns, yielding
  16~possible 2D combinations per block. VVC adds a non-separable secondary
  transform (LFNST) applied to low-frequency coefficients after the primary
  separable transform.}
\label{tab:transforms}
\end{table}

The Asymmetric Discrete Sine Transform (ADST) used in VP9 and AV1 is
particularly effective for intra-predicted residuals, which tend to have
larger magnitudes near the prediction boundary. AV1's FlipADST extends
this to blocks where the large residuals are on the opposite edge, while
the Identity transform (skip) handles blocks where the residual is already
well-distributed in the spatial domain.

The progression from fixed transforms (H.264) to RD-optimized multi-type
selection (AV1, VVC) directly reduces $\Delta_{\text{transform}}$ in
(\ref{eq:delta-decomp}): better energy compaction means fewer bits are
needed to represent the same residual at a given distortion level.

%% ---------------------------------------------------------------------------
\subsection{Motion Estimation and Inter Prediction}
\label{sec:motion}

Inter prediction exploits temporal redundancy by predicting blocks in the
current frame from previously coded reference frames. This is the single
largest source of compression gain---the prediction residual
(eq.~\ref{eq:residual}) typically has 10--20$\times$ less energy than the
original signal.

\subsubsection{Block Matching and Motion Vectors}
\label{sec:block-matching}

The encoder searches for the best-matching region in a reference frame by
minimizing a cost function over candidate displacement vectors
$\mathbf{d} = (d_x, d_y)$:
\begin{equation}
  \text{SAD}(\mathbf{d}) = \sum_{(x,y) \in B}
    \left| f_t(x,y) - f_{t-1}(x + d_x,\, y + d_y) \right|
  \label{eq:sad}
\end{equation}
The Sum of Absolute Differences (SAD) is the simplest distortion measure.
A better rate-distortion proxy is the Sum of Absolute Transformed
Differences (SATD), which applies the Hadamard transform before summation:
\begin{equation}
  \text{SATD}(\mathbf{d}) = \sum \left| \mathbf{H}
    \cdot \left( f_t - f_{t-1}(\mathbf{d}) \right) \right|
  \label{eq:satd}
\end{equation}
SATD approximates the post-transform energy and thus better predicts the
actual coded bit cost. Rate-constrained motion estimation augments the
distortion with a motion vector coding cost:
\begin{equation}
  \mathbf{d}^* = \arg\min_{\mathbf{d}} \left[
    \text{SATD}(\mathbf{d}) + \lambda_{\text{motion}} \cdot R(\mathbf{d})
  \right]
  \label{eq:rcme}
\end{equation}
where $R(\mathbf{d})$ is the bit cost of encoding the motion vector
relative to a predictor.

\subsubsection{Sub-Pixel Interpolation}
\label{sec:subpixel}

Motion in natural video rarely aligns to integer pixel positions. All modern
codecs perform sub-pixel motion compensation by interpolating reference
frame samples at fractional positions.

H.264 uses a 6-tap filter for half-pixel luma interpolation:
\begin{equation}
  h_{\frac{1}{2}} = \{1, -5, 20, 20, -5, 1\} / 32
  \label{eq:h264-halfpel}
\end{equation}
with quarter-pixel positions obtained by bilinear averaging of integer and
half-pixel samples. HEVC extends this to an 8-tap filter for half-pixel and
a 7-tap filter for quarter-pixel positions. AV1 introduces a switchable
filter bank with three 8-tap options---Regular (sharp), Smooth (low-pass),
and Sharp (high-frequency preserving)---selected per block via RD
optimization. VVC adds DMVR (Decoder-side Motion Vector Refinement) and
BDOF (Bi-Directional Optical Flow) for sub-pixel refinement at the decoder
without additional signaling bits.

The progression from H.264's fixed 6-tap filter to AV1's switchable bank
reduces $\Delta_{\text{pred}}$: better interpolation means smaller
prediction residuals, which require fewer bits.

\subsubsection{Advanced Inter Prediction Tools}
\label{sec:advanced-inter}

Beyond basic block matching, each codec generation adds tools to handle
complex motion and occlusion:

\begin{itemize}
\item \textbf{Multiple reference frames.} H.264 introduced up to 16
  reference frames; practical WebRTC encoders use 1--3 for latency reasons.
  More references improve prediction for periodic motion and uncovered
  backgrounds.

\item \textbf{Bi-prediction.} A weighted average of predictions from two
  reference frames, reducing noise and handling gradual scene changes.
  Available from H.264 onward (in B-frames) and in AV1's compound modes.

\item \textbf{Compound prediction (AV1).} Wedge masks and
  difference-weighted blending combine two predictions at the block level,
  handling occlusion boundaries where a single motion model fails.

\item \textbf{Overlapped Block Motion Compensation (AV1).} Blends
  predictions from neighboring blocks at motion boundaries, reducing
  blocking artifacts without post-processing.

\item \textbf{Warped motion (AV1).} An affine model with up to 6
  parameters handles rotation, zoom, and perspective---common in
  screen sharing and camera motion.

\item \textbf{VVC additions.} Geometric Partitioning Mode (non-rectangular
  block splits for motion boundaries), BDOF (optical flow refinement at the
  decoder), and DMVR (decoder-side MV refinement using bilateral matching).
\end{itemize}

\excursus{Excursus: Motion estimation in real-time encoding.}{Full-search
block matching has complexity $O(N^2 \cdot S^2)$ where $N$ is the block
dimension and $S$ is the search range in pixels. For a $64\!\times\!64$
block with $S = 64$, this is ${\sim}17$ million SAD operations per
block---far too expensive for real-time encoding. WebRTC encoders use
hierarchical search patterns (diamond, hexagon) that reduce complexity to
$O(N^2 \cdot \log S)$ with typically $< 5\%$ rate-distortion penalty
relative to full search. This is why real-time encoders use only a subset
of each codec's inter prediction tools: the full tool set is available
for offline encoding, but the real-time constraint forces aggressive
pruning of the mode decision search space.}

%% ---------------------------------------------------------------------------
\subsection{Quantization}
\label{sec:quantization}

Quantization is the sole lossy step in the hybrid coding pipeline. It maps
continuous-valued (or high-precision integer) transform coefficients to a
discrete set of reconstruction levels, introducing controlled information
loss to achieve the target bit rate.

\subsubsection{Scalar Quantization}

Uniform scalar quantization maps coefficient $c$ to a quantization index:
\begin{equation}
  q = \text{round}\!\left(\frac{c}{Q_{\text{step}}}\right), \quad
  \hat{c} = q \cdot Q_{\text{step}}
  \label{eq:uniform-quant}
\end{equation}
Modern codecs use \emph{dead-zone quantization}, where coefficients near
zero are mapped to zero more aggressively than uniform rounding predicts:
\begin{equation}
  q = \text{sign}(c) \cdot \left\lfloor
    \frac{|c| - f \cdot Q_{\text{step}}}{Q_{\text{step}}}
  \right\rfloor
  \label{eq:deadzone}
\end{equation}
where $f \in [0, 0.5)$ controls the dead-zone width. Dead-zone
quantization increases sparsity---the fraction of zero-valued
coefficients---which dramatically improves entropy coding efficiency.

\subsubsection{QP-to-Step-Size Mapping}

The quantization parameter (QP) controls the step size via an exponential
relationship. In the H.264/HEVC/VVC family:
\begin{equation}
  Q_{\text{step}} = Q_0 \cdot 2^{(QP - QP_0)/6}
  \label{eq:qp-mapping}
\end{equation}
A 6-unit increase in QP doubles the step size, approximately doubling the
bit rate. AV1 uses a similar but codec-specific mapping with separate QP
values for DC and AC coefficients, plus per-segment QP adjustment for
region-of-interest coding.

\subsubsection{Quantization Distortion}

For a single coefficient under uniform quantization at high rate, the
distortion is:
\begin{equation}
  D_q = \frac{Q_{\text{step}}^2}{12}
  \label{eq:quant-distortion}
\end{equation}
For a block of $N$ transform coefficients where $K$ survive quantization
(nonzero) and $N - K$ are quantized to zero:
\begin{equation}
  D_{\text{block}} \approx K \cdot \frac{Q_{\text{step}}^2}{12}
    + \sum_{k \notin K} c_k^2
  \label{eq:block-distortion}
\end{equation}
The second term represents the energy of coefficients eliminated by the
dead zone---information permanently discarded. The QP is thus the primary
control knob for navigating the rate-distortion curve in
Figure~\ref{fig:rate-distortion}. The WebCodecs \texttt{VideoEncoder} API
exposes this control via the \texttt{quantizer} configuration parameter,
giving JavaScript applications direct access to this fundamental
trade-off.

%% ---------------------------------------------------------------------------
\subsection{Entropy Coding}
\label{sec:entropy-coding}

After quantization, the surviving coefficients and all side information
(motion vectors, prediction modes, block partitions) must be losslessly
compressed. The entropy coder converts these symbols into a bitstream,
approaching the conditional entropy $H(X|\text{context})$ as closely as
possible. The gap between the entropy coder's output rate and the true
conditional entropy is $\Delta_{\text{entropy}}$ in
(\ref{eq:delta-decomp}).

\subsubsection{Context-Adaptive Binary Arithmetic Coding (CABAC)}
\label{sec:cabac}

CABAC, used in H.264, HEVC, and VVC, encodes each syntax element as a
sequence of binary decisions (bins). For each bin, the arithmetic coder
subdivides the current interval $[L, H)$ based on the estimated
probability $p$:
\begin{equation}
  [L, H) \to \begin{cases}
    [L,\; L + p \cdot (H - L)) & \text{if bin} = 0 \\
    [L + p \cdot (H - L),\; H) & \text{if bin} = 1
  \end{cases}
  \label{eq:arithmetic-coding}
\end{equation}
The theoretical output length approaches the entropy:
\begin{equation}
  \ell \to \sum_{i=1}^{n} -\log_2 p(b_i \mid \text{ctx}_i)
  \label{eq:cabac-length}
\end{equation}
CABAC achieves redundancy below 0.01 bits/symbol relative to the true
conditional entropy---near-optimal compression.

Context modeling is the key to CABAC's effectiveness. Each bin type has an
associated context, whose probability is updated after each observation via
a state-machine-based exponential moving average:
\begin{equation}
  p_{k+1} = p_k + \alpha(s_k) \cdot (b_k - p_k)
  \label{eq:cabac-update}
\end{equation}
where $\alpha(s_k)$ is a state-dependent adaptation rate, indexed by a
lookup table with 64 states (H.264) or 128 states (HEVC/VVC). This design
avoids floating-point arithmetic entirely.

CABAC also provides a \emph{bypass mode} for bins with approximately
uniform probability (e.g., sign bits, high-magnitude coefficient levels),
which skips the context lookup and probability estimation for throughput.

\subsubsection{Multi-Symbol Arithmetic Coding in AV1}
\label{sec:av1-entropy}

AV1 departs from binary arithmetic coding. Its entropy coder operates on
multi-symbol alphabets using CDF (Cumulative Distribution Function) tables:
\begin{equation}
  \text{Encode symbol } s: \quad
  [L, H) \to \left[L + \frac{\text{CDF}(s-1)}{M} \cdot (H-L),\;
    L + \frac{\text{CDF}(s)}{M} \cdot (H-L)\right)
  \label{eq:av1-entropy}
\end{equation}
where $M$ is the total frequency count. CDF tables are adapted after
each symbol:
\begin{equation}
  \text{CDF}_{k+1}(s) = \text{CDF}_k(s) + \left\lfloor
    \frac{\text{target}(s) - \text{CDF}_k(s)}{n}
  \right\rfloor
  \label{eq:cdf-update}
\end{equation}
where $n$ controls the adaptation rate. Multi-symbol coding is more
efficient than CABAC's binarization for syntax elements with moderate
alphabet sizes (e.g., intra prediction modes, transform types), as it
avoids the overhead of converting each symbol into a binary tree of
decisions.

\begin{table}[ht]
\centering
\small
\renewcommand{\arraystretch}{1.4}
\begin{tabularx}{\textwidth}{lXXXXX}
\toprule
\textbf{Feature} & \textbf{H.264} & \textbf{HEVC} & \textbf{VP9}
  & \textbf{AV1} & \textbf{VVC} \\
\midrule
Method &
  Binary AC &
  Binary AC &
  Boolean AC &
  Multi-sym AC &
  Binary AC \\
Contexts &
  ${\sim}400$ &
  ${\sim}200$ &
  ${\sim}900$ &
  ${\sim}300$ &
  ${\sim}700$ \\
Adaptation &
  64-state LUT &
  128-state LUT &
  Counter &
  CDF update &
  128-state LUT \\
Bypass mode &
  Yes &
  Yes &
  N/A &
  N/A &
  Yes \\
Parallelism &
  Slice &
  WPP / Tiles &
  Tiles &
  Tiles &
  WPP / Tiles \\
\bottomrule
\end{tabularx}
\caption{Entropy coding comparison across codecs. HEVC reduced context
  count relative to H.264 for throughput; VVC increased it again for
  improved compression. AV1's multi-symbol approach operates on full
  alphabets rather than binary decompositions.}
\label{tab:entropy-coding}
\end{table}

The frame entropy $\hat{H}(f_t)$ from eq.~(\ref{eq:frame-entropy}) is
ultimately determined by how well the entropy coder's probability model
matches the true symbol distribution. Both CABAC and AV1's multi-symbol
coder achieve near-zero entropy coding gaps; the remaining
$\Delta_{\text{entropy}}$ arises primarily from context model mismatch on
non-stationary content (scene changes, sudden motion onset).

%% ---------------------------------------------------------------------------
\subsection{In-Loop Filtering}
\label{sec:filtering}

Block-based coding introduces visible artifacts at block boundaries:
discontinuities from independent quantization of adjacent blocks, and
ringing from coefficient truncation. In-loop filters reduce these artifacts
\emph{before} the reconstructed frame enters the reference buffer, improving
prediction quality for subsequent frames and thus reducing
$\Delta_{\text{pred}}$.

\begin{itemize}
\item \textbf{Deblocking filter (all codecs).} Smooths discontinuities
  at block edges, with filter strength adapted to the quantization level
  and local gradient. Introduced in H.264; refined in each subsequent
  codec.

\item \textbf{Sample Adaptive Offset---SAO (HEVC).} Classifies
  reconstructed samples by intensity or edge direction and applies
  per-class offsets to reduce systematic bias introduced by quantization.

\item \textbf{Constrained Directional Enhancement Filter---CDEF (AV1).}
  A directional filter that identifies edge orientation per
  $8\!\times\!8$ block and applies filtering only along the edge,
  preserving sharpness while reducing ringing.

\item \textbf{Loop Restoration Filter (AV1).} A switchable per-tile
  filter choosing between a Wiener filter (FIR designed to minimize MSE
  between original and reconstruction) and a Self-Guided filter (non-linear
  denoising based on guided image filtering). This is AV1's most powerful
  quality enhancement tool.

\item \textbf{Adaptive Loop Filter---ALF (VVC).} A diamond-shaped filter
  with coefficients signaled per frame, applied adaptively per
  $4\!\times\!4$ block based on local classification. Subsumes HEVC's
  SAO functionality.
\end{itemize}

The evolution from simple deblocking (H.264) to multi-stage adaptive
filtering (AV1: deblocking $\to$ CDEF $\to$ loop restoration) exemplifies
how each codec generation reduces $\Delta_{\text{pred}}$ by improving
reference frame quality.

%% ---------------------------------------------------------------------------
\subsection{Codec Lineage: MPEG/ITU Family}
\label{sec:mpeg-lineage}

The MPEG/ITU standardization track has produced three codec generations
relevant to the 2021--2031 timeframe of this document.

\subsubsection{H.264/AVC (2003)}
\label{sec:h264}

H.264 defined the modern era of video compression and remains the most
widely deployed codec. Its key innovations over its predecessors (H.263,
MPEG-2):

\begin{itemize}
\item Variable block sizes: $16\!\times\!16$ macroblock partitioned into
  $16\!\times\!8$, $8\!\times\!16$, $8\!\times\!8$, $8\!\times\!4$,
  $4\!\times\!8$, and $4\!\times\!4$ sub-partitions for motion
  compensation.

\item Quarter-pixel motion compensation with 6-tap interpolation
  (eq.~\ref{eq:h264-halfpel}).

\item Context-adaptive intra prediction: 9~angular modes for
  $4\!\times\!4$ blocks, 4~modes for $16\!\times\!16$.

\item Dual entropy coding: CAVLC (simpler, faster) and CABAC
  (Section~\ref{sec:cabac}; 10--15\% better compression).

\item In-loop deblocking filter, mandatory for the decoder.

\item Multiple reference frames (up to 16) and weighted prediction.
\end{itemize}

Three profiles are relevant to WebRTC: \textbf{Constrained Baseline}
(no CABAC, no B-frames, lowest complexity---the mandatory WebRTC profile),
\textbf{Main} (CABAC, B-frames, ${\sim}$10\% better compression), and
\textbf{High} ($8\!\times\!8$ transform, adaptive quantization matrices,
${\sim}$15\% better).

At the 720p/30fps reference point, H.264 Constrained Baseline operates at
approximately 2.0--2.5~Mbps for ``good'' quality (PSNR ${\sim}$37~dB),
corresponding to $\Delta_{\text{H.264}} \approx 2.1$ from
Appendix~\ref{app:info-theory}. H.264 was the only codec Safari supported
in 2021 (Section~3.1), and its higher $\Delta$ is directly responsible
for the bandwidth disparity between Safari-only and VP9-capable sessions
documented in the main text.

\subsubsection{H.265/HEVC (2013)}
\label{sec:hevc}

HEVC replaced H.264's fixed $16\!\times\!16$ macroblock with a recursive
Coding Tree Unit (CTU) of up to $64\!\times\!64$ pixels, using quad-tree
partitioning down to $8\!\times\!8$. Key advances:

\begin{itemize}
\item 35 intra prediction angular directions (vs.\ H.264's 9).
\item Transform sizes from $4\!\times\!4$ to $32\!\times\!32$, with
  DST for $4\!\times\!4$ intra blocks.
\item Advanced Motion Vector Prediction (AMVP): two MVP candidates selected
  from spatial and temporal neighbors, replacing H.264's less structured
  motion vector prediction.
\item Merge mode: inherits the full motion information (MV, reference
  index, prediction direction) from a neighboring block, requiring zero
  additional motion bits.
\item Sample Adaptive Offset (SAO) in-loop filter.
\item Parallel processing: Wavefront Parallel Processing (WPP) and tiles.
\end{itemize}

HEVC achieves approximately 40\% rate reduction over H.264 at equivalent
PSNR. However, HEVC is absent from WebRTC entirely---not due to technical
limitations, but due to patent licensing.

\excursus{Excursus: Why no HEVC in WebRTC?}{HEVC's patent licensing is
fragmented across three separate pools (MPEG-LA, HEVC Advance, Velos Media)
plus independent licensors who are not members of any pool. The resulting
uncertainty about total royalty cost---estimated at \$0.20--\$0.40 per device
for some use cases---led Google, Mozilla, Cisco, and others to form the
Alliance for Open Media (AOM) in 2015 and develop AV1 as a royalty-free
alternative. This patent-driven fork is why the WebRTC ecosystem followed the
VP8$\to$VP9$\to$AV1 path rather than H.264$\to$HEVC$\to$VVC, and why the two
codec families must be understood in parallel.}

\subsubsection{H.266/VVC (2020)}
\label{sec:vvc}

VVC extends HEVC's quad-tree with binary and ternary splits, creating a
\emph{multi-type tree} (MTT) that enables asymmetric block shapes (e.g.,
$32\!\times\!8$, $8\!\times\!32$, $32\!\times\!4$). Additional
innovations:

\begin{itemize}
\item 67 intra prediction directions plus Matrix-based Intra Prediction
  (MIP), which applies a learned linear transform to neighboring samples.

\item Multiple Transform Selection (MTS): explicit signaling of DCT-II or
  DST-VII per block, plus the Low-Frequency Non-Separable Transform
  (LFNST)---a $16\!\times\!16$ or $8\!\times\!8$ non-separable transform
  applied to low-frequency coefficients after the primary separable
  transform.

\item Affine motion model: 4-parameter (translation $+$ rotation/scaling)
  or 6-parameter (full affine) motion, derived from control point MVs at
  block corners.

\item Decoder-side tools: DMVR refines motion vectors at the decoder using
  bilateral template matching; BDOF applies optical flow equations to
  refine bi-predicted samples, both without additional signaling bits.

\item Geometric Partitioning Mode: divides a block diagonally for separate
  motion compensation of each partition, targeting motion boundaries.

\item Adaptive Loop Filter (ALF): a diamond-shaped Wiener filter with
  coefficients signaled per frame and applied adaptively per
  $4\!\times\!4$ block.
\end{itemize}

VVC achieves approximately 30--50\% rate reduction over HEVC, placing it at
$\Delta_{\text{VVC}} \approx 0.3$. Its relevance to this document is
forward-looking: VVC may appear in WebCodecs codec registries by
${\sim}$2030 if licensing concerns are resolved, and its tool set
represents the current efficiency frontier against which AV1's performance
should be measured.

%% ---------------------------------------------------------------------------
\subsection{Codec Lineage: Google/AOM Family}
\label{sec:aom-lineage}

The royalty-free codec family dominates the WebRTC ecosystem. Its
development was driven by the patent licensing concerns described in the
HEVC excursus above.

\subsubsection{VP8 (2010)}
\label{sec:vp8}

VP8 was Google's first open-source video codec, released alongside the
WebM container. Its design is conservative:

\begin{itemize}
\item $16\!\times\!16$ macroblocks with $4\!\times\!4$ subblock partitions.
\item $4\!\times\!4$ Walsh-Hadamard Transform (WHT) for DC coefficients,
  $4\!\times\!4$ DCT for AC coefficients.
\item Boolean arithmetic coder with tree-structured syntax (single-bit
  alphabet).
\item 4 intra prediction modes for $16\!\times\!16$, 10 for
  $4\!\times\!4$ subblocks.
\item Simple loop filter with per-edge strength adaptation.
\item Single reference frame for WebRTC (plus a ``golden'' frame for
  screen content).
\end{itemize}

VP8's compression efficiency is roughly equivalent to H.264 Constrained
Baseline. It served as the baseline WebRTC codec---supported by Chrome and
Firefox since WebRTC~1.0---but has been superseded by VP9 in all modern
deployments.

\subsubsection{VP9 (2013)}
\label{sec:vp9}

VP9 introduced a superblock architecture and several tools that brought it
within range of HEVC:

\begin{itemize}
\item $64\!\times\!64$ superblocks with recursive quad-tree partitioning
  down to $4\!\times\!4$.
\item 10 intra prediction modes.
\item Transforms: DCT and ADST applied independently per row and column,
  at sizes from $4\!\times\!4$ to $32\!\times\!32$.
\item Switchable interpolation filters: 8-tap Regular, Smooth, and
  Sharp, selected per block.
\item Compound prediction: weighted average from two reference frames.
\item Segmentation: per-segment QP adjustment for region-of-interest
  coding.
\end{itemize}

VP9 achieves approximately 35\% rate reduction over H.264
(eq.~\ref{eq:vp9-savings}), placing it at $\Delta_{\text{VP9}} \approx 0.8$.
VP9 is the dominant WebRTC codec in Chrome and Firefox as of March~2026,
and its temporal SVC support (Section~\ref{sec:svc}) is the primary
mechanism enabling adaptive quality in multi-participant calls.

\subsubsection{AV1 (2018)}
\label{sec:av1}

AV1 represents the most ambitious single-generation improvement in the
royalty-free family. Its tool set approaches VVC in breadth while
maintaining royalty-free status:

\begin{itemize}
\item $128\!\times\!128$ superblocks with flexible partitioning (quad-tree
  plus rectangular splits) down to $4\!\times\!4$.

\item 56+ directional intra prediction modes plus non-directional modes
  (Paeth, Smooth variants, DC).

\item Full inter prediction toolkit: OBMC, warped motion (affine),
  compound modes (Wedge, Difference-weighted), inter-intra compound
  prediction.

\item 16 transform type combinations ($4$ 1D types $\times$ rows/columns),
  at sizes up to $64\!\times\!64$.

\item Film Grain Synthesis: the encoder characterizes noise/grain
  parameters and transmits them as metadata; the decoder re-synthesizes
  grain after reconstruction. This removes non-compressible high-frequency
  noise from the coding loop---a major bitrate saving for camera-captured
  content.

\item Multi-stage in-loop filtering: deblocking $\to$ CDEF $\to$ loop
  restoration (Wiener or Self-Guided), as described in
  Section~\ref{sec:filtering}.

\item Multi-symbol arithmetic coding with CDF-based context adaptation
  (Section~\ref{sec:av1-entropy}).
\end{itemize}

AV1 achieves approximately 30\% rate reduction over VP9
(eq.~\ref{eq:av1-savings}), a cumulative ${\sim}$54\% over H.264, placing
it at $\Delta_{\text{AV1}} \approx 0.5$ with current hardware encoders and
approaching 0.4 by 2030 as encoder implementations mature.

AV1's tools directly explain why its $\Delta$ is smaller than VP9's:
more transform types capture residual structure better (reducing
$\Delta_{\text{transform}}$), warped and compound motion reduce prediction
residual energy (reducing $\Delta_{\text{pred}}$), and film grain synthesis
removes non-compressible variance from the coding loop (effectively
reducing the source variance $\sigma^2$ that the codec must represent).

%% ---------------------------------------------------------------------------
\subsection{Block Partitioning: A Generational Comparison}
\label{sec:partitioning}

Block partitioning is the most architecturally consequential difference
between codec generations. Larger and more flexible partitions allow the
encoder to match block boundaries to content structure---smooth regions
get large blocks (few bits for partition overhead), detailed regions get
small blocks (better prediction accuracy).

\begin{figure}[ht]
\centering
\begin{tikzpicture}[>=Stealth, scale=0.55, every node/.style={scale=0.55}]

% --- H.264 panel ---
\begin{scope}[xshift=0cm]
  \node[font=\normalsize\bfseries, anchor=south] at (2, 8.4) {H.264};
  \node[font=\small, anchor=south] at (2, 7.8) {$16\!\times\!16$ max};
  % 4x4 grid of 16x16 macroblocks in a 64x64 region
  \draw[thick] (0,0) rectangle (4,4);
  \draw[thick] (0,0) rectangle (4,4);
  % macroblock grid
  \foreach \x in {0,1,2,3} {
    \foreach \y in {0,1,2,3} {
      \draw[gray] (\x, \y) rectangle (\x+1, \y+1);
    }
  }
  % show some sub-partitions
  \draw[blue!60, thick] (0,3) -- (0.5,3) -- (0.5,4) -- (0,4);
  \draw[blue!60, thick] (0.5,3) -- (0.5,4);
  \draw[blue!60, thick] (0,3.5) -- (1,3.5);
  % 8x8 in another
  \draw[blue!60, thick] (1,2.5) -- (2,2.5);
  \draw[blue!60, thick] (1.5,2) -- (1.5,3);
  % 4x4 in another
  \draw[red!60, thick] (2,0) -- (2,1) -- (3,1) -- (3,0);
  \draw[red!60, thick] (2.5,0) -- (2.5,0.5);
  \draw[red!60, thick] (2,0.5) -- (2.5,0.5);
  \draw[red!60, thick] (2.5,0.5) -- (2.5,1);
  \draw[red!60, thick] (2.5,0) -- (3,0) -- (3,0.5) -- (2.5,0.5);
  % Scale label
  \node[font=\small, anchor=north] at (2, -0.3) {64 pixels};
\end{scope}

% --- VP9/HEVC panel ---
\begin{scope}[xshift=6cm]
  \node[font=\normalsize\bfseries, anchor=south] at (4, 8.4) {VP9 / HEVC};
  \node[font=\small, anchor=south] at (4, 7.8) {$64\!\times\!64$ max};
  % One 64x64 superblock
  \draw[thick] (0,0) rectangle (8,8);
  % quad-tree splits
  \draw[thick] (4,0) -- (4,8);
  \draw[thick] (0,4) -- (8,4);
  % further splits in top-left quadrant
  \draw[blue!60, thick] (2,4) -- (2,8);
  \draw[blue!60, thick] (0,6) -- (4,6);
  % further split in top-left-top-left
  \draw[red!60, thick] (1,6) -- (1,8);
  \draw[red!60, thick] (0,7) -- (2,7);
  % split bottom-right quadrant
  \draw[blue!60, thick] (6,0) -- (6,4);
  \draw[blue!60, thick] (4,2) -- (8,2);
  % deeper in bottom-right-bottom-left
  \draw[red!60, thick] (5,0) -- (5,2);
  \draw[red!60, thick] (4,1) -- (6,1);
  % Scale label
  \node[font=\small, anchor=north] at (4, -0.3) {64 pixels};
\end{scope}

% --- AV1 panel ---
\begin{scope}[xshift=16cm]
  \node[font=\normalsize\bfseries, anchor=south] at (4, 8.4) {AV1};
  \node[font=\small, anchor=south] at (4, 7.8) {$128\!\times\!128$ max};
  % 128x128 shown as 8x8 grid (each cell = 16px)
  \draw[thick] (0,0) rectangle (8,8);
  % quad split: top half / bottom half
  \draw[thick] (4,0) -- (4,8);
  \draw[thick] (0,4) -- (8,4);
  % rectangular splits in top-left
  \draw[blue!60, thick] (0,6) -- (4,6);
  \draw[blue!60, thick] (2,4) -- (2,6);
  % asymmetric in bottom-left (4:1 rect)
  \draw[orange!80, thick] (0,3) -- (4,3);
  \draw[orange!80, thick] (0,2) -- (4,2);
  \draw[orange!80, thick] (0,1) -- (4,1);
  % quad + rect in top-right
  \draw[blue!60, thick] (6,4) -- (6,8);
  \draw[red!60, thick] (4,6) -- (6,6);
  \draw[red!60, thick] (5,4) -- (5,6);
  % bottom-right: large blocks
  \draw[blue!60, thick] (4,2) -- (8,2);
  % Scale label
  \node[font=\small, anchor=north] at (4, -0.3) {128 pixels};
\end{scope}

% --- VVC panel ---
\begin{scope}[xshift=26cm]
  \node[font=\normalsize\bfseries, anchor=south] at (4, 8.4) {VVC};
  \node[font=\small, anchor=south] at (4, 7.8) {$128\!\times\!128$ max};
  \draw[thick] (0,0) rectangle (8,8);
  % quad split
  \draw[thick] (4,0) -- (4,8);
  \draw[thick] (0,4) -- (8,4);
  % binary split (horizontal) in top-left
  \draw[blue!60, thick] (0,6) -- (4,6);
  % ternary split (vertical) in top-left-bottom
  \draw[orange!80, thick] (1,4) -- (1,6);
  \draw[orange!80, thick] (3,4) -- (3,6);
  % binary split (vertical) in top-right
  \draw[blue!60, thick] (6,4) -- (6,8);
  % ternary horizontal in top-right-left
  \draw[orange!80, thick] (4,5.33) -- (6,5.33);
  \draw[orange!80, thick] (4,6.67) -- (6,6.67);
  % bottom-left: binary horizontal
  \draw[blue!60, thick] (0,2) -- (4,2);
  % deeper ternary in bottom-left-top
  \draw[orange!80, thick] (1,2) -- (1,4);
  \draw[orange!80, thick] (3,2) -- (3,4);
  % bottom-right: mostly large
  \draw[blue!60, thick] (6,0) -- (6,4);
  \draw[blue!60, thick] (4,2) -- (8,2);
  % Scale label
  \node[font=\small, anchor=north] at (4, -0.3) {128 pixels};
\end{scope}

\end{tikzpicture}
\caption{Block partitioning comparison across codec generations. H.264 is
  limited to $16\!\times\!16$ macroblocks with fixed sub-partitions (left).
  VP9/HEVC use quad-tree partitioning within $64\!\times\!64$ superblocks
  (center-left). AV1 extends to $128\!\times\!128$ with quad-tree plus
  rectangular splits (center-right). VVC adds binary and ternary splits for
  maximum flexibility (right). Blue lines show first-level splits; red
  shows deeper partitioning; orange shows asymmetric/ternary splits.}
\label{fig:block-partitioning}
\end{figure}

\begin{table}[ht]
\centering
\small
\renewcommand{\arraystretch}{1.4}
\begin{tabularx}{\textwidth}{lXXXXX}
\toprule
\textbf{Feature} & \textbf{H.264} & \textbf{VP9} & \textbf{HEVC}
  & \textbf{AV1} & \textbf{VVC} \\
\midrule
Max block &
  $16\!\times\!16$ &
  $64\!\times\!64$ &
  $64\!\times\!64$ &
  $128\!\times\!128$ &
  $128\!\times\!128$ \\
Min block &
  $4\!\times\!4$ &
  $4\!\times\!4$ &
  $8\!\times\!8$ &
  $4\!\times\!4$ &
  $4\!\times\!4$ \\
Split types &
  Fixed partitions &
  Quad-tree &
  Quad-tree &
  Quad $+$ rect &
  Quad $+$ BT $+$ TT \\
Asymmetric &
  Limited &
  No &
  No &
  4:1 rectangles &
  Full (BT/TT) \\
\bottomrule
\end{tabularx}
\caption{Block partitioning capabilities by codec. BT = binary tree
  (horizontal or vertical), TT = ternary tree (1:2:1 split). Larger maximum
  blocks reduce overhead for homogeneous regions; asymmetric splits improve
  partitioning efficiency at content boundaries.}
\label{tab:partitioning}
\end{table}

%% ---------------------------------------------------------------------------
\subsection{Rate-Distortion Operating Points by Codec}
\label{sec:rd-comparison}

The codec-specific analyses of
Sections~\ref{sec:mpeg-lineage}--\ref{sec:aom-lineage} can be unified
through the Bj\o ntegaard Delta Rate (BD-Rate), the standard metric for
comparing codec efficiency. BD-Rate computes the average percentage rate
difference between two codecs at equivalent distortion by integrating the
difference between their rate-distortion curves fit as cubic polynomials in
PSNR-vs-log-rate space:
\begin{equation}
  \text{BD-Rate} = \frac{1}{\Delta\text{PSNR}}
    \int_{\text{PSNR}_{\min}}^{\text{PSNR}_{\max}}
    \left[ \ln R_A(p) - \ln R_B(p) \right] dp
  \label{eq:bd-rate}
\end{equation}
A negative BD-Rate indicates that codec~$A$ requires fewer bits than
codec~$B$ at the same quality.

\begin{figure}[ht]
\centering
\begin{tikzpicture}[>=Stealth]

% --- Axes ---
\draw[->, thick] (0, 0) -- (11, 0);
\draw[->, thick] (0, 0) -- (0, 6.5);

% --- Axis titles ---
\node[font=\small] at (5.5, -0.85) {Rate [Mbps]};
\node[font=\small, rotate=90] at (-0.9, 3.25) {PSNR [dB]};

% --- X-axis labels ---
\foreach \x/\lab in {0/0, 2/0.5, 4/1.0, 6/1.5, 8/2.0, 10/2.5} {
  \draw (\x, -0.1) -- (\x, 0.1);
  \node[font=\tiny, anchor=north] at (\x, -0.15) {\lab};
}

% --- Y-axis labels ---
\foreach \y/\lab in {0/28, 1/30, 2/32, 3/34, 4/36, 5/38, 6/40} {
  \draw (-0.1, \y) -- (0.1, \y);
  \node[font=\tiny, anchor=east] at (-0.15, \y) {\lab};
}

% --- Grid ---
\foreach \y in {1,2,3,4,5,6} {
  \draw[gray!15] (0, \y) -- (10.5, \y);
}

% --- Shannon R(D) bound (leftmost curve) ---
\draw[very thick, green!60!black]
  (0.4, 0.5) .. controls (0.6, 2) and (0.8, 3.5) ..
  (1.2, 4.5) .. controls (1.8, 5.5) and (2.5, 6.0) ..
  (3.5, 6.3);
\node[font=\scriptsize, green!60!black, anchor=south west] at (2.5, 6.1)
  {$R(D)$};

% --- VVC curve ---
\draw[very thick, violet!70]
  (0.8, 0.3) .. controls (1.2, 2.0) and (1.6, 3.5) ..
  (2.2, 4.5) .. controls (3.0, 5.5) and (4.0, 6.0) ..
  (5.5, 6.3);
\node[font=\scriptsize, violet!70, anchor=south west] at (4.2, 6.1)
  {VVC};

% --- AV1 curve ---
\draw[very thick, orange!80!black]
  (1.0, 0.2) .. controls (1.5, 1.8) and (2.2, 3.3) ..
  (3.0, 4.3) .. controls (4.0, 5.3) and (5.2, 5.9) ..
  (7.0, 6.2);
\node[font=\scriptsize, orange!80!black, anchor=south] at (6.0, 6.1)
  {AV1};

% --- HEVC curve ---
\draw[very thick, teal!70]
  (1.2, 0.1) .. controls (1.8, 1.6) and (2.6, 3.1) ..
  (3.6, 4.1) .. controls (4.8, 5.1) and (6.2, 5.8) ..
  (8.0, 6.1);
\node[font=\scriptsize, teal!70, anchor=south] at (7.2, 6.0)
  {HEVC};

% --- VP9 curve ---
\draw[very thick, blue!70]
  (1.4, 0.0) .. controls (2.0, 1.4) and (3.0, 2.9) ..
  (4.2, 4.0) .. controls (5.5, 5.0) and (7.0, 5.7) ..
  (9.0, 6.0);
\node[font=\scriptsize, blue!70, anchor=south west] at (8.2, 5.9)
  {VP9};

% --- H.264 curve (rightmost) ---
\draw[very thick, red!70!black]
  (2.0, 0.0) .. controls (3.0, 1.2) and (4.2, 2.6) ..
  (5.5, 3.7) .. controls (7.0, 4.7) and (8.5, 5.4) ..
  (10.2, 5.8);
\node[font=\scriptsize, red!70!black, anchor=south west] at (9.5, 5.6)
  {H.264};

% --- 37 dB reference line ---
\draw[thin, gray, dashed] (0, 4.5) -- (10.5, 4.5);
\node[font=\tiny, gray, anchor=west] at (10.6, 4.5) {37\,dB};

% --- Operating points at 37 dB ---
\fill[red!70!black] (5.2, 4.5) circle (2.5pt);
\fill[blue!70] (3.9, 4.5) circle (2.5pt);
\fill[teal!70] (3.3, 4.5) circle (2.5pt);
\fill[orange!80!black] (2.7, 4.5) circle (2.5pt);
\fill[violet!70] (2.0, 4.5) circle (2.5pt);
\fill[green!60!black] (1.1, 4.5) circle (2.5pt);

% --- Annotation ---
\node[font=\scriptsize\itshape, text=black!60, align=center,
  anchor=north] at (5.5, -1.1)
  {At 37~dB PSNR (720p ``good'' quality), each generation
  requires fewer bits.\\
  The gap to the Shannon bound $R(D)$ narrows but never closes.};

\end{tikzpicture}
\caption{Rate-distortion curves for six codecs at 720p/30fps, with the
  Shannon $R(D)$ bound (green). At constant quality (37~dB dashed line),
  the rate progression from H.264 through VVC illustrates the BD-Rate
  savings in Table~\ref{tab:bd-rate}. This figure extends
  Figure~\ref{fig:rate-distortion} from Appendix~\ref{app:info-theory}
  with the full codec lineage.}
\label{fig:rd-curves-detailed}
\end{figure}

\begin{table}[ht]
\centering
\small
\renewcommand{\arraystretch}{1.4}
\begin{tabularx}{\textwidth}{lXXXX}
\toprule
\textbf{Codec} & \textbf{BD-Rate vs.\ H.264}
  & \textbf{$\Delta$ (720p)} & \textbf{Family} & \textbf{Year} \\
\midrule
VP8  & ${\sim}$0\%    & ${\sim}$2.1 & Google    & 2010 \\
VP9  & $-$35\%        & ${\sim}$0.8 & AOM       & 2013 \\
HEVC & $-$40\%        & ${\sim}$0.7 & MPEG/ITU  & 2013 \\
AV1  & $-$54\%        & ${\sim}$0.5 & AOM       & 2018 \\
VVC  & $-$63\%        & ${\sim}$0.3 & MPEG/ITU  & 2020 \\
\bottomrule
\end{tabularx}
\caption{BD-Rate savings and coding efficiency gap $\Delta$
  (eq.~\ref{eq:coding-gap}) for each codec relative to H.264 at 720p/30fps.
  The two families (MPEG/ITU and Google/AOM) achieve comparable efficiency
  within each generation; AV1 slightly trails HEVC, VVC leads AV1 by
  ${\sim}$20\%.}
\label{tab:bd-rate}
\end{table}

Table~\ref{tab:bd-rate} provides the DSP-grounded explanation for the codec
ordering in eq.~(\ref{eq:codec-ordering}) and the savings figures in
eqs.~(\ref{eq:vp9-savings})--(\ref{eq:av1-savings}): each generation adds
tools that reduce $\Delta_{\text{pred}}$, $\Delta_{\text{transform}}$, and
$\Delta_{\text{quant}}$ in (\ref{eq:delta-decomp}), with diminishing
returns as the Shannon bound is approached.

%% ---------------------------------------------------------------------------
\subsection{Scalable Video Coding and Simulcast}
\label{sec:svc-simulcast}

Multi-participant WebRTC sessions require adaptive quality: different
receivers have different bandwidths, and a single encode rate cannot serve
all. Two architectures address this.

\subsubsection{Simulcast}
\label{sec:simulcast}

Simulcast encodes the source at $N$ independent resolution/bitrate
combinations:
\begin{equation}
  R_{\text{simulcast}} = \sum_{i=1}^{N} R_i
  \label{eq:simulcast-rate}
\end{equation}
Each encoding is a fully independent bitstream. A Selective Forwarding
Unit (SFU) selects which stream to forward to each receiver based on
reported bandwidth. The overhead is substantial: a typical 3-layer
simulcast (720p/1.5~Mbps, 360p/500~kbps, 180p/150~kbps) requires
2.15~Mbps total---approximately 1.4$\times$ the cost of the highest
layer alone---with no inter-layer compression benefit.

\subsubsection{Scalable Video Coding (SVC)}
\label{sec:svc}

SVC produces a layered bitstream where a base layer $L_0$ is independently
decodable and enhancement layers $L_1, \ldots, L_K$ refine quality
progressively. Three scalability dimensions exist:

\textbf{Temporal scalability.} The base layer encodes at a reduced frame
rate (e.g., 7.5~fps); enhancement layers add intermediate frames:
\begin{equation}
  R_{\text{temporal}} = R_0 + \sum_{k=1}^{K} \Delta R_k
  \label{eq:temporal-svc}
\end{equation}
where $\Delta R_k \ll R_0$ because enhancement frames reference lower
layers and carry primarily motion-compensated residuals. An SFU drops
enhancement layers for bandwidth-constrained receivers, gracefully
degrading frame rate while preserving spatial quality.

\textbf{Spatial scalability.} The base layer encodes at reduced resolution;
enhancement layers add detail via inter-layer prediction:
\begin{equation}
  R_{\text{spatial}} = R_0 + \Delta R_{\text{upscale}}
  \label{eq:spatial-svc}
\end{equation}
where $\Delta R_{\text{upscale}}$ exploits correlation between the
upsampled base layer and the full-resolution source.

\textbf{Quality (SNR) scalability.} Refinement of quantization at the same
resolution, encoded as the difference between coarse and fine
reconstructions.

SVC achieves 20--30\% rate savings over simulcast for equivalent
adaptivity because enhancement layers exploit inter-layer prediction:
\begin{equation}
  R_{\text{SVC}} < R_{\text{simulcast}} \quad
  \text{(typically by 20--30\%)}
  \label{eq:svc-advantage}
\end{equation}

\begin{figure}[ht]
\centering
\begin{tikzpicture}[>=Stealth, scale=0.85, every node/.style={scale=0.85},
  layer/.style={draw, rounded corners=2pt, minimum width=1.8cm,
    minimum height=0.6cm, font=\scriptsize, align=center, thick}]

% --- Simulcast (left) ---
\node[font=\normalsize\bfseries] at (2, 5.8) {Simulcast};

% Three independent columns
\node[layer, fill=red!15] (s720) at (0.5, 4.5) {720p\\1.5 Mbps};
\node[layer, fill=blue!15] (s360) at (2.0, 4.5) {360p\\500 kbps};
\node[layer, fill=green!15] (s180) at (3.5, 4.5) {180p\\150 kbps};

\node[layer, fill=gray!10] (sfu1) at (2.0, 2.5) {SFU};

\draw[->, thick] (s720) -- (sfu1);
\draw[->, thick] (s360) -- (sfu1);
\draw[->, thick] (s180) -- (sfu1);

\node[font=\scriptsize, text=red!60!black] at (2.0, 1.5)
  {Total: 2.15 Mbps};
\node[font=\tiny\itshape, text=black!50] at (2.0, 1.0)
  {No inter-layer benefit};

% --- SVC (right) ---
\node[font=\normalsize\bfseries] at (8.5, 5.8) {SVC};

% Pyramid
\node[layer, fill=green!15, minimum width=2.0cm] (t0) at (8.5, 2.0)
  {$T_0$: Base\\7.5 fps};
\node[layer, fill=blue!15, minimum width=2.8cm] (t1) at (8.5, 3.2)
  {$T_1$: Enhancement\\15 fps};
\node[layer, fill=red!15, minimum width=3.6cm] (t2) at (8.5, 4.4)
  {$T_2$: Enhancement\\30 fps};

% Dependency arrows
\draw[->, thick, green!50!black] (t0) -- (t1);
\draw[->, thick, blue!50!black] (t1) -- (t2);

% Drop annotations
\draw[thick, dashed, red!40] (6.2, 3.8) -- (10.8, 3.8);
\node[font=\tiny\itshape, text=red!50!black, anchor=west] at (10.9, 3.8)
  {SFU can drop};

\node[font=\scriptsize, text=green!50!black] at (8.5, 1.0)
  {Total: ${\sim}$1.7 Mbps};
\node[font=\tiny\itshape, text=black!50] at (8.5, 0.5)
  {20--30\% savings via};
\node[font=\tiny\itshape, text=black!50] at (8.5, 0.1)
  {inter-layer prediction};

\end{tikzpicture}
\caption{Simulcast (left) encodes three independent bitstreams with no
  inter-layer compression. SVC (right) encodes a layered bitstream where
  temporal enhancement layers ($T_1$, $T_2$) reference the base layer
  ($T_0$), achieving 20--30\% savings. The SFU drops enhancement layers
  for bandwidth-constrained receivers.}
\label{fig:svc-architecture}
\end{figure}

Codec-specific SVC support relevant to WebRTC:

\begin{itemize}
\item \textbf{H.264 SVC (Annex~G).} Defines temporal, spatial, and quality
  scalability. Complex to implement; rarely used in WebRTC deployments.

\item \textbf{VP9 SVC.} Temporal scalability is widely deployed in WebRTC
  via libvpx (up to 3~temporal layers). Spatial SVC is defined but
  implementation-limited. VP9 temporal SVC is the primary mechanism
  enabling adaptive quality in Chrome/Firefox SFU-based calls.

\item \textbf{AV1 SVC.} Temporal scalability supported in both libaom and
  SVT-AV1. A key enabler for AV1 adoption in production WebRTC
  infrastructure, as SFUs require layered bitstreams to serve heterogeneous
  receivers.
\end{itemize}

Safari's lack of simulcast and SVC support in 2021 (Section~3.1) meant
that Safari users in multi-participant calls could not benefit from adaptive
quality---every receiver got the same stream regardless of bandwidth,
degrading the experience for all participants.

%% ---------------------------------------------------------------------------
\subsection{Hardware vs.\ Software Encoding}
\label{sec:hw-sw}

The main paper identifies hardware encoding as the deployment gate for
AV1 (Sections~8.5, 9.5). This section quantifies why.

\subsubsection{Encoding Complexity}

The computational cost of encoding a video frame is dominated by motion
estimation and RD-optimized mode decision. We express complexity as
operations per pixel:
\begin{equation}
  C_{\text{enc}} = C_{\text{ME}} + C_{\text{transform}}
    + C_{\text{quant}} + C_{\text{entropy}}
    + C_{\text{mode}}
  \label{eq:enc-complexity}
\end{equation}

Representative values for real-time presets:

\begin{itemize}
\item H.264 Constrained Baseline (``veryfast''): $C_{\text{enc}} \approx
  200$~ops/pixel
\item VP9 (speed~8, real-time): $C_{\text{enc}} \approx 600$~ops/pixel
\item AV1 (speed~10, real-time): $C_{\text{enc}} \approx 2000$~ops/pixel
\end{itemize}

The real-time constraint requires:
\begin{equation}
  C_{\text{enc}} \times W \times H \times \text{fps}
    \leq F_{\text{available}}
  \label{eq:realtime-constraint}
\end{equation}
where $F_{\text{available}}$ is the available computational throughput
(operations per second). For 720p/30fps AV1 at 2000~ops/pixel:
\begin{equation}
  2000 \times 921{,}600 \times 30 = 55.3 \times 10^9 \;\text{ops/s}
  \label{eq:av1-cost}
\end{equation}
This consumes approximately one full core of a modern desktop CPU. On
mobile devices with thermal and power constraints, software AV1 encoding
at 720p/30fps remains impractical without hardware acceleration.

\subsubsection{Hardware Encoder Architecture}

Hardware encoders implement a fixed subset of the codec's tool set in
dedicated silicon. The trade-off is quality for throughput: hardware
encoders typically produce output 10--20\% larger than software encoders
at equivalent visual quality because they implement fewer mode decision
candidates and simpler motion search:
\begin{equation}
  \Delta_{\text{HW}} \approx (1.1 \text{--} 1.2)
    \cdot \Delta_{\text{SW}}
  \label{eq:hw-quality-gap}
\end{equation}

This quality penalty is acceptable for real-time communication because:
(a)~the alternative---software encoding at a lower speed preset or reduced
resolution---incurs a larger quality penalty, and (b)~the freed CPU
capacity can be used for pre/post-processing, noise reduction, or
additional application logic.

\subsubsection{Hardware Encoder Availability}

\begin{table}[ht]
\centering
\small
\renewcommand{\arraystretch}{1.4}
\begin{tabularx}{\textwidth}{lXXXX}
\toprule
\textbf{Platform} & \textbf{H.264 HW} & \textbf{VP9 HW encode}
  & \textbf{AV1 HW encode} & \textbf{AV1 year} \\
\midrule
Intel (QSV) &
  Sandy Bridge+ (2011) &
  Kaby Lake+ (2017) &
  Arc / Meteor Lake+ &
  2022 \\
NVIDIA (NVENC) &
  Kepler+ (2012) &
  --- &
  Ada Lovelace+ &
  2022 \\
AMD (VCN) &
  VCN~1.0+ (2017) &
  --- &
  VCN~4.0+ (RDNA~3) &
  2023 \\
Apple (VideoToolbox) &
  A4+ (2010) &
  --- &
  M3+ &
  2023 \\
Qualcomm &
  SD~820+ (2016) &
  --- &
  SD~8 Gen~2+ &
  2023 \\
\bottomrule
\end{tabularx}
\caption{Hardware encoder availability timeline. H.264 hardware encoding
  has been ubiquitous for over a decade. AV1 hardware encoding began
  shipping in 2022--2023 across all major platforms; the main paper projects
  it becomes the practical WebRTC default by 2028--2030 as these chipsets
  reach market saturation.}
\label{tab:hw-encoders}
\end{table}

\excursus{Excursus: WebCodecs and hardware acceleration.}{When a WebCodecs
\texttt{VideoEncoder} is configured with \texttt{codec: "av01"}, the browser
selects hardware or software encoding based on platform capability and the
\texttt{hardwareAcceleration} preference (\texttt{"prefer-hardware"},
\texttt{"prefer-software"}, or \texttt{"no-preference"}). The application
receives encoded \texttt{EncodedVideoChunk} objects either way, but latency
and CPU cost differ dramatically: hardware encoding typically adds
$<\!1$~ms per frame and consumes near-zero CPU, while software AV1 at
real-time presets consumes an entire CPU core
(eq.~\ref{eq:av1-cost}). This abstraction is what makes the hardware
timeline in Table~\ref{tab:hw-encoders} directly relevant to the WebCodecs
adoption curve in Section~7.}

The hardware encoding timeline directly underpins the main paper's
projections: the ``AV1 becomes practical'' milestone (Section~8.5) and
``AV1 HW ubiquitous'' milestone (Section~9.5) depend on the chipset
generations in Table~\ref{tab:hw-encoders} reaching market saturation.

%% ---------------------------------------------------------------------------
\subsection{Implications for Real-Time Web Communication}
\label{sec:video-implications}

The DSP analysis of this appendix maps back to concrete implications for
the technologies surveyed in the main text.

\textbf{The hybrid coding framework is universal; innovation is
incremental.} Every codec from VP8 to VVC implements the same
prediction--transform--quantization--entropy framework
(Figure~\ref{fig:hybrid-encoder}). Generational improvements come from
expanding the tool set within this framework---more block sizes
(Table~\ref{tab:partitioning}), more prediction modes, more transform
types (Table~\ref{tab:transforms})---not from fundamentally new
architectures. This means the efficiency gap $\Delta$ approaches but never
reaches zero: each additional tool yields diminishing rate-distortion
improvement at increasing computational cost.

\textbf{Transform and prediction innovations explain the BD-Rate
progression.} The 35\% rate saving from H.264 to VP9 comes primarily from
larger blocks ($64\!\times\!64$ vs.\ $16\!\times\!16$) and adaptive
transforms (DCT$+$ADST vs.\ fixed DCT). The additional 30\% from VP9 to
AV1 comes from compound and warped motion, film grain synthesis, and
16-way transform type selection (Table~\ref{tab:bd-rate}). Each tool adds
encoding complexity proportional to its rate-distortion benefit, which is
why real-time encoders must aggressively prune the tool set.

\textbf{SVC is what makes multi-participant WebRTC work.} Without temporal
SVC, an SFU must forward full-rate streams to all participants regardless
of their bandwidth. With VP9/AV1 temporal SVC
(Section~\ref{sec:svc}), the SFU drops enhancement temporal layers for
bandwidth-constrained receivers, degrading frame rate gracefully while
preserving spatial quality. The 20--30\% savings over simulcast
(eq.~\ref{eq:svc-advantage}) translate directly to reduced server egress
bandwidth in large meetings.

\textbf{Hardware encoding is the deployment gate, not codec specification.}
AV1 was specified in 2018 but becomes a practical WebRTC default only
when hardware encoders are ubiquitous (${\sim}$2028--2030). The DSP tools
that make AV1 efficient---warped motion, OBMC, $128\!\times\!128$
superblocks with flexible partitioning---are precisely the tools that make
it computationally expensive in software (eq.~\ref{eq:av1-cost}). The
${\sim}$10--20\% quality penalty of hardware encoders
(eq.~\ref{eq:hw-quality-gap}) is the price of real-time feasibility, and
it is a price worth paying: $\Delta_{\text{AV1,HW}} \approx 0.6$ is still
far below $\Delta_{\text{VP9,SW}} \approx 0.8$.

\textbf{The convergence toward Shannon is real but asymptotic.}
Cross-referencing with Appendix~\ref{app:info-theory}: H.264's
$\Delta \approx 2.1$ means it operates at $3.1\times$ the theoretical
minimum rate. VP9's $\Delta \approx 0.8$ means $1.8\times$. AV1's
$\Delta \approx 0.5$ means $1.5\times$. VVC's $\Delta \approx 0.3$ means
$1.3\times$. Each generation captures more spatial, temporal, and
perceptual redundancy, but the returns are diminishing: the distance from
H.264 to VP9 ($-$35\%) exceeds VP9 to AV1 ($-$30\%), which exceeds AV1 to
VVC ($-$17\%). The next generation of efficiency gains for real-time web
communication will come not from codecs alone, but from the co-optimization
of transport (QUIC, WebTransport), application architecture (WebCodecs,
unbundled pipelines), and network infrastructure that the main text
describes.

%
% Requires in the preamble: amsmath, amssymb, tabularx, booktabs

\section{Video Compression Fundamentals}
\label{app:video-compression}

Appendix~\ref{app:info-theory} quantified the coding efficiency gap $\Delta$
(eq.~\ref{eq:coding-gap}) for each codec generation, treating codecs as black
boxes characterized by their distance from the Shannon rate-distortion bound.
This appendix opens those black boxes. We examine the signal processing and
coding techniques that determine $\Delta$: how video codecs exploit spatial,
temporal, and perceptual redundancy, what distinguishes one codec generation
from the next, and why computational complexity---not algorithmic
capability---remains the binding constraint for real-time web communication.

%% ---------------------------------------------------------------------------
\subsection{Why Video Is Compressible}
\label{sec:compressible}

A raw digital video signal carries far more data than information. Three
fundamental redundancies make compression ratios of 200:1 to 500:1 achievable
at perceptually acceptable quality.

\subsubsection{Spatial Redundancy}

Natural images exhibit strong local correlation. The autocorrelation function
for a typical image row can be modeled as a first-order Markov process:
\begin{equation}
  R_{xx}(k) = \sigma^2 \rho^{|k|}
  \label{eq:spatial-autocorr}
\end{equation}
where $\sigma^2$ is the pixel variance and $\rho \in [0.90, 0.97]$ for
natural content. The separable two-dimensional extension is:
\begin{equation}
  R_{xx}(k, l) = \sigma^2 \,\rho_h^{|k|}\, \rho_v^{|l|}
  \label{eq:spatial-autocorr-2d}
\end{equation}
The power spectral density---the Fourier transform of
(\ref{eq:spatial-autocorr-2d})---is concentrated at low spatial frequencies,
meaning most of the signal energy occupies a small fraction of the frequency
domain. This concentration is the foundation of transform coding
(Section~\ref{sec:transforms}).

\subsubsection{Temporal Redundancy}

Consecutive video frames are highly correlated. For typical conversational
video at 30~fps, the inter-frame correlation coefficient is:
\begin{equation}
  \text{Corr}(f_t, f_{t-1}) \approx 0.90\text{--}0.98
  \label{eq:temporal-corr}
\end{equation}
Most pixels change little between frames; only regions containing motion,
lighting changes, or scene cuts carry significant new information. Motion
estimation (Section~\ref{sec:motion}) exploits this by predicting each block
from a displaced version of a reference frame, transmitting only the
prediction residual.

\subsubsection{Perceptual Irrelevance}

The human visual system (HVS) has lower spatial acuity for chrominance than
luminance. This asymmetry motivates \emph{chroma subsampling}: representing
color difference channels at reduced resolution relative to the luma channel.

\excursus{Excursus: Why 4:2:0?}{All WebRTC codecs operate on YCbCr 4:2:0
video, where each chroma channel is sampled at half the horizontal and half
the vertical resolution of luma. This reduces the sample count by a factor of
$1 + 0.25 + 0.25 = 1.5$ relative to $1 + 1 + 1 = 3$ for full-resolution
4:4:4 color, a 2$\times$ saving before any coding begins. Perceptual
experiments confirm that 4:2:0 is essentially indistinguishable from 4:4:4 at
typical viewing distances for video conferencing---the HVS cannot resolve
the missing chroma detail.}

\subsubsection{The Compression Imperative}

A raw 720p/30fps 4:2:0 video stream at 8~bits per sample requires:
\begin{equation}
  R_{\text{raw}} = 1280 \times 720 \times 1.5 \times 8 \times 30
    = 331.8 \;\text{Mbps}
  \label{eq:raw-rate}
\end{equation}
This exceeds the 120~Mbps median link capacity cited in
Table~\ref{tab:inet}. The Shannon rate-distortion bound
(eq.~\ref{eq:rate-distortion}) yields approximately 0.7--1.5~Mbps at
perceptually acceptable distortion---a compression ratio of 220:1 to 475:1.
Video compression is not optional for real-time communication; it is a
prerequisite.

The first-order entropy of raw pixel values $H(X)$ greatly exceeds the
conditional entropy given neighboring samples:
\begin{equation}
  H(X) \gg H(X \mid X_{\text{neighbors}})
  \label{eq:conditional-entropy}
\end{equation}
This gap---the mutual information between a pixel and its
neighbors---represents the exploitable redundancy. The entire hybrid coding
framework (Section~\ref{sec:hybrid}) is designed to extract this mutual
information through prediction, leaving only the irreducible residual for
transform coding and quantization.

%% ---------------------------------------------------------------------------
\subsection{The Hybrid Coding Framework}
\label{sec:hybrid}

Every major video codec since H.261 (1988)---including every codec in the
WebRTC stack---is a variation on the \emph{block-based hybrid coding}
framework. The encoder partitions each frame into blocks, predicts each block
from previously coded data, transforms and quantizes the prediction residual,
and entropy codes the result.

\subsubsection{Encoder Signal Flow}

For a block at position $(x,y)$ in frame $f_t$, the encoder computes a
prediction $\hat{f}_{\text{pred}}(x,y)$ from either the current frame
(intra prediction) or a reference frame (inter prediction). The prediction
residual is:
\begin{equation}
  e(x,y) = f_t(x,y) - \hat{f}_{\text{pred}}(x,y)
  \label{eq:residual}
\end{equation}
The residual is transformed, quantized, and entropy coded to produce the
bitstream. The decoder reconstructs the frame as:
\begin{equation}
  \hat{f}(x,y) = \hat{e}(x,y) + \hat{f}_{\text{pred}}(x,y)
  \label{eq:reconstruction}
\end{equation}
where $\hat{e}$ is the inverse-transformed, dequantized residual.

\subsubsection{Rate-Distortion Optimization}

Every coding decision---block partition, prediction mode, motion vector,
transform type, quantization level---is made by minimizing a Lagrangian
rate-distortion cost:
\begin{equation}
  J(m) = D(m) + \lambda \cdot R(m)
  \label{eq:rd-cost}
\end{equation}
where $m \in \mathcal{M}$ is a coding mode, $D(m)$ is the distortion
(typically sum of squared errors), $R(m)$ is the bit cost, and $\lambda$ is a
Lagrange multiplier controlling the rate-distortion trade-off. The optimal
mode satisfies:
\begin{equation}
  m^* = \arg\min_{m \in \mathcal{M}} \left[ D(m) + \lambda \cdot R(m) \right]
  \label{eq:rd-opt}
\end{equation}
The Lagrange multiplier relates to the quantization parameter (QP) via an
empirically validated exponential relationship:
\begin{equation}
  \lambda = c \cdot 2^{(QP - 12)/3}
  \label{eq:lambda-qp}
\end{equation}
where $c$ is a codec-specific constant. This relationship ensures that
finer quantization (lower QP) assigns higher cost to bits, favoring modes
that reduce distortion, while coarser quantization (higher QP) favors
modes that reduce rate.

\begin{figure}[ht]
\centering
\begin{tikzpicture}[>=Stealth, scale=0.85, every node/.style={scale=0.85},
  block/.style={draw, rounded corners=3pt, minimum width=1.6cm,
    minimum height=0.8cm, font=\scriptsize, align=center, thick},
  arrow/.style={->, thick},
  label/.style={font=\tiny, text=black!60}]

% --- Input ---
\node[block, fill=blue!15] (input) at (0, 0) {Input\\frame $f_t$};

% --- Subtract ---
\node[draw, circle, inner sep=1.5pt, thick] (sub) at (2.5, 0) {$-$};

% --- Transform ---
\node[block, fill=orange!20] (transform) at (5, 0) {Transform\\$T$};

% --- Quantize ---
\node[block, fill=orange!20] (quant) at (7.5, 0) {Quantize\\$Q$};

% --- Entropy ---
\node[block, fill=purple!20] (entropy) at (10.5, 0) {Entropy\\coder};

% --- Bitstream ---
\node[font=\scriptsize, right=0.4cm of entropy] (bitstream) {Bitstream};

% --- Inverse Quantize ---
\node[block, fill=orange!20] (iquant) at (7.5, -2) {Inverse\\$Q^{-1}$};

% --- Inverse Transform ---
\node[block, fill=orange!20] (itransform) at (5, -2) {Inverse\\$T^{-1}$};

% --- Add ---
\node[draw, circle, inner sep=1.5pt, thick] (add) at (2.5, -2) {$+$};

% --- In-loop filter ---
\node[block, fill=red!15] (filter) at (2.5, -3.8) {In-loop\\filter};

% --- Reference buffer ---
\node[block, fill=red!15] (refbuf) at (5.5, -3.8) {Reference\\buffer};

% --- Mode decision ---
\node[block, fill=green!15] (mode) at (8.5, -3.8) {Mode\\decision};

% --- Prediction ---
\node[block, fill=green!15] (pred) at (5.5, -5.5) {Prediction\\(intra/inter)};

% --- Forward path ---
\draw[arrow] (input) -- (sub);
\draw[arrow] (sub) -- node[above, label] {$e$} (transform);
\draw[arrow] (transform) -- (quant);
\draw[arrow] (quant) -- (entropy);
\draw[arrow] (entropy) -- (bitstream);

% --- Reconstruction path ---
\draw[arrow] (quant) -- (iquant);
\draw[arrow] (iquant) -- (itransform);
\draw[arrow] (itransform) -- node[above, label] {$\hat{e}$} (add);

% --- Filter and buffer ---
\draw[arrow] (add) -- (filter);
\draw[arrow] (filter) -- (refbuf);

% --- Prediction feedback ---
\draw[arrow] (refbuf) -- (mode);
\draw[arrow] (mode) -- (pred);
\draw[arrow] (pred) -| (sub);
\draw[arrow] (pred) -| (add);

% --- Labels ---
\node[label, anchor=west] at (0, -5.5) {$\hat{f}_{\text{pred}}$};

\end{tikzpicture}
\caption{Block-based hybrid encoder architecture. Every codec from H.264
  through VVC follows this structure: predict each block from previously
  coded data (green), transform and quantize the residual (orange), entropy
  code the result (purple), then reconstruct and filter for future
  reference (red). The coding efficiency gap $\Delta$
  (eq.~\ref{eq:coding-gap}) arises from sub-optimal prediction, imperfect
  energy compaction, quantization beyond the rate-distortion bound, and
  entropy coding overhead.}
\label{fig:hybrid-encoder}
\end{figure}

The four sub-optimal components of Figure~\ref{fig:hybrid-encoder}---prediction,
transform, quantization, and entropy coding---each contribute to $\Delta$:
\begin{equation}
  \Delta \approx \Delta_{\text{pred}} + \Delta_{\text{transform}}
    + \Delta_{\text{quant}} + \Delta_{\text{entropy}}
  \label{eq:delta-decomp}
\end{equation}
Generational codec improvements reduce one or more of these terms. The codec
lineages in Sections~\ref{sec:mpeg-lineage}--\ref{sec:aom-lineage} trace
exactly which tools target which terms.

%% ---------------------------------------------------------------------------
\subsection{Transforms}
\label{sec:transforms}

The transform stage decorrelates the prediction residual, concentrating
energy into a few coefficients that can be efficiently quantized and coded.

\subsubsection{The KLT and Its DCT Approximation}
\label{sec:klt-dct}

The optimal decorrelating transform for a stationary source is the
\emph{Karhunen-Lo\`eve Transform} (KLT), whose basis vectors are the
eigenvectors of the input autocorrelation matrix $\mathbf{R}_{xx}$. For a
first-order Markov source with correlation $\rho$ (eq.~\ref{eq:spatial-autocorr}),
the KLT basis converges to the Type-II DCT as block size grows:
\begin{equation}
  C_k(n) = \alpha_k \cos\!\left[\frac{\pi k (2n+1)}{2N}\right],
  \quad k,n = 0, \ldots, N-1
  \label{eq:dct-ii}
\end{equation}
where $\alpha_0 = \sqrt{1/N}$ and $\alpha_k = \sqrt{2/N}$ for $k > 0$.

The energy compaction efficiency---the fraction of total energy captured by
the first $K$ of $N$ coefficients---is:
\begin{equation}
  \eta(K, N) = \frac{\sum_{k=0}^{K-1} \sigma_k^2}
    {\sum_{k=0}^{N-1} \sigma_k^2}
  \label{eq:energy-compaction}
\end{equation}
where $\sigma_k^2$ is the variance of the $k$-th transform coefficient.
For a typical intra residual with $\rho = 0.95$ and $N = 8$, the first 3
DCT coefficients capture approximately 95\% of the energy. This
concentration enables aggressive quantization of high-frequency coefficients
with minimal perceptual impact.

\subsubsection{Integer Transforms in Practice}
\label{sec:integer-transforms}

Practical codecs use integer approximations of the DCT to ensure bit-exact
encoder-decoder agreement. H.264 introduced a $4 \times 4$ integer
transform:
\begin{equation}
  \mathbf{H}_4 = \begin{pmatrix}
    1 &  1 &  1 &  1 \\
    2 &  1 & -1 & -2 \\
    1 & -1 & -1 &  1 \\
    1 & -2 &  2 & -1
  \end{pmatrix}
  \label{eq:h264-transform}
\end{equation}
The non-unity scaling factors are absorbed into the quantization step,
preserving exact integer arithmetic throughout the encode-decode loop.

Each codec generation expands the transform toolkit:

\begin{table}[ht]
\centering
\small
\renewcommand{\arraystretch}{1.4}
\begin{tabularx}{\textwidth}{lXXX}
\toprule
\textbf{Codec} & \textbf{Transform sizes} & \textbf{Transform types}
  & \textbf{Selection method} \\
\midrule
H.264 &
  $4\!\times\!4$, $8\!\times\!8$ &
  Integer DCT &
  Fixed per block size \\
VP9 &
  $4\!\times\!4$ to $32\!\times\!32$ &
  DCT, ADST &
  Row/column adaptive \\
HEVC &
  $4\!\times\!4$ to $32\!\times\!32$ &
  Integer DCT, DST ($4\!\times\!4$ intra) &
  Mode-dependent \\
AV1 &
  $4\!\times\!4$ to $64\!\times\!64$ &
  DCT, ADST, FlipADST, Identity &
  RD-optimized per block \\
VVC &
  $4\!\times\!4$ to $64\!\times\!64$ &
  DCT-II, DST-VII $+$ LFNST &
  RD-optimized $+$ secondary \\
\bottomrule
\end{tabularx}
\caption{Transform sizes and types across codec generations. AV1 applies
  four 1D transform types independently to rows and columns, yielding
  16~possible 2D combinations per block. VVC adds a non-separable secondary
  transform (LFNST) applied to low-frequency coefficients after the primary
  separable transform.}
\label{tab:transforms}
\end{table}

The Asymmetric Discrete Sine Transform (ADST) used in VP9 and AV1 is
particularly effective for intra-predicted residuals, which tend to have
larger magnitudes near the prediction boundary. AV1's FlipADST extends
this to blocks where the large residuals are on the opposite edge, while
the Identity transform (skip) handles blocks where the residual is already
well-distributed in the spatial domain.

The progression from fixed transforms (H.264) to RD-optimized multi-type
selection (AV1, VVC) directly reduces $\Delta_{\text{transform}}$ in
(\ref{eq:delta-decomp}): better energy compaction means fewer bits are
needed to represent the same residual at a given distortion level.

%% ---------------------------------------------------------------------------
\subsection{Motion Estimation and Inter Prediction}
\label{sec:motion}

Inter prediction exploits temporal redundancy by predicting blocks in the
current frame from previously coded reference frames. This is the single
largest source of compression gain---the prediction residual
(eq.~\ref{eq:residual}) typically has 10--20$\times$ less energy than the
original signal.

\subsubsection{Block Matching and Motion Vectors}
\label{sec:block-matching}

The encoder searches for the best-matching region in a reference frame by
minimizing a cost function over candidate displacement vectors
$\mathbf{d} = (d_x, d_y)$:
\begin{equation}
  \text{SAD}(\mathbf{d}) = \sum_{(x,y) \in B}
    \left| f_t(x,y) - f_{t-1}(x + d_x,\, y + d_y) \right|
  \label{eq:sad}
\end{equation}
The Sum of Absolute Differences (SAD) is the simplest distortion measure.
A better rate-distortion proxy is the Sum of Absolute Transformed
Differences (SATD), which applies the Hadamard transform before summation:
\begin{equation}
  \text{SATD}(\mathbf{d}) = \sum \left| \mathbf{H}
    \cdot \left( f_t - f_{t-1}(\mathbf{d}) \right) \right|
  \label{eq:satd}
\end{equation}
SATD approximates the post-transform energy and thus better predicts the
actual coded bit cost. Rate-constrained motion estimation augments the
distortion with a motion vector coding cost:
\begin{equation}
  \mathbf{d}^* = \arg\min_{\mathbf{d}} \left[
    \text{SATD}(\mathbf{d}) + \lambda_{\text{motion}} \cdot R(\mathbf{d})
  \right]
  \label{eq:rcme}
\end{equation}
where $R(\mathbf{d})$ is the bit cost of encoding the motion vector
relative to a predictor.

\subsubsection{Sub-Pixel Interpolation}
\label{sec:subpixel}

Motion in natural video rarely aligns to integer pixel positions. All modern
codecs perform sub-pixel motion compensation by interpolating reference
frame samples at fractional positions.

H.264 uses a 6-tap filter for half-pixel luma interpolation:
\begin{equation}
  h_{\frac{1}{2}} = \{1, -5, 20, 20, -5, 1\} / 32
  \label{eq:h264-halfpel}
\end{equation}
with quarter-pixel positions obtained by bilinear averaging of integer and
half-pixel samples. HEVC extends this to an 8-tap filter for half-pixel and
a 7-tap filter for quarter-pixel positions. AV1 introduces a switchable
filter bank with three 8-tap options---Regular (sharp), Smooth (low-pass),
and Sharp (high-frequency preserving)---selected per block via RD
optimization. VVC adds DMVR (Decoder-side Motion Vector Refinement) and
BDOF (Bi-Directional Optical Flow) for sub-pixel refinement at the decoder
without additional signaling bits.

The progression from H.264's fixed 6-tap filter to AV1's switchable bank
reduces $\Delta_{\text{pred}}$: better interpolation means smaller
prediction residuals, which require fewer bits.

\subsubsection{Advanced Inter Prediction Tools}
\label{sec:advanced-inter}

Beyond basic block matching, each codec generation adds tools to handle
complex motion and occlusion:

\begin{itemize}
\item \textbf{Multiple reference frames.} H.264 introduced up to 16
  reference frames; practical WebRTC encoders use 1--3 for latency reasons.
  More references improve prediction for periodic motion and uncovered
  backgrounds.

\item \textbf{Bi-prediction.} A weighted average of predictions from two
  reference frames, reducing noise and handling gradual scene changes.
  Available from H.264 onward (in B-frames) and in AV1's compound modes.

\item \textbf{Compound prediction (AV1).} Wedge masks and
  difference-weighted blending combine two predictions at the block level,
  handling occlusion boundaries where a single motion model fails.

\item \textbf{Overlapped Block Motion Compensation (AV1).} Blends
  predictions from neighboring blocks at motion boundaries, reducing
  blocking artifacts without post-processing.

\item \textbf{Warped motion (AV1).} An affine model with up to 6
  parameters handles rotation, zoom, and perspective---common in
  screen sharing and camera motion.

\item \textbf{VVC additions.} Geometric Partitioning Mode (non-rectangular
  block splits for motion boundaries), BDOF (optical flow refinement at the
  decoder), and DMVR (decoder-side MV refinement using bilateral matching).
\end{itemize}

\excursus{Excursus: Motion estimation in real-time encoding.}{Full-search
block matching has complexity $O(N^2 \cdot S^2)$ where $N$ is the block
dimension and $S$ is the search range in pixels. For a $64\!\times\!64$
block with $S = 64$, this is ${\sim}17$ million SAD operations per
block---far too expensive for real-time encoding. WebRTC encoders use
hierarchical search patterns (diamond, hexagon) that reduce complexity to
$O(N^2 \cdot \log S)$ with typically $< 5\%$ rate-distortion penalty
relative to full search. This is why real-time encoders use only a subset
of each codec's inter prediction tools: the full tool set is available
for offline encoding, but the real-time constraint forces aggressive
pruning of the mode decision search space.}

%% ---------------------------------------------------------------------------
\subsection{Quantization}
\label{sec:quantization}

Quantization is the sole lossy step in the hybrid coding pipeline. It maps
continuous-valued (or high-precision integer) transform coefficients to a
discrete set of reconstruction levels, introducing controlled information
loss to achieve the target bit rate.

\subsubsection{Scalar Quantization}

Uniform scalar quantization maps coefficient $c$ to a quantization index:
\begin{equation}
  q = \text{round}\!\left(\frac{c}{Q_{\text{step}}}\right), \quad
  \hat{c} = q \cdot Q_{\text{step}}
  \label{eq:uniform-quant}
\end{equation}
Modern codecs use \emph{dead-zone quantization}, where coefficients near
zero are mapped to zero more aggressively than uniform rounding predicts:
\begin{equation}
  q = \text{sign}(c) \cdot \left\lfloor
    \frac{|c| - f \cdot Q_{\text{step}}}{Q_{\text{step}}}
  \right\rfloor
  \label{eq:deadzone}
\end{equation}
where $f \in [0, 0.5)$ controls the dead-zone width. Dead-zone
quantization increases sparsity---the fraction of zero-valued
coefficients---which dramatically improves entropy coding efficiency.

\subsubsection{QP-to-Step-Size Mapping}

The quantization parameter (QP) controls the step size via an exponential
relationship. In the H.264/HEVC/VVC family:
\begin{equation}
  Q_{\text{step}} = Q_0 \cdot 2^{(QP - QP_0)/6}
  \label{eq:qp-mapping}
\end{equation}
A 6-unit increase in QP doubles the step size, approximately doubling the
bit rate. AV1 uses a similar but codec-specific mapping with separate QP
values for DC and AC coefficients, plus per-segment QP adjustment for
region-of-interest coding.

\subsubsection{Quantization Distortion}

For a single coefficient under uniform quantization at high rate, the
distortion is:
\begin{equation}
  D_q = \frac{Q_{\text{step}}^2}{12}
  \label{eq:quant-distortion}
\end{equation}
For a block of $N$ transform coefficients where $K$ survive quantization
(nonzero) and $N - K$ are quantized to zero:
\begin{equation}
  D_{\text{block}} \approx K \cdot \frac{Q_{\text{step}}^2}{12}
    + \sum_{k \notin K} c_k^2
  \label{eq:block-distortion}
\end{equation}
The second term represents the energy of coefficients eliminated by the
dead zone---information permanently discarded. The QP is thus the primary
control knob for navigating the rate-distortion curve in
Figure~\ref{fig:rate-distortion}. The WebCodecs \texttt{VideoEncoder} API
exposes this control via the \texttt{quantizer} configuration parameter,
giving JavaScript applications direct access to this fundamental
trade-off.

%% ---------------------------------------------------------------------------
\subsection{Entropy Coding}
\label{sec:entropy-coding}

After quantization, the surviving coefficients and all side information
(motion vectors, prediction modes, block partitions) must be losslessly
compressed. The entropy coder converts these symbols into a bitstream,
approaching the conditional entropy $H(X|\text{context})$ as closely as
possible. The gap between the entropy coder's output rate and the true
conditional entropy is $\Delta_{\text{entropy}}$ in
(\ref{eq:delta-decomp}).

\subsubsection{Context-Adaptive Binary Arithmetic Coding (CABAC)}
\label{sec:cabac}

CABAC, used in H.264, HEVC, and VVC, encodes each syntax element as a
sequence of binary decisions (bins). For each bin, the arithmetic coder
subdivides the current interval $[L, H)$ based on the estimated
probability $p$:
\begin{equation}
  [L, H) \to \begin{cases}
    [L,\; L + p \cdot (H - L)) & \text{if bin} = 0 \\
    [L + p \cdot (H - L),\; H) & \text{if bin} = 1
  \end{cases}
  \label{eq:arithmetic-coding}
\end{equation}
The theoretical output length approaches the entropy:
\begin{equation}
  \ell \to \sum_{i=1}^{n} -\log_2 p(b_i \mid \text{ctx}_i)
  \label{eq:cabac-length}
\end{equation}
CABAC achieves redundancy below 0.01 bits/symbol relative to the true
conditional entropy---near-optimal compression.

Context modeling is the key to CABAC's effectiveness. Each bin type has an
associated context, whose probability is updated after each observation via
a state-machine-based exponential moving average:
\begin{equation}
  p_{k+1} = p_k + \alpha(s_k) \cdot (b_k - p_k)
  \label{eq:cabac-update}
\end{equation}
where $\alpha(s_k)$ is a state-dependent adaptation rate, indexed by a
lookup table with 64 states (H.264) or 128 states (HEVC/VVC). This design
avoids floating-point arithmetic entirely.

CABAC also provides a \emph{bypass mode} for bins with approximately
uniform probability (e.g., sign bits, high-magnitude coefficient levels),
which skips the context lookup and probability estimation for throughput.

\subsubsection{Multi-Symbol Arithmetic Coding in AV1}
\label{sec:av1-entropy}

AV1 departs from binary arithmetic coding. Its entropy coder operates on
multi-symbol alphabets using CDF (Cumulative Distribution Function) tables:
\begin{equation}
  \text{Encode symbol } s: \quad
  [L, H) \to \left[L + \frac{\text{CDF}(s-1)}{M} \cdot (H-L),\;
    L + \frac{\text{CDF}(s)}{M} \cdot (H-L)\right)
  \label{eq:av1-entropy}
\end{equation}
where $M$ is the total frequency count. CDF tables are adapted after
each symbol:
\begin{equation}
  \text{CDF}_{k+1}(s) = \text{CDF}_k(s) + \left\lfloor
    \frac{\text{target}(s) - \text{CDF}_k(s)}{n}
  \right\rfloor
  \label{eq:cdf-update}
\end{equation}
where $n$ controls the adaptation rate. Multi-symbol coding is more
efficient than CABAC's binarization for syntax elements with moderate
alphabet sizes (e.g., intra prediction modes, transform types), as it
avoids the overhead of converting each symbol into a binary tree of
decisions.

\begin{table}[ht]
\centering
\small
\renewcommand{\arraystretch}{1.4}
\begin{tabularx}{\textwidth}{lXXXXX}
\toprule
\textbf{Feature} & \textbf{H.264} & \textbf{HEVC} & \textbf{VP9}
  & \textbf{AV1} & \textbf{VVC} \\
\midrule
Method &
  Binary AC &
  Binary AC &
  Boolean AC &
  Multi-sym AC &
  Binary AC \\
Contexts &
  ${\sim}400$ &
  ${\sim}200$ &
  ${\sim}900$ &
  ${\sim}300$ &
  ${\sim}700$ \\
Adaptation &
  64-state LUT &
  128-state LUT &
  Counter &
  CDF update &
  128-state LUT \\
Bypass mode &
  Yes &
  Yes &
  N/A &
  N/A &
  Yes \\
Parallelism &
  Slice &
  WPP / Tiles &
  Tiles &
  Tiles &
  WPP / Tiles \\
\bottomrule
\end{tabularx}
\caption{Entropy coding comparison across codecs. HEVC reduced context
  count relative to H.264 for throughput; VVC increased it again for
  improved compression. AV1's multi-symbol approach operates on full
  alphabets rather than binary decompositions.}
\label{tab:entropy-coding}
\end{table}

The frame entropy $\hat{H}(f_t)$ from eq.~(\ref{eq:frame-entropy}) is
ultimately determined by how well the entropy coder's probability model
matches the true symbol distribution. Both CABAC and AV1's multi-symbol
coder achieve near-zero entropy coding gaps; the remaining
$\Delta_{\text{entropy}}$ arises primarily from context model mismatch on
non-stationary content (scene changes, sudden motion onset).

%% ---------------------------------------------------------------------------
\subsection{In-Loop Filtering}
\label{sec:filtering}

Block-based coding introduces visible artifacts at block boundaries:
discontinuities from independent quantization of adjacent blocks, and
ringing from coefficient truncation. In-loop filters reduce these artifacts
\emph{before} the reconstructed frame enters the reference buffer, improving
prediction quality for subsequent frames and thus reducing
$\Delta_{\text{pred}}$.

\begin{itemize}
\item \textbf{Deblocking filter (all codecs).} Smooths discontinuities
  at block edges, with filter strength adapted to the quantization level
  and local gradient. Introduced in H.264; refined in each subsequent
  codec.

\item \textbf{Sample Adaptive Offset---SAO (HEVC).} Classifies
  reconstructed samples by intensity or edge direction and applies
  per-class offsets to reduce systematic bias introduced by quantization.

\item \textbf{Constrained Directional Enhancement Filter---CDEF (AV1).}
  A directional filter that identifies edge orientation per
  $8\!\times\!8$ block and applies filtering only along the edge,
  preserving sharpness while reducing ringing.

\item \textbf{Loop Restoration Filter (AV1).} A switchable per-tile
  filter choosing between a Wiener filter (FIR designed to minimize MSE
  between original and reconstruction) and a Self-Guided filter (non-linear
  denoising based on guided image filtering). This is AV1's most powerful
  quality enhancement tool.

\item \textbf{Adaptive Loop Filter---ALF (VVC).} A diamond-shaped filter
  with coefficients signaled per frame, applied adaptively per
  $4\!\times\!4$ block based on local classification. Subsumes HEVC's
  SAO functionality.
\end{itemize}

The evolution from simple deblocking (H.264) to multi-stage adaptive
filtering (AV1: deblocking $\to$ CDEF $\to$ loop restoration) exemplifies
how each codec generation reduces $\Delta_{\text{pred}}$ by improving
reference frame quality.

%% ---------------------------------------------------------------------------
\subsection{Codec Lineage: MPEG/ITU Family}
\label{sec:mpeg-lineage}

The MPEG/ITU standardization track has produced three codec generations
relevant to the 2021--2031 timeframe of this document.

\subsubsection{H.264/AVC (2003)}
\label{sec:h264}

H.264 defined the modern era of video compression and remains the most
widely deployed codec. Its key innovations over its predecessors (H.263,
MPEG-2):

\begin{itemize}
\item Variable block sizes: $16\!\times\!16$ macroblock partitioned into
  $16\!\times\!8$, $8\!\times\!16$, $8\!\times\!8$, $8\!\times\!4$,
  $4\!\times\!8$, and $4\!\times\!4$ sub-partitions for motion
  compensation.

\item Quarter-pixel motion compensation with 6-tap interpolation
  (eq.~\ref{eq:h264-halfpel}).

\item Context-adaptive intra prediction: 9~angular modes for
  $4\!\times\!4$ blocks, 4~modes for $16\!\times\!16$.

\item Dual entropy coding: CAVLC (simpler, faster) and CABAC
  (Section~\ref{sec:cabac}; 10--15\% better compression).

\item In-loop deblocking filter, mandatory for the decoder.

\item Multiple reference frames (up to 16) and weighted prediction.
\end{itemize}

Three profiles are relevant to WebRTC: \textbf{Constrained Baseline}
(no CABAC, no B-frames, lowest complexity---the mandatory WebRTC profile),
\textbf{Main} (CABAC, B-frames, ${\sim}$10\% better compression), and
\textbf{High} ($8\!\times\!8$ transform, adaptive quantization matrices,
${\sim}$15\% better).

At the 720p/30fps reference point, H.264 Constrained Baseline operates at
approximately 2.0--2.5~Mbps for ``good'' quality (PSNR ${\sim}$37~dB),
corresponding to $\Delta_{\text{H.264}} \approx 2.1$ from
Appendix~\ref{app:info-theory}. H.264 was the only codec Safari supported
in 2021 (Section~3.1), and its higher $\Delta$ is directly responsible
for the bandwidth disparity between Safari-only and VP9-capable sessions
documented in the main text.

\subsubsection{H.265/HEVC (2013)}
\label{sec:hevc}

HEVC replaced H.264's fixed $16\!\times\!16$ macroblock with a recursive
Coding Tree Unit (CTU) of up to $64\!\times\!64$ pixels, using quad-tree
partitioning down to $8\!\times\!8$. Key advances:

\begin{itemize}
\item 35 intra prediction angular directions (vs.\ H.264's 9).
\item Transform sizes from $4\!\times\!4$ to $32\!\times\!32$, with
  DST for $4\!\times\!4$ intra blocks.
\item Advanced Motion Vector Prediction (AMVP): two MVP candidates selected
  from spatial and temporal neighbors, replacing H.264's less structured
  motion vector prediction.
\item Merge mode: inherits the full motion information (MV, reference
  index, prediction direction) from a neighboring block, requiring zero
  additional motion bits.
\item Sample Adaptive Offset (SAO) in-loop filter.
\item Parallel processing: Wavefront Parallel Processing (WPP) and tiles.
\end{itemize}

HEVC achieves approximately 40\% rate reduction over H.264 at equivalent
PSNR. However, HEVC is absent from WebRTC entirely---not due to technical
limitations, but due to patent licensing.

\excursus{Excursus: Why no HEVC in WebRTC?}{HEVC's patent licensing is
fragmented across three separate pools (MPEG-LA, HEVC Advance, Velos Media)
plus independent licensors who are not members of any pool. The resulting
uncertainty about total royalty cost---estimated at \$0.20--\$0.40 per device
for some use cases---led Google, Mozilla, Cisco, and others to form the
Alliance for Open Media (AOM) in 2015 and develop AV1 as a royalty-free
alternative. This patent-driven fork is why the WebRTC ecosystem followed the
VP8$\to$VP9$\to$AV1 path rather than H.264$\to$HEVC$\to$VVC, and why the two
codec families must be understood in parallel.}

\subsubsection{H.266/VVC (2020)}
\label{sec:vvc}

VVC extends HEVC's quad-tree with binary and ternary splits, creating a
\emph{multi-type tree} (MTT) that enables asymmetric block shapes (e.g.,
$32\!\times\!8$, $8\!\times\!32$, $32\!\times\!4$). Additional
innovations:

\begin{itemize}
\item 67 intra prediction directions plus Matrix-based Intra Prediction
  (MIP), which applies a learned linear transform to neighboring samples.

\item Multiple Transform Selection (MTS): explicit signaling of DCT-II or
  DST-VII per block, plus the Low-Frequency Non-Separable Transform
  (LFNST)---a $16\!\times\!16$ or $8\!\times\!8$ non-separable transform
  applied to low-frequency coefficients after the primary separable
  transform.

\item Affine motion model: 4-parameter (translation $+$ rotation/scaling)
  or 6-parameter (full affine) motion, derived from control point MVs at
  block corners.

\item Decoder-side tools: DMVR refines motion vectors at the decoder using
  bilateral template matching; BDOF applies optical flow equations to
  refine bi-predicted samples, both without additional signaling bits.

\item Geometric Partitioning Mode: divides a block diagonally for separate
  motion compensation of each partition, targeting motion boundaries.

\item Adaptive Loop Filter (ALF): a diamond-shaped Wiener filter with
  coefficients signaled per frame and applied adaptively per
  $4\!\times\!4$ block.
\end{itemize}

VVC achieves approximately 30--50\% rate reduction over HEVC, placing it at
$\Delta_{\text{VVC}} \approx 0.3$. Its relevance to this document is
forward-looking: VVC may appear in WebCodecs codec registries by
${\sim}$2030 if licensing concerns are resolved, and its tool set
represents the current efficiency frontier against which AV1's performance
should be measured.

%% ---------------------------------------------------------------------------
\subsection{Codec Lineage: Google/AOM Family}
\label{sec:aom-lineage}

The royalty-free codec family dominates the WebRTC ecosystem. Its
development was driven by the patent licensing concerns described in the
HEVC excursus above.

\subsubsection{VP8 (2010)}
\label{sec:vp8}

VP8 was Google's first open-source video codec, released alongside the
WebM container. Its design is conservative:

\begin{itemize}
\item $16\!\times\!16$ macroblocks with $4\!\times\!4$ subblock partitions.
\item $4\!\times\!4$ Walsh-Hadamard Transform (WHT) for DC coefficients,
  $4\!\times\!4$ DCT for AC coefficients.
\item Boolean arithmetic coder with tree-structured syntax (single-bit
  alphabet).
\item 4 intra prediction modes for $16\!\times\!16$, 10 for
  $4\!\times\!4$ subblocks.
\item Simple loop filter with per-edge strength adaptation.
\item Single reference frame for WebRTC (plus a ``golden'' frame for
  screen content).
\end{itemize}

VP8's compression efficiency is roughly equivalent to H.264 Constrained
Baseline. It served as the baseline WebRTC codec---supported by Chrome and
Firefox since WebRTC~1.0---but has been superseded by VP9 in all modern
deployments.

\subsubsection{VP9 (2013)}
\label{sec:vp9}

VP9 introduced a superblock architecture and several tools that brought it
within range of HEVC:

\begin{itemize}
\item $64\!\times\!64$ superblocks with recursive quad-tree partitioning
  down to $4\!\times\!4$.
\item 10 intra prediction modes.
\item Transforms: DCT and ADST applied independently per row and column,
  at sizes from $4\!\times\!4$ to $32\!\times\!32$.
\item Switchable interpolation filters: 8-tap Regular, Smooth, and
  Sharp, selected per block.
\item Compound prediction: weighted average from two reference frames.
\item Segmentation: per-segment QP adjustment for region-of-interest
  coding.
\end{itemize}

VP9 achieves approximately 35\% rate reduction over H.264
(eq.~\ref{eq:vp9-savings}), placing it at $\Delta_{\text{VP9}} \approx 0.8$.
VP9 is the dominant WebRTC codec in Chrome and Firefox as of March~2026,
and its temporal SVC support (Section~\ref{sec:svc}) is the primary
mechanism enabling adaptive quality in multi-participant calls.

\subsubsection{AV1 (2018)}
\label{sec:av1}

AV1 represents the most ambitious single-generation improvement in the
royalty-free family. Its tool set approaches VVC in breadth while
maintaining royalty-free status:

\begin{itemize}
\item $128\!\times\!128$ superblocks with flexible partitioning (quad-tree
  plus rectangular splits) down to $4\!\times\!4$.

\item 56+ directional intra prediction modes plus non-directional modes
  (Paeth, Smooth variants, DC).

\item Full inter prediction toolkit: OBMC, warped motion (affine),
  compound modes (Wedge, Difference-weighted), inter-intra compound
  prediction.

\item 16 transform type combinations ($4$ 1D types $\times$ rows/columns),
  at sizes up to $64\!\times\!64$.

\item Film Grain Synthesis: the encoder characterizes noise/grain
  parameters and transmits them as metadata; the decoder re-synthesizes
  grain after reconstruction. This removes non-compressible high-frequency
  noise from the coding loop---a major bitrate saving for camera-captured
  content.

\item Multi-stage in-loop filtering: deblocking $\to$ CDEF $\to$ loop
  restoration (Wiener or Self-Guided), as described in
  Section~\ref{sec:filtering}.

\item Multi-symbol arithmetic coding with CDF-based context adaptation
  (Section~\ref{sec:av1-entropy}).
\end{itemize}

AV1 achieves approximately 30\% rate reduction over VP9
(eq.~\ref{eq:av1-savings}), a cumulative ${\sim}$54\% over H.264, placing
it at $\Delta_{\text{AV1}} \approx 0.5$ with current hardware encoders and
approaching 0.4 by 2030 as encoder implementations mature.

AV1's tools directly explain why its $\Delta$ is smaller than VP9's:
more transform types capture residual structure better (reducing
$\Delta_{\text{transform}}$), warped and compound motion reduce prediction
residual energy (reducing $\Delta_{\text{pred}}$), and film grain synthesis
removes non-compressible variance from the coding loop (effectively
reducing the source variance $\sigma^2$ that the codec must represent).

%% ---------------------------------------------------------------------------
\subsection{Block Partitioning: A Generational Comparison}
\label{sec:partitioning}

Block partitioning is the most architecturally consequential difference
between codec generations. Larger and more flexible partitions allow the
encoder to match block boundaries to content structure---smooth regions
get large blocks (few bits for partition overhead), detailed regions get
small blocks (better prediction accuracy).

\begin{figure}[ht]
\centering
\begin{tikzpicture}[>=Stealth, scale=0.55, every node/.style={scale=0.55}]

% --- H.264 panel ---
\begin{scope}[xshift=0cm]
  \node[font=\normalsize\bfseries, anchor=south] at (2, 8.4) {H.264};
  \node[font=\small, anchor=south] at (2, 7.8) {$16\!\times\!16$ max};
  % 4x4 grid of 16x16 macroblocks in a 64x64 region
  \draw[thick] (0,0) rectangle (4,4);
  \draw[thick] (0,0) rectangle (4,4);
  % macroblock grid
  \foreach \x in {0,1,2,3} {
    \foreach \y in {0,1,2,3} {
      \draw[gray] (\x, \y) rectangle (\x+1, \y+1);
    }
  }
  % show some sub-partitions
  \draw[blue!60, thick] (0,3) -- (0.5,3) -- (0.5,4) -- (0,4);
  \draw[blue!60, thick] (0.5,3) -- (0.5,4);
  \draw[blue!60, thick] (0,3.5) -- (1,3.5);
  % 8x8 in another
  \draw[blue!60, thick] (1,2.5) -- (2,2.5);
  \draw[blue!60, thick] (1.5,2) -- (1.5,3);
  % 4x4 in another
  \draw[red!60, thick] (2,0) -- (2,1) -- (3,1) -- (3,0);
  \draw[red!60, thick] (2.5,0) -- (2.5,0.5);
  \draw[red!60, thick] (2,0.5) -- (2.5,0.5);
  \draw[red!60, thick] (2.5,0.5) -- (2.5,1);
  \draw[red!60, thick] (2.5,0) -- (3,0) -- (3,0.5) -- (2.5,0.5);
  % Scale label
  \node[font=\small, anchor=north] at (2, -0.3) {64 pixels};
\end{scope}

% --- VP9/HEVC panel ---
\begin{scope}[xshift=6cm]
  \node[font=\normalsize\bfseries, anchor=south] at (4, 8.4) {VP9 / HEVC};
  \node[font=\small, anchor=south] at (4, 7.8) {$64\!\times\!64$ max};
  % One 64x64 superblock
  \draw[thick] (0,0) rectangle (8,8);
  % quad-tree splits
  \draw[thick] (4,0) -- (4,8);
  \draw[thick] (0,4) -- (8,4);
  % further splits in top-left quadrant
  \draw[blue!60, thick] (2,4) -- (2,8);
  \draw[blue!60, thick] (0,6) -- (4,6);
  % further split in top-left-top-left
  \draw[red!60, thick] (1,6) -- (1,8);
  \draw[red!60, thick] (0,7) -- (2,7);
  % split bottom-right quadrant
  \draw[blue!60, thick] (6,0) -- (6,4);
  \draw[blue!60, thick] (4,2) -- (8,2);
  % deeper in bottom-right-bottom-left
  \draw[red!60, thick] (5,0) -- (5,2);
  \draw[red!60, thick] (4,1) -- (6,1);
  % Scale label
  \node[font=\small, anchor=north] at (4, -0.3) {64 pixels};
\end{scope}

% --- AV1 panel ---
\begin{scope}[xshift=16cm]
  \node[font=\normalsize\bfseries, anchor=south] at (4, 8.4) {AV1};
  \node[font=\small, anchor=south] at (4, 7.8) {$128\!\times\!128$ max};
  % 128x128 shown as 8x8 grid (each cell = 16px)
  \draw[thick] (0,0) rectangle (8,8);
  % quad split: top half / bottom half
  \draw[thick] (4,0) -- (4,8);
  \draw[thick] (0,4) -- (8,4);
  % rectangular splits in top-left
  \draw[blue!60, thick] (0,6) -- (4,6);
  \draw[blue!60, thick] (2,4) -- (2,6);
  % asymmetric in bottom-left (4:1 rect)
  \draw[orange!80, thick] (0,3) -- (4,3);
  \draw[orange!80, thick] (0,2) -- (4,2);
  \draw[orange!80, thick] (0,1) -- (4,1);
  % quad + rect in top-right
  \draw[blue!60, thick] (6,4) -- (6,8);
  \draw[red!60, thick] (4,6) -- (6,6);
  \draw[red!60, thick] (5,4) -- (5,6);
  % bottom-right: large blocks
  \draw[blue!60, thick] (4,2) -- (8,2);
  % Scale label
  \node[font=\small, anchor=north] at (4, -0.3) {128 pixels};
\end{scope}

% --- VVC panel ---
\begin{scope}[xshift=26cm]
  \node[font=\normalsize\bfseries, anchor=south] at (4, 8.4) {VVC};
  \node[font=\small, anchor=south] at (4, 7.8) {$128\!\times\!128$ max};
  \draw[thick] (0,0) rectangle (8,8);
  % quad split
  \draw[thick] (4,0) -- (4,8);
  \draw[thick] (0,4) -- (8,4);
  % binary split (horizontal) in top-left
  \draw[blue!60, thick] (0,6) -- (4,6);
  % ternary split (vertical) in top-left-bottom
  \draw[orange!80, thick] (1,4) -- (1,6);
  \draw[orange!80, thick] (3,4) -- (3,6);
  % binary split (vertical) in top-right
  \draw[blue!60, thick] (6,4) -- (6,8);
  % ternary horizontal in top-right-left
  \draw[orange!80, thick] (4,5.33) -- (6,5.33);
  \draw[orange!80, thick] (4,6.67) -- (6,6.67);
  % bottom-left: binary horizontal
  \draw[blue!60, thick] (0,2) -- (4,2);
  % deeper ternary in bottom-left-top
  \draw[orange!80, thick] (1,2) -- (1,4);
  \draw[orange!80, thick] (3,2) -- (3,4);
  % bottom-right: mostly large
  \draw[blue!60, thick] (6,0) -- (6,4);
  \draw[blue!60, thick] (4,2) -- (8,2);
  % Scale label
  \node[font=\small, anchor=north] at (4, -0.3) {128 pixels};
\end{scope}

\end{tikzpicture}
\caption{Block partitioning comparison across codec generations. H.264 is
  limited to $16\!\times\!16$ macroblocks with fixed sub-partitions (left).
  VP9/HEVC use quad-tree partitioning within $64\!\times\!64$ superblocks
  (center-left). AV1 extends to $128\!\times\!128$ with quad-tree plus
  rectangular splits (center-right). VVC adds binary and ternary splits for
  maximum flexibility (right). Blue lines show first-level splits; red
  shows deeper partitioning; orange shows asymmetric/ternary splits.}
\label{fig:block-partitioning}
\end{figure}

\begin{table}[ht]
\centering
\small
\renewcommand{\arraystretch}{1.4}
\begin{tabularx}{\textwidth}{lXXXXX}
\toprule
\textbf{Feature} & \textbf{H.264} & \textbf{VP9} & \textbf{HEVC}
  & \textbf{AV1} & \textbf{VVC} \\
\midrule
Max block &
  $16\!\times\!16$ &
  $64\!\times\!64$ &
  $64\!\times\!64$ &
  $128\!\times\!128$ &
  $128\!\times\!128$ \\
Min block &
  $4\!\times\!4$ &
  $4\!\times\!4$ &
  $8\!\times\!8$ &
  $4\!\times\!4$ &
  $4\!\times\!4$ \\
Split types &
  Fixed partitions &
  Quad-tree &
  Quad-tree &
  Quad $+$ rect &
  Quad $+$ BT $+$ TT \\
Asymmetric &
  Limited &
  No &
  No &
  4:1 rectangles &
  Full (BT/TT) \\
\bottomrule
\end{tabularx}
\caption{Block partitioning capabilities by codec. BT = binary tree
  (horizontal or vertical), TT = ternary tree (1:2:1 split). Larger maximum
  blocks reduce overhead for homogeneous regions; asymmetric splits improve
  partitioning efficiency at content boundaries.}
\label{tab:partitioning}
\end{table}

%% ---------------------------------------------------------------------------
\subsection{Rate-Distortion Operating Points by Codec}
\label{sec:rd-comparison}

The codec-specific analyses of
Sections~\ref{sec:mpeg-lineage}--\ref{sec:aom-lineage} can be unified
through the Bj\o ntegaard Delta Rate (BD-Rate), the standard metric for
comparing codec efficiency. BD-Rate computes the average percentage rate
difference between two codecs at equivalent distortion by integrating the
difference between their rate-distortion curves fit as cubic polynomials in
PSNR-vs-log-rate space:
\begin{equation}
  \text{BD-Rate} = \frac{1}{\Delta\text{PSNR}}
    \int_{\text{PSNR}_{\min}}^{\text{PSNR}_{\max}}
    \left[ \ln R_A(p) - \ln R_B(p) \right] dp
  \label{eq:bd-rate}
\end{equation}
A negative BD-Rate indicates that codec~$A$ requires fewer bits than
codec~$B$ at the same quality.

\begin{figure}[ht]
\centering
\begin{tikzpicture}[>=Stealth]

% --- Axes ---
\draw[->, thick] (0, 0) -- (11, 0);
\draw[->, thick] (0, 0) -- (0, 6.5);

% --- Axis titles ---
\node[font=\small] at (5.5, -0.85) {Rate [Mbps]};
\node[font=\small, rotate=90] at (-0.9, 3.25) {PSNR [dB]};

% --- X-axis labels ---
\foreach \x/\lab in {0/0, 2/0.5, 4/1.0, 6/1.5, 8/2.0, 10/2.5} {
  \draw (\x, -0.1) -- (\x, 0.1);
  \node[font=\tiny, anchor=north] at (\x, -0.15) {\lab};
}

% --- Y-axis labels ---
\foreach \y/\lab in {0/28, 1/30, 2/32, 3/34, 4/36, 5/38, 6/40} {
  \draw (-0.1, \y) -- (0.1, \y);
  \node[font=\tiny, anchor=east] at (-0.15, \y) {\lab};
}

% --- Grid ---
\foreach \y in {1,2,3,4,5,6} {
  \draw[gray!15] (0, \y) -- (10.5, \y);
}

% --- Shannon R(D) bound (leftmost curve) ---
\draw[very thick, green!60!black]
  (0.4, 0.5) .. controls (0.6, 2) and (0.8, 3.5) ..
  (1.2, 4.5) .. controls (1.8, 5.5) and (2.5, 6.0) ..
  (3.5, 6.3);
\node[font=\scriptsize, green!60!black, anchor=south west] at (2.5, 6.1)
  {$R(D)$};

% --- VVC curve ---
\draw[very thick, violet!70]
  (0.8, 0.3) .. controls (1.2, 2.0) and (1.6, 3.5) ..
  (2.2, 4.5) .. controls (3.0, 5.5) and (4.0, 6.0) ..
  (5.5, 6.3);
\node[font=\scriptsize, violet!70, anchor=south west] at (4.2, 6.1)
  {VVC};

% --- AV1 curve ---
\draw[very thick, orange!80!black]
  (1.0, 0.2) .. controls (1.5, 1.8) and (2.2, 3.3) ..
  (3.0, 4.3) .. controls (4.0, 5.3) and (5.2, 5.9) ..
  (7.0, 6.2);
\node[font=\scriptsize, orange!80!black, anchor=south] at (6.0, 6.1)
  {AV1};

% --- HEVC curve ---
\draw[very thick, teal!70]
  (1.2, 0.1) .. controls (1.8, 1.6) and (2.6, 3.1) ..
  (3.6, 4.1) .. controls (4.8, 5.1) and (6.2, 5.8) ..
  (8.0, 6.1);
\node[font=\scriptsize, teal!70, anchor=south] at (7.2, 6.0)
  {HEVC};

% --- VP9 curve ---
\draw[very thick, blue!70]
  (1.4, 0.0) .. controls (2.0, 1.4) and (3.0, 2.9) ..
  (4.2, 4.0) .. controls (5.5, 5.0) and (7.0, 5.7) ..
  (9.0, 6.0);
\node[font=\scriptsize, blue!70, anchor=south west] at (8.2, 5.9)
  {VP9};

% --- H.264 curve (rightmost) ---
\draw[very thick, red!70!black]
  (2.0, 0.0) .. controls (3.0, 1.2) and (4.2, 2.6) ..
  (5.5, 3.7) .. controls (7.0, 4.7) and (8.5, 5.4) ..
  (10.2, 5.8);
\node[font=\scriptsize, red!70!black, anchor=south west] at (9.5, 5.6)
  {H.264};

% --- 37 dB reference line ---
\draw[thin, gray, dashed] (0, 4.5) -- (10.5, 4.5);
\node[font=\tiny, gray, anchor=west] at (10.6, 4.5) {37\,dB};

% --- Operating points at 37 dB ---
\fill[red!70!black] (5.2, 4.5) circle (2.5pt);
\fill[blue!70] (3.9, 4.5) circle (2.5pt);
\fill[teal!70] (3.3, 4.5) circle (2.5pt);
\fill[orange!80!black] (2.7, 4.5) circle (2.5pt);
\fill[violet!70] (2.0, 4.5) circle (2.5pt);
\fill[green!60!black] (1.1, 4.5) circle (2.5pt);

% --- Annotation ---
\node[font=\scriptsize\itshape, text=black!60, align=center,
  anchor=north] at (5.5, -1.1)
  {At 37~dB PSNR (720p ``good'' quality), each generation
  requires fewer bits.\\
  The gap to the Shannon bound $R(D)$ narrows but never closes.};

\end{tikzpicture}
\caption{Rate-distortion curves for six codecs at 720p/30fps, with the
  Shannon $R(D)$ bound (green). At constant quality (37~dB dashed line),
  the rate progression from H.264 through VVC illustrates the BD-Rate
  savings in Table~\ref{tab:bd-rate}. This figure extends
  Figure~\ref{fig:rate-distortion} from Appendix~\ref{app:info-theory}
  with the full codec lineage.}
\label{fig:rd-curves-detailed}
\end{figure}

\begin{table}[ht]
\centering
\small
\renewcommand{\arraystretch}{1.4}
\begin{tabularx}{\textwidth}{lXXXX}
\toprule
\textbf{Codec} & \textbf{BD-Rate vs.\ H.264}
  & \textbf{$\Delta$ (720p)} & \textbf{Family} & \textbf{Year} \\
\midrule
VP8  & ${\sim}$0\%    & ${\sim}$2.1 & Google    & 2010 \\
VP9  & $-$35\%        & ${\sim}$0.8 & AOM       & 2013 \\
HEVC & $-$40\%        & ${\sim}$0.7 & MPEG/ITU  & 2013 \\
AV1  & $-$54\%        & ${\sim}$0.5 & AOM       & 2018 \\
VVC  & $-$63\%        & ${\sim}$0.3 & MPEG/ITU  & 2020 \\
\bottomrule
\end{tabularx}
\caption{BD-Rate savings and coding efficiency gap $\Delta$
  (eq.~\ref{eq:coding-gap}) for each codec relative to H.264 at 720p/30fps.
  The two families (MPEG/ITU and Google/AOM) achieve comparable efficiency
  within each generation; AV1 slightly trails HEVC, VVC leads AV1 by
  ${\sim}$20\%.}
\label{tab:bd-rate}
\end{table}

Table~\ref{tab:bd-rate} provides the DSP-grounded explanation for the codec
ordering in eq.~(\ref{eq:codec-ordering}) and the savings figures in
eqs.~(\ref{eq:vp9-savings})--(\ref{eq:av1-savings}): each generation adds
tools that reduce $\Delta_{\text{pred}}$, $\Delta_{\text{transform}}$, and
$\Delta_{\text{quant}}$ in (\ref{eq:delta-decomp}), with diminishing
returns as the Shannon bound is approached.

%% ---------------------------------------------------------------------------
\subsection{Scalable Video Coding and Simulcast}
\label{sec:svc-simulcast}

Multi-participant WebRTC sessions require adaptive quality: different
receivers have different bandwidths, and a single encode rate cannot serve
all. Two architectures address this.

\subsubsection{Simulcast}
\label{sec:simulcast}

Simulcast encodes the source at $N$ independent resolution/bitrate
combinations:
\begin{equation}
  R_{\text{simulcast}} = \sum_{i=1}^{N} R_i
  \label{eq:simulcast-rate}
\end{equation}
Each encoding is a fully independent bitstream. A Selective Forwarding
Unit (SFU) selects which stream to forward to each receiver based on
reported bandwidth. The overhead is substantial: a typical 3-layer
simulcast (720p/1.5~Mbps, 360p/500~kbps, 180p/150~kbps) requires
2.15~Mbps total---approximately 1.4$\times$ the cost of the highest
layer alone---with no inter-layer compression benefit.

\subsubsection{Scalable Video Coding (SVC)}
\label{sec:svc}

SVC produces a layered bitstream where a base layer $L_0$ is independently
decodable and enhancement layers $L_1, \ldots, L_K$ refine quality
progressively. Three scalability dimensions exist:

\textbf{Temporal scalability.} The base layer encodes at a reduced frame
rate (e.g., 7.5~fps); enhancement layers add intermediate frames:
\begin{equation}
  R_{\text{temporal}} = R_0 + \sum_{k=1}^{K} \Delta R_k
  \label{eq:temporal-svc}
\end{equation}
where $\Delta R_k \ll R_0$ because enhancement frames reference lower
layers and carry primarily motion-compensated residuals. An SFU drops
enhancement layers for bandwidth-constrained receivers, gracefully
degrading frame rate while preserving spatial quality.

\textbf{Spatial scalability.} The base layer encodes at reduced resolution;
enhancement layers add detail via inter-layer prediction:
\begin{equation}
  R_{\text{spatial}} = R_0 + \Delta R_{\text{upscale}}
  \label{eq:spatial-svc}
\end{equation}
where $\Delta R_{\text{upscale}}$ exploits correlation between the
upsampled base layer and the full-resolution source.

\textbf{Quality (SNR) scalability.} Refinement of quantization at the same
resolution, encoded as the difference between coarse and fine
reconstructions.

SVC achieves 20--30\% rate savings over simulcast for equivalent
adaptivity because enhancement layers exploit inter-layer prediction:
\begin{equation}
  R_{\text{SVC}} < R_{\text{simulcast}} \quad
  \text{(typically by 20--30\%)}
  \label{eq:svc-advantage}
\end{equation}

\begin{figure}[ht]
\centering
\begin{tikzpicture}[>=Stealth, scale=0.85, every node/.style={scale=0.85},
  layer/.style={draw, rounded corners=2pt, minimum width=1.8cm,
    minimum height=0.6cm, font=\scriptsize, align=center, thick}]

% --- Simulcast (left) ---
\node[font=\normalsize\bfseries] at (2, 5.8) {Simulcast};

% Three independent columns
\node[layer, fill=red!15] (s720) at (0.5, 4.5) {720p\\1.5 Mbps};
\node[layer, fill=blue!15] (s360) at (2.0, 4.5) {360p\\500 kbps};
\node[layer, fill=green!15] (s180) at (3.5, 4.5) {180p\\150 kbps};

\node[layer, fill=gray!10] (sfu1) at (2.0, 2.5) {SFU};

\draw[->, thick] (s720) -- (sfu1);
\draw[->, thick] (s360) -- (sfu1);
\draw[->, thick] (s180) -- (sfu1);

\node[font=\scriptsize, text=red!60!black] at (2.0, 1.5)
  {Total: 2.15 Mbps};
\node[font=\tiny\itshape, text=black!50] at (2.0, 1.0)
  {No inter-layer benefit};

% --- SVC (right) ---
\node[font=\normalsize\bfseries] at (8.5, 5.8) {SVC};

% Pyramid
\node[layer, fill=green!15, minimum width=2.0cm] (t0) at (8.5, 2.0)
  {$T_0$: Base\\7.5 fps};
\node[layer, fill=blue!15, minimum width=2.8cm] (t1) at (8.5, 3.2)
  {$T_1$: Enhancement\\15 fps};
\node[layer, fill=red!15, minimum width=3.6cm] (t2) at (8.5, 4.4)
  {$T_2$: Enhancement\\30 fps};

% Dependency arrows
\draw[->, thick, green!50!black] (t0) -- (t1);
\draw[->, thick, blue!50!black] (t1) -- (t2);

% Drop annotations
\draw[thick, dashed, red!40] (6.2, 3.8) -- (10.8, 3.8);
\node[font=\tiny\itshape, text=red!50!black, anchor=west] at (10.9, 3.8)
  {SFU can drop};

\node[font=\scriptsize, text=green!50!black] at (8.5, 1.0)
  {Total: ${\sim}$1.7 Mbps};
\node[font=\tiny\itshape, text=black!50] at (8.5, 0.5)
  {20--30\% savings via};
\node[font=\tiny\itshape, text=black!50] at (8.5, 0.1)
  {inter-layer prediction};

\end{tikzpicture}
\caption{Simulcast (left) encodes three independent bitstreams with no
  inter-layer compression. SVC (right) encodes a layered bitstream where
  temporal enhancement layers ($T_1$, $T_2$) reference the base layer
  ($T_0$), achieving 20--30\% savings. The SFU drops enhancement layers
  for bandwidth-constrained receivers.}
\label{fig:svc-architecture}
\end{figure}

Codec-specific SVC support relevant to WebRTC:

\begin{itemize}
\item \textbf{H.264 SVC (Annex~G).} Defines temporal, spatial, and quality
  scalability. Complex to implement; rarely used in WebRTC deployments.

\item \textbf{VP9 SVC.} Temporal scalability is widely deployed in WebRTC
  via libvpx (up to 3~temporal layers). Spatial SVC is defined but
  implementation-limited. VP9 temporal SVC is the primary mechanism
  enabling adaptive quality in Chrome/Firefox SFU-based calls.

\item \textbf{AV1 SVC.} Temporal scalability supported in both libaom and
  SVT-AV1. A key enabler for AV1 adoption in production WebRTC
  infrastructure, as SFUs require layered bitstreams to serve heterogeneous
  receivers.
\end{itemize}

Safari's lack of simulcast and SVC support in 2021 (Section~3.1) meant
that Safari users in multi-participant calls could not benefit from adaptive
quality---every receiver got the same stream regardless of bandwidth,
degrading the experience for all participants.

%% ---------------------------------------------------------------------------
\subsection{Hardware vs.\ Software Encoding}
\label{sec:hw-sw}

The main paper identifies hardware encoding as the deployment gate for
AV1 (Sections~8.5, 9.5). This section quantifies why.

\subsubsection{Encoding Complexity}

The computational cost of encoding a video frame is dominated by motion
estimation and RD-optimized mode decision. We express complexity as
operations per pixel:
\begin{equation}
  C_{\text{enc}} = C_{\text{ME}} + C_{\text{transform}}
    + C_{\text{quant}} + C_{\text{entropy}}
    + C_{\text{mode}}
  \label{eq:enc-complexity}
\end{equation}

Representative values for real-time presets:

\begin{itemize}
\item H.264 Constrained Baseline (``veryfast''): $C_{\text{enc}} \approx
  200$~ops/pixel
\item VP9 (speed~8, real-time): $C_{\text{enc}} \approx 600$~ops/pixel
\item AV1 (speed~10, real-time): $C_{\text{enc}} \approx 2000$~ops/pixel
\end{itemize}

The real-time constraint requires:
\begin{equation}
  C_{\text{enc}} \times W \times H \times \text{fps}
    \leq F_{\text{available}}
  \label{eq:realtime-constraint}
\end{equation}
where $F_{\text{available}}$ is the available computational throughput
(operations per second). For 720p/30fps AV1 at 2000~ops/pixel:
\begin{equation}
  2000 \times 921{,}600 \times 30 = 55.3 \times 10^9 \;\text{ops/s}
  \label{eq:av1-cost}
\end{equation}
This consumes approximately one full core of a modern desktop CPU. On
mobile devices with thermal and power constraints, software AV1 encoding
at 720p/30fps remains impractical without hardware acceleration.

\subsubsection{Hardware Encoder Architecture}

Hardware encoders implement a fixed subset of the codec's tool set in
dedicated silicon. The trade-off is quality for throughput: hardware
encoders typically produce output 10--20\% larger than software encoders
at equivalent visual quality because they implement fewer mode decision
candidates and simpler motion search:
\begin{equation}
  \Delta_{\text{HW}} \approx (1.1 \text{--} 1.2)
    \cdot \Delta_{\text{SW}}
  \label{eq:hw-quality-gap}
\end{equation}

This quality penalty is acceptable for real-time communication because:
(a)~the alternative---software encoding at a lower speed preset or reduced
resolution---incurs a larger quality penalty, and (b)~the freed CPU
capacity can be used for pre/post-processing, noise reduction, or
additional application logic.

\subsubsection{Hardware Encoder Availability}

\begin{table}[ht]
\centering
\small
\renewcommand{\arraystretch}{1.4}
\begin{tabularx}{\textwidth}{lXXXX}
\toprule
\textbf{Platform} & \textbf{H.264 HW} & \textbf{VP9 HW encode}
  & \textbf{AV1 HW encode} & \textbf{AV1 year} \\
\midrule
Intel (QSV) &
  Sandy Bridge+ (2011) &
  Kaby Lake+ (2017) &
  Arc / Meteor Lake+ &
  2022 \\
NVIDIA (NVENC) &
  Kepler+ (2012) &
  --- &
  Ada Lovelace+ &
  2022 \\
AMD (VCN) &
  VCN~1.0+ (2017) &
  --- &
  VCN~4.0+ (RDNA~3) &
  2023 \\
Apple (VideoToolbox) &
  A4+ (2010) &
  --- &
  M3+ &
  2023 \\
Qualcomm &
  SD~820+ (2016) &
  --- &
  SD~8 Gen~2+ &
  2023 \\
\bottomrule
\end{tabularx}
\caption{Hardware encoder availability timeline. H.264 hardware encoding
  has been ubiquitous for over a decade. AV1 hardware encoding began
  shipping in 2022--2023 across all major platforms; the main paper projects
  it becomes the practical WebRTC default by 2028--2030 as these chipsets
  reach market saturation.}
\label{tab:hw-encoders}
\end{table}

\excursus{Excursus: WebCodecs and hardware acceleration.}{When a WebCodecs
\texttt{VideoEncoder} is configured with \texttt{codec: "av01"}, the browser
selects hardware or software encoding based on platform capability and the
\texttt{hardwareAcceleration} preference (\texttt{"prefer-hardware"},
\texttt{"prefer-software"}, or \texttt{"no-preference"}). The application
receives encoded \texttt{EncodedVideoChunk} objects either way, but latency
and CPU cost differ dramatically: hardware encoding typically adds
$<\!1$~ms per frame and consumes near-zero CPU, while software AV1 at
real-time presets consumes an entire CPU core
(eq.~\ref{eq:av1-cost}). This abstraction is what makes the hardware
timeline in Table~\ref{tab:hw-encoders} directly relevant to the WebCodecs
adoption curve in Section~7.}

The hardware encoding timeline directly underpins the main paper's
projections: the ``AV1 becomes practical'' milestone (Section~8.5) and
``AV1 HW ubiquitous'' milestone (Section~9.5) depend on the chipset
generations in Table~\ref{tab:hw-encoders} reaching market saturation.

%% ---------------------------------------------------------------------------
\subsection{Implications for Real-Time Web Communication}
\label{sec:video-implications}

The DSP analysis of this appendix maps back to concrete implications for
the technologies surveyed in the main text.

\textbf{The hybrid coding framework is universal; innovation is
incremental.} Every codec from VP8 to VVC implements the same
prediction--transform--quantization--entropy framework
(Figure~\ref{fig:hybrid-encoder}). Generational improvements come from
expanding the tool set within this framework---more block sizes
(Table~\ref{tab:partitioning}), more prediction modes, more transform
types (Table~\ref{tab:transforms})---not from fundamentally new
architectures. This means the efficiency gap $\Delta$ approaches but never
reaches zero: each additional tool yields diminishing rate-distortion
improvement at increasing computational cost.

\textbf{Transform and prediction innovations explain the BD-Rate
progression.} The 35\% rate saving from H.264 to VP9 comes primarily from
larger blocks ($64\!\times\!64$ vs.\ $16\!\times\!16$) and adaptive
transforms (DCT$+$ADST vs.\ fixed DCT). The additional 30\% from VP9 to
AV1 comes from compound and warped motion, film grain synthesis, and
16-way transform type selection (Table~\ref{tab:bd-rate}). Each tool adds
encoding complexity proportional to its rate-distortion benefit, which is
why real-time encoders must aggressively prune the tool set.

\textbf{SVC is what makes multi-participant WebRTC work.} Without temporal
SVC, an SFU must forward full-rate streams to all participants regardless
of their bandwidth. With VP9/AV1 temporal SVC
(Section~\ref{sec:svc}), the SFU drops enhancement temporal layers for
bandwidth-constrained receivers, degrading frame rate gracefully while
preserving spatial quality. The 20--30\% savings over simulcast
(eq.~\ref{eq:svc-advantage}) translate directly to reduced server egress
bandwidth in large meetings.

\textbf{Hardware encoding is the deployment gate, not codec specification.}
AV1 was specified in 2018 but becomes a practical WebRTC default only
when hardware encoders are ubiquitous (${\sim}$2028--2030). The DSP tools
that make AV1 efficient---warped motion, OBMC, $128\!\times\!128$
superblocks with flexible partitioning---are precisely the tools that make
it computationally expensive in software (eq.~\ref{eq:av1-cost}). The
${\sim}$10--20\% quality penalty of hardware encoders
(eq.~\ref{eq:hw-quality-gap}) is the price of real-time feasibility, and
it is a price worth paying: $\Delta_{\text{AV1,HW}} \approx 0.6$ is still
far below $\Delta_{\text{VP9,SW}} \approx 0.8$.

\textbf{The convergence toward Shannon is real but asymptotic.}
Cross-referencing with Appendix~\ref{app:info-theory}: H.264's
$\Delta \approx 2.1$ means it operates at $3.1\times$ the theoretical
minimum rate. VP9's $\Delta \approx 0.8$ means $1.8\times$. AV1's
$\Delta \approx 0.5$ means $1.5\times$. VVC's $\Delta \approx 0.3$ means
$1.3\times$. Each generation captures more spatial, temporal, and
perceptual redundancy, but the returns are diminishing: the distance from
H.264 to VP9 ($-$35\%) exceeds VP9 to AV1 ($-$30\%), which exceeds AV1 to
VVC ($-$17\%). The next generation of efficiency gains for real-time web
communication will come not from codecs alone, but from the co-optimization
of transport (QUIC, WebTransport), application architecture (WebCodecs,
unbundled pipelines), and network infrastructure that the main text
describes.

%
% Requires in the preamble: amsmath, amssymb, tabularx, booktabs

\section{Video Compression Fundamentals}
\label{app:video-compression}

Appendix~\ref{app:info-theory} quantified the coding efficiency gap $\Delta$
(eq.~\ref{eq:coding-gap}) for each codec generation, treating codecs as black
boxes characterized by their distance from the Shannon rate-distortion bound.
This appendix opens those black boxes. We examine the signal processing and
coding techniques that determine $\Delta$: how video codecs exploit spatial,
temporal, and perceptual redundancy, what distinguishes one codec generation
from the next, and why computational complexity---not algorithmic
capability---remains the binding constraint for real-time web communication.

%% ---------------------------------------------------------------------------
\subsection{Why Video Is Compressible}
\label{sec:compressible}

A raw digital video signal carries far more data than information. Three
fundamental redundancies make compression ratios of 200:1 to 500:1 achievable
at perceptually acceptable quality.

\subsubsection{Spatial Redundancy}

Natural images exhibit strong local correlation. The autocorrelation function
for a typical image row can be modeled as a first-order Markov process:
\begin{equation}
  R_{xx}(k) = \sigma^2 \rho^{|k|}
  \label{eq:spatial-autocorr}
\end{equation}
where $\sigma^2$ is the pixel variance and $\rho \in [0.90, 0.97]$ for
natural content. The separable two-dimensional extension is:
\begin{equation}
  R_{xx}(k, l) = \sigma^2 \,\rho_h^{|k|}\, \rho_v^{|l|}
  \label{eq:spatial-autocorr-2d}
\end{equation}
The power spectral density---the Fourier transform of
(\ref{eq:spatial-autocorr-2d})---is concentrated at low spatial frequencies,
meaning most of the signal energy occupies a small fraction of the frequency
domain. This concentration is the foundation of transform coding
(Section~\ref{sec:transforms}).

\subsubsection{Temporal Redundancy}

Consecutive video frames are highly correlated. For typical conversational
video at 30~fps, the inter-frame correlation coefficient is:
\begin{equation}
  \text{Corr}(f_t, f_{t-1}) \approx 0.90\text{--}0.98
  \label{eq:temporal-corr}
\end{equation}
Most pixels change little between frames; only regions containing motion,
lighting changes, or scene cuts carry significant new information. Motion
estimation (Section~\ref{sec:motion}) exploits this by predicting each block
from a displaced version of a reference frame, transmitting only the
prediction residual.

\subsubsection{Perceptual Irrelevance}

The human visual system (HVS) has lower spatial acuity for chrominance than
luminance. This asymmetry motivates \emph{chroma subsampling}: representing
color difference channels at reduced resolution relative to the luma channel.

\excursus{Excursus: Why 4:2:0?}{All WebRTC codecs operate on YCbCr 4:2:0
video, where each chroma channel is sampled at half the horizontal and half
the vertical resolution of luma. This reduces the sample count by a factor of
$1 + 0.25 + 0.25 = 1.5$ relative to $1 + 1 + 1 = 3$ for full-resolution
4:4:4 color, a 2$\times$ saving before any coding begins. Perceptual
experiments confirm that 4:2:0 is essentially indistinguishable from 4:4:4 at
typical viewing distances for video conferencing---the HVS cannot resolve
the missing chroma detail.}

\subsubsection{The Compression Imperative}

A raw 720p/30fps 4:2:0 video stream at 8~bits per sample requires:
\begin{equation}
  R_{\text{raw}} = 1280 \times 720 \times 1.5 \times 8 \times 30
    = 331.8 \;\text{Mbps}
  \label{eq:raw-rate}
\end{equation}
This exceeds the 120~Mbps median link capacity cited in
Table~\ref{tab:inet}. The Shannon rate-distortion bound
(eq.~\ref{eq:rate-distortion}) yields approximately 0.7--1.5~Mbps at
perceptually acceptable distortion---a compression ratio of 220:1 to 475:1.
Video compression is not optional for real-time communication; it is a
prerequisite.

The first-order entropy of raw pixel values $H(X)$ greatly exceeds the
conditional entropy given neighboring samples:
\begin{equation}
  H(X) \gg H(X \mid X_{\text{neighbors}})
  \label{eq:conditional-entropy}
\end{equation}
This gap---the mutual information between a pixel and its
neighbors---represents the exploitable redundancy. The entire hybrid coding
framework (Section~\ref{sec:hybrid}) is designed to extract this mutual
information through prediction, leaving only the irreducible residual for
transform coding and quantization.

%% ---------------------------------------------------------------------------
\subsection{The Hybrid Coding Framework}
\label{sec:hybrid}

Every major video codec since H.261 (1988)---including every codec in the
WebRTC stack---is a variation on the \emph{block-based hybrid coding}
framework. The encoder partitions each frame into blocks, predicts each block
from previously coded data, transforms and quantizes the prediction residual,
and entropy codes the result.

\subsubsection{Encoder Signal Flow}

For a block at position $(x,y)$ in frame $f_t$, the encoder computes a
prediction $\hat{f}_{\text{pred}}(x,y)$ from either the current frame
(intra prediction) or a reference frame (inter prediction). The prediction
residual is:
\begin{equation}
  e(x,y) = f_t(x,y) - \hat{f}_{\text{pred}}(x,y)
  \label{eq:residual}
\end{equation}
The residual is transformed, quantized, and entropy coded to produce the
bitstream. The decoder reconstructs the frame as:
\begin{equation}
  \hat{f}(x,y) = \hat{e}(x,y) + \hat{f}_{\text{pred}}(x,y)
  \label{eq:reconstruction}
\end{equation}
where $\hat{e}$ is the inverse-transformed, dequantized residual.

\subsubsection{Rate-Distortion Optimization}

Every coding decision---block partition, prediction mode, motion vector,
transform type, quantization level---is made by minimizing a Lagrangian
rate-distortion cost:
\begin{equation}
  J(m) = D(m) + \lambda \cdot R(m)
  \label{eq:rd-cost}
\end{equation}
where $m \in \mathcal{M}$ is a coding mode, $D(m)$ is the distortion
(typically sum of squared errors), $R(m)$ is the bit cost, and $\lambda$ is a
Lagrange multiplier controlling the rate-distortion trade-off. The optimal
mode satisfies:
\begin{equation}
  m^* = \arg\min_{m \in \mathcal{M}} \left[ D(m) + \lambda \cdot R(m) \right]
  \label{eq:rd-opt}
\end{equation}
The Lagrange multiplier relates to the quantization parameter (QP) via an
empirically validated exponential relationship:
\begin{equation}
  \lambda = c \cdot 2^{(QP - 12)/3}
  \label{eq:lambda-qp}
\end{equation}
where $c$ is a codec-specific constant. This relationship ensures that
finer quantization (lower QP) assigns higher cost to bits, favoring modes
that reduce distortion, while coarser quantization (higher QP) favors
modes that reduce rate.

\begin{figure}[ht]
\centering
\begin{tikzpicture}[>=Stealth, scale=0.85, every node/.style={scale=0.85},
  block/.style={draw, rounded corners=3pt, minimum width=1.6cm,
    minimum height=0.8cm, font=\scriptsize, align=center, thick},
  arrow/.style={->, thick},
  label/.style={font=\tiny, text=black!60}]

% --- Input ---
\node[block, fill=blue!15] (input) at (0, 0) {Input\\frame $f_t$};

% --- Subtract ---
\node[draw, circle, inner sep=1.5pt, thick] (sub) at (2.5, 0) {$-$};

% --- Transform ---
\node[block, fill=orange!20] (transform) at (5, 0) {Transform\\$T$};

% --- Quantize ---
\node[block, fill=orange!20] (quant) at (7.5, 0) {Quantize\\$Q$};

% --- Entropy ---
\node[block, fill=purple!20] (entropy) at (10.5, 0) {Entropy\\coder};

% --- Bitstream ---
\node[font=\scriptsize, right=0.4cm of entropy] (bitstream) {Bitstream};

% --- Inverse Quantize ---
\node[block, fill=orange!20] (iquant) at (7.5, -2) {Inverse\\$Q^{-1}$};

% --- Inverse Transform ---
\node[block, fill=orange!20] (itransform) at (5, -2) {Inverse\\$T^{-1}$};

% --- Add ---
\node[draw, circle, inner sep=1.5pt, thick] (add) at (2.5, -2) {$+$};

% --- In-loop filter ---
\node[block, fill=red!15] (filter) at (2.5, -3.8) {In-loop\\filter};

% --- Reference buffer ---
\node[block, fill=red!15] (refbuf) at (5.5, -3.8) {Reference\\buffer};

% --- Mode decision ---
\node[block, fill=green!15] (mode) at (8.5, -3.8) {Mode\\decision};

% --- Prediction ---
\node[block, fill=green!15] (pred) at (5.5, -5.5) {Prediction\\(intra/inter)};

% --- Forward path ---
\draw[arrow] (input) -- (sub);
\draw[arrow] (sub) -- node[above, label] {$e$} (transform);
\draw[arrow] (transform) -- (quant);
\draw[arrow] (quant) -- (entropy);
\draw[arrow] (entropy) -- (bitstream);

% --- Reconstruction path ---
\draw[arrow] (quant) -- (iquant);
\draw[arrow] (iquant) -- (itransform);
\draw[arrow] (itransform) -- node[above, label] {$\hat{e}$} (add);

% --- Filter and buffer ---
\draw[arrow] (add) -- (filter);
\draw[arrow] (filter) -- (refbuf);

% --- Prediction feedback ---
\draw[arrow] (refbuf) -- (mode);
\draw[arrow] (mode) -- (pred);
\draw[arrow] (pred) -| (sub);
\draw[arrow] (pred) -| (add);

% --- Labels ---
\node[label, anchor=west] at (0, -5.5) {$\hat{f}_{\text{pred}}$};

\end{tikzpicture}
\caption{Block-based hybrid encoder architecture. Every codec from H.264
  through VVC follows this structure: predict each block from previously
  coded data (green), transform and quantize the residual (orange), entropy
  code the result (purple), then reconstruct and filter for future
  reference (red). The coding efficiency gap $\Delta$
  (eq.~\ref{eq:coding-gap}) arises from sub-optimal prediction, imperfect
  energy compaction, quantization beyond the rate-distortion bound, and
  entropy coding overhead.}
\label{fig:hybrid-encoder}
\end{figure}

The four sub-optimal components of Figure~\ref{fig:hybrid-encoder}---prediction,
transform, quantization, and entropy coding---each contribute to $\Delta$:
\begin{equation}
  \Delta \approx \Delta_{\text{pred}} + \Delta_{\text{transform}}
    + \Delta_{\text{quant}} + \Delta_{\text{entropy}}
  \label{eq:delta-decomp}
\end{equation}
Generational codec improvements reduce one or more of these terms. The codec
lineages in Sections~\ref{sec:mpeg-lineage}--\ref{sec:aom-lineage} trace
exactly which tools target which terms.

%% ---------------------------------------------------------------------------
\subsection{Transforms}
\label{sec:transforms}

The transform stage decorrelates the prediction residual, concentrating
energy into a few coefficients that can be efficiently quantized and coded.

\subsubsection{The KLT and Its DCT Approximation}
\label{sec:klt-dct}

The optimal decorrelating transform for a stationary source is the
\emph{Karhunen-Lo\`eve Transform} (KLT), whose basis vectors are the
eigenvectors of the input autocorrelation matrix $\mathbf{R}_{xx}$. For a
first-order Markov source with correlation $\rho$ (eq.~\ref{eq:spatial-autocorr}),
the KLT basis converges to the Type-II DCT as block size grows:
\begin{equation}
  C_k(n) = \alpha_k \cos\!\left[\frac{\pi k (2n+1)}{2N}\right],
  \quad k,n = 0, \ldots, N-1
  \label{eq:dct-ii}
\end{equation}
where $\alpha_0 = \sqrt{1/N}$ and $\alpha_k = \sqrt{2/N}$ for $k > 0$.

The energy compaction efficiency---the fraction of total energy captured by
the first $K$ of $N$ coefficients---is:
\begin{equation}
  \eta(K, N) = \frac{\sum_{k=0}^{K-1} \sigma_k^2}
    {\sum_{k=0}^{N-1} \sigma_k^2}
  \label{eq:energy-compaction}
\end{equation}
where $\sigma_k^2$ is the variance of the $k$-th transform coefficient.
For a typical intra residual with $\rho = 0.95$ and $N = 8$, the first 3
DCT coefficients capture approximately 95\% of the energy. This
concentration enables aggressive quantization of high-frequency coefficients
with minimal perceptual impact.

\subsubsection{Integer Transforms in Practice}
\label{sec:integer-transforms}

Practical codecs use integer approximations of the DCT to ensure bit-exact
encoder-decoder agreement. H.264 introduced a $4 \times 4$ integer
transform:
\begin{equation}
  \mathbf{H}_4 = \begin{pmatrix}
    1 &  1 &  1 &  1 \\
    2 &  1 & -1 & -2 \\
    1 & -1 & -1 &  1 \\
    1 & -2 &  2 & -1
  \end{pmatrix}
  \label{eq:h264-transform}
\end{equation}
The non-unity scaling factors are absorbed into the quantization step,
preserving exact integer arithmetic throughout the encode-decode loop.

Each codec generation expands the transform toolkit:

\begin{table}[ht]
\centering
\small
\renewcommand{\arraystretch}{1.4}
\begin{tabularx}{\textwidth}{lXXX}
\toprule
\textbf{Codec} & \textbf{Transform sizes} & \textbf{Transform types}
  & \textbf{Selection method} \\
\midrule
H.264 &
  $4\!\times\!4$, $8\!\times\!8$ &
  Integer DCT &
  Fixed per block size \\
VP9 &
  $4\!\times\!4$ to $32\!\times\!32$ &
  DCT, ADST &
  Row/column adaptive \\
HEVC &
  $4\!\times\!4$ to $32\!\times\!32$ &
  Integer DCT, DST ($4\!\times\!4$ intra) &
  Mode-dependent \\
AV1 &
  $4\!\times\!4$ to $64\!\times\!64$ &
  DCT, ADST, FlipADST, Identity &
  RD-optimized per block \\
VVC &
  $4\!\times\!4$ to $64\!\times\!64$ &
  DCT-II, DST-VII $+$ LFNST &
  RD-optimized $+$ secondary \\
\bottomrule
\end{tabularx}
\caption{Transform sizes and types across codec generations. AV1 applies
  four 1D transform types independently to rows and columns, yielding
  16~possible 2D combinations per block. VVC adds a non-separable secondary
  transform (LFNST) applied to low-frequency coefficients after the primary
  separable transform.}
\label{tab:transforms}
\end{table}

The Asymmetric Discrete Sine Transform (ADST) used in VP9 and AV1 is
particularly effective for intra-predicted residuals, which tend to have
larger magnitudes near the prediction boundary. AV1's FlipADST extends
this to blocks where the large residuals are on the opposite edge, while
the Identity transform (skip) handles blocks where the residual is already
well-distributed in the spatial domain.

The progression from fixed transforms (H.264) to RD-optimized multi-type
selection (AV1, VVC) directly reduces $\Delta_{\text{transform}}$ in
(\ref{eq:delta-decomp}): better energy compaction means fewer bits are
needed to represent the same residual at a given distortion level.

%% ---------------------------------------------------------------------------
\subsection{Motion Estimation and Inter Prediction}
\label{sec:motion}

Inter prediction exploits temporal redundancy by predicting blocks in the
current frame from previously coded reference frames. This is the single
largest source of compression gain---the prediction residual
(eq.~\ref{eq:residual}) typically has 10--20$\times$ less energy than the
original signal.

\subsubsection{Block Matching and Motion Vectors}
\label{sec:block-matching}

The encoder searches for the best-matching region in a reference frame by
minimizing a cost function over candidate displacement vectors
$\mathbf{d} = (d_x, d_y)$:
\begin{equation}
  \text{SAD}(\mathbf{d}) = \sum_{(x,y) \in B}
    \left| f_t(x,y) - f_{t-1}(x + d_x,\, y + d_y) \right|
  \label{eq:sad}
\end{equation}
The Sum of Absolute Differences (SAD) is the simplest distortion measure.
A better rate-distortion proxy is the Sum of Absolute Transformed
Differences (SATD), which applies the Hadamard transform before summation:
\begin{equation}
  \text{SATD}(\mathbf{d}) = \sum \left| \mathbf{H}
    \cdot \left( f_t - f_{t-1}(\mathbf{d}) \right) \right|
  \label{eq:satd}
\end{equation}
SATD approximates the post-transform energy and thus better predicts the
actual coded bit cost. Rate-constrained motion estimation augments the
distortion with a motion vector coding cost:
\begin{equation}
  \mathbf{d}^* = \arg\min_{\mathbf{d}} \left[
    \text{SATD}(\mathbf{d}) + \lambda_{\text{motion}} \cdot R(\mathbf{d})
  \right]
  \label{eq:rcme}
\end{equation}
where $R(\mathbf{d})$ is the bit cost of encoding the motion vector
relative to a predictor.

\subsubsection{Sub-Pixel Interpolation}
\label{sec:subpixel}

Motion in natural video rarely aligns to integer pixel positions. All modern
codecs perform sub-pixel motion compensation by interpolating reference
frame samples at fractional positions.

H.264 uses a 6-tap filter for half-pixel luma interpolation:
\begin{equation}
  h_{\frac{1}{2}} = \{1, -5, 20, 20, -5, 1\} / 32
  \label{eq:h264-halfpel}
\end{equation}
with quarter-pixel positions obtained by bilinear averaging of integer and
half-pixel samples. HEVC extends this to an 8-tap filter for half-pixel and
a 7-tap filter for quarter-pixel positions. AV1 introduces a switchable
filter bank with three 8-tap options---Regular (sharp), Smooth (low-pass),
and Sharp (high-frequency preserving)---selected per block via RD
optimization. VVC adds DMVR (Decoder-side Motion Vector Refinement) and
BDOF (Bi-Directional Optical Flow) for sub-pixel refinement at the decoder
without additional signaling bits.

The progression from H.264's fixed 6-tap filter to AV1's switchable bank
reduces $\Delta_{\text{pred}}$: better interpolation means smaller
prediction residuals, which require fewer bits.

\subsubsection{Advanced Inter Prediction Tools}
\label{sec:advanced-inter}

Beyond basic block matching, each codec generation adds tools to handle
complex motion and occlusion:

\begin{itemize}
\item \textbf{Multiple reference frames.} H.264 introduced up to 16
  reference frames; practical WebRTC encoders use 1--3 for latency reasons.
  More references improve prediction for periodic motion and uncovered
  backgrounds.

\item \textbf{Bi-prediction.} A weighted average of predictions from two
  reference frames, reducing noise and handling gradual scene changes.
  Available from H.264 onward (in B-frames) and in AV1's compound modes.

\item \textbf{Compound prediction (AV1).} Wedge masks and
  difference-weighted blending combine two predictions at the block level,
  handling occlusion boundaries where a single motion model fails.

\item \textbf{Overlapped Block Motion Compensation (AV1).} Blends
  predictions from neighboring blocks at motion boundaries, reducing
  blocking artifacts without post-processing.

\item \textbf{Warped motion (AV1).} An affine model with up to 6
  parameters handles rotation, zoom, and perspective---common in
  screen sharing and camera motion.

\item \textbf{VVC additions.} Geometric Partitioning Mode (non-rectangular
  block splits for motion boundaries), BDOF (optical flow refinement at the
  decoder), and DMVR (decoder-side MV refinement using bilateral matching).
\end{itemize}

\excursus{Excursus: Motion estimation in real-time encoding.}{Full-search
block matching has complexity $O(N^2 \cdot S^2)$ where $N$ is the block
dimension and $S$ is the search range in pixels. For a $64\!\times\!64$
block with $S = 64$, this is ${\sim}17$ million SAD operations per
block---far too expensive for real-time encoding. WebRTC encoders use
hierarchical search patterns (diamond, hexagon) that reduce complexity to
$O(N^2 \cdot \log S)$ with typically $< 5\%$ rate-distortion penalty
relative to full search. This is why real-time encoders use only a subset
of each codec's inter prediction tools: the full tool set is available
for offline encoding, but the real-time constraint forces aggressive
pruning of the mode decision search space.}

%% ---------------------------------------------------------------------------
\subsection{Quantization}
\label{sec:quantization}

Quantization is the sole lossy step in the hybrid coding pipeline. It maps
continuous-valued (or high-precision integer) transform coefficients to a
discrete set of reconstruction levels, introducing controlled information
loss to achieve the target bit rate.

\subsubsection{Scalar Quantization}

Uniform scalar quantization maps coefficient $c$ to a quantization index:
\begin{equation}
  q = \text{round}\!\left(\frac{c}{Q_{\text{step}}}\right), \quad
  \hat{c} = q \cdot Q_{\text{step}}
  \label{eq:uniform-quant}
\end{equation}
Modern codecs use \emph{dead-zone quantization}, where coefficients near
zero are mapped to zero more aggressively than uniform rounding predicts:
\begin{equation}
  q = \text{sign}(c) \cdot \left\lfloor
    \frac{|c| - f \cdot Q_{\text{step}}}{Q_{\text{step}}}
  \right\rfloor
  \label{eq:deadzone}
\end{equation}
where $f \in [0, 0.5)$ controls the dead-zone width. Dead-zone
quantization increases sparsity---the fraction of zero-valued
coefficients---which dramatically improves entropy coding efficiency.

\subsubsection{QP-to-Step-Size Mapping}

The quantization parameter (QP) controls the step size via an exponential
relationship. In the H.264/HEVC/VVC family:
\begin{equation}
  Q_{\text{step}} = Q_0 \cdot 2^{(QP - QP_0)/6}
  \label{eq:qp-mapping}
\end{equation}
A 6-unit increase in QP doubles the step size, approximately doubling the
bit rate. AV1 uses a similar but codec-specific mapping with separate QP
values for DC and AC coefficients, plus per-segment QP adjustment for
region-of-interest coding.

\subsubsection{Quantization Distortion}

For a single coefficient under uniform quantization at high rate, the
distortion is:
\begin{equation}
  D_q = \frac{Q_{\text{step}}^2}{12}
  \label{eq:quant-distortion}
\end{equation}
For a block of $N$ transform coefficients where $K$ survive quantization
(nonzero) and $N - K$ are quantized to zero:
\begin{equation}
  D_{\text{block}} \approx K \cdot \frac{Q_{\text{step}}^2}{12}
    + \sum_{k \notin K} c_k^2
  \label{eq:block-distortion}
\end{equation}
The second term represents the energy of coefficients eliminated by the
dead zone---information permanently discarded. The QP is thus the primary
control knob for navigating the rate-distortion curve in
Figure~\ref{fig:rate-distortion}. The WebCodecs \texttt{VideoEncoder} API
exposes this control via the \texttt{quantizer} configuration parameter,
giving JavaScript applications direct access to this fundamental
trade-off.

%% ---------------------------------------------------------------------------
\subsection{Entropy Coding}
\label{sec:entropy-coding}

After quantization, the surviving coefficients and all side information
(motion vectors, prediction modes, block partitions) must be losslessly
compressed. The entropy coder converts these symbols into a bitstream,
approaching the conditional entropy $H(X|\text{context})$ as closely as
possible. The gap between the entropy coder's output rate and the true
conditional entropy is $\Delta_{\text{entropy}}$ in
(\ref{eq:delta-decomp}).

\subsubsection{Context-Adaptive Binary Arithmetic Coding (CABAC)}
\label{sec:cabac}

CABAC, used in H.264, HEVC, and VVC, encodes each syntax element as a
sequence of binary decisions (bins). For each bin, the arithmetic coder
subdivides the current interval $[L, H)$ based on the estimated
probability $p$:
\begin{equation}
  [L, H) \to \begin{cases}
    [L,\; L + p \cdot (H - L)) & \text{if bin} = 0 \\
    [L + p \cdot (H - L),\; H) & \text{if bin} = 1
  \end{cases}
  \label{eq:arithmetic-coding}
\end{equation}
The theoretical output length approaches the entropy:
\begin{equation}
  \ell \to \sum_{i=1}^{n} -\log_2 p(b_i \mid \text{ctx}_i)
  \label{eq:cabac-length}
\end{equation}
CABAC achieves redundancy below 0.01 bits/symbol relative to the true
conditional entropy---near-optimal compression.

Context modeling is the key to CABAC's effectiveness. Each bin type has an
associated context, whose probability is updated after each observation via
a state-machine-based exponential moving average:
\begin{equation}
  p_{k+1} = p_k + \alpha(s_k) \cdot (b_k - p_k)
  \label{eq:cabac-update}
\end{equation}
where $\alpha(s_k)$ is a state-dependent adaptation rate, indexed by a
lookup table with 64 states (H.264) or 128 states (HEVC/VVC). This design
avoids floating-point arithmetic entirely.

CABAC also provides a \emph{bypass mode} for bins with approximately
uniform probability (e.g., sign bits, high-magnitude coefficient levels),
which skips the context lookup and probability estimation for throughput.

\subsubsection{Multi-Symbol Arithmetic Coding in AV1}
\label{sec:av1-entropy}

AV1 departs from binary arithmetic coding. Its entropy coder operates on
multi-symbol alphabets using CDF (Cumulative Distribution Function) tables:
\begin{equation}
  \text{Encode symbol } s: \quad
  [L, H) \to \left[L + \frac{\text{CDF}(s-1)}{M} \cdot (H-L),\;
    L + \frac{\text{CDF}(s)}{M} \cdot (H-L)\right)
  \label{eq:av1-entropy}
\end{equation}
where $M$ is the total frequency count. CDF tables are adapted after
each symbol:
\begin{equation}
  \text{CDF}_{k+1}(s) = \text{CDF}_k(s) + \left\lfloor
    \frac{\text{target}(s) - \text{CDF}_k(s)}{n}
  \right\rfloor
  \label{eq:cdf-update}
\end{equation}
where $n$ controls the adaptation rate. Multi-symbol coding is more
efficient than CABAC's binarization for syntax elements with moderate
alphabet sizes (e.g., intra prediction modes, transform types), as it
avoids the overhead of converting each symbol into a binary tree of
decisions.

\begin{table}[ht]
\centering
\small
\renewcommand{\arraystretch}{1.4}
\begin{tabularx}{\textwidth}{lXXXXX}
\toprule
\textbf{Feature} & \textbf{H.264} & \textbf{HEVC} & \textbf{VP9}
  & \textbf{AV1} & \textbf{VVC} \\
\midrule
Method &
  Binary AC &
  Binary AC &
  Boolean AC &
  Multi-sym AC &
  Binary AC \\
Contexts &
  ${\sim}400$ &
  ${\sim}200$ &
  ${\sim}900$ &
  ${\sim}300$ &
  ${\sim}700$ \\
Adaptation &
  64-state LUT &
  128-state LUT &
  Counter &
  CDF update &
  128-state LUT \\
Bypass mode &
  Yes &
  Yes &
  N/A &
  N/A &
  Yes \\
Parallelism &
  Slice &
  WPP / Tiles &
  Tiles &
  Tiles &
  WPP / Tiles \\
\bottomrule
\end{tabularx}
\caption{Entropy coding comparison across codecs. HEVC reduced context
  count relative to H.264 for throughput; VVC increased it again for
  improved compression. AV1's multi-symbol approach operates on full
  alphabets rather than binary decompositions.}
\label{tab:entropy-coding}
\end{table}

The frame entropy $\hat{H}(f_t)$ from eq.~(\ref{eq:frame-entropy}) is
ultimately determined by how well the entropy coder's probability model
matches the true symbol distribution. Both CABAC and AV1's multi-symbol
coder achieve near-zero entropy coding gaps; the remaining
$\Delta_{\text{entropy}}$ arises primarily from context model mismatch on
non-stationary content (scene changes, sudden motion onset).

%% ---------------------------------------------------------------------------
\subsection{In-Loop Filtering}
\label{sec:filtering}

Block-based coding introduces visible artifacts at block boundaries:
discontinuities from independent quantization of adjacent blocks, and
ringing from coefficient truncation. In-loop filters reduce these artifacts
\emph{before} the reconstructed frame enters the reference buffer, improving
prediction quality for subsequent frames and thus reducing
$\Delta_{\text{pred}}$.

\begin{itemize}
\item \textbf{Deblocking filter (all codecs).} Smooths discontinuities
  at block edges, with filter strength adapted to the quantization level
  and local gradient. Introduced in H.264; refined in each subsequent
  codec.

\item \textbf{Sample Adaptive Offset---SAO (HEVC).} Classifies
  reconstructed samples by intensity or edge direction and applies
  per-class offsets to reduce systematic bias introduced by quantization.

\item \textbf{Constrained Directional Enhancement Filter---CDEF (AV1).}
  A directional filter that identifies edge orientation per
  $8\!\times\!8$ block and applies filtering only along the edge,
  preserving sharpness while reducing ringing.

\item \textbf{Loop Restoration Filter (AV1).} A switchable per-tile
  filter choosing between a Wiener filter (FIR designed to minimize MSE
  between original and reconstruction) and a Self-Guided filter (non-linear
  denoising based on guided image filtering). This is AV1's most powerful
  quality enhancement tool.

\item \textbf{Adaptive Loop Filter---ALF (VVC).} A diamond-shaped filter
  with coefficients signaled per frame, applied adaptively per
  $4\!\times\!4$ block based on local classification. Subsumes HEVC's
  SAO functionality.
\end{itemize}

The evolution from simple deblocking (H.264) to multi-stage adaptive
filtering (AV1: deblocking $\to$ CDEF $\to$ loop restoration) exemplifies
how each codec generation reduces $\Delta_{\text{pred}}$ by improving
reference frame quality.

%% ---------------------------------------------------------------------------
\subsection{Codec Lineage: MPEG/ITU Family}
\label{sec:mpeg-lineage}

The MPEG/ITU standardization track has produced three codec generations
relevant to the 2021--2031 timeframe of this document.

\subsubsection{H.264/AVC (2003)}
\label{sec:h264}

H.264 defined the modern era of video compression and remains the most
widely deployed codec. Its key innovations over its predecessors (H.263,
MPEG-2):

\begin{itemize}
\item Variable block sizes: $16\!\times\!16$ macroblock partitioned into
  $16\!\times\!8$, $8\!\times\!16$, $8\!\times\!8$, $8\!\times\!4$,
  $4\!\times\!8$, and $4\!\times\!4$ sub-partitions for motion
  compensation.

\item Quarter-pixel motion compensation with 6-tap interpolation
  (eq.~\ref{eq:h264-halfpel}).

\item Context-adaptive intra prediction: 9~angular modes for
  $4\!\times\!4$ blocks, 4~modes for $16\!\times\!16$.

\item Dual entropy coding: CAVLC (simpler, faster) and CABAC
  (Section~\ref{sec:cabac}; 10--15\% better compression).

\item In-loop deblocking filter, mandatory for the decoder.

\item Multiple reference frames (up to 16) and weighted prediction.
\end{itemize}

Three profiles are relevant to WebRTC: \textbf{Constrained Baseline}
(no CABAC, no B-frames, lowest complexity---the mandatory WebRTC profile),
\textbf{Main} (CABAC, B-frames, ${\sim}$10\% better compression), and
\textbf{High} ($8\!\times\!8$ transform, adaptive quantization matrices,
${\sim}$15\% better).

At the 720p/30fps reference point, H.264 Constrained Baseline operates at
approximately 2.0--2.5~Mbps for ``good'' quality (PSNR ${\sim}$37~dB),
corresponding to $\Delta_{\text{H.264}} \approx 2.1$ from
Appendix~\ref{app:info-theory}. H.264 was the only codec Safari supported
in 2021 (Section~3.1), and its higher $\Delta$ is directly responsible
for the bandwidth disparity between Safari-only and VP9-capable sessions
documented in the main text.

\subsubsection{H.265/HEVC (2013)}
\label{sec:hevc}

HEVC replaced H.264's fixed $16\!\times\!16$ macroblock with a recursive
Coding Tree Unit (CTU) of up to $64\!\times\!64$ pixels, using quad-tree
partitioning down to $8\!\times\!8$. Key advances:

\begin{itemize}
\item 35 intra prediction angular directions (vs.\ H.264's 9).
\item Transform sizes from $4\!\times\!4$ to $32\!\times\!32$, with
  DST for $4\!\times\!4$ intra blocks.
\item Advanced Motion Vector Prediction (AMVP): two MVP candidates selected
  from spatial and temporal neighbors, replacing H.264's less structured
  motion vector prediction.
\item Merge mode: inherits the full motion information (MV, reference
  index, prediction direction) from a neighboring block, requiring zero
  additional motion bits.
\item Sample Adaptive Offset (SAO) in-loop filter.
\item Parallel processing: Wavefront Parallel Processing (WPP) and tiles.
\end{itemize}

HEVC achieves approximately 40\% rate reduction over H.264 at equivalent
PSNR. However, HEVC is absent from WebRTC entirely---not due to technical
limitations, but due to patent licensing.

\excursus{Excursus: Why no HEVC in WebRTC?}{HEVC's patent licensing is
fragmented across three separate pools (MPEG-LA, HEVC Advance, Velos Media)
plus independent licensors who are not members of any pool. The resulting
uncertainty about total royalty cost---estimated at \$0.20--\$0.40 per device
for some use cases---led Google, Mozilla, Cisco, and others to form the
Alliance for Open Media (AOM) in 2015 and develop AV1 as a royalty-free
alternative. This patent-driven fork is why the WebRTC ecosystem followed the
VP8$\to$VP9$\to$AV1 path rather than H.264$\to$HEVC$\to$VVC, and why the two
codec families must be understood in parallel.}

\subsubsection{H.266/VVC (2020)}
\label{sec:vvc}

VVC extends HEVC's quad-tree with binary and ternary splits, creating a
\emph{multi-type tree} (MTT) that enables asymmetric block shapes (e.g.,
$32\!\times\!8$, $8\!\times\!32$, $32\!\times\!4$). Additional
innovations:

\begin{itemize}
\item 67 intra prediction directions plus Matrix-based Intra Prediction
  (MIP), which applies a learned linear transform to neighboring samples.

\item Multiple Transform Selection (MTS): explicit signaling of DCT-II or
  DST-VII per block, plus the Low-Frequency Non-Separable Transform
  (LFNST)---a $16\!\times\!16$ or $8\!\times\!8$ non-separable transform
  applied to low-frequency coefficients after the primary separable
  transform.

\item Affine motion model: 4-parameter (translation $+$ rotation/scaling)
  or 6-parameter (full affine) motion, derived from control point MVs at
  block corners.

\item Decoder-side tools: DMVR refines motion vectors at the decoder using
  bilateral template matching; BDOF applies optical flow equations to
  refine bi-predicted samples, both without additional signaling bits.

\item Geometric Partitioning Mode: divides a block diagonally for separate
  motion compensation of each partition, targeting motion boundaries.

\item Adaptive Loop Filter (ALF): a diamond-shaped Wiener filter with
  coefficients signaled per frame and applied adaptively per
  $4\!\times\!4$ block.
\end{itemize}

VVC achieves approximately 30--50\% rate reduction over HEVC, placing it at
$\Delta_{\text{VVC}} \approx 0.3$. Its relevance to this document is
forward-looking: VVC may appear in WebCodecs codec registries by
${\sim}$2030 if licensing concerns are resolved, and its tool set
represents the current efficiency frontier against which AV1's performance
should be measured.

%% ---------------------------------------------------------------------------
\subsection{Codec Lineage: Google/AOM Family}
\label{sec:aom-lineage}

The royalty-free codec family dominates the WebRTC ecosystem. Its
development was driven by the patent licensing concerns described in the
HEVC excursus above.

\subsubsection{VP8 (2010)}
\label{sec:vp8}

VP8 was Google's first open-source video codec, released alongside the
WebM container. Its design is conservative:

\begin{itemize}
\item $16\!\times\!16$ macroblocks with $4\!\times\!4$ subblock partitions.
\item $4\!\times\!4$ Walsh-Hadamard Transform (WHT) for DC coefficients,
  $4\!\times\!4$ DCT for AC coefficients.
\item Boolean arithmetic coder with tree-structured syntax (single-bit
  alphabet).
\item 4 intra prediction modes for $16\!\times\!16$, 10 for
  $4\!\times\!4$ subblocks.
\item Simple loop filter with per-edge strength adaptation.
\item Single reference frame for WebRTC (plus a ``golden'' frame for
  screen content).
\end{itemize}

VP8's compression efficiency is roughly equivalent to H.264 Constrained
Baseline. It served as the baseline WebRTC codec---supported by Chrome and
Firefox since WebRTC~1.0---but has been superseded by VP9 in all modern
deployments.

\subsubsection{VP9 (2013)}
\label{sec:vp9}

VP9 introduced a superblock architecture and several tools that brought it
within range of HEVC:

\begin{itemize}
\item $64\!\times\!64$ superblocks with recursive quad-tree partitioning
  down to $4\!\times\!4$.
\item 10 intra prediction modes.
\item Transforms: DCT and ADST applied independently per row and column,
  at sizes from $4\!\times\!4$ to $32\!\times\!32$.
\item Switchable interpolation filters: 8-tap Regular, Smooth, and
  Sharp, selected per block.
\item Compound prediction: weighted average from two reference frames.
\item Segmentation: per-segment QP adjustment for region-of-interest
  coding.
\end{itemize}

VP9 achieves approximately 35\% rate reduction over H.264
(eq.~\ref{eq:vp9-savings}), placing it at $\Delta_{\text{VP9}} \approx 0.8$.
VP9 is the dominant WebRTC codec in Chrome and Firefox as of March~2026,
and its temporal SVC support (Section~\ref{sec:svc}) is the primary
mechanism enabling adaptive quality in multi-participant calls.

\subsubsection{AV1 (2018)}
\label{sec:av1}

AV1 represents the most ambitious single-generation improvement in the
royalty-free family. Its tool set approaches VVC in breadth while
maintaining royalty-free status:

\begin{itemize}
\item $128\!\times\!128$ superblocks with flexible partitioning (quad-tree
  plus rectangular splits) down to $4\!\times\!4$.

\item 56+ directional intra prediction modes plus non-directional modes
  (Paeth, Smooth variants, DC).

\item Full inter prediction toolkit: OBMC, warped motion (affine),
  compound modes (Wedge, Difference-weighted), inter-intra compound
  prediction.

\item 16 transform type combinations ($4$ 1D types $\times$ rows/columns),
  at sizes up to $64\!\times\!64$.

\item Film Grain Synthesis: the encoder characterizes noise/grain
  parameters and transmits them as metadata; the decoder re-synthesizes
  grain after reconstruction. This removes non-compressible high-frequency
  noise from the coding loop---a major bitrate saving for camera-captured
  content.

\item Multi-stage in-loop filtering: deblocking $\to$ CDEF $\to$ loop
  restoration (Wiener or Self-Guided), as described in
  Section~\ref{sec:filtering}.

\item Multi-symbol arithmetic coding with CDF-based context adaptation
  (Section~\ref{sec:av1-entropy}).
\end{itemize}

AV1 achieves approximately 30\% rate reduction over VP9
(eq.~\ref{eq:av1-savings}), a cumulative ${\sim}$54\% over H.264, placing
it at $\Delta_{\text{AV1}} \approx 0.5$ with current hardware encoders and
approaching 0.4 by 2030 as encoder implementations mature.

AV1's tools directly explain why its $\Delta$ is smaller than VP9's:
more transform types capture residual structure better (reducing
$\Delta_{\text{transform}}$), warped and compound motion reduce prediction
residual energy (reducing $\Delta_{\text{pred}}$), and film grain synthesis
removes non-compressible variance from the coding loop (effectively
reducing the source variance $\sigma^2$ that the codec must represent).

%% ---------------------------------------------------------------------------
\subsection{Block Partitioning: A Generational Comparison}
\label{sec:partitioning}

Block partitioning is the most architecturally consequential difference
between codec generations. Larger and more flexible partitions allow the
encoder to match block boundaries to content structure---smooth regions
get large blocks (few bits for partition overhead), detailed regions get
small blocks (better prediction accuracy).

\begin{figure}[ht]
\centering
\begin{tikzpicture}[>=Stealth, scale=0.55, every node/.style={scale=0.55}]

% --- H.264 panel ---
\begin{scope}[xshift=0cm]
  \node[font=\normalsize\bfseries, anchor=south] at (2, 8.4) {H.264};
  \node[font=\small, anchor=south] at (2, 7.8) {$16\!\times\!16$ max};
  % 4x4 grid of 16x16 macroblocks in a 64x64 region
  \draw[thick] (0,0) rectangle (4,4);
  \draw[thick] (0,0) rectangle (4,4);
  % macroblock grid
  \foreach \x in {0,1,2,3} {
    \foreach \y in {0,1,2,3} {
      \draw[gray] (\x, \y) rectangle (\x+1, \y+1);
    }
  }
  % show some sub-partitions
  \draw[blue!60, thick] (0,3) -- (0.5,3) -- (0.5,4) -- (0,4);
  \draw[blue!60, thick] (0.5,3) -- (0.5,4);
  \draw[blue!60, thick] (0,3.5) -- (1,3.5);
  % 8x8 in another
  \draw[blue!60, thick] (1,2.5) -- (2,2.5);
  \draw[blue!60, thick] (1.5,2) -- (1.5,3);
  % 4x4 in another
  \draw[red!60, thick] (2,0) -- (2,1) -- (3,1) -- (3,0);
  \draw[red!60, thick] (2.5,0) -- (2.5,0.5);
  \draw[red!60, thick] (2,0.5) -- (2.5,0.5);
  \draw[red!60, thick] (2.5,0.5) -- (2.5,1);
  \draw[red!60, thick] (2.5,0) -- (3,0) -- (3,0.5) -- (2.5,0.5);
  % Scale label
  \node[font=\small, anchor=north] at (2, -0.3) {64 pixels};
\end{scope}

% --- VP9/HEVC panel ---
\begin{scope}[xshift=6cm]
  \node[font=\normalsize\bfseries, anchor=south] at (4, 8.4) {VP9 / HEVC};
  \node[font=\small, anchor=south] at (4, 7.8) {$64\!\times\!64$ max};
  % One 64x64 superblock
  \draw[thick] (0,0) rectangle (8,8);
  % quad-tree splits
  \draw[thick] (4,0) -- (4,8);
  \draw[thick] (0,4) -- (8,4);
  % further splits in top-left quadrant
  \draw[blue!60, thick] (2,4) -- (2,8);
  \draw[blue!60, thick] (0,6) -- (4,6);
  % further split in top-left-top-left
  \draw[red!60, thick] (1,6) -- (1,8);
  \draw[red!60, thick] (0,7) -- (2,7);
  % split bottom-right quadrant
  \draw[blue!60, thick] (6,0) -- (6,4);
  \draw[blue!60, thick] (4,2) -- (8,2);
  % deeper in bottom-right-bottom-left
  \draw[red!60, thick] (5,0) -- (5,2);
  \draw[red!60, thick] (4,1) -- (6,1);
  % Scale label
  \node[font=\small, anchor=north] at (4, -0.3) {64 pixels};
\end{scope}

% --- AV1 panel ---
\begin{scope}[xshift=16cm]
  \node[font=\normalsize\bfseries, anchor=south] at (4, 8.4) {AV1};
  \node[font=\small, anchor=south] at (4, 7.8) {$128\!\times\!128$ max};
  % 128x128 shown as 8x8 grid (each cell = 16px)
  \draw[thick] (0,0) rectangle (8,8);
  % quad split: top half / bottom half
  \draw[thick] (4,0) -- (4,8);
  \draw[thick] (0,4) -- (8,4);
  % rectangular splits in top-left
  \draw[blue!60, thick] (0,6) -- (4,6);
  \draw[blue!60, thick] (2,4) -- (2,6);
  % asymmetric in bottom-left (4:1 rect)
  \draw[orange!80, thick] (0,3) -- (4,3);
  \draw[orange!80, thick] (0,2) -- (4,2);
  \draw[orange!80, thick] (0,1) -- (4,1);
  % quad + rect in top-right
  \draw[blue!60, thick] (6,4) -- (6,8);
  \draw[red!60, thick] (4,6) -- (6,6);
  \draw[red!60, thick] (5,4) -- (5,6);
  % bottom-right: large blocks
  \draw[blue!60, thick] (4,2) -- (8,2);
  % Scale label
  \node[font=\small, anchor=north] at (4, -0.3) {128 pixels};
\end{scope}

% --- VVC panel ---
\begin{scope}[xshift=26cm]
  \node[font=\normalsize\bfseries, anchor=south] at (4, 8.4) {VVC};
  \node[font=\small, anchor=south] at (4, 7.8) {$128\!\times\!128$ max};
  \draw[thick] (0,0) rectangle (8,8);
  % quad split
  \draw[thick] (4,0) -- (4,8);
  \draw[thick] (0,4) -- (8,4);
  % binary split (horizontal) in top-left
  \draw[blue!60, thick] (0,6) -- (4,6);
  % ternary split (vertical) in top-left-bottom
  \draw[orange!80, thick] (1,4) -- (1,6);
  \draw[orange!80, thick] (3,4) -- (3,6);
  % binary split (vertical) in top-right
  \draw[blue!60, thick] (6,4) -- (6,8);
  % ternary horizontal in top-right-left
  \draw[orange!80, thick] (4,5.33) -- (6,5.33);
  \draw[orange!80, thick] (4,6.67) -- (6,6.67);
  % bottom-left: binary horizontal
  \draw[blue!60, thick] (0,2) -- (4,2);
  % deeper ternary in bottom-left-top
  \draw[orange!80, thick] (1,2) -- (1,4);
  \draw[orange!80, thick] (3,2) -- (3,4);
  % bottom-right: mostly large
  \draw[blue!60, thick] (6,0) -- (6,4);
  \draw[blue!60, thick] (4,2) -- (8,2);
  % Scale label
  \node[font=\small, anchor=north] at (4, -0.3) {128 pixels};
\end{scope}

\end{tikzpicture}
\caption{Block partitioning comparison across codec generations. H.264 is
  limited to $16\!\times\!16$ macroblocks with fixed sub-partitions (left).
  VP9/HEVC use quad-tree partitioning within $64\!\times\!64$ superblocks
  (center-left). AV1 extends to $128\!\times\!128$ with quad-tree plus
  rectangular splits (center-right). VVC adds binary and ternary splits for
  maximum flexibility (right). Blue lines show first-level splits; red
  shows deeper partitioning; orange shows asymmetric/ternary splits.}
\label{fig:block-partitioning}
\end{figure}

\begin{table}[ht]
\centering
\small
\renewcommand{\arraystretch}{1.4}
\begin{tabularx}{\textwidth}{lXXXXX}
\toprule
\textbf{Feature} & \textbf{H.264} & \textbf{VP9} & \textbf{HEVC}
  & \textbf{AV1} & \textbf{VVC} \\
\midrule
Max block &
  $16\!\times\!16$ &
  $64\!\times\!64$ &
  $64\!\times\!64$ &
  $128\!\times\!128$ &
  $128\!\times\!128$ \\
Min block &
  $4\!\times\!4$ &
  $4\!\times\!4$ &
  $8\!\times\!8$ &
  $4\!\times\!4$ &
  $4\!\times\!4$ \\
Split types &
  Fixed partitions &
  Quad-tree &
  Quad-tree &
  Quad $+$ rect &
  Quad $+$ BT $+$ TT \\
Asymmetric &
  Limited &
  No &
  No &
  4:1 rectangles &
  Full (BT/TT) \\
\bottomrule
\end{tabularx}
\caption{Block partitioning capabilities by codec. BT = binary tree
  (horizontal or vertical), TT = ternary tree (1:2:1 split). Larger maximum
  blocks reduce overhead for homogeneous regions; asymmetric splits improve
  partitioning efficiency at content boundaries.}
\label{tab:partitioning}
\end{table}

%% ---------------------------------------------------------------------------
\subsection{Rate-Distortion Operating Points by Codec}
\label{sec:rd-comparison}

The codec-specific analyses of
Sections~\ref{sec:mpeg-lineage}--\ref{sec:aom-lineage} can be unified
through the Bj\o ntegaard Delta Rate (BD-Rate), the standard metric for
comparing codec efficiency. BD-Rate computes the average percentage rate
difference between two codecs at equivalent distortion by integrating the
difference between their rate-distortion curves fit as cubic polynomials in
PSNR-vs-log-rate space:
\begin{equation}
  \text{BD-Rate} = \frac{1}{\Delta\text{PSNR}}
    \int_{\text{PSNR}_{\min}}^{\text{PSNR}_{\max}}
    \left[ \ln R_A(p) - \ln R_B(p) \right] dp
  \label{eq:bd-rate}
\end{equation}
A negative BD-Rate indicates that codec~$A$ requires fewer bits than
codec~$B$ at the same quality.

\begin{figure}[ht]
\centering
\begin{tikzpicture}[>=Stealth]

% --- Axes ---
\draw[->, thick] (0, 0) -- (11, 0);
\draw[->, thick] (0, 0) -- (0, 6.5);

% --- Axis titles ---
\node[font=\small] at (5.5, -0.85) {Rate [Mbps]};
\node[font=\small, rotate=90] at (-0.9, 3.25) {PSNR [dB]};

% --- X-axis labels ---
\foreach \x/\lab in {0/0, 2/0.5, 4/1.0, 6/1.5, 8/2.0, 10/2.5} {
  \draw (\x, -0.1) -- (\x, 0.1);
  \node[font=\tiny, anchor=north] at (\x, -0.15) {\lab};
}

% --- Y-axis labels ---
\foreach \y/\lab in {0/28, 1/30, 2/32, 3/34, 4/36, 5/38, 6/40} {
  \draw (-0.1, \y) -- (0.1, \y);
  \node[font=\tiny, anchor=east] at (-0.15, \y) {\lab};
}

% --- Grid ---
\foreach \y in {1,2,3,4,5,6} {
  \draw[gray!15] (0, \y) -- (10.5, \y);
}

% --- Shannon R(D) bound (leftmost curve) ---
\draw[very thick, green!60!black]
  (0.4, 0.5) .. controls (0.6, 2) and (0.8, 3.5) ..
  (1.2, 4.5) .. controls (1.8, 5.5) and (2.5, 6.0) ..
  (3.5, 6.3);
\node[font=\scriptsize, green!60!black, anchor=south west] at (2.5, 6.1)
  {$R(D)$};

% --- VVC curve ---
\draw[very thick, violet!70]
  (0.8, 0.3) .. controls (1.2, 2.0) and (1.6, 3.5) ..
  (2.2, 4.5) .. controls (3.0, 5.5) and (4.0, 6.0) ..
  (5.5, 6.3);
\node[font=\scriptsize, violet!70, anchor=south west] at (4.2, 6.1)
  {VVC};

% --- AV1 curve ---
\draw[very thick, orange!80!black]
  (1.0, 0.2) .. controls (1.5, 1.8) and (2.2, 3.3) ..
  (3.0, 4.3) .. controls (4.0, 5.3) and (5.2, 5.9) ..
  (7.0, 6.2);
\node[font=\scriptsize, orange!80!black, anchor=south] at (6.0, 6.1)
  {AV1};

% --- HEVC curve ---
\draw[very thick, teal!70]
  (1.2, 0.1) .. controls (1.8, 1.6) and (2.6, 3.1) ..
  (3.6, 4.1) .. controls (4.8, 5.1) and (6.2, 5.8) ..
  (8.0, 6.1);
\node[font=\scriptsize, teal!70, anchor=south] at (7.2, 6.0)
  {HEVC};

% --- VP9 curve ---
\draw[very thick, blue!70]
  (1.4, 0.0) .. controls (2.0, 1.4) and (3.0, 2.9) ..
  (4.2, 4.0) .. controls (5.5, 5.0) and (7.0, 5.7) ..
  (9.0, 6.0);
\node[font=\scriptsize, blue!70, anchor=south west] at (8.2, 5.9)
  {VP9};

% --- H.264 curve (rightmost) ---
\draw[very thick, red!70!black]
  (2.0, 0.0) .. controls (3.0, 1.2) and (4.2, 2.6) ..
  (5.5, 3.7) .. controls (7.0, 4.7) and (8.5, 5.4) ..
  (10.2, 5.8);
\node[font=\scriptsize, red!70!black, anchor=south west] at (9.5, 5.6)
  {H.264};

% --- 37 dB reference line ---
\draw[thin, gray, dashed] (0, 4.5) -- (10.5, 4.5);
\node[font=\tiny, gray, anchor=west] at (10.6, 4.5) {37\,dB};

% --- Operating points at 37 dB ---
\fill[red!70!black] (5.2, 4.5) circle (2.5pt);
\fill[blue!70] (3.9, 4.5) circle (2.5pt);
\fill[teal!70] (3.3, 4.5) circle (2.5pt);
\fill[orange!80!black] (2.7, 4.5) circle (2.5pt);
\fill[violet!70] (2.0, 4.5) circle (2.5pt);
\fill[green!60!black] (1.1, 4.5) circle (2.5pt);

% --- Annotation ---
\node[font=\scriptsize\itshape, text=black!60, align=center,
  anchor=north] at (5.5, -1.1)
  {At 37~dB PSNR (720p ``good'' quality), each generation
  requires fewer bits.\\
  The gap to the Shannon bound $R(D)$ narrows but never closes.};

\end{tikzpicture}
\caption{Rate-distortion curves for six codecs at 720p/30fps, with the
  Shannon $R(D)$ bound (green). At constant quality (37~dB dashed line),
  the rate progression from H.264 through VVC illustrates the BD-Rate
  savings in Table~\ref{tab:bd-rate}. This figure extends
  Figure~\ref{fig:rate-distortion} from Appendix~\ref{app:info-theory}
  with the full codec lineage.}
\label{fig:rd-curves-detailed}
\end{figure}

\begin{table}[ht]
\centering
\small
\renewcommand{\arraystretch}{1.4}
\begin{tabularx}{\textwidth}{lXXXX}
\toprule
\textbf{Codec} & \textbf{BD-Rate vs.\ H.264}
  & \textbf{$\Delta$ (720p)} & \textbf{Family} & \textbf{Year} \\
\midrule
VP8  & ${\sim}$0\%    & ${\sim}$2.1 & Google    & 2010 \\
VP9  & $-$35\%        & ${\sim}$0.8 & AOM       & 2013 \\
HEVC & $-$40\%        & ${\sim}$0.7 & MPEG/ITU  & 2013 \\
AV1  & $-$54\%        & ${\sim}$0.5 & AOM       & 2018 \\
VVC  & $-$63\%        & ${\sim}$0.3 & MPEG/ITU  & 2020 \\
\bottomrule
\end{tabularx}
\caption{BD-Rate savings and coding efficiency gap $\Delta$
  (eq.~\ref{eq:coding-gap}) for each codec relative to H.264 at 720p/30fps.
  The two families (MPEG/ITU and Google/AOM) achieve comparable efficiency
  within each generation; AV1 slightly trails HEVC, VVC leads AV1 by
  ${\sim}$20\%.}
\label{tab:bd-rate}
\end{table}

Table~\ref{tab:bd-rate} provides the DSP-grounded explanation for the codec
ordering in eq.~(\ref{eq:codec-ordering}) and the savings figures in
eqs.~(\ref{eq:vp9-savings})--(\ref{eq:av1-savings}): each generation adds
tools that reduce $\Delta_{\text{pred}}$, $\Delta_{\text{transform}}$, and
$\Delta_{\text{quant}}$ in (\ref{eq:delta-decomp}), with diminishing
returns as the Shannon bound is approached.

%% ---------------------------------------------------------------------------
\subsection{Scalable Video Coding and Simulcast}
\label{sec:svc-simulcast}

Multi-participant WebRTC sessions require adaptive quality: different
receivers have different bandwidths, and a single encode rate cannot serve
all. Two architectures address this.

\subsubsection{Simulcast}
\label{sec:simulcast}

Simulcast encodes the source at $N$ independent resolution/bitrate
combinations:
\begin{equation}
  R_{\text{simulcast}} = \sum_{i=1}^{N} R_i
  \label{eq:simulcast-rate}
\end{equation}
Each encoding is a fully independent bitstream. A Selective Forwarding
Unit (SFU) selects which stream to forward to each receiver based on
reported bandwidth. The overhead is substantial: a typical 3-layer
simulcast (720p/1.5~Mbps, 360p/500~kbps, 180p/150~kbps) requires
2.15~Mbps total---approximately 1.4$\times$ the cost of the highest
layer alone---with no inter-layer compression benefit.

\subsubsection{Scalable Video Coding (SVC)}
\label{sec:svc}

SVC produces a layered bitstream where a base layer $L_0$ is independently
decodable and enhancement layers $L_1, \ldots, L_K$ refine quality
progressively. Three scalability dimensions exist:

\textbf{Temporal scalability.} The base layer encodes at a reduced frame
rate (e.g., 7.5~fps); enhancement layers add intermediate frames:
\begin{equation}
  R_{\text{temporal}} = R_0 + \sum_{k=1}^{K} \Delta R_k
  \label{eq:temporal-svc}
\end{equation}
where $\Delta R_k \ll R_0$ because enhancement frames reference lower
layers and carry primarily motion-compensated residuals. An SFU drops
enhancement layers for bandwidth-constrained receivers, gracefully
degrading frame rate while preserving spatial quality.

\textbf{Spatial scalability.} The base layer encodes at reduced resolution;
enhancement layers add detail via inter-layer prediction:
\begin{equation}
  R_{\text{spatial}} = R_0 + \Delta R_{\text{upscale}}
  \label{eq:spatial-svc}
\end{equation}
where $\Delta R_{\text{upscale}}$ exploits correlation between the
upsampled base layer and the full-resolution source.

\textbf{Quality (SNR) scalability.} Refinement of quantization at the same
resolution, encoded as the difference between coarse and fine
reconstructions.

SVC achieves 20--30\% rate savings over simulcast for equivalent
adaptivity because enhancement layers exploit inter-layer prediction:
\begin{equation}
  R_{\text{SVC}} < R_{\text{simulcast}} \quad
  \text{(typically by 20--30\%)}
  \label{eq:svc-advantage}
\end{equation}

\begin{figure}[ht]
\centering
\begin{tikzpicture}[>=Stealth, scale=0.85, every node/.style={scale=0.85},
  layer/.style={draw, rounded corners=2pt, minimum width=1.8cm,
    minimum height=0.6cm, font=\scriptsize, align=center, thick}]

% --- Simulcast (left) ---
\node[font=\normalsize\bfseries] at (2, 5.8) {Simulcast};

% Three independent columns
\node[layer, fill=red!15] (s720) at (0.5, 4.5) {720p\\1.5 Mbps};
\node[layer, fill=blue!15] (s360) at (2.0, 4.5) {360p\\500 kbps};
\node[layer, fill=green!15] (s180) at (3.5, 4.5) {180p\\150 kbps};

\node[layer, fill=gray!10] (sfu1) at (2.0, 2.5) {SFU};

\draw[->, thick] (s720) -- (sfu1);
\draw[->, thick] (s360) -- (sfu1);
\draw[->, thick] (s180) -- (sfu1);

\node[font=\scriptsize, text=red!60!black] at (2.0, 1.5)
  {Total: 2.15 Mbps};
\node[font=\tiny\itshape, text=black!50] at (2.0, 1.0)
  {No inter-layer benefit};

% --- SVC (right) ---
\node[font=\normalsize\bfseries] at (8.5, 5.8) {SVC};

% Pyramid
\node[layer, fill=green!15, minimum width=2.0cm] (t0) at (8.5, 2.0)
  {$T_0$: Base\\7.5 fps};
\node[layer, fill=blue!15, minimum width=2.8cm] (t1) at (8.5, 3.2)
  {$T_1$: Enhancement\\15 fps};
\node[layer, fill=red!15, minimum width=3.6cm] (t2) at (8.5, 4.4)
  {$T_2$: Enhancement\\30 fps};

% Dependency arrows
\draw[->, thick, green!50!black] (t0) -- (t1);
\draw[->, thick, blue!50!black] (t1) -- (t2);

% Drop annotations
\draw[thick, dashed, red!40] (6.2, 3.8) -- (10.8, 3.8);
\node[font=\tiny\itshape, text=red!50!black, anchor=west] at (10.9, 3.8)
  {SFU can drop};

\node[font=\scriptsize, text=green!50!black] at (8.5, 1.0)
  {Total: ${\sim}$1.7 Mbps};
\node[font=\tiny\itshape, text=black!50] at (8.5, 0.5)
  {20--30\% savings via};
\node[font=\tiny\itshape, text=black!50] at (8.5, 0.1)
  {inter-layer prediction};

\end{tikzpicture}
\caption{Simulcast (left) encodes three independent bitstreams with no
  inter-layer compression. SVC (right) encodes a layered bitstream where
  temporal enhancement layers ($T_1$, $T_2$) reference the base layer
  ($T_0$), achieving 20--30\% savings. The SFU drops enhancement layers
  for bandwidth-constrained receivers.}
\label{fig:svc-architecture}
\end{figure}

Codec-specific SVC support relevant to WebRTC:

\begin{itemize}
\item \textbf{H.264 SVC (Annex~G).} Defines temporal, spatial, and quality
  scalability. Complex to implement; rarely used in WebRTC deployments.

\item \textbf{VP9 SVC.} Temporal scalability is widely deployed in WebRTC
  via libvpx (up to 3~temporal layers). Spatial SVC is defined but
  implementation-limited. VP9 temporal SVC is the primary mechanism
  enabling adaptive quality in Chrome/Firefox SFU-based calls.

\item \textbf{AV1 SVC.} Temporal scalability supported in both libaom and
  SVT-AV1. A key enabler for AV1 adoption in production WebRTC
  infrastructure, as SFUs require layered bitstreams to serve heterogeneous
  receivers.
\end{itemize}

Safari's lack of simulcast and SVC support in 2021 (Section~3.1) meant
that Safari users in multi-participant calls could not benefit from adaptive
quality---every receiver got the same stream regardless of bandwidth,
degrading the experience for all participants.

%% ---------------------------------------------------------------------------
\subsection{Hardware vs.\ Software Encoding}
\label{sec:hw-sw}

The main paper identifies hardware encoding as the deployment gate for
AV1 (Sections~8.5, 9.5). This section quantifies why.

\subsubsection{Encoding Complexity}

The computational cost of encoding a video frame is dominated by motion
estimation and RD-optimized mode decision. We express complexity as
operations per pixel:
\begin{equation}
  C_{\text{enc}} = C_{\text{ME}} + C_{\text{transform}}
    + C_{\text{quant}} + C_{\text{entropy}}
    + C_{\text{mode}}
  \label{eq:enc-complexity}
\end{equation}

Representative values for real-time presets:

\begin{itemize}
\item H.264 Constrained Baseline (``veryfast''): $C_{\text{enc}} \approx
  200$~ops/pixel
\item VP9 (speed~8, real-time): $C_{\text{enc}} \approx 600$~ops/pixel
\item AV1 (speed~10, real-time): $C_{\text{enc}} \approx 2000$~ops/pixel
\end{itemize}

The real-time constraint requires:
\begin{equation}
  C_{\text{enc}} \times W \times H \times \text{fps}
    \leq F_{\text{available}}
  \label{eq:realtime-constraint}
\end{equation}
where $F_{\text{available}}$ is the available computational throughput
(operations per second). For 720p/30fps AV1 at 2000~ops/pixel:
\begin{equation}
  2000 \times 921{,}600 \times 30 = 55.3 \times 10^9 \;\text{ops/s}
  \label{eq:av1-cost}
\end{equation}
This consumes approximately one full core of a modern desktop CPU. On
mobile devices with thermal and power constraints, software AV1 encoding
at 720p/30fps remains impractical without hardware acceleration.

\subsubsection{Hardware Encoder Architecture}

Hardware encoders implement a fixed subset of the codec's tool set in
dedicated silicon. The trade-off is quality for throughput: hardware
encoders typically produce output 10--20\% larger than software encoders
at equivalent visual quality because they implement fewer mode decision
candidates and simpler motion search:
\begin{equation}
  \Delta_{\text{HW}} \approx (1.1 \text{--} 1.2)
    \cdot \Delta_{\text{SW}}
  \label{eq:hw-quality-gap}
\end{equation}

This quality penalty is acceptable for real-time communication because:
(a)~the alternative---software encoding at a lower speed preset or reduced
resolution---incurs a larger quality penalty, and (b)~the freed CPU
capacity can be used for pre/post-processing, noise reduction, or
additional application logic.

\subsubsection{Hardware Encoder Availability}

\begin{table}[ht]
\centering
\small
\renewcommand{\arraystretch}{1.4}
\begin{tabularx}{\textwidth}{lXXXX}
\toprule
\textbf{Platform} & \textbf{H.264 HW} & \textbf{VP9 HW encode}
  & \textbf{AV1 HW encode} & \textbf{AV1 year} \\
\midrule
Intel (QSV) &
  Sandy Bridge+ (2011) &
  Kaby Lake+ (2017) &
  Arc / Meteor Lake+ &
  2022 \\
NVIDIA (NVENC) &
  Kepler+ (2012) &
  --- &
  Ada Lovelace+ &
  2022 \\
AMD (VCN) &
  VCN~1.0+ (2017) &
  --- &
  VCN~4.0+ (RDNA~3) &
  2023 \\
Apple (VideoToolbox) &
  A4+ (2010) &
  --- &
  M3+ &
  2023 \\
Qualcomm &
  SD~820+ (2016) &
  --- &
  SD~8 Gen~2+ &
  2023 \\
\bottomrule
\end{tabularx}
\caption{Hardware encoder availability timeline. H.264 hardware encoding
  has been ubiquitous for over a decade. AV1 hardware encoding began
  shipping in 2022--2023 across all major platforms; the main paper projects
  it becomes the practical WebRTC default by 2028--2030 as these chipsets
  reach market saturation.}
\label{tab:hw-encoders}
\end{table}

\excursus{Excursus: WebCodecs and hardware acceleration.}{When a WebCodecs
\texttt{VideoEncoder} is configured with \texttt{codec: "av01"}, the browser
selects hardware or software encoding based on platform capability and the
\texttt{hardwareAcceleration} preference (\texttt{"prefer-hardware"},
\texttt{"prefer-software"}, or \texttt{"no-preference"}). The application
receives encoded \texttt{EncodedVideoChunk} objects either way, but latency
and CPU cost differ dramatically: hardware encoding typically adds
$<\!1$~ms per frame and consumes near-zero CPU, while software AV1 at
real-time presets consumes an entire CPU core
(eq.~\ref{eq:av1-cost}). This abstraction is what makes the hardware
timeline in Table~\ref{tab:hw-encoders} directly relevant to the WebCodecs
adoption curve in Section~7.}

The hardware encoding timeline directly underpins the main paper's
projections: the ``AV1 becomes practical'' milestone (Section~8.5) and
``AV1 HW ubiquitous'' milestone (Section~9.5) depend on the chipset
generations in Table~\ref{tab:hw-encoders} reaching market saturation.

%% ---------------------------------------------------------------------------
\subsection{Implications for Real-Time Web Communication}
\label{sec:video-implications}

The DSP analysis of this appendix maps back to concrete implications for
the technologies surveyed in the main text.

\textbf{The hybrid coding framework is universal; innovation is
incremental.} Every codec from VP8 to VVC implements the same
prediction--transform--quantization--entropy framework
(Figure~\ref{fig:hybrid-encoder}). Generational improvements come from
expanding the tool set within this framework---more block sizes
(Table~\ref{tab:partitioning}), more prediction modes, more transform
types (Table~\ref{tab:transforms})---not from fundamentally new
architectures. This means the efficiency gap $\Delta$ approaches but never
reaches zero: each additional tool yields diminishing rate-distortion
improvement at increasing computational cost.

\textbf{Transform and prediction innovations explain the BD-Rate
progression.} The 35\% rate saving from H.264 to VP9 comes primarily from
larger blocks ($64\!\times\!64$ vs.\ $16\!\times\!16$) and adaptive
transforms (DCT$+$ADST vs.\ fixed DCT). The additional 30\% from VP9 to
AV1 comes from compound and warped motion, film grain synthesis, and
16-way transform type selection (Table~\ref{tab:bd-rate}). Each tool adds
encoding complexity proportional to its rate-distortion benefit, which is
why real-time encoders must aggressively prune the tool set.

\textbf{SVC is what makes multi-participant WebRTC work.} Without temporal
SVC, an SFU must forward full-rate streams to all participants regardless
of their bandwidth. With VP9/AV1 temporal SVC
(Section~\ref{sec:svc}), the SFU drops enhancement temporal layers for
bandwidth-constrained receivers, degrading frame rate gracefully while
preserving spatial quality. The 20--30\% savings over simulcast
(eq.~\ref{eq:svc-advantage}) translate directly to reduced server egress
bandwidth in large meetings.

\textbf{Hardware encoding is the deployment gate, not codec specification.}
AV1 was specified in 2018 but becomes a practical WebRTC default only
when hardware encoders are ubiquitous (${\sim}$2028--2030). The DSP tools
that make AV1 efficient---warped motion, OBMC, $128\!\times\!128$
superblocks with flexible partitioning---are precisely the tools that make
it computationally expensive in software (eq.~\ref{eq:av1-cost}). The
${\sim}$10--20\% quality penalty of hardware encoders
(eq.~\ref{eq:hw-quality-gap}) is the price of real-time feasibility, and
it is a price worth paying: $\Delta_{\text{AV1,HW}} \approx 0.6$ is still
far below $\Delta_{\text{VP9,SW}} \approx 0.8$.

\textbf{The convergence toward Shannon is real but asymptotic.}
Cross-referencing with Appendix~\ref{app:info-theory}: H.264's
$\Delta \approx 2.1$ means it operates at $3.1\times$ the theoretical
minimum rate. VP9's $\Delta \approx 0.8$ means $1.8\times$. AV1's
$\Delta \approx 0.5$ means $1.5\times$. VVC's $\Delta \approx 0.3$ means
$1.3\times$. Each generation captures more spatial, temporal, and
perceptual redundancy, but the returns are diminishing: the distance from
H.264 to VP9 ($-$35\%) exceeds VP9 to AV1 ($-$30\%), which exceeds AV1 to
VVC ($-$17\%). The next generation of efficiency gains for real-time web
communication will come not from codecs alone, but from the co-optimization
of transport (QUIC, WebTransport), application architecture (WebCodecs,
unbundled pipelines), and network infrastructure that the main text
describes.

%
% Requires in the preamble: amsmath, amssymb, tabularx, booktabs

\section{Video Compression Fundamentals}
\label{app:video-compression}

Appendix~\ref{app:info-theory} quantified the coding efficiency gap $\Delta$
(eq.~\ref{eq:coding-gap}) for each codec generation, treating codecs as black
boxes characterized by their distance from the Shannon rate-distortion bound.
This appendix opens those black boxes. We examine the signal processing and
coding techniques that determine $\Delta$: how video codecs exploit spatial,
temporal, and perceptual redundancy, what distinguishes one codec generation
from the next, and why computational complexity---not algorithmic
capability---remains the binding constraint for real-time web communication.

%% ---------------------------------------------------------------------------
\subsection{Why Video Is Compressible}
\label{sec:compressible}

A raw digital video signal carries far more data than information. Three
fundamental redundancies make compression ratios of 200:1 to 500:1 achievable
at perceptually acceptable quality.

\subsubsection{Spatial Redundancy}

Natural images exhibit strong local correlation. The autocorrelation function
for a typical image row can be modeled as a first-order Markov process:
\begin{equation}
  R_{xx}(k) = \sigma^2 \rho^{|k|}
  \label{eq:spatial-autocorr}
\end{equation}
where $\sigma^2$ is the pixel variance and $\rho \in [0.90, 0.97]$ for
natural content. The separable two-dimensional extension is:
\begin{equation}
  R_{xx}(k, l) = \sigma^2 \,\rho_h^{|k|}\, \rho_v^{|l|}
  \label{eq:spatial-autocorr-2d}
\end{equation}
The power spectral density---the Fourier transform of
(\ref{eq:spatial-autocorr-2d})---is concentrated at low spatial frequencies,
meaning most of the signal energy occupies a small fraction of the frequency
domain. This concentration is the foundation of transform coding
(Section~\ref{sec:transforms}).

\subsubsection{Temporal Redundancy}

Consecutive video frames are highly correlated. For typical conversational
video at 30~fps, the inter-frame correlation coefficient is:
\begin{equation}
  \text{Corr}(f_t, f_{t-1}) \approx 0.90\text{--}0.98
  \label{eq:temporal-corr}
\end{equation}
Most pixels change little between frames; only regions containing motion,
lighting changes, or scene cuts carry significant new information. Motion
estimation (Section~\ref{sec:motion}) exploits this by predicting each block
from a displaced version of a reference frame, transmitting only the
prediction residual.

\subsubsection{Perceptual Irrelevance}

The human visual system (HVS) has lower spatial acuity for chrominance than
luminance. This asymmetry motivates \emph{chroma subsampling}: representing
color difference channels at reduced resolution relative to the luma channel.

\excursus{Excursus: Why 4:2:0?}{All WebRTC codecs operate on YCbCr 4:2:0
video, where each chroma channel is sampled at half the horizontal and half
the vertical resolution of luma. This reduces the sample count by a factor of
$1 + 0.25 + 0.25 = 1.5$ relative to $1 + 1 + 1 = 3$ for full-resolution
4:4:4 color, a 2$\times$ saving before any coding begins. Perceptual
experiments confirm that 4:2:0 is essentially indistinguishable from 4:4:4 at
typical viewing distances for video conferencing---the HVS cannot resolve
the missing chroma detail.}

\subsubsection{The Compression Imperative}

A raw 720p/30fps 4:2:0 video stream at 8~bits per sample requires:
\begin{equation}
  R_{\text{raw}} = 1280 \times 720 \times 1.5 \times 8 \times 30
    = 331.8 \;\text{Mbps}
  \label{eq:raw-rate}
\end{equation}
This exceeds the 120~Mbps median link capacity cited in
Table~\ref{tab:inet}. The Shannon rate-distortion bound
(eq.~\ref{eq:rate-distortion}) yields approximately 0.7--1.5~Mbps at
perceptually acceptable distortion---a compression ratio of 220:1 to 475:1.
Video compression is not optional for real-time communication; it is a
prerequisite.

The first-order entropy of raw pixel values $H(X)$ greatly exceeds the
conditional entropy given neighboring samples:
\begin{equation}
  H(X) \gg H(X \mid X_{\text{neighbors}})
  \label{eq:conditional-entropy}
\end{equation}
This gap---the mutual information between a pixel and its
neighbors---represents the exploitable redundancy. The entire hybrid coding
framework (Section~\ref{sec:hybrid}) is designed to extract this mutual
information through prediction, leaving only the irreducible residual for
transform coding and quantization.

%% ---------------------------------------------------------------------------
\subsection{The Hybrid Coding Framework}
\label{sec:hybrid}

Every major video codec since H.261 (1988)---including every codec in the
WebRTC stack---is a variation on the \emph{block-based hybrid coding}
framework. The encoder partitions each frame into blocks, predicts each block
from previously coded data, transforms and quantizes the prediction residual,
and entropy codes the result.

\subsubsection{Encoder Signal Flow}

For a block at position $(x,y)$ in frame $f_t$, the encoder computes a
prediction $\hat{f}_{\text{pred}}(x,y)$ from either the current frame
(intra prediction) or a reference frame (inter prediction). The prediction
residual is:
\begin{equation}
  e(x,y) = f_t(x,y) - \hat{f}_{\text{pred}}(x,y)
  \label{eq:residual}
\end{equation}
The residual is transformed, quantized, and entropy coded to produce the
bitstream. The decoder reconstructs the frame as:
\begin{equation}
  \hat{f}(x,y) = \hat{e}(x,y) + \hat{f}_{\text{pred}}(x,y)
  \label{eq:reconstruction}
\end{equation}
where $\hat{e}$ is the inverse-transformed, dequantized residual.

\subsubsection{Rate-Distortion Optimization}

Every coding decision---block partition, prediction mode, motion vector,
transform type, quantization level---is made by minimizing a Lagrangian
rate-distortion cost:
\begin{equation}
  J(m) = D(m) + \lambda \cdot R(m)
  \label{eq:rd-cost}
\end{equation}
where $m \in \mathcal{M}$ is a coding mode, $D(m)$ is the distortion
(typically sum of squared errors), $R(m)$ is the bit cost, and $\lambda$ is a
Lagrange multiplier controlling the rate-distortion trade-off. The optimal
mode satisfies:
\begin{equation}
  m^* = \arg\min_{m \in \mathcal{M}} \left[ D(m) + \lambda \cdot R(m) \right]
  \label{eq:rd-opt}
\end{equation}
The Lagrange multiplier relates to the quantization parameter (QP) via an
empirically validated exponential relationship:
\begin{equation}
  \lambda = c \cdot 2^{(QP - 12)/3}
  \label{eq:lambda-qp}
\end{equation}
where $c$ is a codec-specific constant. This relationship ensures that
finer quantization (lower QP) assigns higher cost to bits, favoring modes
that reduce distortion, while coarser quantization (higher QP) favors
modes that reduce rate.

\begin{figure}[ht]
\centering
\begin{tikzpicture}[>=Stealth, scale=0.85, every node/.style={scale=0.85},
  block/.style={draw, rounded corners=3pt, minimum width=1.6cm,
    minimum height=0.8cm, font=\scriptsize, align=center, thick},
  arrow/.style={->, thick},
  label/.style={font=\tiny, text=black!60}]

% --- Input ---
\node[block, fill=blue!15] (input) at (0, 0) {Input\\frame $f_t$};

% --- Subtract ---
\node[draw, circle, inner sep=1.5pt, thick] (sub) at (2.5, 0) {$-$};

% --- Transform ---
\node[block, fill=orange!20] (transform) at (5, 0) {Transform\\$T$};

% --- Quantize ---
\node[block, fill=orange!20] (quant) at (7.5, 0) {Quantize\\$Q$};

% --- Entropy ---
\node[block, fill=purple!20] (entropy) at (10.5, 0) {Entropy\\coder};

% --- Bitstream ---
\node[font=\scriptsize, right=0.4cm of entropy] (bitstream) {Bitstream};

% --- Inverse Quantize ---
\node[block, fill=orange!20] (iquant) at (7.5, -2) {Inverse\\$Q^{-1}$};

% --- Inverse Transform ---
\node[block, fill=orange!20] (itransform) at (5, -2) {Inverse\\$T^{-1}$};

% --- Add ---
\node[draw, circle, inner sep=1.5pt, thick] (add) at (2.5, -2) {$+$};

% --- In-loop filter ---
\node[block, fill=red!15] (filter) at (2.5, -3.8) {In-loop\\filter};

% --- Reference buffer ---
\node[block, fill=red!15] (refbuf) at (5.5, -3.8) {Reference\\buffer};

% --- Mode decision ---
\node[block, fill=green!15] (mode) at (8.5, -3.8) {Mode\\decision};

% --- Prediction ---
\node[block, fill=green!15] (pred) at (5.5, -5.5) {Prediction\\(intra/inter)};

% --- Forward path ---
\draw[arrow] (input) -- (sub);
\draw[arrow] (sub) -- node[above, label] {$e$} (transform);
\draw[arrow] (transform) -- (quant);
\draw[arrow] (quant) -- (entropy);
\draw[arrow] (entropy) -- (bitstream);

% --- Reconstruction path ---
\draw[arrow] (quant) -- (iquant);
\draw[arrow] (iquant) -- (itransform);
\draw[arrow] (itransform) -- node[above, label] {$\hat{e}$} (add);

% --- Filter and buffer ---
\draw[arrow] (add) -- (filter);
\draw[arrow] (filter) -- (refbuf);

% --- Prediction feedback ---
\draw[arrow] (refbuf) -- (mode);
\draw[arrow] (mode) -- (pred);
\draw[arrow] (pred) -| (sub);
\draw[arrow] (pred) -| (add);

% --- Labels ---
\node[label, anchor=west] at (0, -5.5) {$\hat{f}_{\text{pred}}$};

\end{tikzpicture}
\caption{Block-based hybrid encoder architecture. Every codec from H.264
  through VVC follows this structure: predict each block from previously
  coded data (green), transform and quantize the residual (orange), entropy
  code the result (purple), then reconstruct and filter for future
  reference (red). The coding efficiency gap $\Delta$
  (eq.~\ref{eq:coding-gap}) arises from sub-optimal prediction, imperfect
  energy compaction, quantization beyond the rate-distortion bound, and
  entropy coding overhead.}
\label{fig:hybrid-encoder}
\end{figure}

The four sub-optimal components of Figure~\ref{fig:hybrid-encoder}---prediction,
transform, quantization, and entropy coding---each contribute to $\Delta$:
\begin{equation}
  \Delta \approx \Delta_{\text{pred}} + \Delta_{\text{transform}}
    + \Delta_{\text{quant}} + \Delta_{\text{entropy}}
  \label{eq:delta-decomp}
\end{equation}
Generational codec improvements reduce one or more of these terms. The codec
lineages in Sections~\ref{sec:mpeg-lineage}--\ref{sec:aom-lineage} trace
exactly which tools target which terms.

%% ---------------------------------------------------------------------------
\subsection{Transforms}
\label{sec:transforms}

The transform stage decorrelates the prediction residual, concentrating
energy into a few coefficients that can be efficiently quantized and coded.

\subsubsection{The KLT and Its DCT Approximation}
\label{sec:klt-dct}

The optimal decorrelating transform for a stationary source is the
\emph{Karhunen-Lo\`eve Transform} (KLT), whose basis vectors are the
eigenvectors of the input autocorrelation matrix $\mathbf{R}_{xx}$. For a
first-order Markov source with correlation $\rho$ (eq.~\ref{eq:spatial-autocorr}),
the KLT basis converges to the Type-II DCT as block size grows:
\begin{equation}
  C_k(n) = \alpha_k \cos\!\left[\frac{\pi k (2n+1)}{2N}\right],
  \quad k,n = 0, \ldots, N-1
  \label{eq:dct-ii}
\end{equation}
where $\alpha_0 = \sqrt{1/N}$ and $\alpha_k = \sqrt{2/N}$ for $k > 0$.

The energy compaction efficiency---the fraction of total energy captured by
the first $K$ of $N$ coefficients---is:
\begin{equation}
  \eta(K, N) = \frac{\sum_{k=0}^{K-1} \sigma_k^2}
    {\sum_{k=0}^{N-1} \sigma_k^2}
  \label{eq:energy-compaction}
\end{equation}
where $\sigma_k^2$ is the variance of the $k$-th transform coefficient.
For a typical intra residual with $\rho = 0.95$ and $N = 8$, the first 3
DCT coefficients capture approximately 95\% of the energy. This
concentration enables aggressive quantization of high-frequency coefficients
with minimal perceptual impact.

\subsubsection{Integer Transforms in Practice}
\label{sec:integer-transforms}

Practical codecs use integer approximations of the DCT to ensure bit-exact
encoder-decoder agreement. H.264 introduced a $4 \times 4$ integer
transform:
\begin{equation}
  \mathbf{H}_4 = \begin{pmatrix}
    1 &  1 &  1 &  1 \\
    2 &  1 & -1 & -2 \\
    1 & -1 & -1 &  1 \\
    1 & -2 &  2 & -1
  \end{pmatrix}
  \label{eq:h264-transform}
\end{equation}
The non-unity scaling factors are absorbed into the quantization step,
preserving exact integer arithmetic throughout the encode-decode loop.

Each codec generation expands the transform toolkit:

\begin{table}[ht]
\centering
\small
\renewcommand{\arraystretch}{1.4}
\begin{tabularx}{\textwidth}{lXXX}
\toprule
\textbf{Codec} & \textbf{Transform sizes} & \textbf{Transform types}
  & \textbf{Selection method} \\
\midrule
H.264 &
  $4\!\times\!4$, $8\!\times\!8$ &
  Integer DCT &
  Fixed per block size \\
VP9 &
  $4\!\times\!4$ to $32\!\times\!32$ &
  DCT, ADST &
  Row/column adaptive \\
HEVC &
  $4\!\times\!4$ to $32\!\times\!32$ &
  Integer DCT, DST ($4\!\times\!4$ intra) &
  Mode-dependent \\
AV1 &
  $4\!\times\!4$ to $64\!\times\!64$ &
  DCT, ADST, FlipADST, Identity &
  RD-optimized per block \\
VVC &
  $4\!\times\!4$ to $64\!\times\!64$ &
  DCT-II, DST-VII $+$ LFNST &
  RD-optimized $+$ secondary \\
\bottomrule
\end{tabularx}
\caption{Transform sizes and types across codec generations. AV1 applies
  four 1D transform types independently to rows and columns, yielding
  16~possible 2D combinations per block. VVC adds a non-separable secondary
  transform (LFNST) applied to low-frequency coefficients after the primary
  separable transform.}
\label{tab:transforms}
\end{table}

The Asymmetric Discrete Sine Transform (ADST) used in VP9 and AV1 is
particularly effective for intra-predicted residuals, which tend to have
larger magnitudes near the prediction boundary. AV1's FlipADST extends
this to blocks where the large residuals are on the opposite edge, while
the Identity transform (skip) handles blocks where the residual is already
well-distributed in the spatial domain.

The progression from fixed transforms (H.264) to RD-optimized multi-type
selection (AV1, VVC) directly reduces $\Delta_{\text{transform}}$ in
(\ref{eq:delta-decomp}): better energy compaction means fewer bits are
needed to represent the same residual at a given distortion level.

%% ---------------------------------------------------------------------------
\subsection{Motion Estimation and Inter Prediction}
\label{sec:motion}

Inter prediction exploits temporal redundancy by predicting blocks in the
current frame from previously coded reference frames. This is the single
largest source of compression gain---the prediction residual
(eq.~\ref{eq:residual}) typically has 10--20$\times$ less energy than the
original signal.

\subsubsection{Block Matching and Motion Vectors}
\label{sec:block-matching}

The encoder searches for the best-matching region in a reference frame by
minimizing a cost function over candidate displacement vectors
$\mathbf{d} = (d_x, d_y)$:
\begin{equation}
  \text{SAD}(\mathbf{d}) = \sum_{(x,y) \in B}
    \left| f_t(x,y) - f_{t-1}(x + d_x,\, y + d_y) \right|
  \label{eq:sad}
\end{equation}
The Sum of Absolute Differences (SAD) is the simplest distortion measure.
A better rate-distortion proxy is the Sum of Absolute Transformed
Differences (SATD), which applies the Hadamard transform before summation:
\begin{equation}
  \text{SATD}(\mathbf{d}) = \sum \left| \mathbf{H}
    \cdot \left( f_t - f_{t-1}(\mathbf{d}) \right) \right|
  \label{eq:satd}
\end{equation}
SATD approximates the post-transform energy and thus better predicts the
actual coded bit cost. Rate-constrained motion estimation augments the
distortion with a motion vector coding cost:
\begin{equation}
  \mathbf{d}^* = \arg\min_{\mathbf{d}} \left[
    \text{SATD}(\mathbf{d}) + \lambda_{\text{motion}} \cdot R(\mathbf{d})
  \right]
  \label{eq:rcme}
\end{equation}
where $R(\mathbf{d})$ is the bit cost of encoding the motion vector
relative to a predictor.

\subsubsection{Sub-Pixel Interpolation}
\label{sec:subpixel}

Motion in natural video rarely aligns to integer pixel positions. All modern
codecs perform sub-pixel motion compensation by interpolating reference
frame samples at fractional positions.

H.264 uses a 6-tap filter for half-pixel luma interpolation:
\begin{equation}
  h_{\frac{1}{2}} = \{1, -5, 20, 20, -5, 1\} / 32
  \label{eq:h264-halfpel}
\end{equation}
with quarter-pixel positions obtained by bilinear averaging of integer and
half-pixel samples. HEVC extends this to an 8-tap filter for half-pixel and
a 7-tap filter for quarter-pixel positions. AV1 introduces a switchable
filter bank with three 8-tap options---Regular (sharp), Smooth (low-pass),
and Sharp (high-frequency preserving)---selected per block via RD
optimization. VVC adds DMVR (Decoder-side Motion Vector Refinement) and
BDOF (Bi-Directional Optical Flow) for sub-pixel refinement at the decoder
without additional signaling bits.

The progression from H.264's fixed 6-tap filter to AV1's switchable bank
reduces $\Delta_{\text{pred}}$: better interpolation means smaller
prediction residuals, which require fewer bits.

\subsubsection{Advanced Inter Prediction Tools}
\label{sec:advanced-inter}

Beyond basic block matching, each codec generation adds tools to handle
complex motion and occlusion:

\begin{itemize}
\item \textbf{Multiple reference frames.} H.264 introduced up to 16
  reference frames; practical WebRTC encoders use 1--3 for latency reasons.
  More references improve prediction for periodic motion and uncovered
  backgrounds.

\item \textbf{Bi-prediction.} A weighted average of predictions from two
  reference frames, reducing noise and handling gradual scene changes.
  Available from H.264 onward (in B-frames) and in AV1's compound modes.

\item \textbf{Compound prediction (AV1).} Wedge masks and
  difference-weighted blending combine two predictions at the block level,
  handling occlusion boundaries where a single motion model fails.

\item \textbf{Overlapped Block Motion Compensation (AV1).} Blends
  predictions from neighboring blocks at motion boundaries, reducing
  blocking artifacts without post-processing.

\item \textbf{Warped motion (AV1).} An affine model with up to 6
  parameters handles rotation, zoom, and perspective---common in
  screen sharing and camera motion.

\item \textbf{VVC additions.} Geometric Partitioning Mode (non-rectangular
  block splits for motion boundaries), BDOF (optical flow refinement at the
  decoder), and DMVR (decoder-side MV refinement using bilateral matching).
\end{itemize}

\excursus{Excursus: Motion estimation in real-time encoding.}{Full-search
block matching has complexity $O(N^2 \cdot S^2)$ where $N$ is the block
dimension and $S$ is the search range in pixels. For a $64\!\times\!64$
block with $S = 64$, this is ${\sim}17$ million SAD operations per
block---far too expensive for real-time encoding. WebRTC encoders use
hierarchical search patterns (diamond, hexagon) that reduce complexity to
$O(N^2 \cdot \log S)$ with typically $< 5\%$ rate-distortion penalty
relative to full search. This is why real-time encoders use only a subset
of each codec's inter prediction tools: the full tool set is available
for offline encoding, but the real-time constraint forces aggressive
pruning of the mode decision search space.}

%% ---------------------------------------------------------------------------
\subsection{Quantization}
\label{sec:quantization}

Quantization is the sole lossy step in the hybrid coding pipeline. It maps
continuous-valued (or high-precision integer) transform coefficients to a
discrete set of reconstruction levels, introducing controlled information
loss to achieve the target bit rate.

\subsubsection{Scalar Quantization}

Uniform scalar quantization maps coefficient $c$ to a quantization index:
\begin{equation}
  q = \text{round}\!\left(\frac{c}{Q_{\text{step}}}\right), \quad
  \hat{c} = q \cdot Q_{\text{step}}
  \label{eq:uniform-quant}
\end{equation}
Modern codecs use \emph{dead-zone quantization}, where coefficients near
zero are mapped to zero more aggressively than uniform rounding predicts:
\begin{equation}
  q = \text{sign}(c) \cdot \left\lfloor
    \frac{|c| - f \cdot Q_{\text{step}}}{Q_{\text{step}}}
  \right\rfloor
  \label{eq:deadzone}
\end{equation}
where $f \in [0, 0.5)$ controls the dead-zone width. Dead-zone
quantization increases sparsity---the fraction of zero-valued
coefficients---which dramatically improves entropy coding efficiency.

\subsubsection{QP-to-Step-Size Mapping}

The quantization parameter (QP) controls the step size via an exponential
relationship. In the H.264/HEVC/VVC family:
\begin{equation}
  Q_{\text{step}} = Q_0 \cdot 2^{(QP - QP_0)/6}
  \label{eq:qp-mapping}
\end{equation}
A 6-unit increase in QP doubles the step size, approximately doubling the
bit rate. AV1 uses a similar but codec-specific mapping with separate QP
values for DC and AC coefficients, plus per-segment QP adjustment for
region-of-interest coding.

\subsubsection{Quantization Distortion}

For a single coefficient under uniform quantization at high rate, the
distortion is:
\begin{equation}
  D_q = \frac{Q_{\text{step}}^2}{12}
  \label{eq:quant-distortion}
\end{equation}
For a block of $N$ transform coefficients where $K$ survive quantization
(nonzero) and $N - K$ are quantized to zero:
\begin{equation}
  D_{\text{block}} \approx K \cdot \frac{Q_{\text{step}}^2}{12}
    + \sum_{k \notin K} c_k^2
  \label{eq:block-distortion}
\end{equation}
The second term represents the energy of coefficients eliminated by the
dead zone---information permanently discarded. The QP is thus the primary
control knob for navigating the rate-distortion curve in
Figure~\ref{fig:rate-distortion}. The WebCodecs \texttt{VideoEncoder} API
exposes this control via the \texttt{quantizer} configuration parameter,
giving JavaScript applications direct access to this fundamental
trade-off.

%% ---------------------------------------------------------------------------
\subsection{Entropy Coding}
\label{sec:entropy-coding}

After quantization, the surviving coefficients and all side information
(motion vectors, prediction modes, block partitions) must be losslessly
compressed. The entropy coder converts these symbols into a bitstream,
approaching the conditional entropy $H(X|\text{context})$ as closely as
possible. The gap between the entropy coder's output rate and the true
conditional entropy is $\Delta_{\text{entropy}}$ in
(\ref{eq:delta-decomp}).

\subsubsection{Context-Adaptive Binary Arithmetic Coding (CABAC)}
\label{sec:cabac}

CABAC, used in H.264, HEVC, and VVC, encodes each syntax element as a
sequence of binary decisions (bins). For each bin, the arithmetic coder
subdivides the current interval $[L, H)$ based on the estimated
probability $p$:
\begin{equation}
  [L, H) \to \begin{cases}
    [L,\; L + p \cdot (H - L)) & \text{if bin} = 0 \\
    [L + p \cdot (H - L),\; H) & \text{if bin} = 1
  \end{cases}
  \label{eq:arithmetic-coding}
\end{equation}
The theoretical output length approaches the entropy:
\begin{equation}
  \ell \to \sum_{i=1}^{n} -\log_2 p(b_i \mid \text{ctx}_i)
  \label{eq:cabac-length}
\end{equation}
CABAC achieves redundancy below 0.01 bits/symbol relative to the true
conditional entropy---near-optimal compression.

Context modeling is the key to CABAC's effectiveness. Each bin type has an
associated context, whose probability is updated after each observation via
a state-machine-based exponential moving average:
\begin{equation}
  p_{k+1} = p_k + \alpha(s_k) \cdot (b_k - p_k)
  \label{eq:cabac-update}
\end{equation}
where $\alpha(s_k)$ is a state-dependent adaptation rate, indexed by a
lookup table with 64 states (H.264) or 128 states (HEVC/VVC). This design
avoids floating-point arithmetic entirely.

CABAC also provides a \emph{bypass mode} for bins with approximately
uniform probability (e.g., sign bits, high-magnitude coefficient levels),
which skips the context lookup and probability estimation for throughput.

\subsubsection{Multi-Symbol Arithmetic Coding in AV1}
\label{sec:av1-entropy}

AV1 departs from binary arithmetic coding. Its entropy coder operates on
multi-symbol alphabets using CDF (Cumulative Distribution Function) tables:
\begin{equation}
  \text{Encode symbol } s: \quad
  [L, H) \to \left[L + \frac{\text{CDF}(s-1)}{M} \cdot (H-L),\;
    L + \frac{\text{CDF}(s)}{M} \cdot (H-L)\right)
  \label{eq:av1-entropy}
\end{equation}
where $M$ is the total frequency count. CDF tables are adapted after
each symbol:
\begin{equation}
  \text{CDF}_{k+1}(s) = \text{CDF}_k(s) + \left\lfloor
    \frac{\text{target}(s) - \text{CDF}_k(s)}{n}
  \right\rfloor
  \label{eq:cdf-update}
\end{equation}
where $n$ controls the adaptation rate. Multi-symbol coding is more
efficient than CABAC's binarization for syntax elements with moderate
alphabet sizes (e.g., intra prediction modes, transform types), as it
avoids the overhead of converting each symbol into a binary tree of
decisions.

\begin{table}[ht]
\centering
\small
\renewcommand{\arraystretch}{1.4}
\begin{tabularx}{\textwidth}{lXXXXX}
\toprule
\textbf{Feature} & \textbf{H.264} & \textbf{HEVC} & \textbf{VP9}
  & \textbf{AV1} & \textbf{VVC} \\
\midrule
Method &
  Binary AC &
  Binary AC &
  Boolean AC &
  Multi-sym AC &
  Binary AC \\
Contexts &
  ${\sim}400$ &
  ${\sim}200$ &
  ${\sim}900$ &
  ${\sim}300$ &
  ${\sim}700$ \\
Adaptation &
  64-state LUT &
  128-state LUT &
  Counter &
  CDF update &
  128-state LUT \\
Bypass mode &
  Yes &
  Yes &
  N/A &
  N/A &
  Yes \\
Parallelism &
  Slice &
  WPP / Tiles &
  Tiles &
  Tiles &
  WPP / Tiles \\
\bottomrule
\end{tabularx}
\caption{Entropy coding comparison across codecs. HEVC reduced context
  count relative to H.264 for throughput; VVC increased it again for
  improved compression. AV1's multi-symbol approach operates on full
  alphabets rather than binary decompositions.}
\label{tab:entropy-coding}
\end{table}

The frame entropy $\hat{H}(f_t)$ from eq.~(\ref{eq:frame-entropy}) is
ultimately determined by how well the entropy coder's probability model
matches the true symbol distribution. Both CABAC and AV1's multi-symbol
coder achieve near-zero entropy coding gaps; the remaining
$\Delta_{\text{entropy}}$ arises primarily from context model mismatch on
non-stationary content (scene changes, sudden motion onset).

%% ---------------------------------------------------------------------------
\subsection{In-Loop Filtering}
\label{sec:filtering}

Block-based coding introduces visible artifacts at block boundaries:
discontinuities from independent quantization of adjacent blocks, and
ringing from coefficient truncation. In-loop filters reduce these artifacts
\emph{before} the reconstructed frame enters the reference buffer, improving
prediction quality for subsequent frames and thus reducing
$\Delta_{\text{pred}}$.

\begin{itemize}
\item \textbf{Deblocking filter (all codecs).} Smooths discontinuities
  at block edges, with filter strength adapted to the quantization level
  and local gradient. Introduced in H.264; refined in each subsequent
  codec.

\item \textbf{Sample Adaptive Offset---SAO (HEVC).} Classifies
  reconstructed samples by intensity or edge direction and applies
  per-class offsets to reduce systematic bias introduced by quantization.

\item \textbf{Constrained Directional Enhancement Filter---CDEF (AV1).}
  A directional filter that identifies edge orientation per
  $8\!\times\!8$ block and applies filtering only along the edge,
  preserving sharpness while reducing ringing.

\item \textbf{Loop Restoration Filter (AV1).} A switchable per-tile
  filter choosing between a Wiener filter (FIR designed to minimize MSE
  between original and reconstruction) and a Self-Guided filter (non-linear
  denoising based on guided image filtering). This is AV1's most powerful
  quality enhancement tool.

\item \textbf{Adaptive Loop Filter---ALF (VVC).} A diamond-shaped filter
  with coefficients signaled per frame, applied adaptively per
  $4\!\times\!4$ block based on local classification. Subsumes HEVC's
  SAO functionality.
\end{itemize}

The evolution from simple deblocking (H.264) to multi-stage adaptive
filtering (AV1: deblocking $\to$ CDEF $\to$ loop restoration) exemplifies
how each codec generation reduces $\Delta_{\text{pred}}$ by improving
reference frame quality.

%% ---------------------------------------------------------------------------
\subsection{Codec Lineage: MPEG/ITU Family}
\label{sec:mpeg-lineage}

The MPEG/ITU standardization track has produced three codec generations
relevant to the 2021--2031 timeframe of this document.

\subsubsection{H.264/AVC (2003)}
\label{sec:h264}

H.264 defined the modern era of video compression and remains the most
widely deployed codec. Its key innovations over its predecessors (H.263,
MPEG-2):

\begin{itemize}
\item Variable block sizes: $16\!\times\!16$ macroblock partitioned into
  $16\!\times\!8$, $8\!\times\!16$, $8\!\times\!8$, $8\!\times\!4$,
  $4\!\times\!8$, and $4\!\times\!4$ sub-partitions for motion
  compensation.

\item Quarter-pixel motion compensation with 6-tap interpolation
  (eq.~\ref{eq:h264-halfpel}).

\item Context-adaptive intra prediction: 9~angular modes for
  $4\!\times\!4$ blocks, 4~modes for $16\!\times\!16$.

\item Dual entropy coding: CAVLC (simpler, faster) and CABAC
  (Section~\ref{sec:cabac}; 10--15\% better compression).

\item In-loop deblocking filter, mandatory for the decoder.

\item Multiple reference frames (up to 16) and weighted prediction.
\end{itemize}

Three profiles are relevant to WebRTC: \textbf{Constrained Baseline}
(no CABAC, no B-frames, lowest complexity---the mandatory WebRTC profile),
\textbf{Main} (CABAC, B-frames, ${\sim}$10\% better compression), and
\textbf{High} ($8\!\times\!8$ transform, adaptive quantization matrices,
${\sim}$15\% better).

At the 720p/30fps reference point, H.264 Constrained Baseline operates at
approximately 2.0--2.5~Mbps for ``good'' quality (PSNR ${\sim}$37~dB),
corresponding to $\Delta_{\text{H.264}} \approx 2.1$ from
Appendix~\ref{app:info-theory}. H.264 was the only codec Safari supported
in 2021 (Section~3.1), and its higher $\Delta$ is directly responsible
for the bandwidth disparity between Safari-only and VP9-capable sessions
documented in the main text.

\subsubsection{H.265/HEVC (2013)}
\label{sec:hevc}

HEVC replaced H.264's fixed $16\!\times\!16$ macroblock with a recursive
Coding Tree Unit (CTU) of up to $64\!\times\!64$ pixels, using quad-tree
partitioning down to $8\!\times\!8$. Key advances:

\begin{itemize}
\item 35 intra prediction angular directions (vs.\ H.264's 9).
\item Transform sizes from $4\!\times\!4$ to $32\!\times\!32$, with
  DST for $4\!\times\!4$ intra blocks.
\item Advanced Motion Vector Prediction (AMVP): two MVP candidates selected
  from spatial and temporal neighbors, replacing H.264's less structured
  motion vector prediction.
\item Merge mode: inherits the full motion information (MV, reference
  index, prediction direction) from a neighboring block, requiring zero
  additional motion bits.
\item Sample Adaptive Offset (SAO) in-loop filter.
\item Parallel processing: Wavefront Parallel Processing (WPP) and tiles.
\end{itemize}

HEVC achieves approximately 40\% rate reduction over H.264 at equivalent
PSNR. However, HEVC is absent from WebRTC entirely---not due to technical
limitations, but due to patent licensing.

\excursus{Excursus: Why no HEVC in WebRTC?}{HEVC's patent licensing is
fragmented across three separate pools (MPEG-LA, HEVC Advance, Velos Media)
plus independent licensors who are not members of any pool. The resulting
uncertainty about total royalty cost---estimated at \$0.20--\$0.40 per device
for some use cases---led Google, Mozilla, Cisco, and others to form the
Alliance for Open Media (AOM) in 2015 and develop AV1 as a royalty-free
alternative. This patent-driven fork is why the WebRTC ecosystem followed the
VP8$\to$VP9$\to$AV1 path rather than H.264$\to$HEVC$\to$VVC, and why the two
codec families must be understood in parallel.}

\subsubsection{H.266/VVC (2020)}
\label{sec:vvc}

VVC extends HEVC's quad-tree with binary and ternary splits, creating a
\emph{multi-type tree} (MTT) that enables asymmetric block shapes (e.g.,
$32\!\times\!8$, $8\!\times\!32$, $32\!\times\!4$). Additional
innovations:

\begin{itemize}
\item 67 intra prediction directions plus Matrix-based Intra Prediction
  (MIP), which applies a learned linear transform to neighboring samples.

\item Multiple Transform Selection (MTS): explicit signaling of DCT-II or
  DST-VII per block, plus the Low-Frequency Non-Separable Transform
  (LFNST)---a $16\!\times\!16$ or $8\!\times\!8$ non-separable transform
  applied to low-frequency coefficients after the primary separable
  transform.

\item Affine motion model: 4-parameter (translation $+$ rotation/scaling)
  or 6-parameter (full affine) motion, derived from control point MVs at
  block corners.

\item Decoder-side tools: DMVR refines motion vectors at the decoder using
  bilateral template matching; BDOF applies optical flow equations to
  refine bi-predicted samples, both without additional signaling bits.

\item Geometric Partitioning Mode: divides a block diagonally for separate
  motion compensation of each partition, targeting motion boundaries.

\item Adaptive Loop Filter (ALF): a diamond-shaped Wiener filter with
  coefficients signaled per frame and applied adaptively per
  $4\!\times\!4$ block.
\end{itemize}

VVC achieves approximately 30--50\% rate reduction over HEVC, placing it at
$\Delta_{\text{VVC}} \approx 0.3$. Its relevance to this document is
forward-looking: VVC may appear in WebCodecs codec registries by
${\sim}$2030 if licensing concerns are resolved, and its tool set
represents the current efficiency frontier against which AV1's performance
should be measured.

%% ---------------------------------------------------------------------------
\subsection{Codec Lineage: Google/AOM Family}
\label{sec:aom-lineage}

The royalty-free codec family dominates the WebRTC ecosystem. Its
development was driven by the patent licensing concerns described in the
HEVC excursus above.

\subsubsection{VP8 (2010)}
\label{sec:vp8}

VP8 was Google's first open-source video codec, released alongside the
WebM container. Its design is conservative:

\begin{itemize}
\item $16\!\times\!16$ macroblocks with $4\!\times\!4$ subblock partitions.
\item $4\!\times\!4$ Walsh-Hadamard Transform (WHT) for DC coefficients,
  $4\!\times\!4$ DCT for AC coefficients.
\item Boolean arithmetic coder with tree-structured syntax (single-bit
  alphabet).
\item 4 intra prediction modes for $16\!\times\!16$, 10 for
  $4\!\times\!4$ subblocks.
\item Simple loop filter with per-edge strength adaptation.
\item Single reference frame for WebRTC (plus a ``golden'' frame for
  screen content).
\end{itemize}

VP8's compression efficiency is roughly equivalent to H.264 Constrained
Baseline. It served as the baseline WebRTC codec---supported by Chrome and
Firefox since WebRTC~1.0---but has been superseded by VP9 in all modern
deployments.

\subsubsection{VP9 (2013)}
\label{sec:vp9}

VP9 introduced a superblock architecture and several tools that brought it
within range of HEVC:

\begin{itemize}
\item $64\!\times\!64$ superblocks with recursive quad-tree partitioning
  down to $4\!\times\!4$.
\item 10 intra prediction modes.
\item Transforms: DCT and ADST applied independently per row and column,
  at sizes from $4\!\times\!4$ to $32\!\times\!32$.
\item Switchable interpolation filters: 8-tap Regular, Smooth, and
  Sharp, selected per block.
\item Compound prediction: weighted average from two reference frames.
\item Segmentation: per-segment QP adjustment for region-of-interest
  coding.
\end{itemize}

VP9 achieves approximately 35\% rate reduction over H.264
(eq.~\ref{eq:vp9-savings}), placing it at $\Delta_{\text{VP9}} \approx 0.8$.
VP9 is the dominant WebRTC codec in Chrome and Firefox as of March~2026,
and its temporal SVC support (Section~\ref{sec:svc}) is the primary
mechanism enabling adaptive quality in multi-participant calls.

\subsubsection{AV1 (2018)}
\label{sec:av1}

AV1 represents the most ambitious single-generation improvement in the
royalty-free family. Its tool set approaches VVC in breadth while
maintaining royalty-free status:

\begin{itemize}
\item $128\!\times\!128$ superblocks with flexible partitioning (quad-tree
  plus rectangular splits) down to $4\!\times\!4$.

\item 56+ directional intra prediction modes plus non-directional modes
  (Paeth, Smooth variants, DC).

\item Full inter prediction toolkit: OBMC, warped motion (affine),
  compound modes (Wedge, Difference-weighted), inter-intra compound
  prediction.

\item 16 transform type combinations ($4$ 1D types $\times$ rows/columns),
  at sizes up to $64\!\times\!64$.

\item Film Grain Synthesis: the encoder characterizes noise/grain
  parameters and transmits them as metadata; the decoder re-synthesizes
  grain after reconstruction. This removes non-compressible high-frequency
  noise from the coding loop---a major bitrate saving for camera-captured
  content.

\item Multi-stage in-loop filtering: deblocking $\to$ CDEF $\to$ loop
  restoration (Wiener or Self-Guided), as described in
  Section~\ref{sec:filtering}.

\item Multi-symbol arithmetic coding with CDF-based context adaptation
  (Section~\ref{sec:av1-entropy}).
\end{itemize}

AV1 achieves approximately 30\% rate reduction over VP9
(eq.~\ref{eq:av1-savings}), a cumulative ${\sim}$54\% over H.264, placing
it at $\Delta_{\text{AV1}} \approx 0.5$ with current hardware encoders and
approaching 0.4 by 2030 as encoder implementations mature.

AV1's tools directly explain why its $\Delta$ is smaller than VP9's:
more transform types capture residual structure better (reducing
$\Delta_{\text{transform}}$), warped and compound motion reduce prediction
residual energy (reducing $\Delta_{\text{pred}}$), and film grain synthesis
removes non-compressible variance from the coding loop (effectively
reducing the source variance $\sigma^2$ that the codec must represent).

%% ---------------------------------------------------------------------------
\subsection{Block Partitioning: A Generational Comparison}
\label{sec:partitioning}

Block partitioning is the most architecturally consequential difference
between codec generations. Larger and more flexible partitions allow the
encoder to match block boundaries to content structure---smooth regions
get large blocks (few bits for partition overhead), detailed regions get
small blocks (better prediction accuracy).

\begin{figure}[ht]
\centering
\begin{tikzpicture}[>=Stealth, scale=0.55, every node/.style={scale=0.55}]

% --- H.264 panel ---
\begin{scope}[xshift=0cm]
  \node[font=\normalsize\bfseries, anchor=south] at (2, 8.4) {H.264};
  \node[font=\small, anchor=south] at (2, 7.8) {$16\!\times\!16$ max};
  % 4x4 grid of 16x16 macroblocks in a 64x64 region
  \draw[thick] (0,0) rectangle (4,4);
  \draw[thick] (0,0) rectangle (4,4);
  % macroblock grid
  \foreach \x in {0,1,2,3} {
    \foreach \y in {0,1,2,3} {
      \draw[gray] (\x, \y) rectangle (\x+1, \y+1);
    }
  }
  % show some sub-partitions
  \draw[blue!60, thick] (0,3) -- (0.5,3) -- (0.5,4) -- (0,4);
  \draw[blue!60, thick] (0.5,3) -- (0.5,4);
  \draw[blue!60, thick] (0,3.5) -- (1,3.5);
  % 8x8 in another
  \draw[blue!60, thick] (1,2.5) -- (2,2.5);
  \draw[blue!60, thick] (1.5,2) -- (1.5,3);
  % 4x4 in another
  \draw[red!60, thick] (2,0) -- (2,1) -- (3,1) -- (3,0);
  \draw[red!60, thick] (2.5,0) -- (2.5,0.5);
  \draw[red!60, thick] (2,0.5) -- (2.5,0.5);
  \draw[red!60, thick] (2.5,0.5) -- (2.5,1);
  \draw[red!60, thick] (2.5,0) -- (3,0) -- (3,0.5) -- (2.5,0.5);
  % Scale label
  \node[font=\small, anchor=north] at (2, -0.3) {64 pixels};
\end{scope}

% --- VP9/HEVC panel ---
\begin{scope}[xshift=6cm]
  \node[font=\normalsize\bfseries, anchor=south] at (4, 8.4) {VP9 / HEVC};
  \node[font=\small, anchor=south] at (4, 7.8) {$64\!\times\!64$ max};
  % One 64x64 superblock
  \draw[thick] (0,0) rectangle (8,8);
  % quad-tree splits
  \draw[thick] (4,0) -- (4,8);
  \draw[thick] (0,4) -- (8,4);
  % further splits in top-left quadrant
  \draw[blue!60, thick] (2,4) -- (2,8);
  \draw[blue!60, thick] (0,6) -- (4,6);
  % further split in top-left-top-left
  \draw[red!60, thick] (1,6) -- (1,8);
  \draw[red!60, thick] (0,7) -- (2,7);
  % split bottom-right quadrant
  \draw[blue!60, thick] (6,0) -- (6,4);
  \draw[blue!60, thick] (4,2) -- (8,2);
  % deeper in bottom-right-bottom-left
  \draw[red!60, thick] (5,0) -- (5,2);
  \draw[red!60, thick] (4,1) -- (6,1);
  % Scale label
  \node[font=\small, anchor=north] at (4, -0.3) {64 pixels};
\end{scope}

% --- AV1 panel ---
\begin{scope}[xshift=16cm]
  \node[font=\normalsize\bfseries, anchor=south] at (4, 8.4) {AV1};
  \node[font=\small, anchor=south] at (4, 7.8) {$128\!\times\!128$ max};
  % 128x128 shown as 8x8 grid (each cell = 16px)
  \draw[thick] (0,0) rectangle (8,8);
  % quad split: top half / bottom half
  \draw[thick] (4,0) -- (4,8);
  \draw[thick] (0,4) -- (8,4);
  % rectangular splits in top-left
  \draw[blue!60, thick] (0,6) -- (4,6);
  \draw[blue!60, thick] (2,4) -- (2,6);
  % asymmetric in bottom-left (4:1 rect)
  \draw[orange!80, thick] (0,3) -- (4,3);
  \draw[orange!80, thick] (0,2) -- (4,2);
  \draw[orange!80, thick] (0,1) -- (4,1);
  % quad + rect in top-right
  \draw[blue!60, thick] (6,4) -- (6,8);
  \draw[red!60, thick] (4,6) -- (6,6);
  \draw[red!60, thick] (5,4) -- (5,6);
  % bottom-right: large blocks
  \draw[blue!60, thick] (4,2) -- (8,2);
  % Scale label
  \node[font=\small, anchor=north] at (4, -0.3) {128 pixels};
\end{scope}

% --- VVC panel ---
\begin{scope}[xshift=26cm]
  \node[font=\normalsize\bfseries, anchor=south] at (4, 8.4) {VVC};
  \node[font=\small, anchor=south] at (4, 7.8) {$128\!\times\!128$ max};
  \draw[thick] (0,0) rectangle (8,8);
  % quad split
  \draw[thick] (4,0) -- (4,8);
  \draw[thick] (0,4) -- (8,4);
  % binary split (horizontal) in top-left
  \draw[blue!60, thick] (0,6) -- (4,6);
  % ternary split (vertical) in top-left-bottom
  \draw[orange!80, thick] (1,4) -- (1,6);
  \draw[orange!80, thick] (3,4) -- (3,6);
  % binary split (vertical) in top-right
  \draw[blue!60, thick] (6,4) -- (6,8);
  % ternary horizontal in top-right-left
  \draw[orange!80, thick] (4,5.33) -- (6,5.33);
  \draw[orange!80, thick] (4,6.67) -- (6,6.67);
  % bottom-left: binary horizontal
  \draw[blue!60, thick] (0,2) -- (4,2);
  % deeper ternary in bottom-left-top
  \draw[orange!80, thick] (1,2) -- (1,4);
  \draw[orange!80, thick] (3,2) -- (3,4);
  % bottom-right: mostly large
  \draw[blue!60, thick] (6,0) -- (6,4);
  \draw[blue!60, thick] (4,2) -- (8,2);
  % Scale label
  \node[font=\small, anchor=north] at (4, -0.3) {128 pixels};
\end{scope}

\end{tikzpicture}
\caption{Block partitioning comparison across codec generations. H.264 is
  limited to $16\!\times\!16$ macroblocks with fixed sub-partitions (left).
  VP9/HEVC use quad-tree partitioning within $64\!\times\!64$ superblocks
  (center-left). AV1 extends to $128\!\times\!128$ with quad-tree plus
  rectangular splits (center-right). VVC adds binary and ternary splits for
  maximum flexibility (right). Blue lines show first-level splits; red
  shows deeper partitioning; orange shows asymmetric/ternary splits.}
\label{fig:block-partitioning}
\end{figure}

\begin{table}[ht]
\centering
\small
\renewcommand{\arraystretch}{1.4}
\begin{tabularx}{\textwidth}{lXXXXX}
\toprule
\textbf{Feature} & \textbf{H.264} & \textbf{VP9} & \textbf{HEVC}
  & \textbf{AV1} & \textbf{VVC} \\
\midrule
Max block &
  $16\!\times\!16$ &
  $64\!\times\!64$ &
  $64\!\times\!64$ &
  $128\!\times\!128$ &
  $128\!\times\!128$ \\
Min block &
  $4\!\times\!4$ &
  $4\!\times\!4$ &
  $8\!\times\!8$ &
  $4\!\times\!4$ &
  $4\!\times\!4$ \\
Split types &
  Fixed partitions &
  Quad-tree &
  Quad-tree &
  Quad $+$ rect &
  Quad $+$ BT $+$ TT \\
Asymmetric &
  Limited &
  No &
  No &
  4:1 rectangles &
  Full (BT/TT) \\
\bottomrule
\end{tabularx}
\caption{Block partitioning capabilities by codec. BT = binary tree
  (horizontal or vertical), TT = ternary tree (1:2:1 split). Larger maximum
  blocks reduce overhead for homogeneous regions; asymmetric splits improve
  partitioning efficiency at content boundaries.}
\label{tab:partitioning}
\end{table}

%% ---------------------------------------------------------------------------
\subsection{Rate-Distortion Operating Points by Codec}
\label{sec:rd-comparison}

The codec-specific analyses of
Sections~\ref{sec:mpeg-lineage}--\ref{sec:aom-lineage} can be unified
through the Bj\o ntegaard Delta Rate (BD-Rate), the standard metric for
comparing codec efficiency. BD-Rate computes the average percentage rate
difference between two codecs at equivalent distortion by integrating the
difference between their rate-distortion curves fit as cubic polynomials in
PSNR-vs-log-rate space:
\begin{equation}
  \text{BD-Rate} = \frac{1}{\Delta\text{PSNR}}
    \int_{\text{PSNR}_{\min}}^{\text{PSNR}_{\max}}
    \left[ \ln R_A(p) - \ln R_B(p) \right] dp
  \label{eq:bd-rate}
\end{equation}
A negative BD-Rate indicates that codec~$A$ requires fewer bits than
codec~$B$ at the same quality.

\begin{figure}[ht]
\centering
\begin{tikzpicture}[>=Stealth]

% --- Axes ---
\draw[->, thick] (0, 0) -- (11, 0);
\draw[->, thick] (0, 0) -- (0, 6.5);

% --- Axis titles ---
\node[font=\small] at (5.5, -0.85) {Rate [Mbps]};
\node[font=\small, rotate=90] at (-0.9, 3.25) {PSNR [dB]};

% --- X-axis labels ---
\foreach \x/\lab in {0/0, 2/0.5, 4/1.0, 6/1.5, 8/2.0, 10/2.5} {
  \draw (\x, -0.1) -- (\x, 0.1);
  \node[font=\tiny, anchor=north] at (\x, -0.15) {\lab};
}

% --- Y-axis labels ---
\foreach \y/\lab in {0/28, 1/30, 2/32, 3/34, 4/36, 5/38, 6/40} {
  \draw (-0.1, \y) -- (0.1, \y);
  \node[font=\tiny, anchor=east] at (-0.15, \y) {\lab};
}

% --- Grid ---
\foreach \y in {1,2,3,4,5,6} {
  \draw[gray!15] (0, \y) -- (10.5, \y);
}

% --- Shannon R(D) bound (leftmost curve) ---
\draw[very thick, green!60!black]
  (0.4, 0.5) .. controls (0.6, 2) and (0.8, 3.5) ..
  (1.2, 4.5) .. controls (1.8, 5.5) and (2.5, 6.0) ..
  (3.5, 6.3);
\node[font=\scriptsize, green!60!black, anchor=south west] at (2.5, 6.1)
  {$R(D)$};

% --- VVC curve ---
\draw[very thick, violet!70]
  (0.8, 0.3) .. controls (1.2, 2.0) and (1.6, 3.5) ..
  (2.2, 4.5) .. controls (3.0, 5.5) and (4.0, 6.0) ..
  (5.5, 6.3);
\node[font=\scriptsize, violet!70, anchor=south west] at (4.2, 6.1)
  {VVC};

% --- AV1 curve ---
\draw[very thick, orange!80!black]
  (1.0, 0.2) .. controls (1.5, 1.8) and (2.2, 3.3) ..
  (3.0, 4.3) .. controls (4.0, 5.3) and (5.2, 5.9) ..
  (7.0, 6.2);
\node[font=\scriptsize, orange!80!black, anchor=south] at (6.0, 6.1)
  {AV1};

% --- HEVC curve ---
\draw[very thick, teal!70]
  (1.2, 0.1) .. controls (1.8, 1.6) and (2.6, 3.1) ..
  (3.6, 4.1) .. controls (4.8, 5.1) and (6.2, 5.8) ..
  (8.0, 6.1);
\node[font=\scriptsize, teal!70, anchor=south] at (7.2, 6.0)
  {HEVC};

% --- VP9 curve ---
\draw[very thick, blue!70]
  (1.4, 0.0) .. controls (2.0, 1.4) and (3.0, 2.9) ..
  (4.2, 4.0) .. controls (5.5, 5.0) and (7.0, 5.7) ..
  (9.0, 6.0);
\node[font=\scriptsize, blue!70, anchor=south west] at (8.2, 5.9)
  {VP9};

% --- H.264 curve (rightmost) ---
\draw[very thick, red!70!black]
  (2.0, 0.0) .. controls (3.0, 1.2) and (4.2, 2.6) ..
  (5.5, 3.7) .. controls (7.0, 4.7) and (8.5, 5.4) ..
  (10.2, 5.8);
\node[font=\scriptsize, red!70!black, anchor=south west] at (9.5, 5.6)
  {H.264};

% --- 37 dB reference line ---
\draw[thin, gray, dashed] (0, 4.5) -- (10.5, 4.5);
\node[font=\tiny, gray, anchor=west] at (10.6, 4.5) {37\,dB};

% --- Operating points at 37 dB ---
\fill[red!70!black] (5.2, 4.5) circle (2.5pt);
\fill[blue!70] (3.9, 4.5) circle (2.5pt);
\fill[teal!70] (3.3, 4.5) circle (2.5pt);
\fill[orange!80!black] (2.7, 4.5) circle (2.5pt);
\fill[violet!70] (2.0, 4.5) circle (2.5pt);
\fill[green!60!black] (1.1, 4.5) circle (2.5pt);

% --- Annotation ---
\node[font=\scriptsize\itshape, text=black!60, align=center,
  anchor=north] at (5.5, -1.1)
  {At 37~dB PSNR (720p ``good'' quality), each generation
  requires fewer bits.\\
  The gap to the Shannon bound $R(D)$ narrows but never closes.};

\end{tikzpicture}
\caption{Rate-distortion curves for six codecs at 720p/30fps, with the
  Shannon $R(D)$ bound (green). At constant quality (37~dB dashed line),
  the rate progression from H.264 through VVC illustrates the BD-Rate
  savings in Table~\ref{tab:bd-rate}. This figure extends
  Figure~\ref{fig:rate-distortion} from Appendix~\ref{app:info-theory}
  with the full codec lineage.}
\label{fig:rd-curves-detailed}
\end{figure}

\begin{table}[ht]
\centering
\small
\renewcommand{\arraystretch}{1.4}
\begin{tabularx}{\textwidth}{lXXXX}
\toprule
\textbf{Codec} & \textbf{BD-Rate vs.\ H.264}
  & \textbf{$\Delta$ (720p)} & \textbf{Family} & \textbf{Year} \\
\midrule
VP8  & ${\sim}$0\%    & ${\sim}$2.1 & Google    & 2010 \\
VP9  & $-$35\%        & ${\sim}$0.8 & AOM       & 2013 \\
HEVC & $-$40\%        & ${\sim}$0.7 & MPEG/ITU  & 2013 \\
AV1  & $-$54\%        & ${\sim}$0.5 & AOM       & 2018 \\
VVC  & $-$63\%        & ${\sim}$0.3 & MPEG/ITU  & 2020 \\
\bottomrule
\end{tabularx}
\caption{BD-Rate savings and coding efficiency gap $\Delta$
  (eq.~\ref{eq:coding-gap}) for each codec relative to H.264 at 720p/30fps.
  The two families (MPEG/ITU and Google/AOM) achieve comparable efficiency
  within each generation; AV1 slightly trails HEVC, VVC leads AV1 by
  ${\sim}$20\%.}
\label{tab:bd-rate}
\end{table}

Table~\ref{tab:bd-rate} provides the DSP-grounded explanation for the codec
ordering in eq.~(\ref{eq:codec-ordering}) and the savings figures in
eqs.~(\ref{eq:vp9-savings})--(\ref{eq:av1-savings}): each generation adds
tools that reduce $\Delta_{\text{pred}}$, $\Delta_{\text{transform}}$, and
$\Delta_{\text{quant}}$ in (\ref{eq:delta-decomp}), with diminishing
returns as the Shannon bound is approached.

%% ---------------------------------------------------------------------------
\subsection{Scalable Video Coding and Simulcast}
\label{sec:svc-simulcast}

Multi-participant WebRTC sessions require adaptive quality: different
receivers have different bandwidths, and a single encode rate cannot serve
all. Two architectures address this.

\subsubsection{Simulcast}
\label{sec:simulcast}

Simulcast encodes the source at $N$ independent resolution/bitrate
combinations:
\begin{equation}
  R_{\text{simulcast}} = \sum_{i=1}^{N} R_i
  \label{eq:simulcast-rate}
\end{equation}
Each encoding is a fully independent bitstream. A Selective Forwarding
Unit (SFU) selects which stream to forward to each receiver based on
reported bandwidth. The overhead is substantial: a typical 3-layer
simulcast (720p/1.5~Mbps, 360p/500~kbps, 180p/150~kbps) requires
2.15~Mbps total---approximately 1.4$\times$ the cost of the highest
layer alone---with no inter-layer compression benefit.

\subsubsection{Scalable Video Coding (SVC)}
\label{sec:svc}

SVC produces a layered bitstream where a base layer $L_0$ is independently
decodable and enhancement layers $L_1, \ldots, L_K$ refine quality
progressively. Three scalability dimensions exist:

\textbf{Temporal scalability.} The base layer encodes at a reduced frame
rate (e.g., 7.5~fps); enhancement layers add intermediate frames:
\begin{equation}
  R_{\text{temporal}} = R_0 + \sum_{k=1}^{K} \Delta R_k
  \label{eq:temporal-svc}
\end{equation}
where $\Delta R_k \ll R_0$ because enhancement frames reference lower
layers and carry primarily motion-compensated residuals. An SFU drops
enhancement layers for bandwidth-constrained receivers, gracefully
degrading frame rate while preserving spatial quality.

\textbf{Spatial scalability.} The base layer encodes at reduced resolution;
enhancement layers add detail via inter-layer prediction:
\begin{equation}
  R_{\text{spatial}} = R_0 + \Delta R_{\text{upscale}}
  \label{eq:spatial-svc}
\end{equation}
where $\Delta R_{\text{upscale}}$ exploits correlation between the
upsampled base layer and the full-resolution source.

\textbf{Quality (SNR) scalability.} Refinement of quantization at the same
resolution, encoded as the difference between coarse and fine
reconstructions.

SVC achieves 20--30\% rate savings over simulcast for equivalent
adaptivity because enhancement layers exploit inter-layer prediction:
\begin{equation}
  R_{\text{SVC}} < R_{\text{simulcast}} \quad
  \text{(typically by 20--30\%)}
  \label{eq:svc-advantage}
\end{equation}

\begin{figure}[ht]
\centering
\begin{tikzpicture}[>=Stealth, scale=0.85, every node/.style={scale=0.85},
  layer/.style={draw, rounded corners=2pt, minimum width=1.8cm,
    minimum height=0.6cm, font=\scriptsize, align=center, thick}]

% --- Simulcast (left) ---
\node[font=\normalsize\bfseries] at (2, 5.8) {Simulcast};

% Three independent columns
\node[layer, fill=red!15] (s720) at (0.5, 4.5) {720p\\1.5 Mbps};
\node[layer, fill=blue!15] (s360) at (2.0, 4.5) {360p\\500 kbps};
\node[layer, fill=green!15] (s180) at (3.5, 4.5) {180p\\150 kbps};

\node[layer, fill=gray!10] (sfu1) at (2.0, 2.5) {SFU};

\draw[->, thick] (s720) -- (sfu1);
\draw[->, thick] (s360) -- (sfu1);
\draw[->, thick] (s180) -- (sfu1);

\node[font=\scriptsize, text=red!60!black] at (2.0, 1.5)
  {Total: 2.15 Mbps};
\node[font=\tiny\itshape, text=black!50] at (2.0, 1.0)
  {No inter-layer benefit};

% --- SVC (right) ---
\node[font=\normalsize\bfseries] at (8.5, 5.8) {SVC};

% Pyramid
\node[layer, fill=green!15, minimum width=2.0cm] (t0) at (8.5, 2.0)
  {$T_0$: Base\\7.5 fps};
\node[layer, fill=blue!15, minimum width=2.8cm] (t1) at (8.5, 3.2)
  {$T_1$: Enhancement\\15 fps};
\node[layer, fill=red!15, minimum width=3.6cm] (t2) at (8.5, 4.4)
  {$T_2$: Enhancement\\30 fps};

% Dependency arrows
\draw[->, thick, green!50!black] (t0) -- (t1);
\draw[->, thick, blue!50!black] (t1) -- (t2);

% Drop annotations
\draw[thick, dashed, red!40] (6.2, 3.8) -- (10.8, 3.8);
\node[font=\tiny\itshape, text=red!50!black, anchor=west] at (10.9, 3.8)
  {SFU can drop};

\node[font=\scriptsize, text=green!50!black] at (8.5, 1.0)
  {Total: ${\sim}$1.7 Mbps};
\node[font=\tiny\itshape, text=black!50] at (8.5, 0.5)
  {20--30\% savings via};
\node[font=\tiny\itshape, text=black!50] at (8.5, 0.1)
  {inter-layer prediction};

\end{tikzpicture}
\caption{Simulcast (left) encodes three independent bitstreams with no
  inter-layer compression. SVC (right) encodes a layered bitstream where
  temporal enhancement layers ($T_1$, $T_2$) reference the base layer
  ($T_0$), achieving 20--30\% savings. The SFU drops enhancement layers
  for bandwidth-constrained receivers.}
\label{fig:svc-architecture}
\end{figure}

Codec-specific SVC support relevant to WebRTC:

\begin{itemize}
\item \textbf{H.264 SVC (Annex~G).} Defines temporal, spatial, and quality
  scalability. Complex to implement; rarely used in WebRTC deployments.

\item \textbf{VP9 SVC.} Temporal scalability is widely deployed in WebRTC
  via libvpx (up to 3~temporal layers). Spatial SVC is defined but
  implementation-limited. VP9 temporal SVC is the primary mechanism
  enabling adaptive quality in Chrome/Firefox SFU-based calls.

\item \textbf{AV1 SVC.} Temporal scalability supported in both libaom and
  SVT-AV1. A key enabler for AV1 adoption in production WebRTC
  infrastructure, as SFUs require layered bitstreams to serve heterogeneous
  receivers.
\end{itemize}

Safari's lack of simulcast and SVC support in 2021 (Section~3.1) meant
that Safari users in multi-participant calls could not benefit from adaptive
quality---every receiver got the same stream regardless of bandwidth,
degrading the experience for all participants.

%% ---------------------------------------------------------------------------
\subsection{Hardware vs.\ Software Encoding}
\label{sec:hw-sw}

The main paper identifies hardware encoding as the deployment gate for
AV1 (Sections~8.5, 9.5). This section quantifies why.

\subsubsection{Encoding Complexity}

The computational cost of encoding a video frame is dominated by motion
estimation and RD-optimized mode decision. We express complexity as
operations per pixel:
\begin{equation}
  C_{\text{enc}} = C_{\text{ME}} + C_{\text{transform}}
    + C_{\text{quant}} + C_{\text{entropy}}
    + C_{\text{mode}}
  \label{eq:enc-complexity}
\end{equation}

Representative values for real-time presets:

\begin{itemize}
\item H.264 Constrained Baseline (``veryfast''): $C_{\text{enc}} \approx
  200$~ops/pixel
\item VP9 (speed~8, real-time): $C_{\text{enc}} \approx 600$~ops/pixel
\item AV1 (speed~10, real-time): $C_{\text{enc}} \approx 2000$~ops/pixel
\end{itemize}

The real-time constraint requires:
\begin{equation}
  C_{\text{enc}} \times W \times H \times \text{fps}
    \leq F_{\text{available}}
  \label{eq:realtime-constraint}
\end{equation}
where $F_{\text{available}}$ is the available computational throughput
(operations per second). For 720p/30fps AV1 at 2000~ops/pixel:
\begin{equation}
  2000 \times 921{,}600 \times 30 = 55.3 \times 10^9 \;\text{ops/s}
  \label{eq:av1-cost}
\end{equation}
This consumes approximately one full core of a modern desktop CPU. On
mobile devices with thermal and power constraints, software AV1 encoding
at 720p/30fps remains impractical without hardware acceleration.

\subsubsection{Hardware Encoder Architecture}

Hardware encoders implement a fixed subset of the codec's tool set in
dedicated silicon. The trade-off is quality for throughput: hardware
encoders typically produce output 10--20\% larger than software encoders
at equivalent visual quality because they implement fewer mode decision
candidates and simpler motion search:
\begin{equation}
  \Delta_{\text{HW}} \approx (1.1 \text{--} 1.2)
    \cdot \Delta_{\text{SW}}
  \label{eq:hw-quality-gap}
\end{equation}

This quality penalty is acceptable for real-time communication because:
(a)~the alternative---software encoding at a lower speed preset or reduced
resolution---incurs a larger quality penalty, and (b)~the freed CPU
capacity can be used for pre/post-processing, noise reduction, or
additional application logic.

\subsubsection{Hardware Encoder Availability}

\begin{table}[ht]
\centering
\small
\renewcommand{\arraystretch}{1.4}
\begin{tabularx}{\textwidth}{lXXXX}
\toprule
\textbf{Platform} & \textbf{H.264 HW} & \textbf{VP9 HW encode}
  & \textbf{AV1 HW encode} & \textbf{AV1 year} \\
\midrule
Intel (QSV) &
  Sandy Bridge+ (2011) &
  Kaby Lake+ (2017) &
  Arc / Meteor Lake+ &
  2022 \\
NVIDIA (NVENC) &
  Kepler+ (2012) &
  --- &
  Ada Lovelace+ &
  2022 \\
AMD (VCN) &
  VCN~1.0+ (2017) &
  --- &
  VCN~4.0+ (RDNA~3) &
  2023 \\
Apple (VideoToolbox) &
  A4+ (2010) &
  --- &
  M3+ &
  2023 \\
Qualcomm &
  SD~820+ (2016) &
  --- &
  SD~8 Gen~2+ &
  2023 \\
\bottomrule
\end{tabularx}
\caption{Hardware encoder availability timeline. H.264 hardware encoding
  has been ubiquitous for over a decade. AV1 hardware encoding began
  shipping in 2022--2023 across all major platforms; the main paper projects
  it becomes the practical WebRTC default by 2028--2030 as these chipsets
  reach market saturation.}
\label{tab:hw-encoders}
\end{table}

\excursus{Excursus: WebCodecs and hardware acceleration.}{When a WebCodecs
\texttt{VideoEncoder} is configured with \texttt{codec: "av01"}, the browser
selects hardware or software encoding based on platform capability and the
\texttt{hardwareAcceleration} preference (\texttt{"prefer-hardware"},
\texttt{"prefer-software"}, or \texttt{"no-preference"}). The application
receives encoded \texttt{EncodedVideoChunk} objects either way, but latency
and CPU cost differ dramatically: hardware encoding typically adds
$<\!1$~ms per frame and consumes near-zero CPU, while software AV1 at
real-time presets consumes an entire CPU core
(eq.~\ref{eq:av1-cost}). This abstraction is what makes the hardware
timeline in Table~\ref{tab:hw-encoders} directly relevant to the WebCodecs
adoption curve in Section~7.}

The hardware encoding timeline directly underpins the main paper's
projections: the ``AV1 becomes practical'' milestone (Section~8.5) and
``AV1 HW ubiquitous'' milestone (Section~9.5) depend on the chipset
generations in Table~\ref{tab:hw-encoders} reaching market saturation.

%% ---------------------------------------------------------------------------
\subsection{Implications for Real-Time Web Communication}
\label{sec:video-implications}

The DSP analysis of this appendix maps back to concrete implications for
the technologies surveyed in the main text.

\textbf{The hybrid coding framework is universal; innovation is
incremental.} Every codec from VP8 to VVC implements the same
prediction--transform--quantization--entropy framework
(Figure~\ref{fig:hybrid-encoder}). Generational improvements come from
expanding the tool set within this framework---more block sizes
(Table~\ref{tab:partitioning}), more prediction modes, more transform
types (Table~\ref{tab:transforms})---not from fundamentally new
architectures. This means the efficiency gap $\Delta$ approaches but never
reaches zero: each additional tool yields diminishing rate-distortion
improvement at increasing computational cost.

\textbf{Transform and prediction innovations explain the BD-Rate
progression.} The 35\% rate saving from H.264 to VP9 comes primarily from
larger blocks ($64\!\times\!64$ vs.\ $16\!\times\!16$) and adaptive
transforms (DCT$+$ADST vs.\ fixed DCT). The additional 30\% from VP9 to
AV1 comes from compound and warped motion, film grain synthesis, and
16-way transform type selection (Table~\ref{tab:bd-rate}). Each tool adds
encoding complexity proportional to its rate-distortion benefit, which is
why real-time encoders must aggressively prune the tool set.

\textbf{SVC is what makes multi-participant WebRTC work.} Without temporal
SVC, an SFU must forward full-rate streams to all participants regardless
of their bandwidth. With VP9/AV1 temporal SVC
(Section~\ref{sec:svc}), the SFU drops enhancement temporal layers for
bandwidth-constrained receivers, degrading frame rate gracefully while
preserving spatial quality. The 20--30\% savings over simulcast
(eq.~\ref{eq:svc-advantage}) translate directly to reduced server egress
bandwidth in large meetings.

\textbf{Hardware encoding is the deployment gate, not codec specification.}
AV1 was specified in 2018 but becomes a practical WebRTC default only
when hardware encoders are ubiquitous (${\sim}$2028--2030). The DSP tools
that make AV1 efficient---warped motion, OBMC, $128\!\times\!128$
superblocks with flexible partitioning---are precisely the tools that make
it computationally expensive in software (eq.~\ref{eq:av1-cost}). The
${\sim}$10--20\% quality penalty of hardware encoders
(eq.~\ref{eq:hw-quality-gap}) is the price of real-time feasibility, and
it is a price worth paying: $\Delta_{\text{AV1,HW}} \approx 0.6$ is still
far below $\Delta_{\text{VP9,SW}} \approx 0.8$.

\textbf{The convergence toward Shannon is real but asymptotic.}
Cross-referencing with Appendix~\ref{app:info-theory}: H.264's
$\Delta \approx 2.1$ means it operates at $3.1\times$ the theoretical
minimum rate. VP9's $\Delta \approx 0.8$ means $1.8\times$. AV1's
$\Delta \approx 0.5$ means $1.5\times$. VVC's $\Delta \approx 0.3$ means
$1.3\times$. Each generation captures more spatial, temporal, and
perceptual redundancy, but the returns are diminishing: the distance from
H.264 to VP9 ($-$35\%) exceeds VP9 to AV1 ($-$30\%), which exceeds AV1 to
VVC ($-$17\%). The next generation of efficiency gains for real-time web
communication will come not from codecs alone, but from the co-optimization
of transport (QUIC, WebTransport), application architecture (WebCodecs,
unbundled pipelines), and network infrastructure that the main text
describes.
